\documentclass[10pt, letterpaper]{article}

% Inhaltsverzeichnis für Pakettypen (nur für Übersicht im Header, wird nicht im Dokument angezeigt)
% 1. Seitenlayout und Ränder
% 2. Sprache und Zeichensatz
% 3. Mathematik und Theorem-Umgebungen
% 4. Eigene Makros
% 5. Diagramme und Grafiken
% 6. Tabellen und Aufzählungen
% 7. Inhaltsverzeichnis
% 8. Abschnittsüberschriften
% 9. Abstrakt-Umgebung
% 10. Todos/Notizen
% 11. Rahmen/Box-Umgebungen
% 12. Python-Integration
% 13. Literaturverwaltung
% 14. Hyperlinks
% 15. Absatzeinstellungen
% 16. Umgebungen
% 17  Graphik
% 18  Extra
% 00. Titel und Autor

\usepackage{slashed}

% --- 1. Seitenlayout und Ränder ---
\usepackage[margin=3cm]{geometry}

% --- 2. Sprache und Zeichensatz ---
\usepackage[english]{babel}
\usepackage[T1]{fontenc}
\usepackage[utf8]{inputenc}

% --- 3. Mathematik und Theorem-Umgebungen ---
\usepackage{amsmath, amssymb, amsthm}
\usepackage{mathrsfs}
\DeclareMathOperator{\WF}{WF}

% --- 4. Eigene Makros ---
\usepackage{xcolor}
\newcommand{\SKP}{\langle\cdot,\cdot\rangle}
\newcommand{\id}{\operatorname{id}}
\newcommand{\sgn}{\operatorname{sgn}}
\newcommand{\Ad}{\operatorname{Ad}}
\newcommand{\ad}{\operatorname{ad}}
\newcommand{\R}{\mathbb{R}}
\newcommand{\N}{\mathbb{N}}
\newcommand{\Q}{\mathbb{Q}}
\newcommand{\Z}{\mathbb{Z}}
\newcommand{\C}{\mathbb{C}}
\newcommand{\QU}{\mathbb{H}}
\newcommand{\Cinf}{C^{\infty}}
\newcommand{\Mod}{\mathrm{Mod}}
\newcommand{\Alg}{\mathrm{Alg}}
\newcommand{\Diff}{\mathrm{Diff}}
\newcommand{\Iso}{\mathrm{Iso}}
\newcommand{\Aut}{\mathrm{Aut}}
\newcommand{\End}{\mathrm{End}}
\newcommand{\Hom}{\mathrm{Hom}}
\newcommand{\Cl}{\mathrm{Cl}}
\newcommand{\entwurf}[1]{\textcolor{red}{#1}}
\newcommand{\GL}{\mathrm{GL}}
\newcommand{\Mat}{\mathrm{Mat}}
\newcommand{\SL}{\mathrm{SL}}
\newcommand{\SO}{\mathrm{SO}}
\newcommand{\SU}{\mathrm{SU}}
\newcommand{\Sp}{\mathrm{Sp}}
\newcommand{\Spin}{\mathrm{Spin}}
\newcommand{\Pin}{\mathrm{Pin}}
\newcommand{\Lor}{\mathrm{Lor}}
\newcommand{\Ogroup}{\mathrm{O}}
\newcommand{\U}{\mathrm{U}}

% --- 5. Diagramme und Grafiken ---
\usepackage{graphicx}
\graphicspath{{../../Library/Images/}}
\usepackage[export]{adjustbox}
\usepackage{tikz}
\usetikzlibrary{decorations.pathreplacing, arrows.meta, positioning}
\usepackage{tikz-cd}

% --- 6. Tabellen und Aufzählungen ---
\usepackage{enumitem}
\setlist[itemize]{left=0.5cm}

\newenvironment{romanenum}[1][]
  {%
    \ifx&#1&
    \else
      \textbf{#1}\quad
    \fi
    \begin{enumerate}[label=\roman*)]
  }
  {%
    \end{enumerate}%
  }

% --- 7. Inhaltsverzeichnis ---
\usepackage{tocloft}
\renewcommand{\cftsecfont}{\footnotesize}
\renewcommand{\cftsubsecfont}{\footnotesize}
\renewcommand{\cftsubsubsecfont}{\footnotesize}
\renewcommand{\cftsecpagefont}{\footnotesize}
\renewcommand{\cftsubsecpagefont}{\footnotesize}
\renewcommand{\cftsubsubsecpagefont}{\footnotesize}
\usepackage{etoc}

% --- 8. Abschnittsüberschriften ---
\usepackage{titlesec}
\titleformat{\section}{\normalfont\large\bfseries}{\thesection}{1em}{}
\titleformat{\subsection}{\normalfont\normalsize\bfseries}{\thesubsection}{0.5em}{}
\titleformat{\subsubsection}{\normalfont\normalsize\bfseries}{\thesubsubsection}{0.5em}{}
\setcounter{secnumdepth}{4}

% --- 9. Abstrakt-Umgebung ---
\usepackage{changepage}
\renewenvironment{abstract}
  {
    \begin{adjustwidth}{1.5cm}{1.5cm}
    \small
    \textsc{Abstract. –}%
  }
  {
    \end{adjustwidth}
  }

% --- 10. Todos/Notizen ---
\usepackage{todonotes}
% Hellblau definieren
\definecolor{hellblau}{rgb}{0.3,0.6,1}
\newenvironment{note}{\par\begingroup\color{hellblau}\itshape}{\par\endgroup}

% --- 11. Rahmen/Box-Umgebungen ---
\usepackage{mdframed}
\usepackage{tcolorbox}
\colorlet{shadecolor}{gray!25}

\newenvironment{customTheorem}
  {\vspace{10pt}%
   \begin{mdframed}[
     backgroundcolor=gray!20,
     linewidth=0pt,
     innertopmargin=10pt,
     innerbottommargin=10pt,
     skipabove=\dimexpr\topsep+\ht\strutbox\relax,
     skipbelow=\topsep,
   ]}
  {\end{mdframed}
   \vspace{10pt}%
  }

% --- 12. Python-Integration ---
% (Deaktiviert in dieser Version, aktiviere bei Bedarf)
% \usepackage{pythontex}
% \usepackage[makestderr]{pythontex}

% --- 13. Literaturverwaltung ---
\usepackage{csquotes}
\usepackage[backend=biber, style=alphabetic, citestyle=alphabetic]{biblatex}
\addbibresource{bibliography.bib}

% --- 14. Hyperlinks ---
\usepackage{hyperref}
\hypersetup{
  colorlinks   = true,
  urlcolor     = blue,
  linkcolor    = blue,
  citecolor    = blue,
  frenchlinks  = true
}

% --- 15. Absatzeinstellungen ---
\usepackage[parfill]{parskip}
\sloppy

% --- 16. Umgebungen ---
\usepackage{thmtools}

\newcommand{\CustomHeading}[3]{%
  \par\medskip\noindent%
  \textbf{#1 #2} \textnormal{(#3)}.\enskip%
}

\newenvironment{DEF}[2]{\begin{unitbox}\CustomHeading{Definition}{#1}{#2}}{\end{unitbox}}
\newenvironment{PROP}[2]{\begin{unitbox}\CustomHeading{Proposition}{#1}{#2}}{\end{unitbox}}
\newenvironment{THEO}[2]{\begin{unitbox}\CustomHeading{Theorem}{#1}{#2}}{\end{unitbox}}
\newenvironment{LEM}[2]{\begin{unitbox}\CustomHeading{Lemma}{#1}{#2}}{\end{unitbox}}
\newenvironment{KORO}[2]{\begin{unitbox}\CustomHeading{Corollar}{#1}{#2}}{\end{unitbox}}
\newenvironment{REM}[2]{\begin{unitbox}\CustomHeading{Remark}{#1}{#2}}{\end{unitbox}}
\newenvironment{EXA}[2]{\begin{unitbox}\CustomHeading{Example}{#1}{#2}}{\end{unitbox}}
\newenvironment{STUD}[2]{\begin{unitbox}\CustomHeading{Study}{#1}{#2}}{\end{unitbox}}
\newenvironment{CONC}[2]{\begin{unitbox}\CustomHeading{Concept}{#1}{#2}}{\end{unitbox}}
\newenvironment{OTH}[2]{\begin{unitbox}\CustomHeading{Other}{#1}{#2}}{\end{unitbox}}
\newenvironment{EXE}[2]{\begin{unitbox}\CustomHeading{Exercise}{#1}{#2}}{\end{unitbox}}
\newenvironment{MOT}[2]{\begin{unitbox}\CustomHeading{Motivation}{#1}{#2}}{\end{unitbox}}
\newenvironment{PROOF}[2]{\begin{unitbox}\CustomHeading{Proof}{#1}{#2}}{\end{unitbox}}

% --- Unit Umgebung für Source-Inhalte ---
\usepackage{mdframed}
\newmdenv[
  linewidth=1pt,
  topline=false,
  bottomline=false,
  rightline=false,
  leftmargin=0cm,
  rightmargin=0cm,
  skipabove=10pt,
  skipbelow=10pt,
  innertopmargin=0.5\baselineskip,
  innerbottommargin=0.5\baselineskip,
  backgroundcolor=gray!10,
  linecolor=gray
]{unitbox}

\newenvironment{unit}[1]
  {\begin{unitbox}\textbf{Unit #1}\par\smallskip}
  {\end{unitbox}}

% --- 17. ... ---

% --- 18. Extras ---
\usepackage{stmaryrd}
\usepackage{bbold}  % falls du \mathbb{1} nutzen willst

% --- 00. Titel und Autor ---
\title{Mein Titel}
\author{Tim Jaschik}
\date{\today}

\begin{document}

\maketitle
\rule{\textwidth}{0.5pt}
\begin{abstract}
Kurze Beschreibung …
\end{abstract}
\rule{\textwidth}{0.5pt}
\vspace{0.5cm}

\tableofcontents

\pagebreak

\section*{Prologue}
\begin{displayquote}
Men can do nothing without the make-believe of a beginning. Even Science, the strict measurer, is obliged to start with a make-believe unit, and must fix on a point in the stars' unceasing journey when his sidereal clock shall pretend time is at Nought. His less accurate grandmother Poetry has always been understood to start in the middle, but on reflection it appears that her proceeding is not very different from his; since Science, too, reckons backwards as well as forwards, divides his unit into billions, and with his clock-finger at Nought really sets off in medias res. No retrospect will take us to the true beginning; and whether our prologue be in heaven or on earth, it is but a fraction of the all-presupposing fact with which our story sets out.\\
-George Eliot, Daniel Deronda (epigraph of Chapter 1)
\end{displayquote}

The purpose of this preliminary chapter is to present some notation and terminology, discuss a few matters on which mathematicians' and physics usages differ, and provide some basic definitions and formulas from physics and Lie theory that will be needed throughout the book.

\subsection*{Linguistic prologue: notation and terminology}
Hilbert space and its operators. Quantum mechanicians and mathematical analysts both spend a lot of time in Hilbert space, but they have different ways of speaking about it. In this book we generally follow physicists' conventions, so it will be well to explain how to translate from one dialect to the other. (The physicists' dialect is largely due to Dirac, and its ubiquity is a testimony to the profound influence that his book [22] had on several generations of physicists.)

To begin with, complex conjugates are denoted by asterisks rather than overlines:

$$
x-i y=(x+i y)^{*} \text { rather than } \overline{x+i y} .
$$

The inner product on a Hilbert space is denoted by $\langle\cdot \mid \cdot\rangle$ and is taken to be linear in the second variable and conjugate-linear in the first:

$$
\langle u \mid a v+b w\rangle=a\langle u \mid v\rangle+b\langle u \mid w\rangle, \quad\langle v \mid u\rangle=\langle u \mid v\rangle^{*}
$$

The Hermitian adjoint (conjugate transpose) of a matrix or an operator is denoted by a dagger rather than an asterisk:

$$
\text { if } A=\left(x_{j k}+i y_{j k}\right) \text {, then }\left(x_{k j}-i y_{k j}\right)=A^{\dagger} \text { rather than } A^{*} \text {. }
$$

(More about adjoints below.) These shifted uses of * and ${ }^{\dagger}$ may seem annoying at first, but one gets used to them - and there is another consideration. The overline is employed in physics not for conjugation but to denote the "Dirac adjoint" of a Dirac spinor, as we shall explain in §4.1. This usage turns up sufficiently frequently\\
that any attempt to permute the physicists' meanings of *, ${ }^{\dagger}$, and - would threaten mass confusion.

We shall, however, employ the standard mathematical notation for the norm on the Hilbert space: $\|u\|=\langle u \mid u\rangle^{1 / 2}$. We also denote the transpose of a complex matrix $A$ by $A^{T}$; for real matrices we shall generally use $A^{\dagger}$ instead.

Now, some more subtle matters. In mathematicians' dialect, $\langle u \mid v\rangle$ is the inner product of two vectors $u$ and $v$ in the Hilbert space $\mathcal{H}$. Moreover, if $u \in \mathcal{H}$, the map $\phi_{u}(v)=\langle u \mid v\rangle$ is a bounded linear functional on $\mathcal{H}$; the correspondence $u \leftrightarrow \phi_{u}$ gives a conjugate-linear identification of $\mathcal{H}$ with its dual $\mathcal{H}^{\prime}$, which mathematicians generally take for granted without employing any special notation for it. Physicists, on the other hand, distinguish between elements of $\mathcal{H}$ and elements of $\mathcal{H}^{\prime}$, which they respectively call ket vectors and bra vectors and denote by symbols of the form $|u\rangle$ and $\langle u|$. If $|u\rangle$ is a ket vector (what mathematicians might call simply $u$ ), the corresponding bra vector $\langle u|$ is the linear functional denoted by $\phi_{u}$ above. The action of the bra vector (linear functional) $\langle u|$ on the ket vector (element of $\mathcal{H})|v\rangle$ is the inner product (i.e., "bracket" or "bra-ket") $\langle u \mid v\rangle$.

In this system " $u$ " can be any sort of convenient label to identify the vector: either a name for the vector itself, like the mathematicians' $u$, or an index or set of indices that specify the vector within an indexed family. For example, if one is working with an operator with a set of multiplicity-one eigenvalues, one might denote a unit eigenvector for the eigenvalue $\lambda$ by $|\lambda\rangle$, or an eigenvector for the $n$th eigenvalue simply by $|n\rangle$. Likewise, joint eigenvectors for a pair of operators with eigenvalues $\lambda_{n}$ and $\mu_{j}$ might be denoted by $|n, j\rangle$; and so forth. (In most situations the ambiguity of a scalar factor of modulus one in the choice of eigenvector is of no importance, for reasons that will become clear in Chapter 3.) Mathematicians may find this convention uncomfortably informal, but its virtue lies in its flexibility and its ability to strip away inessential symbols. Mathematicians would typically denote the $n$th eigenvector by something like $u_{n}$, but the symbol $u$ is just a placeholder; it is the $n$ that carries the useful information, and the physicists' notation gives it the starring role it deserves.

Next, let $A$ be a linear operator ${ }^{1}$ on $\mathcal{H}$ : in mathematicians' notation, $A$ maps a vector $v$ to $A v$; and in physicists' notation, it maps $|v\rangle$ to $A|v\rangle$ : no surprises here. But the physicists' notation for the scalar product of $A|v\rangle$ with the bra vector $\langle u|$ is $\langle u| A|v\rangle$ :

$$
\text { physicists' }\langle u| A|v\rangle=\text { mathematicians' }\langle u \mid A v\rangle \text {. }
$$

Now, $\langle u| A|v\rangle$ can also be considered as the scalar product of a bra $\langle u| A$ with the ket $|v\rangle$, and in this way $A$ defines a linear operator on bra vectors. This operator is what the mathematicians call the adjoint of $A$ when $\mathcal{H}^{\prime}$ is identified with $\mathcal{H}$ - a point that can lead to some confusion if one does not remain alert. Recalling that the adjoint of $A$ is denoted by $A^{\dagger}$ and the conjugate of $a \in \mathbb{C}$ is denoted by $a^{*}$, we have

$$
\langle u| A|v\rangle=\langle v| A^{\dagger}|u\rangle^{*}
$$

On the theory that flexibility and concision are more important than consistency, we shall feel free to denote elements of the Hilbert space by either $u$ or $|u\rangle$, and we shall write either $\langle u| A|v\rangle$ or $\langle u \mid A v\rangle$ as convenience dictates.

\footnotetext{${ }^{1} A$ need not be bounded. The notion of adjoint for unbounded operators involves some technicalities, but they are beside the point here.
}By the way, we will often write a scalar multiple of the identity operator, $\lambda I$, simply as $\lambda$. This commonly used bit of shorthand will cause no confusion if readers will keep in mind that when they encounter a scalar where an operator or a matrix seems to be needed, they should tacitly insert an $I$.

Vectors, tensors, and derivatives. In this book we reserve boldface type almost exclusively to denote 3 -dimensional vectors related to $\mathbb{R}^{3}$ as a model for the physical space in which we live. (Exception: In §4.5 it is used to denote elements of tensor products of Hilbert spaces.) Elements of $\mathbb{R}^{n}$ for general $n$ are usually denoted by lower-case italic letters $(x, p, \ldots)$. We denote the $n$-tuple of partial derivatives $\left(\partial_{1}, \ldots, \partial_{n}\right)$ on $\mathbb{R}^{n}$ by $\nabla$, and we denote the Laplacian $\nabla \cdot \nabla=\sum \partial_{j}^{2}$ by $\nabla^{2}$ (rather than $\Delta$, for which we will have other uses).

When $n=4$, a variant of this notation will be used for calculations arising from relativistic mechanics. Four-dimensional space-time is taken to be $\mathbb{R}^{4}$ with coordinates

$$
x^{0}=c t, \quad x^{1}=x, \quad x^{2}=y, \quad x^{3}=z
$$

on $\mathbb{R}^{4}$, where $t$ represents time, $c$ the speed of light, and $(x, y, z)$ a set of Cartesian coordinates on physical space $\mathbb{R}^{3}$. Thus, a point $x \in \mathbb{R}^{4}$ may be written as $\left(x^{0}, \mathbf{x}\right)$ when it is important to separate the space and time components. The Lorentz inner product on $\mathbb{R}^{4}$ is the bilinear form $\Lambda$ defined by

$$
\Lambda(x, y)=x^{0} y^{0}-x^{1} y^{1}-x^{2} y^{2}-x^{3} y^{3}
$$

and $\mathbb{R}^{4}$ equipped with this form is called Minkowski space. It is common in the physics literature to denote the vector whose components are $x^{0}, \ldots, x^{3}$ by $x^{\mu}$ rather than $x$. This is the same sort of harmless abuse of language that is involved in speaking of "the sequence $a_{n}$ " or "the function $f(x)$ "; we shall use it when it seems convenient.

We shall generally use classical tensor notation for vectors and tensors associated to Minkowski space. In particular, the Lorentz form is defined by the matrix

$$
g_{\mu \nu}=g^{\mu \nu}=\operatorname{diag}(1,-1,-1,-1)
$$

We employ the Einstein summation convention: in any product of vectors and tensors in which an index appears once as a subscript and once as a superscript, that index is to be summed from 0 to 3 . Thus, for example,

$$
\Lambda(x, y)=g_{\mu \nu} x^{\mu} y^{\nu}
$$

We shall also adopt the convention of using the matrix $g$ to "raise or lower indices":

$$
x_{\mu}=g_{\mu \nu} x^{\nu}, \quad p^{\mu}=g^{\mu \nu} p_{\nu}
$$

the practical effect of which is to change the sign of the last three components. (Strictly speaking, vectors whose components are denoted by subscripts should be construed as elements of the dual space $\left(\mathbb{R}^{4}\right)^{\prime}$, and the map $x^{\mu} \mapsto x_{\mu}$ is the isomorphism of $\mathbb{R}^{4}$ with $\left(\mathbb{R}^{4}\right)^{\prime}$ induced by the Lorentz form. But we shall not attempt to distinguish between $\mathbb{R}^{4}$ and its dual.) Formula (1.2) can then be written as

$$
\Lambda(x, y)=x^{\mu} y_{\mu}=x_{\mu} y^{\mu}
$$

The Lorentz inner product of a vector $x$ with itself is denoted by $x^{2}$ and its Euclidean norm by $|x|$ :

$$
x^{2}=x_{\mu} x^{\mu}, \quad|x|^{2}=x_{0}^{2}+x_{1}^{2}+x_{2}^{2}+x_{3}^{2}
$$

Vectors $x$ such that $x^{2}>0, x^{2}=0$, or $x^{2}<0$ are called timelike, lightlike, or spacelike, respectively.

In this framework, the notation for derivatives on $\mathbb{R}^{4}$ is

$$
\partial_{\mu}=\frac{\partial}{\partial x^{\mu}}
$$

so that, with respect to traditional space-time coordinates $\mathbf{x}$ and $t$,

$$
\left(\partial_{0}, \ldots, \partial_{3}\right)=\left(c^{-1} \partial_{t}, \nabla_{\mathbf{x}}\right), \quad\left(\partial^{0}, \ldots, \partial^{3}\right)=\left(c^{-1} \partial_{t},-\nabla_{\mathbf{x}}\right),
$$

and $\partial^{2}$ is the wave operator or d'Alembertian:

$$
\partial^{2}=\partial_{0}^{2}-\partial_{1}^{2}-\partial_{2}^{2}-\partial_{3}^{2}=c^{-2} \partial_{t}^{2}-\nabla_{\mathbf{x}}^{2}
$$

Integrals. Physicists like to indicate the dimensionality of their integrals explicitly by writing the volume element on $\mathbb{R}^{n}$ as $d^{n} x$. This convention is occasionally an aid to clarity, and we shall generally follow it. Another convention for integrals more commonly used in physics than in mathematics is to put the differential next to the integral sign: $\int d x f(x)$ rather than $\int f(x) d x$. This can also be an aid to clarity, particularly in multiple integrals where writing $\int_{a}^{b} d x \int_{c}^{d} d y f(x, y)$ in preference to $\int_{a}^{b} \int_{c}^{d} f(x, y) d y d x$ makes it easier to see which variables go with which limits of integration. However, for reasons of inertia more than anything else, we shall not adopt this convention.

Fourier transforms. If we are considering functions on the space $\mathbb{R}^{n}$ equipped with the Euclidean inner product $(x, y) \mapsto x \cdot y$, the Fourier transform $f \mapsto \widehat{f}$ and its inverse $f \mapsto f^{\vee}$ will be defined by

$$
\widehat{f}(y)=\int e^{-i x \cdot y} f(x) d^{n} x, \quad f^{\vee}(x)=\int e^{i x \cdot y} f(y) \frac{d^{n} y}{(2 \pi)^{n}}
$$

The Parseval identity is then

$$
\int|f(x)|^{2} d^{n} x=\int|\widehat{f}(y)|^{2} \frac{d^{n} y}{(2 \pi)^{n}}
$$

However, most of the time the Fourier transform will pertain to functions on $\mathbb{R}^{4}$ equipped with the Lorentz metric, in which case we define

$$
\widehat{f}(p)=\int e^{i p_{\mu} x^{\mu}} f(x) d^{4} x, \quad f^{\vee}(x)=\int e^{-i p_{\mu} x^{\mu}} f(p) \frac{d^{4} p}{(2 \pi)^{4}}
$$

(The Parseval identity is again (1.4), with $n=4$.) Thus the Lorentz sign convention in the exponent agrees with the Euclidean one as far as the space variables are concerned, but it is reversed for the time variable. In this situation, generally $x$ is in "position space" and $p$ is in "momentum space," in which case the $p_{\mu} x^{\mu}$ in the exponent really needs to be divided by a normalization factor so that the argument of the exponent is a pure number, independent of the units of measurement. In quantum mechanics this factor is Planck's constant $\hbar$, but it will usually be suppressed since we will be using units in which $\hbar=1$.

Lord Kelvin once quipped that a mathematician is someone to whom it is obvious that $\int_{-\infty}^{\infty} e^{-x^{2}} d x=\sqrt{\pi}$. In the same spirit, I submit that a mathematical physicist is someone to whom it is obvious that

$$
\delta(x)=\int e^{i x \cdot y} \frac{d^{n} y}{(2 \pi)^{n}}
$$

where $\delta$ denotes the delta-function (point mass at the origin) on $\mathbb{R}^{n}$. This is the Fourier inversion formula, stated in the informal language of distributions. Readers who do not yet qualify as mathematical physicists by this criterion are advised to spend some time playing with this formula until they understand how to make sense of it and are convinced of its validity, for it will be used on many occasions in this book. (The necessary tools from distribution theory can be found, for example, in Folland [48].)

It is useful to remember that the factors of $2 \pi$ always go with the measure on momentum space; that is, the standard measure on momentum space (or, more abstractly, "Fourier space" as opposed to "configuration space") $\mathbb{R}^{n}$ is $d^{n} p /(2 \pi)^{n}$. A corollary of this is that delta-functions on momentum space are normally accompanied by a factor of $(2 \pi)^{n}$ so that their integral against this measure is still 1.

Symbolic homonyms. It is an unfortunate fact of life that some letters of the alphabet have more than one conventional meaning, while others are used sometimes conventionally and sometimes just as a convenient label for a variable. In such cases one must rely on context for clarification. For example, all of us have surely seen uses of the letter $\pi$ that have nothing to do with the constant $3.14 \ldots$; there are some in this book. Other notable examples: (1) $e$ is both $\exp (1)$ and the electric charge that functions as the coupling constant in quantum electrodynamics (typically the charge of the electron or its absolute value); occasionally it is also the symbol for the electron. (Fortunately, we never need it to be the eccentricity of an ellipse!) (2) $q$ is sometimes an electric charge and sometimes a momentum. (3) $\alpha^{j}$ is a type of Dirac matrix, but $\alpha$ is the fine structure constant. (4) $\gamma^{\mu}$ is another type of Dirac matrix, but $\gamma$ is the Euler-Mascheroni constant and occasionally the symbol for the photon. (5) $Z$ is sometimes the number of protons in a nucleus, sometimes a field renormalization constant, sometimes the generating functional in the functional integral approach to field theory, and sometimes the symbol for the neutral vector Boson of the weak interaction.

Caveat lector.

\subsection*{Physical prologue: dimensions, units, constants, and particles}
The fundamental aspects of the universe to which all physical measurements relate are mass $[m]$, length or distance $[l]$, and time $[t]$. All quantities in physics come with "dimensions" that can be expressed in terms of these three basic ones. For example, velocity $\mathbf{v}=d \mathbf{x} / d t$ has dimensions of distance per unit time, or $\left[l t^{-1}\right]$. The dimensions of some of the basic quantities of mechanics are summarized in Table 1.1. The fact that angular momentum and action both have dimensions $\left[m l^{2} t^{-1}\right]$ seems a mere coincidence in classical mechanics, but it acquires a deeper significance in quantum mechanics, because [ $m l^{2} t^{-1}$ ] is the dimensions of Planck's constant.

Units for electromagnetic quantities are related to those of mechanics by taking the constant of proportionality in some basic law to be equal to one. For example, the ampere is the unit of current defined as follows: if two infinite, straight, perfectly conducting wires are placed parallel to each other one meter apart and a current of one ampere flows in each of them, the (magnetic) force between them per unit length is $2 \times 10^{-7}$ newton per meter. The other everyday units of electromagnetism (coulomb, volt, ohm, etc.) are then defined in terms of the ampere and the MKS

\begin{table}[h]
\begin{center}
\begin{tabular}{|l|l|l|}
\hline
Quantity & Dimensions & Defining relation \\
\hline
Momentum & [ $m l t^{-1}$ ] & $\mathrm{p}=\mathrm{m} \mathbf{v}$ \\
\hline
Angular momentum & $\left[m l^{2} t^{-1}\right]$ & $\mathbf{l}=\mathbf{x} \times \mathbf{p}$ \\
\hline
Force & [ $m l t^{-2}$ ] & $\mathbf{F}=m(d \mathbf{v} / d t)$ \\
\hline
Energy & $\left[m l^{2} t^{-2}\right]$ & $E=\frac{1}{2} m|\mathbf{v}|^{2}$ or $\int \mathbf{F} \cdot d \mathbf{x}$ or $m c^{2}$ or $\ldots$ \\
\hline
Action & $\left[m l^{2} t^{-1}\right]$ & $S=\int\left(E_{\text {kinetic }}-E_{\text {potential }}\right) d t$ \\
\hline
\end{tabular}

\caption{Table 1.1. Some basic quantities of mechanics.}
\end{center}
\end{table}

mechanical units. However, for purposes of fundamental physics it is better to take electric charge as the basic quantity. Naively, the simplest procedure ("Gaussian units") is to define the unit of charge so that the constant in Coulomb's law is equal to one. That is, having chosen a system of units for mechanics, one sets the unit of charge so that the force between two point charges is equal in magnitude to the product of the charges divided by the square of the distance between them. From a slightly more sophisticated point of view, however, a better procedure ("HeavisideLorentz units") is to take the constant in Coulomb's law to be $1 / 4 \pi$, essentially because $-1 / 4 \pi|\mathbf{x}|$ is the fundamental solution of the Laplacian $\left(\nabla^{2}[-1 / 4 \pi|\mathbf{x}|]=\delta\right)$, and this is the procedure we shall follow. Its practical advantage is that it makes all the $4 \pi$ 's disappear from Maxwell's equations. Either way, the dimension $[q]$ of charge is related to mechanical dimensions by $\left[q^{2} l^{-2}\right]=[$ force $]=\left[m l t^{-2}\right]$, or $[q]=\left[m^{1 / 2} l^{3 / 2} t^{-1}\right]$.

The idea of setting proportionality constants equal to one can be carried further. In relativistic physics it is natural to relate the units of length and time so that the speed of light $c$ is equal to one; doing this makes $[l]$ equivalent to $[t]$. Moreover, in quantum mechanics it is natural to choose units so that Planck's constant $\hbar$ is equal to one. Since $\hbar$ has dimensions $\left[m l^{2} t^{-1}\right]$, in conjunction with the condition $c=1$ this makes $[m]$ equivalent to $\left[l^{-1}\right]$ or $\left[t^{-1}\right]$. We have $c=299792458 \mathrm{~m} / \mathrm{s}$ (an exact equality according to the current official definition of the meter) and $\hbar=1.054589 \times 10^{-34} \mathrm{~kg} \cdot \mathrm{~m}^{2} / \mathrm{s}$, so the equivalences of the basic MKS units are as follows:

$$
1 \text { second } \cong 299792458 \text { meter } \cong 8.522668 \times 10^{50} \text { kilogram }^{-1} .
$$

These large numbers make the MKS system awkward for the world of particle physics. Seconds and centimeters (accompanied by large negative powers of 10) are still generally used for times and lengths, but the other commonly used unit is the electronvolt ( eV ), which is the amount of energy gained by an electron when passing through an electrostatic potential of one volt, or its larger relatives the mega-electronvolt ( MeV ) or giga-electronvolt ( GeV ):

$$
1 \mathrm{eV}=1.602176 \times 10^{-19} \mathrm{~kg} \cdot \mathrm{~m}^{2} / \mathrm{s}^{2}, \quad 1 \mathrm{MeV}=10^{6} \mathrm{eV}, \quad 1 \mathrm{GeV}=10^{9} \mathrm{eV}
$$

The eV is a unit of energy, but setting $c=1$ makes Einstein's $E=m c^{2}$ into an equality of mass and energy, and in the subatomic world this is not just a formal equivalence but an everyday fact of life. The eV , or more commonly the MeV or GeV , is therefore used as a unit of mass:

$$
1 \mathrm{MeV}=1.782661 \times 10^{-27} \text { gram. }
$$

The equivalence of units of mass, length, and time can be restated as follows:

$$
1 \mathrm{MeV} \cong\left(6.581966 \times 10^{-22} \text { second }\right)^{-1} \cong\left(1.973224 \times 10^{-11} \text { centimeter }^{-1}\right.
$$

Again, in the subatomic world these are more than formal equivalences. The precise constants are usually not important, but the orders of magnitude express relations among the characteristic energies, distances, and times of elementary particle interactions.

When we first discuss relativity in Chapter 2 and quantum theory in Chapter 3 , we will include the factors of $c$ and $\hbar$ explicitly for the sake of clarity, but starting with Chapter 4 we will always use "natural" units in which $c$ and $\hbar$ are both equal to 1 . We can still quote lengths in centimeters, times in seconds, and masses in MeV , but having chosen one of these units, we must relate the other two to it by (1.6).

Once one has made $[l]=[t]=\left[m^{-1}\right]$, the dimensions $\left[m^{1 / 2} l^{3 / 2} t^{-1}\right]$ of electric charge cancel out, so that charge is dimensionless, and the unit of charge in the Heaviside-Lorentz system is simply the number 1. The basic physical constant is then the fundamental unit of charge, the charge of a proton or the absolute value of the charge of an electron, which we denote here by $e$. The related quantity $e^{2} / 4 \pi$ (or $e^{2} / 4 \pi \hbar c$ if one does not use "natural" units) is the fine structure constant $\alpha$, which is nearly equal to $1 / 137$ and is often quoted that way:

$$
\alpha=\frac{e^{2}}{4 \pi}=0.0072973525=\frac{1}{137.035}
$$

which gives

$$
e=0.30282212
$$

(In some parts of this book, the letter $e$ may denote the charge of whatever particle is in question at the time; this is always an integer multiple $\left[\frac{1}{3}\right.$-integer in the case of quarks] of the $e$ here.)

One can carry the reduction of different dimensions one final step by including gravity. Having set $\hbar$ and $c$ equal to one, one can set the coefficient $G$ in Newton's law of gravity $F=G m_{1} m_{2} / r^{2}$ equal to one too. This yields an absolute scale of length, time, and mass, called the Planck scale. Since $G=6.674 \times 10^{-11} \mathrm{~m}^{3} / \mathrm{kg} \cdot \mathrm{s}^{2}$, we have

$$
\begin{aligned}
\text { Planck length } & =1.616 \times 10^{-33} \text { centimeter } \\
\text { Planck time } & =5.390 \times 10^{-44} \text { second } \\
\text { Planck mass } & =2.176 \times 10^{-5} \text { gram }=1.221 \times 10^{19} \mathrm{GeV}
\end{aligned}
$$

The Planck length and time are ridiculously small, and the Planck mass ridiculously large, on the scale of ordinary particle physics. There is much speculation about what particle physics on the Planck scale might look like, none of which will be discussed in this book.

Here is a quick review of the terminology concerning subatomic particles. To begin with, there are two basic dichotomies: all particles are either Bosons, with integer spin, or Fermions, with half-integer spin; and all particles are either hadrons, which participate in the strong interaction, or non-hadrons. The fundamental Fermions are quarks, which are hadrons, and leptons, which are not. Quarks combine in triplets to make baryons, of which the most familiar are the proton and\\
neutron, and in quark-antiquark pairs to make mesons, of which the most important are the pions; baryons are Fermions, whereas mesons are Bosons. The leptons comprise electrons, muons, ${ }^{2}$ tauons, and their associated neutrinos. The fundamental Bosons are the mediating particles of the various interactions: photons (electromagnetism), $W^{ \pm}$and $Z$ particles (weak interaction), gluons (strong interaction), and the Higgs Boson (the field that generates particle masses). There is also presumably a graviton for gravity, but it has not been observed as of this writing. For each of the particles mentioned above there is also an antiparticle; in some cases (photons, gravitons, Higgs Bosons, $Z$ particles, and some gluons and mesons) the particle and antiparticle coincide. Griffiths [59] is a good reference for the physics of elementary particles.

The rest masses of the particles that we will encounter most frequently in this book are as follows:

$$
\begin{gathered}
m_{\gamma}=0, \quad m_{e}=0.511 \mathrm{MeV}, \quad m_{p}=938.280 \mathrm{MeV}, \quad m_{n}=939.573 \mathrm{MeV} \\
m_{\mu}=105.659 \mathrm{MeV}, \quad m_{\pi^{ \pm}}=139.569 \mathrm{MeV}, \quad m_{\pi^{0}}=134.964 \mathrm{MeV}
\end{gathered}
$$

where the subscripts denote photons ( $\gamma$ ), electrons, protons, neutrons, muons, charged pions, and neutral pions. The mass of a proton in grams is $1.672635 \times 10^{-24}$; the reciprocal of this number, approximately $6 \times 10^{23}$, is essentially Avogadro's number, which to this one-significant-figure accuracy can be considered as the number of nucleons (protons or neutrons) in one gram of matter. ${ }^{3}$

A few characteristic lengths: The diameter of a proton is about $10^{-13} \mathrm{~cm}$. Diameters of atoms are in the range $1-5 \times 10^{-8} \mathrm{~cm}$. (Heavy atoms are about the same size as light ones, because although the former have more electrons, the increased charge of the nucleus causes them to be packed in more tightly.) Electrons, as far as we know, are point particles, but the Compton radius or classical radius of an electron is the number $r_{0}$ such that the mass of the electron is equal to the electrostatic energy of a solid ball of radius $r_{0}$ with a uniform charge distribution of total charge $e$, which in natural units is $e^{2} / 4 \pi m_{e}$ :

$$
r_{0}=\frac{e^{2}}{4 \pi m_{e}}=2.8 \times 10^{-13} \mathrm{~cm}
$$

\subsection*{Mathematical prologue: some Lie groups and Lie algebras}
In this section we review the basic facts and terminology concerning the Lie groups and Lie algebras that play a central role in quantum mechanics and relativity.

The Lorentz and orthogonal groups and their Lie algebras. The Lorentz group $O(1,3)$ is the group of all linear transformations of $\mathbb{R}^{4}$ (or $4 \times 4$ real matrices) that preserve the Lorentz inner product: ${ }^{4}$

$$
A \in O(1,3) \Longleftrightarrow(A x)_{\mu}(A y)^{\mu}=x_{\mu} y^{\mu} \text { for all } x, y \quad \Longleftrightarrow A^{\dagger} g A=g .
$$

We employ the notation discussed in §1.1; in particular, $g$ is the matrix (1.1).

\footnotetext{${ }^{2}$ Muons were originally called " $\mu$-mesons." This is a misuse of the word "meson" according to modern usage.\\
${ }^{3}$ The actual mass of an atom is generally less than the sum of the masses of its nucleons, even after adding in the mass of its electrons. One has to subtract the binding energy of the nucleus.\\
${ }^{4} O(1,3)$ is more commonly called $O(3,1)$, but I prefer to keep the arguments of $O(\cdot, \cdot)$ in the same order as the coordinates on $\mathbb{R}^{4}$.
}The group $O(1,3)$ has four connected components, which are determined by the values of two homomorphisms from $O(1,3)$ onto the two-element group $\{ \pm 1\}$ :

$$
A \mapsto \operatorname{det} A \quad \text { and } \quad A \mapsto \operatorname{sgn}\left(A_{0}^{0}\right)(\text { the sign of the }(0,0) \text { entry of } A) .
$$

The kernel of the first of these is the special Lorentz group $S O(1,3)$, and the kernel of the second one is the orthochronous Lorentz group $O^{\uparrow}(1,3)$, the subgroup of $O(1,3)$ that preserves the direction of time. The intersection

$$
S O^{\uparrow}(1,3)=S O(1,3) \cap O^{\uparrow}(1,3)
$$

sometimes called the restricted Lorentz group or proper Lorentz group, is the connected component of the identity in $O(1,3)$. The linear isometry group or orthogonal group $O(3)$ of $\mathbb{R}^{3}$ sits inside $O^{\uparrow}(1,3)$ as the subgroup that fixes the point $(1,0,0,0)$, and the rotation group $S O(3)$ is $O(3) \cap S O^{\uparrow}(1,3)$.

The Lie algebra $\mathfrak{s} \mathfrak{o}(1,3)$ of the Lorentz group consists of the $4 \times 4$ real matrices $X$ that satisfy

$$
X^{\dagger} g+g X=0
$$

A convenient basis for it may be described as follows. Let $e_{\mu \nu}$ be the matrix whose $(\mu, \nu)$ entry is 1 and whose other entries are 0 , and define

$$
X_{j k}=e_{k j}-e_{j k}, \quad X_{k 0}=-X_{0 k}=e_{k 0}+e_{0 k} \quad(j, k=1,2,3)
$$

Then $\left\{X_{\mu \nu}: \mu<\nu\right\}$ is a basis for $\mathfrak{s} \mathfrak{o}(1,3)$, and of course one can replace any $X_{\mu \nu}$ by $X_{\nu \mu}=-X_{\mu \nu}$. The commutation relations are given by

$$
\begin{aligned}
{\left[X_{\mu \nu}, X_{\rho \sigma}\right] } & =0 \text { if }\{\mu, \nu\}=\{\rho, \sigma\} \text { or }\{\mu, \nu\} \cap\{\rho, \sigma\}=\varnothing \\
{\left[X_{\mu \nu}, X_{\nu \rho}\right] } & =g_{\nu \nu} X_{\mu \rho}
\end{aligned}
$$

The Lie algebra $\mathfrak{s} \mathfrak{o}(3)$ of the spatial rotation group is the span of $\left\{X_{j k}: j, k>0\right\}$.\\
In detail, the relation between these basis elements and the Lie group is as follows. If $(j, k, l)$ is a cyclic permutation of $(1,2,3), \exp \left(s X_{j k}\right)$ is the rotation through angle $s$ about the $x^{l}$-axis (counterclockwise, as viewed from the positive $x^{l}$-axis). On the other hand, $\exp \left(s X_{0 k}\right)$ is a so-called boost along the $x^{k}$-axis that changes from the initial reference frame "at rest" to one moving with velocity $\tanh s$ along the $x^{k}$-axis. Thus, for example,

$$
\exp \left(s X_{23}\right)=\left(\begin{array}{cc}
I & 0 \\
0 & R(s)
\end{array}\right), \quad R(s)=\left(\begin{array}{cc}
\cos s & -\sin s \\
\sin s & \cos s
\end{array}\right)
$$

and

$$
\exp \left(s X_{01}\right)=\left(\begin{array}{cc}
B(s) & 0 \\
0 & I
\end{array}\right), \quad B(s)=\left(\begin{array}{cc}
\cosh s & \sinh s \\
\sinh s & \cosh s
\end{array}\right)
$$

Orbits and invariant measures. We need to consider the geometry of the action of $O(1,3)$ on $\mathbb{R}^{4}$. Actually, there are two natural actions of $O(1,3)$ on $\mathbb{R}^{4}$ : the identity representation of $O(1,3)$ and its contragredient $A \mapsto A^{\dagger-1}$. Strictly speaking, the latter is the action of $O(1,3)$ on the dual space $\left(\mathbb{R}^{4}\right)^{\prime}$; physically, it is the action on momentum space rather than position space. Since the map $A \mapsto A^{\dagger-1}$ is an automorphism of $O(1,3)$, however, the orbits in $\mathbb{R}^{4}$ under the two actions are identical, and we need not distinguish them. In what follows we shall think of $\mathbb{R}^{4}$ as momentum space, since this is the context in which the orbits usually appear naturally.

The orbits of the restricted Lorentz group and the orthochronous Lorentz group are identical. Parametrized by a nonnegative real number $m$ that generally has a physical interpretation as a mass, they are as follows:

$$
\begin{gathered}
X_{m}^{+}=\left\{p: p^{2}=m^{2}, p_{0}>0\right\}, \quad X_{m}^{-}=\left\{p: p^{2}=m^{2}, p_{0}<0\right\} \\
Y_{m}=\left\{p: p^{2}=-m^{2}\right\} \quad(m>0) \\
\{0\}
\end{gathered}
$$

(Recall that $p^{2}=p_{\mu} p^{\mu}$.) The orbits $X_{m}^{+}$(and sometimes also $X_{m}^{-}$) are known as mass shells, and the orbits $X_{0}^{+}$and $X_{0}^{-}$are called the forward and backward light cones. (Under the action of the full group $O(1,3)$, the two orbits $X_{m}^{ \pm}$coalesce into one.)

Observe that

$$
p \in X_{m}^{+} \Longleftrightarrow p=\left(\omega_{\mathbf{p}}, \mathbf{p}\right), \text { where } \omega_{\mathbf{p}}=\sqrt{m^{2}+|\mathbf{p}|^{2}} .
$$

The symbol $\omega_{\mathbf{p}}$ will carry this meaning throughout the book.\\
Each of the orbits (1.9) has an $O^{\uparrow}(1,3)$-invariant measure, which is known on abstract grounds to be unique up to scalar multiples. The invariant measure on the mass shell $X_{m}^{+}$with $m>0$ may be derived as follows. Let $V=\bigcup_{m>0} X_{m}^{+}=\{p$ : $\left.p^{2}>0, p_{0}>0\right\}$ be the region inside the forward light cone. Then $V$ is $O^{\uparrow}(1,3)$ invariant, and since $|\operatorname{det} T|=1$ for $T \in O^{\uparrow}(1,3)$, the restriction of Lebesgue measure $d^{4} p$ to $V$ is an $O^{\uparrow}(1,3)$-invariant measure on $V$; hence so is $f\left(p^{2}\right) d^{4} p$ for any nonnegative continuous $f$ with support in ( $0, \infty$ ). One obtains the invariant measure on $X_{m}^{+}$by letting $f$ turn into a delta-function with pole at $m^{2}: \delta\left(p^{2}-m^{2}\right) d^{4} p$. The result often appears in precisely this way in the physics literature, but in order to avoid possible pitfalls in using delta-functions with nonlinear arguments it is best to take a little more care.

To wit, consider the map $\phi:(0, \infty) \times \mathbb{R}^{3} \rightarrow V$ defined by

$$
\phi(y, \mathbf{p})=\left(\sqrt{y+|\mathbf{p}|^{2}}, \mathbf{p}\right)
$$

so that $q^{2}=y$ when $q=\phi(y, \mathbf{p})$. $\phi$ is a diffeomorphism, and its Jacobian is $1 / 2 \sqrt{y+|\mathbf{p}|^{2}}$, so for $f \in C_{c}(0, \infty)$ we have

$$
f\left(p^{2}\right) d^{4} p=\frac{f(y) d y d^{3} \mathbf{p}}{2 \sqrt{y+|\mathbf{p}|^{2}}}
$$

Now let $f$ approach the delta-function with pole at $y=m^{2}$ : if we write points in $X_{m}^{+}$as $\left(\omega_{\mathbf{p}}, \mathbf{p}\right)$ where $\omega_{\mathbf{p}}=\sqrt{m^{2}+|\mathbf{p}|^{2}}$, we obtain the invariant measure

$$
\frac{d^{3} \mathbf{p}}{2 \sqrt{m^{2}+|\mathbf{p}|^{2}}}=\frac{d^{3} \mathbf{p}}{2 \omega_{\mathbf{p}}}
$$

A limiting argument then shows that $d^{3} \mathbf{p} / 2|\mathbf{p}|$ is an invariant measure on $X_{0}^{+}$. The invariant measure on $X_{m}^{-}$, of course, is also given by (1.11). We leave the calculation of the invariant measure on $Y_{m}$, for which we shall have no use, to the reader.\\
$S L(2, \mathbb{C}), S U(2)$, and the Pauli matrices. The three Pauli matrices are

$$
\sigma_{1}=\left(\begin{array}{cc}
0 & 1 \\
1 & 0
\end{array}\right), \quad \sigma_{2}=\left(\begin{array}{cc}
0 & -i \\
i & 0
\end{array}\right), \quad \sigma_{3}=\left(\begin{array}{cc}
1 & 0 \\
0 & -1
\end{array}\right)
$$

We shall denote the triple ( $\sigma_{1}, \sigma_{2}, \sigma_{3}$ ) by $\boldsymbol{\sigma}$. For any cyclic permutation ( $j, k, l$ ) of $(1,2,3)$, we have

$$
\sigma_{j} \sigma_{k}=-\sigma_{k} \sigma_{j}=i \sigma_{l}, \quad \text { and hence } \quad\left[\sigma_{j}, \sigma_{k}\right]=2 i \sigma_{l}
$$

$\frac{1}{2 i} \sigma_{1}, \frac{1}{2 i} \sigma_{2}$, and $\frac{1}{2 i} \sigma_{3}$ are a basis for the Lie algebra $\mathfrak{s u}(2)$ of skew-Hermitian $2 \times 2$ matrices of trace zero, and together with $\frac{1}{2} \sigma_{1}, \frac{1}{2} \sigma_{2}$, and $\frac{1}{2} \sigma_{3}$, they are a basis (over $\mathbb{R}$ ) for the Lie algebra $\mathfrak{s l}(2, \mathbb{C})$ of all $2 \times 2$ complex matrices of trace zero. In view of (1.8) and (1.13), the linear map $\kappa^{\prime}: \mathfrak{s} \mathfrak{l}(2, \mathbb{C}) \rightarrow \mathfrak{s} \mathfrak{o}(1,3)$ defined on this basis by (1.14)\\
$\kappa^{\prime}\left(\frac{1}{2 i} \sigma_{j}\right)=X_{k l}$ for $(j, k, l)$ a cyclic permutation of $(1,2,3), \quad \kappa^{\prime}\left(\frac{1}{2} \sigma_{k}\right)=X_{0 k}$\\
is an isomorphism of Lie algebras, and its restriction to $\mathfrak{s u}(2)$ is an isomorphism from $\mathfrak{s u}(2)$ to $\mathfrak{s} \mathfrak{o}(3)$.

The corresponding homomorphism on the group level may be described as follows. Let us add the "fourth Pauli matrix"

$$
\sigma_{0}=I=\left(\begin{array}{ll}
1 & 0 \\
0 & 1
\end{array}\right)
$$

to obtain a basis $\sigma_{0}, \ldots, \sigma_{3}$ for the space $H$ of $2 \times 2$ Hermitian matrices. We then identify $\mathbb{R}^{4}$ with $H$ by the correspondence

$$
x \in \mathbb{R}^{4} \longleftrightarrow M(x)=x^{\mu} \sigma_{\mu}=\left(\begin{array}{cc}
x^{0}+x^{3} & x^{1}-i x^{2} \\
x^{1}+i x^{2} & x^{0}-x^{3}
\end{array}\right)
$$

The crucial feature of this correspondence is that

$$
\operatorname{det} M(x)=\left(x^{0}\right)^{2}-\left(x^{1}\right)^{2}-\left(x^{2}\right)^{2}-\left(x^{3}\right)^{2}=x^{2}
$$

Every $A \in S L(2, \mathbb{C})$ acts on $H$ by the map $X \mapsto A X A^{\dagger}$, and we denote the corresponding action on $\mathbb{R}^{4}$ by $\kappa(A)$ :

$$
M[\kappa(A) x]=A M(x) A^{\dagger}
$$

Since $\operatorname{det} A=1$, we have $\operatorname{det} A M(x) A^{\dagger}=\operatorname{det} M(x)$, and by (1.16) this means that $\kappa$ maps $\operatorname{SL}(2, \mathbb{C})$ into $O(1,3)$. It is easily verified that the differential of $\kappa$ at the identity, which takes $S \in \mathfrak{s l}(2, \mathbb{C})$ to the map $M(x) \mapsto S M(x)+M(x) S^{\dagger}$, is precisely the map $\kappa^{\prime}$ defined in (1.14). Thus $\kappa$ is a local isomorphism, and since $\operatorname{SL}(2, \mathbb{C})$ is connected, its image is the connected component of the identity in $O(1,3)$, namely, $S O^{\uparrow}(1,3)$. Finally, it is easy to verify that the kernel of $\kappa$ is $\pm I$. In short, we have proved:

The map $\kappa$ is a double covering of $S O^{\uparrow}(1,3)$ by $S L(2, \mathbb{C})$.\\
As we observed earlier, $S O(3)$ can be identified with the subgroup of $S O^{\uparrow}(1,3)$ that fixes the point $(1,0,0,0)$. The inverse image of $S O(3)$ in $S L(2, \mathbb{C})$ is therefore the set of all $A \in S L(2, \mathbb{C})$ that fix the point $I=M(1,0,0,0)$, i.e., that satisfy $A A^{\dagger}=I$. This is the group $S U(2)$ of $2 \times 2$ unitary matrices. It is easily verified that $\operatorname{SU}(2)$ is the set of matrices of the form $\left(\begin{array}{cc}a & -\bar{b} \\ b & \bar{a}\end{array}\right)$ where $|a|^{2}+|b|^{2}=1$, so that $\operatorname{SU}(2)$ is homeomorphic to the unit sphere in $\mathbb{C}^{2}$ and in particular is simply connected. Hence:\\
$S U(2)$ is the universal cover of $S O(3)$, and the covering map is the restriction of $\kappa$ to $\operatorname{SU}(2)$.

There is another way to look at this. The map $\mathbf{x} \mapsto i M(0, \mathbf{x})=i \mathbf{x} \cdot \boldsymbol{\sigma}$ is an isomorphism from $\mathbb{R}^{3}$ to the space of $2 \times 2$ skew-Hermitian matrices of trace zero, which is the Lie algebra $\mathfrak{s u}(2)$. Since $A^{\dagger}=A^{-1}$ for $A \in S U(2)$, the action\\
$\left.A \mapsto \kappa(A)\right|_{\{0\} \times \mathbb{R}^{3}}$ of $S U(2)$ on $\mathbb{R}^{3}$ is essentially the adjoint action of $S U(2)$ on its Lie algebra.

The map $\kappa$ has one more important property that is not quite obvious: it respects adjoints, i.e.,

$$
\kappa\left(A^{\dagger}\right)=\kappa(A)^{\dagger}
$$

This is most easily seen on the Lie algebra level: $\kappa^{\prime}$ takes the Hermitian matrices $\frac{1}{2} \sigma_{j}$ to the symmetric matrices $X_{j 0}$ and the skew-Hermitian matrices $\frac{1}{2 i} \sigma_{j}$ to the skewsymmetric matrices $X_{k l}$. Thus $\kappa^{\prime}\left(X^{\dagger}\right)=\kappa^{\prime}(X)^{\dagger}$, so since $\kappa(\exp X)=\exp \kappa^{\prime}(X)$ and exp preserves adjoints, the same is true of $\kappa$.

One can form a group that doubly covers the whole Lorentz group $O(1,3)$ : it is a semidirect product of $\operatorname{SL}(2, \mathbb{C})$ with the group $\mathbb{Z}_{2} \times \mathbb{Z}_{2}$. We leave the details to the reader.

The Poincaré group. The Poincaré group or inhomogeneous Lorentz group $\mathcal{P}$ is the group of transformations of $\mathbb{R}^{4}$ generated by $O(1,3)$ and the group of translations (isomorphic to $\mathbb{R}^{4}$ itself). That is, $\mathcal{P}$ is the semi-direct product of $\mathbb{R}^{4}$ and $O(1,3)$,

$$
\mathcal{P}=\mathbb{R}^{4} \ltimes O(1,3),
$$

whose underlying set is $\mathbb{R}^{4} \times O(1,3)$ with group law given by

$$
(a, S)(b, T)=(a+S b, S T), \quad(a, S)^{-1}=\left(-S^{-1} a, S^{-1}\right)
$$

The action of $(a, A) \in \mathcal{P}$ on $\mathbb{R}^{4}$ is

$$
(a, S) x=S x+a .
$$

Like $O(1,3), \mathcal{P}$ has four connected components, and the component of the identity is $\mathcal{P}_{0}=\mathbb{R}^{4} \ltimes S O^{\uparrow}(3,1)$. The covering map $\kappa$ of $O(1,3)$ by $S L(2, \mathbb{C})$ induces a double covering $(a, A) \mapsto(a, \kappa(A))$ of $\mathcal{P}_{0}$ by the group $\mathbb{R}^{4} \ltimes S L(2, \mathbb{C})$, whose group law is

$$
(a, A)(b, B)=(a+\kappa(A) b, A B) .
$$

\section*{CHAPTER 2}
\section*{Review of Pre-quantum Physics}
\begin{displayquote}
I was lucky enough to attend a few lectures of S. S. Chern just before he retired from Berkeley in which he said that the cotangent bundle (differential forms) is the feminine side of analysis on manifolds, and the tangent bundle (vector fields) is the masculine side. From this perspective, Hamiltonian mechanics is the feminine side of classical physics, [and] its masculine side is Lagrangian mechanics.
\end{displayquote}

-Richard Montgomery ([84], p. 352)\\[0pt]
This chapter is devoted to a brief review of classical mechanics, special relativity, and electromagnetic theory. Among the many available texts on these subjects, the Feynman Lectures [42] are excellent for Newtonian mechanics, relativity, and electromagnetism, and Purcell [91] perhaps even better for the latter. The classic physics text on Hamiltonian and Lagrangian mechanics is Goldstein [56]; Abraham and Marsden [1] and Arnold [4] are good treatments of this material by and for mathematicians.

\subsection*{Mechanics according to Newton and Hamilton}
We may as well begin at the beginning with Newton's law:

$$
F=m a .
$$

Familiar as this is, it needs some explanation. We consider a system of $k$ particles with fixed masses $m_{1}, \ldots, m_{k}$, located at positions $\mathbf{x}_{1}, \ldots, \mathbf{x}_{k} \in \mathbb{R}^{3}$ at time $t \in \mathbb{R}$. The $j$ th particle is acted upon by a force $\mathbf{F}_{j}$ that depends on the positions $\mathbf{x}_{1}, \ldots, \mathbf{x}_{k}$ and the time $t$ (and perhaps on other parameters such as the masses $m_{j}$ ). We then concatenate the positions and forces into $3 k$-vectors,

$$
F=\left(\mathbf{F}_{1}, \ldots, \mathbf{F}_{k}\right) \in \mathbb{R}^{3 k}, \quad x=\left(\mathbf{x}_{1}, \ldots, \mathbf{x}_{k}\right) \in \mathbb{R}^{3 k}
$$

and the masses into a $3 k \times 3 k$ matrix,

$$
m=\left(\begin{array}{cccc}
m_{1} I_{3} & 0 & \ldots & 0 \\
0 & m_{2} I_{3} & \ldots & 0 \\
\vdots & \vdots & & \vdots \\
0 & 0 & \ldots & m_{k} I_{3}
\end{array}\right) \quad\left(I_{3}=3 \times 3 \text { identity matrix }\right),
$$

and set

$$
v=\frac{d x}{d t}, \quad a=\frac{d v}{d t}
$$

Then (2.1) is a second-order ordinary differential equation for $x$ when the force $F$ is known, and a solution is determined by initial values $x_{0}=x\left(t_{0}\right)$ and $v_{0}=v\left(t_{0}\right)$.

Mathematicians look at physical laws such as (2.1) with the expectation that they will contain some universal truth about the world. They do, but not in the\\
absolute way that the formula $A=\pi r^{2}$ is a universal truth about circles; they are always burdened with a certain amount of fine print. The law (2.1) is valid in all places at all times, as far as we know, but it has to be modified in relativistic situations where the masses cannot be treated as constants, and it fades into the background in the study of the submicroscopic world where the whole concept of force loses much of its utility. Moreover, (2.1) is valid only in inertial coordinate systems, so it gives the right answers about what goes on in a laboratory on the surface of the earth only to the extent that the rotation of the earth and its motion about the sun can be neglected. It is important to keep in mind that all physical theories have their limitations; the quantum field theories toward which we are heading are no exception.

We shall consider only autonomous and conservative forces, that is, those $F$ that depend only on the positions $x$ (and not explicitly on $t$ ) and for which the line integral $\int_{C} F \cdot d x$ vanishes for every closed curve $C$ in $\mathbb{R}^{3 n}$. In this case, $F$ is the gradient of a function on $\mathbb{R}^{3 n}$ denoted by $-V$ and called the potential energy. We also have the kinetic energy $T$ of the system,

$$
T=\frac{1}{2} m v \cdot v=\sum \frac{1}{2} m_{j}\left|\mathbf{v}_{j}\right|^{2}
$$

and the total energy,

$$
E=T+V
$$

It is to be noted that $V$, and hence $E$, is well defined only up to an additive constant. Total energy is conserved:

$$
\frac{d E}{d t}=m v \cdot \frac{d v}{d t}+\nabla V \cdot \frac{d x}{d t}=(m a-F) \cdot v=0
$$

The Hamiltonian reformulation of Newtonian mechanics allows more flexibility and reveals some important mathematical structures. For reasons that will become clearer below, it takes the momentum

$$
p=m v
$$

instead of the velocity $v$ as a primary object. The total energy, considered as a function of position and momentum, is called the Hamiltonian and is denoted by $H$ :

$$
H(x, p)=\frac{1}{2} m^{-1} p \cdot p+V(x)
$$

We can rewrite Newton's law $F=m a$ as a first-order system in the variables $x$ and $p$ :

$$
\frac{d x}{d t}=v=m^{-1} p, \quad \frac{d p}{d t}=m a=-\nabla V(x)
$$

The key point is that the quantities on the right are derivatives of the Hamiltonian:

$$
\frac{d x}{d t}=\nabla_{p} H, \quad \frac{d p}{d t}=-\nabla_{x} H
$$

These are Hamilton's equations, which tell how position and momentum evolve with time. From them one easily obtains the evolution equations for any function of position and momentum, $f(x, p)$ :

$$
\frac{d f}{d t}=\nabla_{x} f \cdot \frac{d x}{d t}+\nabla_{p} f \cdot \frac{d p}{d t}=\{f, H\}
$$

where the Poisson bracket $\{f, g\}$ of any functions $f$ and $g$ is defined by

$$
\{f, g\}=\nabla_{x} f \cdot \nabla_{p} g-\nabla_{p} f \cdot \nabla_{x} g
$$

Of particular importance are the Poisson brackets of the coordinate functions themselves:

$$
\left\{x_{i}, p_{j}\right\}=-\left\{p_{j}, x_{i}\right\}=\delta_{i j}, \quad\left\{x_{i}, x_{j}\right\}=\left\{p_{i}, p_{j}\right\}=0
$$

At this point a fundamental mathematical structure comes clearly into view (from the modern perspective): that of a symplectic manifold.

Let us pause for a brief review of this concept. A symplectic manifold is a $C^{\infty}$ manifold $M$ equipped with a differential 2-form $\Omega$ that is closed ( $d \Omega=0$ ) and pointwise nondegenerate, i.e., for each $a \in M$ and $v \in T_{a} M, \Omega_{a}(v, w)=0$ for all $w \in T_{a} M$ only when $v=0$. Nondegeneracy forces the dimension of $M$ to be even, $\operatorname{dim} M=2 n$, and it is equivalent to the condition that the $n$th exterior power of $\Omega$ is everywhere nonvanishing.

The form $\Omega$ gives an identification of 1 -forms with vector fields. Namely, if $\phi$ is a 1 -form, the corresponding vector field $X_{\phi}$ is defined by

$$
\Omega\left(Y, X_{\phi}\right)=\phi(Y)
$$

for all vector fields $Y$. In particular, if $f$ is a smooth function on $M, X_{d f}$ is a smooth vector field on $M$ called the Hamiltonian vector field of $f$; by a small abuse of notation, we denote it by $X_{f}$ rather than $X_{d f}$. Thus, for any vector field $Y$,

$$
\Omega\left(Y, X_{f}\right)=d f(Y)=Y f
$$

The Poisson bracket of two smooth functions $f$ and $g$ is

$$
\{f, g\}=\Omega\left(X_{f}, X_{g}\right)=X_{f} g=-X_{g} f
$$

The Poisson bracket makes $C^{\infty}(M)$ into a Lie algebra; the Jacobi identity is a consequence of the fact that $d \Omega=0$ and the formula for the action of the exterior derivative of a form on vector fields. Moreover, the correspondence $f \mapsto X_{f}$ is a Lie algebra homomorphism:

$$
\begin{aligned}
X_{\{f, g\}} h=-\{h,\{f, g\}\} & =\{g,\{h, f\}\}+\{f,\{g, h\}\} \\
& =-X_{g} X_{f} h+X_{f} X_{g} h=\left[X_{f}, X_{g}\right] h
\end{aligned}
$$

A celebrated theorem of Darboux states that for any point $a$ in a symplectic manifold $M$ there is a system of local coordinates $x_{1}, \ldots, x_{n}, p_{1} \ldots, p_{n}$ on a neighborhood of $a$ such that $\Omega=\sum d x_{j} \wedge d p_{j}$. A coordinate system with this property is called canonical, and the coordinates $x_{j}$ and $p_{j}$ are said to be canonically conjugate. In canonical coordinates, Hamiltonian vector fields and Poisson brackets are given by

$$
X_{f}=\sum\left(\frac{\partial f}{\partial p_{j}} \frac{\partial}{\partial x_{j}}-\frac{\partial f}{\partial x_{j}} \frac{\partial}{\partial p_{j}}\right), \quad\{f, g\}=\sum\left(\frac{\partial f}{\partial x_{j}} \frac{\partial g}{\partial p_{j}}-\frac{\partial f}{\partial p_{j}} \frac{\partial g}{\partial x_{j}}\right) .
$$

A diffeomorphism $\Phi: M \rightarrow M$ that preserves the symplectic structure is called a canonical transformation (by physicists) or a symplectomorphism (by geometers). Suppose $\left\{\Phi_{t}: t \in \mathbb{R}\right\}$ is a one-parameter group of canonical transformations; then its infinitesimal generator, the vector field $X$ defined by $X f=\left.(d / d t) f \circ \Phi_{t}\right|_{t=0}$, satisfies $L_{X} \Omega=0$, where $L$ denotes the Lie derivative. Conversely, if $X$ is a vector field such that $L_{X} \Omega=0$, the flow it generates consists of (local) canonical transformations. But $L_{X} \Omega=i_{X}(d \Omega)+d\left(i_{X} \Omega\right)=d\left(i_{X} \Omega\right)$, where $i_{X}$ denotes contraction with $X$, and $i_{X} \Omega$ is the 1 -form associated to $X$ by $\Omega$, so $X$ is an infinitesimal canonical transformation if and only if its associated 1 -form is closed. But this means that $i_{X} \Omega=d f$ for some function $f$ (perhaps well-defined only on a covering\\
space of $M$ if $M$ is not simply connected), or in other words, that $X=X_{f}$. Thus, with the understanding that functions and transformations may be only locally well-defined, Hamiltonian vector fields are precisely the infinitesimal generators of canonical transformations.

For Hamiltonian mechanics, the symplectic manifolds of primary importance are cotangent bundles. Let $N$ be an $n$-dimensional manifold, $T^{*} N$ its cotangent bundle, and $\pi: T^{*} N \rightarrow N$ the natural projection. There is a canonical 1-form $\omega$ on $T^{*} N$, defined by $\omega(v)=\phi\left(\pi_{*} v\right)$ for $\phi \in T^{*} N$ and $v \in T_{\phi}\left(T^{*} N\right)$, and $\Omega=-d \omega$ is a symplectic form on $T^{*} N$. (The minus sign is inserted to make this consistent with the preceding discussion and some standard conventions in the literature.) Indeed, any system $\left\{x_{1}, \ldots, x_{n}\right\}$ of local coordinates on an open set $U \subset N$ induces a frame $\left\{d x_{1}, \ldots, d x_{n}\right\}$ on $T^{*} U$ and hence a system of linear coordinates $\left\{p_{1}, \ldots, p_{n}\right\}$ on each fiber of $T^{*} U$ (that is, if $\phi$ is a 1 -form on $U, \phi=\sum p_{j}(\phi) d x_{j}$ ), and $\left\{x_{1}, \ldots, x_{n}, p_{1}, \ldots, p_{n}\right\}$ is a system of local coordinates on $T^{*} U$. In these coordinates the canonical 1-form is given by $\omega=\sum p_{j} d x_{j}$, so $\Omega=\sum d x_{j} \wedge d p_{j}$. Thus $\Omega$ is indeed a symplectic form, and the coordinates $x_{j}, p_{j}$ are canonical.

We return to physics. According to Newton's law, once the forces are given, the motion of the system is completely determined by (i) the position $x$ and (ii) the velocity $v$ or the momentum $p$ at an initial time $t_{0}$. We therefore take this data as a complete description of the state of the system. To build a general mathematical framework for dealing with these matters, we start with a configuration space $N$, which is taken to be a manifold and is supposed to be a description of the possible "positions" of the system. There is quite a lot of flexibility here. For example, if the system consists of $k$ particles moving in $\mathbb{R}^{3}$ as discussed previously, $N$ will be $\mathbb{R}^{3 k}$, or perhaps $\left\{\left(\mathbf{x}_{1}, \ldots, \mathbf{x}_{k}\right) \in \mathbb{R}^{3 k}: \mathbf{x}_{i} \neq \mathbf{x}_{j}\right.$ for $\left.i \neq j\right\}$. If the motion is subject to some constraints, $N$ could be a submanifold of $\mathbb{R}^{3 k}$ instead. For an asymmetric rigid body, however, the appropriate configuration space is $\mathbb{R}^{3} \times S O(3)$ : three linear coordinates to give the location of the center of mass (or some other convenient point in the body), and three angular coordinates to give the body's orientation in space.

Velocities are taken to be tangent vectors to the configuration space, so the position-velocity state space is $T N$. On the other hand, the appearance of Poisson brackets in the Hamiltonian formalism leads us to take the position-momentum state space to be the symplectic manifold $T^{*} N$. That is, momenta should be considered as cotangent vectors, and in the relation $p=m v$ the mass matrix $m$ should be interpreted as a Riemannian metric on $N$ that mediates between vectors and covectors. Thus in the setting of $N=\mathbb{R}^{3 k}$ considered above, the inner product on tangent vectors is given by $\langle v, w\rangle_{m}=m v \cdot w$; the corresponding inner product on cotangent vectors is given by $\langle p, q\rangle^{m}=m^{-1} p \cdot q$, and the cotangent vector corresponding to the tangent vector $v$ is $p=m v . T^{*} N$ is traditionally called the phase space of the system. (The origin of the name lies in statistical mechanics.)

As a function on $T^{*} N$, the Hamiltonian is $H(\phi)=\frac{1}{2}\langle\phi, \phi\rangle^{m}+V(\pi(\phi))$, where $\langle\cdot, \cdot\rangle^{m}$ is the Riemannian inner product just described and $\pi: T^{*} N \rightarrow N$ is the natural projection. Hamilton's equations (2.2) say that the Hamiltonian vector field $X_{H}$ is the infinitesimal generator of the one-parameter group $\Phi_{t}$ of time translations of the system. (If $(x, p)$ is the state at time $t_{0}, \Phi_{t}(x, p)$ is the state at time $t+t_{0}$. Actually $\Phi_{t}$ may be only a flow rather than a full one-parameter group, as particles might collide or escape to infinity in a finite time.) Consequently, the time\\
translations are canonical transformations: the symplectic structure is an invariant of the dynamics.

In this setting it is easy to derive the connection between symmetries and conserved quantities generally known as Noether's theorem. To wit, suppose $\left\{\Psi_{s}\right.$ : $s \in \mathbb{R}\}$ is a one-parameter group of canonical transformations of phase space, and suppose the Hamiltonian $H$ is invariant under these transformations. The infinitesimal generator of $\Psi_{s}$ is a Hamiltonian vector field $X_{f}$, and invariance means that $X_{f} H=0$. But $X_{f} H=\{H, f\}=-\{f, H\}=-d f / d t$, so $f$ is invariant under time translations, i.e., $f$ is a conserved quantity. For example, if $N=\mathbb{R}^{3 k}$ and $\Psi_{s}$ is spatial translation in the direction of a unit vector $\mathbf{u} \in \mathbb{R}^{3}$, i.e., $\Psi_{s}(x, p)= \left(\mathbf{x}_{1}+s \mathbf{u}, \ldots, \mathbf{x}_{k}+s \mathbf{u}, p\right)$, then $f$ is the $\mathbf{u}$-component of the total linear momentum, $f(x, p)=\left(\mathbf{p}_{1}+\ldots+\mathbf{p}_{k}\right) \cdot \mathbf{u}$. If $N=\mathbb{R}^{3}$ and $\Psi_{s}$ is rotation through the angle $s$ about the $\mathbf{u}$ axis, then $f$ is the $\mathbf{u}$-component of the angular momentum, $f(\mathbf{x}, \mathbf{p})= (\mathbf{x} \times \mathbf{p}) \cdot \mathbf{u}$.

One of the advantages of the Hamiltonian formulation of mechanics is that it allows the possibility of simplifying the problem by performing arbitrary canonical transformations, including ones that mix up the position and momentum variables. As an illustration, we consider a simple but important problem whose quantum analogue will be of fundamental importance later on: the one-dimensional harmonic oscillator. This is a single particle of mass $m$, moving along the real line subject to a linear restoring force $F(x)=-\kappa x(\kappa>0)$. The associated potential is $V(x)= \frac{1}{2} \kappa x^{2}$. Of course, Newton's equation $m x^{\prime \prime}=-\kappa x$ is easy to solve directly, but the Hamiltonian method allows us to transform this simple differential equation into a completely trivial one.

Let us first observe that canonical transformations on $\mathbb{R}^{2}$ are precisely those transformations that preserve orientation and area, as the symplectic form $d x \wedge d p$ is just the element of oriented area.

Let $\omega=\sqrt{\kappa / m}$. The Hamiltonian is

$$
H=\frac{1}{2 m} p^{2}+\frac{\kappa}{2} x^{2}=\frac{\omega}{2}\left(\frac{1}{\sqrt{\kappa m}} p^{2}+\sqrt{\kappa m} x^{2}\right) .
$$

As a first step, we make the canonical transformation $\widetilde{x}=(\kappa m)^{1 / 4} x, \widetilde{p}=p /(\kappa m)^{1 / 4}$, which makes $H=(\omega / 2)\left(\widetilde{x}^{2}+\widetilde{p}^{2}\right)$. Now, the polar-coordinate map $(\widetilde{x}, \widetilde{p}) \mapsto(r, \theta)$ is not area-preserving, but its close relative $(\widetilde{x}, \widetilde{p}) \mapsto\left(\frac{1}{2} r^{2}, \theta\right)$ is, since $r d r=d\left(\frac{1}{2} r^{2}\right)$. Thus, the map $(\widetilde{x}, \widetilde{p}) \mapsto(s, \theta)$ is canonical if we take $s=\frac{1}{2}\left(\widetilde{x}^{2}+\widetilde{p}^{2}\right)$ and $\theta= \arctan (\widetilde{p} / \widetilde{x})$, and in the new coordinates, $H$ is simply $\omega s$. Hamilton's equations therefore boil down to

$$
\frac{d s}{d t}=\frac{\partial H}{\partial \theta}=0, \quad \frac{d \theta}{d t}=-\frac{\partial H}{\partial s}=-\omega .
$$

Thus $s$ is a constant - namely, $s=E / \omega$ where $E$ is the total energy of the system - and $\theta=\theta_{0}-\omega t$. Finally, $\widetilde{x}=\sqrt{2 s} \cos \theta=\sqrt{2 E / \omega} \cos \left(\omega t-\theta_{0}\right)$, so

$$
x=\sqrt{\frac{2 E}{\kappa}} \cos \left(\omega t-\theta_{0}\right) .
$$

The constant $\theta_{0}$ can be chosen to make the initial values $x\left(t_{0}\right)$ and $p\left(t_{0}\right)$ whatever we wish, subject to the condition that $H\left(x_{0}, p_{0}\right)=E$.

Another canonical transformation that will be useful later is the one that simplifies the two-body problem. Suppose that two particles with masses $m_{1}, m_{2}$, positions $\mathbf{x}_{1}, \mathbf{x}_{2}$, and momenta $\mathbf{p}_{1}, \mathbf{p}_{2}$ are subject to a conservative force with potential $V$ that depends only on their relative displacement $\mathbf{x}_{1}-\mathbf{x}_{2}$, so the Hamiltonian is

$$
H=\frac{\left|\mathbf{p}_{1}\right|^{2}}{2 m_{1}}+\frac{\left|\mathbf{p}_{2}\right|^{2}}{2 m_{2}}+V\left(\mathbf{x}_{1}-\mathbf{x}_{2}\right) .
$$

Let

$$
\begin{array}{cl}
M=m_{1}+m_{2}, \quad m=\frac{m_{1} m_{2}}{m_{1}+m_{2}}, & \mathbf{X}=\frac{m_{1} \mathbf{x}_{1}+m_{2} \mathbf{x}_{2}}{m_{1}+m_{2}}, \quad \mathbf{x}=\mathbf{x}_{1}-\mathbf{x}_{2} \\
\mathbf{P}=M \frac{d \mathbf{X}}{d t}=\mathbf{p}_{1}+\mathbf{p}_{2}, & \mathbf{p}=m \frac{d \mathbf{x}}{d t}=\frac{m_{2} \mathbf{p}_{1}-m_{1} \mathbf{p}_{2}}{m_{1}+m_{2}}
\end{array}
$$

The transformation $\left(\mathbf{x}_{1}, \mathbf{x}_{2}, \mathbf{p}_{1}, \mathbf{p}_{2}\right) \mapsto(\mathbf{X}, \mathbf{x}, \mathbf{P}, \mathbf{p})$ is canonical, and in the new coordinates the Hamiltonian becomes

$$
H=\frac{|\mathbf{P}|^{2}}{2 M}+\frac{|\mathbf{p}|^{2}}{2 m}+V(\mathbf{x})
$$

in which the coordinates are decoupled. Hamilton's equations then give $d \mathbf{X} / d t= \mathbf{P} / M, d \mathbf{P} / d t=0$, so that $\mathbf{X}$, the center of mass of the two-body system, moves with constant velocity; and $d \mathbf{x} / d t=\mathbf{p} / m, d \mathbf{p} / d t=-\nabla V(\mathbf{x})$, so that $m \mathbf{x}^{\prime \prime}(t)=-\nabla V(\mathbf{x})$. In short, after removing the uniform motion of the center of mass, the problem reduces to that of a single particle of mass $m$, the reduced mass of the system, moving in the potential $V(\mathbf{x})$.

\subsection*{Mechanics according to Lagrange}
We now turn to the other major reworking of Newtonian mechanics, the Lagrangian formulation. On the simplest level of a system of particles moving in a potential $V$, whose Hamiltonian is $H(x, p)=\frac{1}{2} m^{-1} p \cdot p+V(x)$, the Lagrangian is the function of position and velocity given by

$$
L(x, v)=m v \cdot v-H(x, m v)=\frac{1}{2} m v \cdot v-V(x) .
$$

Stated in a more intrinsic fashion, if we regard $H$ as a function on the cotangent bundle $T^{*} N$, the Lagrangian is the function on the tangent bundle $T N$ given by $L(\xi)=\left\langle j_{m} \xi, \xi\right\rangle-H\left(j_{m} \xi\right)$, where $j_{m}: T N \rightarrow T^{*} N$ is the isomorphism provided by the inner product $m$, and $\langle\cdot, \cdot\rangle$ is the pairing of cotangent vectors with tangent vectors. Returning to (2.5), we have

$$
\nabla_{x} L=-\nabla_{x} H=-\nabla V, \quad \nabla_{v} L=m v,
$$

so Newton's law becomes Lagrange's equation

$$
\frac{d}{d t} \nabla_{v} L=\nabla_{x} L
$$

The significance of this is that it is the Euler-Lagrange equation for a certain problem in the calculus of variations. Namely, given a path $t \mapsto x(t), t_{0} \leq t \leq t_{1}$, in the configuration space $N$, we consider the action

$$
S=\int_{t_{0}}^{t_{1}} L\left(x(t), x^{\prime}(t)\right) d t
$$

and the problem is to minimize the action over all paths $x(t)$ that begin and end at two given points $x_{0}=x\left(t_{0}\right)$ and $x_{1}=x\left(t_{1}\right)$. Indeed, suppose we replace $x(t)$ by a\\
slightly different path $x(t)+\delta x(t)$, where $\delta x\left(t_{0}\right)=\delta x\left(t_{1}\right)=0$. To first order in $\delta x$, the change in the Lagrangian is

$$
\delta L=L\left(x+\delta x, x^{\prime}+\delta x^{\prime}\right)-L\left(x, x^{\prime}\right)=\nabla_{x} L \cdot \delta x+\nabla_{v} L \cdot \delta x^{\prime}
$$

so the change in the action is

$$
\delta S=\int_{t_{0}}^{t_{1}} \delta L d t=\int_{t_{0}}^{t_{1}}\left[\nabla_{x} L-\frac{d}{d t} \nabla_{v} L\right] \cdot \delta x d t
$$

where the endpoint terms in the integration by parts vanish since $\delta x$ vanishes at the endpoints. At a minimum, $\delta S$ must vanish, and $\delta x$ is arbitrary, so Lagrange's equation must hold.

This result is commonly known as the principle of least action; it will come back to haunt us in Chapter 8. Of course, this name is somewhat inaccurate: a solution of Newton's equations is a path that is a critical point, not necessarily a minimizer, of the action.

The calculation that led us from Hamiltonian mechanics to Lagrangian mechanics is easily reversible. Indeed, suppose we start with the Lagrangian $L(x, v)= \frac{1}{2} m v \cdot v-V(x)$. We define the Hamiltonian $H$ by

$$
H(x, p)=m^{-1} p \cdot p-L\left(x, m^{-1} p\right)=\frac{1}{2} m^{-1} p \cdot p+V(x)
$$

Then $\nabla_{p} H=m^{-1} p=v=d x / d t$, and $\nabla_{v} L=m v=p$, so by Lagrange's equation, $-\nabla_{x} H=\nabla_{x} L=(d / d t)\left(\nabla_{v} L\right)=d p / d t$, and we have Hamilton's equations.

The Lagrangian and Hamiltonian machinery can be used in many situations other than the simple mechanical systems we have mentioned. The most common paradigm is to formulate a physical problem as a variational problem for a Lagrangian $L$ defined on the tangent bundle of some configuration space $N$. The general procedure for converting such a problem to Hamiltonian form is as follows.

Let $x_{1}, \ldots, x_{n}$ be local coordinates on $N$, and let $v_{1}, \ldots, v_{n}$ be the corresponding linear coordinates on the fibers of $T N$; thus $L$ is a function of $x$ and $v$. For the moment, we fix a point $a \in N$ and think of $L$ as a function on $T_{a} N$. Its differential $d L=\sum\left(\partial L / \partial v_{j}\right) d v_{j}$ is a 1 -form on $T_{a} N$. If we identify the vector space $T_{a} N$ with its own tangent spaces, we can think of $d L$ as a map that assigns to each $v \in T_{a} N$ a linear functional on $T_{a} N$, i.e., a map from $T_{a} N$ to $T_{a}^{*} N$. Now letting $a$ vary too, we obtain a map $W: T N \rightarrow T^{*} N$.

Let us write this out in terms of local coordinates. Given coordinates $x_{1}, \ldots, x_{n}$ on $U \subset N$, we identify $T U$ and $T^{*} U$ both with $U \times \mathbb{R}^{n}$; then $L$ is a function of $x$ and the fiber coordinates $v$ on $T U$, and $W(x, v)=\left(x, \nabla_{v} L\right)$. The component $p_{j}=\partial L / \partial v_{j}$ of $\nabla_{v} L$ is called the conjugate momentum to the position variable $x_{j}$ with respect to $L$.

Now suppose that $W$ is invertible. From the Lagrangian $L$ on $T N$ we construct the Hamiltonian $H$ on $T^{*} N$ by

$$
H(\xi)=\left\langle\xi, W^{-1} \xi\right\rangle-L\left(W^{-1} \xi\right)
$$

where the first term on the right is the pairing of a covector with a vector. A calculation similar to the one we performed above then shows that the Euler-Lagrange equation for $L$ is equivalent to Hamilton's equations for $H$. Exactly the same recipe can be used to go from $H$ back to $L$; it is called the Legendre transformation.

In the general situation, this procedure for going from $H$ to $L$ and back may be problematic because of the question of the invertibility of $W$. However, it clearly\\
coincides with the transformation we considered earlier for the case where $L(x, v)= \frac{1}{2} m v \cdot v-V(x)$. More generally, suppose that for each $a \in N$ the restriction of $L$ to $T_{a} N$ is a quadratic function whose pure second-order part is positive definite. Then $W$ is indeed invertible: it is an affine-linear isomorphism of $T_{a} N$ and $T_{a}^{*} N$ for each $a$. This is the most important case in practice, and the only one that will concern us in the sequel.

The Lagrangian formulation of Noether's theorem about symmetries and conserved quantities (which is the original version) is as follows: Suppose $\left\{\psi_{s}: s \in \mathbb{R}\right\}$ is a one-parameter group of diffeomorphisms of the configuration space $N$ that preserve the Lagrangian:

$$
L\left(\psi_{s}(x),\left(\psi_{s}\right)_{*}(v)\right)=L(x, v)
$$

Then the quantity $I(x, v)=\nabla_{v} L(x, v) \cdot d \psi_{s}(x) / d s$ is a constant of the motion. (See Arnold [4] or Goldstein [56].)

Our real interest is not in mechanical systems with finitely many degrees of freedom but in fields, which can be regarded as continuum mechanical systems with infinitely many degrees of freedom. Let us illustrate the transition to this situation with a simple example.

Consider a long elastic rod that is subject to longitudinal vibrations (or, if you prefer, a long air column such as an organ pipe). As a simplified model, we chop the rod up into small bits of equal size, consider each bit as a point particle, and replace the elastic forces within the rod by ideal springs connecting the particles. More precisely, we take each particle to have mass $\Delta m$ and the separation between adjacent particles when the system is at rest to be $\Delta x$. Let $u_{j}$ be the displacement of the $j$ th particle from its position at rest. Then Newton's law gives the following system of equations for the $u_{j}$ 's:

$$
\Delta m \frac{d^{2} u_{j}}{d t^{2}}=k\left(u_{j+1}-u_{j}\right)-k\left(u_{j}-u_{j-1}\right)=k\left(u_{j+1}-2 u_{j}+u_{j-1}\right)
$$

where $k$ is the common spring constant, and the Lagrangian is

$$
L=T-V=\frac{1}{2} \sum \Delta m\left(\frac{d u_{j}}{d t}\right)^{2}-\frac{1}{2} \sum k\left(u_{j+1}-u_{j}\right)^{2} .
$$

Let us divide (2.6) by the equilibrium separation $\Delta x$,

$$
\frac{\Delta m}{\Delta x} \frac{d^{2} u_{j}}{d t^{2}}=(k \Delta x) \frac{u_{j+1}-2 u_{j}+u_{j-1}}{\Delta x^{2}},
$$

and rewrite (2.7) as

$$
L=\frac{1}{2} \sum\left[\frac{\Delta m}{\Delta x}\left(\frac{d u_{j}}{d t}\right)^{2}-(k \Delta x)\left(\frac{u_{j+1}-u_{j}}{\Delta x}\right)^{2}\right] \Delta x .
$$

Now pass to the continuum limit: $\Delta m / \Delta x$ becomes the mass density $\mu, k \Delta x$ becomes the Young's modulus $Y$, the displacements $u_{j}$ become a displacement function $u(x)$, and we obtain

$$
\mu \frac{\partial^{2} u}{\partial t^{2}}=Y \frac{\partial^{2} u}{\partial x^{2}}
$$

from (2.6) and

$$
L\left(u, \partial_{t} u, \partial_{x} u\right)=\frac{1}{2} \int_{a}^{b}\left[\mu\left(\frac{\partial u}{\partial t}\right)^{2}-Y\left(\frac{\partial u}{\partial x}\right)^{2}\right] d x
$$

from (2.7), where $[a, b]$ is the interval occupied by the rod. Equation (2.8) is the familiar wave equation which is indeed the standard model for elastic vibrations, and (2.9) expresses the Lagrangian as the integral of the Lagrangian density

$$
\mathcal{L}\left(u, \partial_{t} u, \partial_{x} u\right)=\frac{1}{2}\left[\mu\left(\frac{\partial u}{\partial t}\right)^{2}-Y\left(\frac{\partial u}{\partial x}\right)^{2}\right] .
$$

The wave equation (2.8) can be derived from the Lagrangian (2.9) by the calculus of variations just as before. That is, one considers the action

$$
S(u)=\int_{t_{0}}^{t_{1}} L d t=\int_{t_{0}}^{t_{1}} \int_{a}^{b} \mathcal{L}\left(u, \partial_{x} u, \partial_{t} u\right) d x d t
$$

and requires $\delta S=S(u+\delta u)-S(u)$ to vanish to first order in $\delta u$ for appropriate $\delta u$. To make this precise, one must consider what happens at the ends of the rod. One possibility is to consider an infinitely long rod, $[a, b]=\mathbb{R}$. The vibrations $u(x, t)$ are then assumed to have compact support in $x$ or vanish rapidly as $x \rightarrow \pm \infty$, so that the Lagrangian makes sense, and one imposes the same requirement on $\delta u$, together with $\delta u=0$ at $t=t_{0}, t_{1}$. Likewise, for a rod with fixed ends or an air column with closed ends, one requires $u$ and $\delta u$ to vanish at $x=a, b$.

A little more interesting is the case of a rod with free ends or an air column with open ends. Here there is no a priori restriction on $u$ and $\delta u$ at the endpoints; let us see what happens. With $L$ given by (2.9), to first order in $\delta u$ we have

$$
\delta S=\int_{t_{0}}^{t_{1}} \int_{a}^{b}\left[\mu \frac{\partial u}{\partial t} \frac{\partial \delta u}{\partial t}-Y \frac{\partial u}{\partial x} \frac{\partial \delta u}{\partial x}\right] d x d t
$$

Integration by parts yields

$$
\delta S=\int_{t_{0}}^{t_{1}} \int_{a}^{b}\left[-\mu \frac{\partial^{2} u}{\partial t^{2}}+Y \frac{\partial^{2} u}{\partial x^{2}}\right] \delta u d x d t-Y \int_{t_{0}}^{t_{1}}\left[\frac{\partial u}{\partial x} \delta u\right]_{x=a}^{b} d t
$$

(We always assume $\delta u$ vanishes at $t=t_{0}, t_{1}$, so the terms with $t$-derivatives contribute no boundary terms.) Thus if we want $\delta S=0$ for arbitrary $\delta u, u$ must satisfy not only the wave equation but the boundary conditions $\partial u / \partial x=0$ at $x=a, b$. These are in fact the physically correct boundary conditions for a rod with free ends or an air column with open ends, and they yield a well-posed boundary value problem for the wave equation.

Let us be clear about the subtle shift of notation that has taken place here. For a system with finitely many degrees of freedom, we had position variables $x_{j}$ and velocity variables $v_{j}$. Here the index $j$ has become the continuous variable $x, x_{j}$ has become $u(x, t)$, and $v_{j}$ has become $\partial_{t} u$. It is not worthwhile to try to maintain the geometric language of tangent and cotangent vectors at this point, but we can still speak of the canonically conjugate momentum density to the field $u$,

$$
p=\frac{\partial \mathcal{L}}{\partial\left(\partial_{t} u\right)}=\mu \frac{\partial u}{\partial t}
$$

and the Hamiltonian density,

$$
\mathcal{H}=p \partial_{t} u-\mathcal{L}=\frac{1}{2}\left[\left(\frac{\partial u}{\partial t}\right)^{2}+Y\left(\frac{\partial u}{\partial x}\right)^{2}\right] .
$$

The Hamiltonian $H=\int \mathcal{H} d x$ still represents the total energy, and it is a conserved quantity.

Similar ideas pertain to the wave equation in higher dimensions. Beginning with the Lagrangian density

$$
\mathcal{L}=\frac{1}{2}\left[c^{-2}\left(\partial_{t} u\right)^{2}-\left|\nabla_{x} u\right|^{2}\right]
$$

for a function $u(x, t)$ with $x \in \Omega$, an open subset of $\mathbb{R}^{n}$, and $t \in \mathbb{R}$, one forms the action integral

$$
S=\int_{t_{0}}^{t_{1}} \int_{\Omega} \mathcal{L} d x d t
$$

Setting the first variation of $S$ equal to zero yields the wave equation

$$
c^{-2} \partial_{t}^{2} u-\nabla^{2} u=0
$$

If $\Omega$ is a bounded domain and one imposes no restrictions on the perturbation $\delta u$ on the boundary $\partial \Omega$, one obtains the Neumann boundary conditions $\partial u / \partial n=0$ for free. Alternatively, one can take $\Omega=\mathbb{R}^{n}, t_{0}=-\infty$, and $t_{1}=+\infty$, with some implicit or explicit assumptions on the vanishing of the fields at infinity, and obtain the wave equation on $\mathbb{R}^{n}$.

This last point leads to a glimpse of the importance of Lagrangians for quantum field theory. As we shall see, the transition from classical mechanics to quantum mechanics proceeds most easily by adopting the Hamiltonian point of view. However, since this approach singles out the time axis for special attention, it lacks relativistic invariance. A Lagrangian density, on the other hand, or the action which is its integral over space-time, employs the space and time variables on an equal footing, and it may have relativistic invariance built in - as (2.11) does, for example. It therefore serves as a useful starting point for building relativistic field theories.

When we return to these ideas in later chapters, we shall often be cavalier about dropping the word "density"; that is, we shall often refer to $\mathcal{L}$ as the Lagrangian and $\mathcal{H}$ as the Hamiltonian; the meaning will always be clear from the context.

\subsection*{Special relativity}
It is taken as axiomatic that the fundamental laws of physics should be invariant under the Euclidean group of translations and rotations of space, as well as under translations of time. Newtonian physics and its immediate descendants are also invariant under the Galilean transformations

$$
(t, \mathbf{x}) \mapsto(t, \mathbf{x}+t \mathbf{v})
$$

for any $\mathbf{v} \in \mathbb{R}^{3}$. The fundamental equations governing electromagnetism, however, are not, as we shall see in the next section. Rather, they are invariant under translations and under the Lorentz group $O(1,3)$ (see § 1.3). For the time being we adopt notation that displays the speed of light $c$ explicitly, so $O(1,3)$ is the group of linear transformations of $\mathbb{R}^{4}$ that preserve the Lorentz bilinear form

$$
\Lambda\left((t, \mathbf{x}),\left(t^{\prime}, \mathbf{x}^{\prime}\right)\right)=c^{2} t t^{\prime}-\mathbf{x} \cdot \mathbf{x}^{\prime}=c^{2} t t^{\prime}-x x^{\prime}-y y^{\prime}-z z^{\prime}
$$

Einstein's great insight was that the laws of mechanics should be modified so as to be invariant under this group too.

A typical Lorentz transformation that mixes space and time variables is the boost along the $x$-axis with velocity $v(|v|<c)$ :

$$
(t, x, y, z) \mapsto\left(\frac{t+\left(v x / c^{2}\right)}{\sqrt{1-(v / c)^{2}}}, \frac{x+v t}{\sqrt{1-(v / c)^{2}}}, y, z\right)
$$

(This description of the boost is related to the one in §1.3 by the substitution $v / c=\tanh s$.) The factors of $\sqrt{1-(v / c)^{2}}$ in (2.12) are what account for the "Fitzgerald contraction" and "time dilation" effects in relativistic motion. (The fact that the word "contraction" is used in one case and "dilation" is used in the other is purely a matter of psychology. The factor $\sqrt{1-(v / c)^{2}}$ is in the denominator for both $x$ and $t$.) The feature of these phenomena that is perplexing to the intuition is that the inverse transformation of (2.12) is of exactly the same form, with $v$ replaced by $-v$, so the same factors of $\sqrt{1-(v / c)^{2}}$ appear in the inverse, not their reciprocals. Nonetheless, everything works out consistently.

For example, suppose an observer $A$ fires a bomb at an object $O$ that is at rest with respect to $A$, and an observer $B$ rides along with the bomb like Slim Pickens in Dr. Strangelove. $A$ and $B$ synchronize their clocks so that the bomb is launched at time 0 . Suppose the bomb is set to go off 1 second after launch, the speed of the bomb is $0.6 c$, and the distance from $A$ to $O$ is $2.25 \times 10^{8}$ meters ( $=0.75$ lightsecond). Result: The bomb goes off when it reaches $O$. From $A$ 's point of view, the bomb takes 1.25 seconds to reach $O$, and it takes that long to go off because the clock on board runs slow by a factor of $\sqrt{1-(v / c)^{2}}=0.8$. From $B$ 's point of view, the bomb takes 1 second to reach $O$, but the distance from $A$ to $O$ is only $1.8 \times 10^{8}$ meters because it is shortened by a factor of 0.8 . (This is why muons created by cosmic rays in the upper atmosphere, around $10^{-4}$ light-second from earth, manage to reach the earth's surface by traveling at almost the speed of light, even though the half-life of a muon is only about $10^{-6}$ second. We think the muons' clocks run slow; they think the distance is short.) On the other hand, from $B$ 's point of view, $A$ 's clock runs slow, so (according to $B$ ) $A$ 's clock only reads 0.8 second when the bomb goes off. The point here is that the events of the explosion and the reading of 0.8 second on $A$ 's clock are simultaneous for $B$ but not for $A$. This is possible because the two events happen at different places: the transformation (2.12) does not take equal times to equal times when $x \neq 0$. In constrast, $A$ and $B$ agree that $B$ 's clock reads 1 second when the bomb goes off; these two events happen at the same place, namely $O$, so their simultaneity is independent of the reference frame.

An elementary derivation of the relativistic formulas for momentum, energy, etc., can be found in Feynman [42]. Here we take a different approach, via Lagrangian mechanics, as a simple exercise in the sort of thought processes that go into finding the laws of quantum field theory.

We consider a free particle and wish to derive its laws of motion from an action functional, $S=\int_{t_{0}}^{t_{1}} L(\mathbf{x}, \mathbf{v}) d t$, where $\mathbf{x}=\mathbf{x}(t)$ is a path with given endpoints $\mathbf{x}\left(t_{0}\right)=\mathbf{x}_{0}$ and $\mathbf{x}\left(t_{1}\right)=\mathbf{x}_{1}$, and $\mathbf{v}(t)=\mathbf{x}^{\prime}(t)$. (We implicitly assume that $|\mathbf{v}(t)|<c$ and that $c^{2}\left(t_{1}-t_{0}\right)^{2}>\left|\mathbf{x}_{1}-\mathbf{x}_{0}\right|^{2}$.) We wish the action to be Lorentz invariant, i.e., to be unchanged if ( $\mathbf{x}_{0}, t_{0}$ ) and ( $\mathbf{x}_{1}, t_{1}$ ) are subjected to the same Lorentz transformation, and we also wish it to yield Newtonian mechanics for a free particle in the limit of small velocities. The first requirement narrows down the possibilities\\
enormously; the simplest one is to take the integrand to be a constant multiple of the Lorentz analogue of the differential of arc length,

$$
a d s=a \sqrt{c^{2} d t^{2}-d x^{2}-d y^{2}-d z^{2}}
$$

which in a given reference frame becomes

$$
a \sqrt{1-\frac{|\mathbf{v}|^{2}}{c^{2}}} c d t
$$

Thus we try the Lagrangian

$$
L(\mathbf{x}, \mathbf{v})=L(\mathbf{v})=a c \sqrt{1-\frac{|\mathbf{v}|^{2}}{c^{2}}}
$$

To figure out the right value of $a$, observe that when $|\mathbf{v}| \ll c$ we have

$$
L(\mathbf{v}) \approx a c-\frac{a|\mathbf{v}|^{2}}{2 c}
$$

Adding or subtracting a constant such as $a c$ to the Lagrangian does not affect the equations of motion, and the second term $-a|\mathbf{v}|^{2} / 2 c$ becomes the classical (nonrelativistic) Lagrangian $m|\mathbf{v}|^{2} / 2$ for a free particle of mass $m$ if we take $a= -m c$. Thus we are led to the Lagrangian

$$
L=-m c^{2} \sqrt{1-\frac{|\mathbf{v}|^{2}}{c^{2}}}
$$

The corresponding conjugate momentum is

$$
\mathbf{p}=\nabla_{\mathbf{v}} L=\frac{m \mathbf{v}}{\sqrt{1-\left|\mathbf{v}^{2}\right| / c^{2}}}
$$

Again, this is approximately the classical momentum $m \mathbf{v}$ when $\mathbf{v}$ is small. In general, it is $M \mathbf{v}$ where

$$
M=M(\mathbf{v})=\frac{m}{\sqrt{1-|\mathbf{v}|^{2} / c^{2}}}
$$

We are led to the conclusion that the effective mass of a particle of (rest) mass $m$ moving with velocity $\mathbf{v}$ is $M(\mathbf{v})$. Thus mass is subject to the same dilation effect as time and distance.

The Euler-Lagrange equation for the Lagrangian (2.14) is

$$
0=\frac{d}{d t} \nabla_{\mathbf{v}} L-\nabla_{\mathbf{x}} L=\frac{d \mathbf{p}}{d t}
$$

which gives motion in a straight line with constant speed as expected. (In general, if forces are present, Newton's law remains valid if " $m \mathbf{a}$ " is reinterpreted as $d \mathbf{p} / d t$. But the question of a proper relativistic interpretation of forces is in general problematic.)

Next, the quantity

$$
E=\mathbf{p} \cdot \mathbf{v}-L
$$

considered as a function of $\mathbf{v}$, is a constant of the motion that constitutes the total energy in nonrelativistic mechanics. In our situation,

$$
E=\frac{m|\mathbf{v}|^{2}}{\sqrt{1-|\mathbf{v}|^{2} / c^{2}}}+m c^{2} \sqrt{1-|\mathbf{v}|^{2} / c^{2}}=\frac{m c^{2}}{\sqrt{1-|\mathbf{v}|^{2} / c^{2}}}=M c^{2}
$$

which is Einstein's formula for the energy of a free particle in relativistic mechanics. When $\mathbf{v}$ is small,

$$
E \approx m c^{2}+\frac{1}{2} m|\mathbf{v}|^{2}
$$

the sum of the "rest energy" $m c^{2}$ and the classical kinetic energy. We observe that

$$
\frac{E^{2}}{c^{2}}-|\mathbf{p}|^{2}=\frac{m^{2} c^{2}}{1-|\mathbf{v}|^{2} / c^{2}}-\frac{m^{2}|\mathbf{v}|^{2}}{1-|\mathbf{v}|^{2} / c^{2}}=m^{2} c^{2}
$$

so that

$$
E=c \sqrt{|\mathbf{p}|^{2}+m^{2} c^{2}}
$$

Expressed in this way as a function of $\mathbf{p}, E$ is the Hamiltonian of the free particle. Again, for small velocities we have $|\mathbf{p}| \ll m c$ and hence

$$
E \approx m c^{2}+\frac{|\mathbf{p}|^{2}}{2 m}
$$

the rest energy plus the classical Hamiltonian.\\
The momentum $\mathbf{p}$ and the energy $E$ are different in different frames of reference related by Lorentz transformations. However, it is not hard to check that $\mathbf{p}$ and $E / c$ transform as the space and time components of a 4 -vector, called the (relativistic) 4 -momentum or energy-momentum vector. The relation (2.15) expresses the fact that the Lorentz inner product of this vector with itself is $m^{2} c^{2}$. It is important to note that in the tensor notation explained in §1.1, this vector is $p^{\mu}$ rather than $p_{\mu}$; otherwise the sign of $\mathbf{p}$ comes out wrong.

\subsection*{Electromagnetism}
Classically, the electric field $\mathbf{E}$ and the magnetic field $\mathbf{B}$ are vector-valued functions of position $\mathbf{x} \in \mathbb{R}^{3}$ and time $t$ that may be operationally defined by the Lorentz force law: The force on a particle with charge $q$ moving with velocity $\mathbf{v}$ is

$$
\mathbf{F}=q \mathbf{E}+\frac{q}{c} \mathbf{v} \times \mathbf{B}
$$

where $c$ is the speed of light, and $\mathbf{E}$ and $\mathbf{B}$ are evaluated at the location of the particle. ${ }^{1}$

The behavior of $\mathbf{E}$ and $\mathbf{B}$ is governed by Maxwell's equations. Using the Heaviside-Lorentz convention on the scale of electric charge (see §1.2), they are ${ }^{2}$

$$
\begin{aligned}
\operatorname{div} \mathbf{E} & =\rho \\
\operatorname{div} \mathbf{B} & =0 \\
\operatorname{curl} \mathbf{E}+\frac{1}{c} \frac{\partial \mathbf{B}}{\partial t} & =0 \\
\operatorname{curl} \mathbf{B}-\frac{1}{c} \frac{\partial \mathbf{E}}{\partial t} & =\frac{1}{c} \mathbf{j}
\end{aligned}
$$

Here $\rho$ is the charge density and $\mathbf{j}$ is the current density, both of which are functions of position and time. Everything here may be interpreted in the sense of

\footnotetext{${ }^{1}$ In some texts the factor of $c$ is omitted, which amounts to a redefinition of $\mathbf{B}$. The point of including it is to make $\mathbf{E}$ and $\mathbf{B}$ have the same dimensions, namely, force per unit charge.\\
${ }^{2}$ Some texts write Maxwell's equations in a way that involves two additional vector fields called $\mathbf{D}$ and $\mathbf{H}$. This is appropriate for the study of electromagnetism in bulk matter, but on the level of fundamental physics, $\mathbf{D = E}$ and $\mathbf{H}=\mathbf{B}$.
}
distributions, so that one may allow point charges, charge distributions on curves or surfaces, etc.

The quantities $\rho$ and $\mathbf{j}$ are not independent; indeed, by (2.18) and (2.21),

$$
\frac{\partial \rho}{\partial t}=\operatorname{div} \frac{\partial \mathbf{E}}{\partial t}=c \operatorname{div}(\operatorname{curl} \mathbf{B})-\operatorname{div} \mathbf{j}=-\operatorname{div} \mathbf{j}
$$

since div curl $=0$. This so-called continuity equation,

$$
\frac{\partial \rho}{\partial t}+\operatorname{div} \mathbf{j}=0
$$

is a strong form of the law of conservation of charge. More precisely, the total charge in a region $R \subset \mathbb{R}^{3}$ is the integral of $\rho$ over $R$, and by the divergence theorem,

$$
\frac{d}{d t} \iiint_{R} \rho d V=-\iiint \operatorname{div} \mathbf{j} d V=-\iint_{\partial R} \mathbf{j} \cdot \mathbf{n} d S
$$

where $\mathbf{n}$ is the unit outward normal to $\partial R$. Thus the total charge in $R$ can change only to the extent that charge enters and leaves through $\partial R$.

On the one hand, electromagnetic fields are produced by charges and currents, in accordance with Maxwell's equations; on the other, they act on charged bodies in accordance with the Lorentz force law. If one combines these things with a specification of the relation between charge and mass distributions, one obtains from (2.17)-(2.21) a system of differential equations that, in principle, completely describes the evolution of the system. It is nonlinear because of the products in (2.17). However, in practice one often considers electromagnetic fields produced by a given system of charges and currents or the motion of some charged particles induced by a given electromagnetic field (produced by some laboratory apparatus, say) without worrying about feedback.

If $\rho$ and $\mathbf{j}$ are taken as given (subject to (2.22)), Maxwell's equations alone provide a linear system of differential equations for the fields. The two equations (2.20) and (2.21) that involve time derivatives form a symmetric hyperbolic system,

$$
\begin{aligned}
& \frac{\partial \mathbf{E}}{\partial t}=c \operatorname{curl} \mathbf{B}-\mathbf{j} \\
& \frac{\partial \mathbf{B}}{\partial t}=-c \operatorname{curl} \mathbf{E}
\end{aligned}
$$

The mathematical theory of such systems is well understood. In particular, the Cauchy problem is well posed: there is a family of theorems that say that if $\mathbf{j}$ is in some nice space $X$ of functions or distributions on $\mathbb{R}^{4}$ and the initial data $\mathbf{E}\left(t_{0}, \cdot\right)$ and $\mathbf{B}\left(t_{0}, \cdot\right)$ are in some related space $y$ of functions or distributions on $\mathbb{R}^{3}$, there is a unique solution ( $\mathbf{E}, \mathbf{B}$ ) of this system in another related space $z$. Versions of this result can be found, for example, in Taylor $[\mathbf{1 1 7}], § 6.5$, and Treves $[\mathbf{1 2 2}], § 15$. As for the other two Maxwell equations, observe that

$$
\begin{gathered}
\frac{\partial}{\partial t} \operatorname{div} \mathbf{B}=\operatorname{div} \frac{\partial \mathbf{B}}{\partial t}=-c \operatorname{div} \operatorname{curl} \mathbf{E}=0 \\
\frac{\partial}{\partial t} \operatorname{div} \mathbf{E}=\operatorname{div} \frac{\partial \mathbf{E}}{\partial t}=c \operatorname{div} \operatorname{curl} \mathbf{B}-\operatorname{div} \mathbf{j}=\frac{\partial \rho}{\partial t}
\end{gathered}
$$

Hence if the equations (2.18) and (2.19) are satisfied at the initial time $t_{0}$, they are satisfied at all other times too, so they are merely restrictions on the initial values of $\mathbf{E}$ and $\mathbf{B}$.

In what follows we assume that the quantities in Maxwell's equations are defined on all of space-time $\mathbb{R}^{4}$, although they may have singularities, in which case the equations are to be interpreted in the sense of distributions. (However, one should not think of the spatial component $\mathbb{R}^{3}$ as representing the entire universe. Rather, one should regard it as a model for some region of which the phenomena in which one is interested are largely concentrated in some compact subregion. Depending on the phenomena one is studying, the actual size of such a region could be anything from a fraction of a meter to many light-years. But one almost always has in mind that the fields and charge distributions vanish at infinity.)

By (2.19), $\mathbf{B}$ is the curl of a vector field $\mathbf{A}$ known as the vector potential. The equation (2.20) then can be rewritten as $\operatorname{curl}\left(\mathbf{E}+c^{-1} \partial_{t} \mathbf{A}\right)=0$, so $\mathbf{E}+c^{-1} \partial_{t} \mathbf{A}$ is the gradient of a function $-\phi$, the scalar potential. Of course $\mathbf{A}$ and $\phi$ are far from unique: given any (smooth) function $\chi$ on $\mathbb{R}^{4}$, one can replace $\mathbf{A}$ by $\mathbf{A}-\nabla_{\mathbf{x}} \chi$ and $\phi$ by $\phi+c^{-1} \partial_{t} \chi$ without changing $\mathbf{E}$ and $\mathbf{B}$. Such adjustments to $\mathbf{A}$ and $\phi$ are called gauge transformations (for reasons that will be explained in §9.1), and the choice of a particular $\chi$ to make $\mathbf{A}$ and $\phi$ satisfy some desired condition is called a choice of gauge.

One of the commonly imposed gauge conditions is

$$
\operatorname{div} \mathbf{A}+\frac{1}{c} \frac{\partial \phi}{\partial t}=0
$$

which can be achieved by starting with any $\mathbf{A}$ and $\phi$ and replacing them with $\mathbf{A}-\nabla \chi$ and $\phi+c^{-1} \partial_{t} \chi$ where $\chi$ satisfies the inhomogeneous wave equation $\nabla^{2} \chi-c^{-2} \partial_{t}^{2} \chi= \operatorname{div} \mathbf{A}+c^{-1} \partial_{t} \phi$. Potentials $\mathbf{A}$ and $\phi$ that satisfy (2.23) are said to be in the Lorentz gauge or Landau gauge. In terms of such potentials, Maxwell's equations take a particularly appealing form. The equations (2.19) and (2.20) are already embodied in the definition of $\mathbf{A}$ and $\phi$. For (2.18), we have

$$
\rho=\operatorname{div} \mathbf{E}=-\nabla^{2} \phi-\frac{1}{c} \frac{\partial}{\partial t} \operatorname{div} \mathbf{A}=-\nabla^{2} \phi+\frac{1}{c^{2}} \frac{\partial^{2} \phi}{\partial t^{2}}
$$

and for (2.21), since $\operatorname{curl}(\operatorname{curl} \mathbf{A})=\nabla(\operatorname{div} \mathbf{A})-\nabla^{2} \mathbf{A}$ we have

$$
\begin{aligned}
\frac{1}{c} \mathbf{j} & =\operatorname{curl}(\operatorname{curl} \mathbf{A})-\frac{1}{c} \frac{\partial}{\partial t}\left(-\nabla \phi-\frac{1}{c} \frac{\partial \mathbf{A}}{\partial t}\right) \\
& =\frac{1}{c^{2}} \frac{\partial^{2} \mathbf{A}}{\partial t^{2}}-\nabla^{2} \mathbf{A}
\end{aligned}
$$
In short, denoting the wave operator or d'Alembertian by,

$$
\square=\frac{1}{c^{2}} \frac{\partial^{2}}{\partial t^{2}}-\nabla^{2}
$$

we see that Maxwell's equations boil down to

$$
\square \phi=\rho, \quad \square \mathbf{A}=\mathbf{j} / c
$$

in the presence of the Lorentz gauge condition (2.23).\\
The situation here is best understood by thinking relativistically and putting space and time on an equal footing. Henceforth we adopt the space-time coordinates

$$
x^{0}=c t, \quad\left(x^{1}, x^{2}, x^{3}\right)=\mathbf{x}
$$

to get rid of the factors of $c$, and we adopt the tensor notation explained in §1.1. We define the electromagnetic 4-potential by

$$
A^{\mu}=(\phi, \mathbf{A}), \quad \text { or } \quad A_{\mu}=(\phi,-\mathbf{A}) .
$$

Gauge transformations now take the form

$$
A^{\mu} \mapsto A^{\mu}+\partial^{\mu} \chi, \quad \text { or } \quad A_{\mu} \mapsto A_{\mu}+\partial_{\mu} \chi
$$

The electric and magnetic fields are combined into the electromagnetic field tensor

$$
F_{\mu \nu}=\partial_{\mu} A_{\nu}-\partial_{\nu} A_{\mu}, \quad \text { or } \quad F^{\mu \nu}=\partial^{\mu} A^{\nu}-\partial^{\nu} A^{\mu}
$$

that is,

$$
F_{\mu \nu}=\left(\begin{array}{cccc}
0 & E_{x} & E_{y} & E_{z} \\
-E_{x} & 0 & -B_{z} & B_{y} \\
-E_{y} & B_{z} & 0 & -B_{x} \\
-E_{z} & -B_{y} & B_{x} & 0
\end{array}\right), \quad F^{\mu \nu}=\left(\begin{array}{cccc}
0 & -E_{x} & -E_{y} & -E_{z} \\
E_{x} & 0 & -B_{z} & B_{y} \\
E_{y} & B_{z} & 0 & -B_{x} \\
E_{z} & -B_{y} & B_{x} & 0
\end{array}\right) .
$$

The Maxwell equations (2.19) and (2.20) become

$$
\partial_{\kappa} F_{\mu \nu}+\partial_{\mu} F_{\nu \kappa}+\partial_{\nu} F_{\kappa \mu}=0
$$

which is an immediate consequence of the form (2.26) of $F_{\mu \nu}$, and the equations (2.18) and (2.21) become

$$
\partial_{\mu} F^{\mu \nu}=j^{\nu}
$$

where the 4 -current density $j^{\mu}$ is defined by

$$
j^{\mu}=(\rho, \mathbf{j} / c)
$$

The Lorentz gauge condition (2.23) becomes

$$
\partial_{\mu} A^{\mu}=0
$$

and since

$$
\partial_{\mu} F^{\mu \nu}=\left(\partial_{\mu} \partial^{\mu}\right) A^{\nu}-\partial^{\nu}\left(\partial_{\mu} A^{\mu}\right)
$$

the Lorentz condition together with (2.28) yields the wave equations (2.24):

$$
\partial^{2} A^{\mu}=j^{\mu}
$$

Moreover, the continuity equation (2.22) becomes

$$
\partial_{\mu} j^{\mu}=0
$$

Fans of the language of differential geometry may prefer to see this restated in terms of differential forms. The potential and current density are expressed as 1-forms:\\
$A=A_{\mu} d x^{\mu}=\phi d t-A_{x} d x-A_{y} d y-A_{z} d z, \quad j=j_{\mu} d x^{\mu}=\rho d t-j_{x} d x-j_{y} d y-j_{z} d z$ (with $t$ measured in units so that $c=1$ ). The electromagnetic field tensor $F$ is the negative of the exterior derivative of $A$ :\\
$F=-d A=\left(E_{x} d x+E_{y} d y+E_{z} d z\right) \wedge d t+B_{x} d y \wedge d z+B_{y} d z \wedge d x+B_{z} d x \wedge d y$.\\
We have

$$
d F=-d d A=0
$$

which is equivalent to (2.27). The other Maxwell equation (2.28) and the continuity equation (2.29) become

$$
* d * F=j, \quad d * j=0
$$

where $*$ is the Hodge star operator relative to the Lorentz form. (It is defined by

$$
\alpha \wedge * \beta=\Lambda(\alpha, \beta) d t \wedge d x \wedge d y \wedge d z
$$

where $\Lambda(\alpha, \beta)$ denotes the bilinear functional on differential forms induced by the Lorentz form.)

In terms of the potential $A$, we have

$$
-* d * d A=j
$$

In the presence of the Lorentz gauge condition

$$
d * A=0
$$

this is equivalent to

$$
-(* d * d A+d * d *) A=j
$$

But $-(* d * d+d * d *)$ is the "Lorentz Laplacian" - that is, the d'Alembertian $\partial^{2}$, acting componentwise; hence we recover (2.24) in the form

$$
\partial^{2} A=j \quad(d * A=0)
$$

We now turn to the Lagrangian formulation of the laws of electromagnetism. We first derive the Lagrangian for a charged particle with rest mass $m$ and charge ${ }^{3} q$, moving in a given electromagnetic field with potential $A^{\mu}=(\phi, \mathbf{A})$. Relativistic invariance severly restricts the possibilities for a Lorentz-invariant action functional, and the simplest one is

$$
S=\int_{\left(t_{0}, \mathbf{x}_{0}\right)}^{\left(t_{1}, \mathbf{x}_{1}\right)}\left(-m d s-q A_{\mu} d x^{\mu}\right)
$$

where $d s$ is as in (2.13) and the integral is taken over a path in spacetime with the given endpoints. The first term gives the action for a free particle; the second one is meant to represent the interaction with the electromagnetic field; and as with the coefficient $-m$ for the free particle, the coefficient $-q$ for the interaction is dictated by the need to obtain the correct equations of motion. Notice that the gauge ambiguity in $A_{\mu}$ is irrelevant: replacing $A_{\mu}$ by $A_{\mu}+\partial_{\mu} \chi$ subtracts the exact differential $d \chi$ from the integrand and the constant $\chi\left(t_{1}, \mathbf{x}_{1}\right)-\chi\left(t_{0}, \mathbf{x}_{0}\right)$ from the action, which does not affect the dynamics.

Choosing a particular reference frame with coordinates $(t, \mathbf{x})$ (and units chosen so that $c=1$ ) and rewriting the action as an integral in $t$ with $\mathbf{v}=d \mathbf{x} / d t$ gives

$$
S=\int_{t_{0}}^{t_{1}}\left(-m \sqrt{1-|\mathbf{v}|^{2}}+q \mathbf{A} \cdot \mathbf{v}-q \phi\right) d t
$$

so the Lagrangian is

$$
L=-m \sqrt{1-|\mathbf{v}|^{2}}+q \mathbf{A} \cdot \mathbf{v}-q \phi
$$

At this point there are two different momenta to be considered. We shall refer to the ordinary mass-times-velocity momentum $m \mathbf{v} / \sqrt{1-|\mathbf{v}|^{2}}$ as the kinematic or mechanical momentum and denote it by $\mathbf{p}_{\kappa}$. On the other hand, $\mathbf{p}$ will denote the canonical momentum associated to the Lagrangian $L$ :

$$
\mathbf{p}=\nabla_{\mathbf{v}} L=\frac{m \mathbf{v}}{\sqrt{1-|\mathbf{v}|^{2}}}+q \mathbf{A}=\mathbf{p}_{\kappa}+q \mathbf{A}
$$

\footnotetext{${ }^{3}$ It is an experimental fact that charge, unlike mass, does not depend on velocity.
}
(This is the momentum that needs to be considered when we pass to quantum mechanics.) The Euler-Lagrange equation is
$$
\frac{d \mathbf{p}}{d t}=\nabla_{\mathbf{x}} L=q \nabla_{\mathbf{x}}(\mathbf{A} \cdot \mathbf{v})-q \nabla \phi
$$

Here $L$ is considered as a function of the independent variables $\mathbf{x}$ and $\mathbf{v}$, so $\mathbf{v}$ is treated as a constant in the expression $\nabla_{\mathbf{x}}(\mathbf{A} \cdot \mathbf{v})$; hence,

$$
\nabla_{\mathbf{x}}(\mathbf{A} \cdot \mathbf{v})=\left(\mathbf{v} \cdot \nabla_{\mathbf{x}}\right) \mathbf{A}+\mathbf{v} \times \operatorname{curl} \mathbf{A} .
$$

Moreover,

$$
\frac{d \mathbf{A}}{d t}=\frac{\partial \mathbf{A}}{\partial t}+\left(\mathbf{v} \cdot \nabla_{\mathbf{x}}\right) \mathbf{A}
$$

Hence the Euler-Lagrange equation is equivalent to

$$
\begin{aligned}
\frac{d \mathbf{p}_{\kappa}}{d t}=\frac{d}{d t}(\mathbf{p}-q \mathbf{A}) & =q\left(-\nabla \phi-\frac{\partial \mathbf{A}}{\partial t}\right)+q \mathbf{v} \times \operatorname{curl} \mathbf{A} \\
& =q \mathbf{E}+q \mathbf{v} \times \mathbf{B}
\end{aligned}
$$

This is the Lorentz force law (2.17) (with $c=1$ ), so our choice (2.30) for the action is indeed correct.

The corresponding Hamiltonian $H=\mathbf{p} \cdot \mathbf{v}-L$ is, in terms of the velocity $\mathbf{v}$,

$$
H=\frac{m|\mathbf{v}|^{2}}{\sqrt{1-|\mathbf{v}|^{2}}} \cdot \mathbf{v}+m \sqrt{1-|\mathbf{v}|^{2}}+q \phi=\frac{m}{\sqrt{1-|\mathbf{v}|^{2}}}+q \phi
$$

the sum of the free energy and the Coulomb potential energy. By the same calculation as in the case of a free particle, we find that $(H-q \phi)^{2}-\left|\mathbf{p}_{\kappa}\right|^{2}=m^{2}$, so in terms of the canonical momentum $\mathbf{p}$,

$$
H=\sqrt{m^{2}+|\mathbf{p}-q \mathbf{A}|^{2}}+q \phi
$$

When $\mathbf{v}$ is small we have

$$
L \approx-m+\frac{1}{2} m|\mathbf{v}|^{2}+q \mathbf{A} \cdot \mathbf{v}-q \phi, \quad H \approx m+\frac{1}{2 m}|\mathbf{p}-q \mathbf{A}|^{2}+q \phi
$$

which yield the Lagrangian and Hamiltonian for nonrelativistic motion of a particle in an electromagnetic field after subtraction of the rest mass $m$.

We observed earlier that the energy and momentum ( $E, \mathbf{p}$ ) of a free particle make up a 4 -vector $p^{\mu}$. In the presence of an electromagnetic potential $A^{\mu}$, so do the total energy given by (2.31) and the canonical momentum $\mathbf{p}$, and we still denote this 4 -vector by $p^{\mu}$. Equation (2.31) implies that

$$
(E-q \phi)^{2}=m^{2}+|\mathbf{p}-q \mathbf{A}|^{2}
$$

which says that the Lorentz inner product of $p^{\mu}-q A^{\mu}$ with itself is $m^{2}:(p-q A)^{2}= m^{2}$. Comparing this with the formula (2.16) for the free energy, we obtain a rule for incorporating the electromagnetic field that will be of basic importance in the quantum theory: To obtain the energy-momentum vector for a particle with charge $q$ in an electromagnetic field $A^{\mu}$ from that for a free particle, simply replace $p^{\mu}$ by $p^{\mu}-q A^{\mu}$.

Next we consider the Lagrangian form of the field equations for electromagnetism. Here, since we are dealing with fields instead of particles, the action will be a 4-dimensional integral of a Lagrangian density that is a function of the fields and their derivatives. Moreover, the basic field is taken to be the potential $A^{\mu}$ rather than the electromagnetic field $F^{\mu \nu}$; the justification for this is that it is the only\\
thing that works at the quantum level, in spite of the ambiguity in the definition of $A^{\mu}$. Moreover, there is an implicit assumption that all fields and charges vanish at spatial infinity in such a way that these 4 -dimensional integrals converge and all boundary terms in spatial integrations by parts vanish. In the preceding paragraphs we developed the Lagrangian for a particle moving in a given electromagnetic field; it was the sum of the Lagrangian for a free particle and a Lagrangian for the interaction. Similarly, we now develop the Lagrangian for a field in the presence of a given system of charges and currents; it will be the sum of a Lagrangian involving only the field and a Lagrangian for the interaction.

The form of the interaction term is suggested by the $-q A_{\mu} d x^{\mu} / d t=-q \phi+q \mathbf{A} \cdot \mathbf{v}$ that we found before. That was for a single charge $q$; if instead we consider a continuous distribution of charge with charge density $\rho$ and current density $\mathbf{j}$, the obvious analogue is $-A_{\mu} j^{\mu}=-\rho \phi+\mathbf{A} \cdot \mathbf{j}$ where $j^{\mu}=(\rho, \mathbf{j})$ is the 4-current density. Thus, the interaction part of the action functional will be

$$
S_{\mathrm{int}}=-\int_{\mathbb{R}^{4}} A_{\mu} j^{\mu} d^{4} x
$$

As before, the ambiguity in $A^{\mu}$ is irrelevant: If we change $A_{\mu}$ by adding $\partial_{\mu} \chi$, where $\chi$ vanishes at infinity, the change in the action is

$$
-\int\left(\partial_{\mu} \chi\right) j^{\mu} d^{4} x=\int \chi \partial_{\mu} j^{\mu} d^{4} x=0
$$

because of (2.29).\\
Now, what about the free-field Lagrangian? It should be Lorentz invariant; it should be unaffected by gauge transformations; it should be quadratic so that the field equations will turn out to be linear (Maxwell's equations). The only possibility is a constant multiple of $F_{\mu \nu} F^{\mu \nu}$ (recall that $F_{\mu \nu}=\partial_{\mu} A_{\nu}-\partial_{\nu} A_{\mu}$ ), and the correct constant turns out to be $-\frac{1}{4}$. Thus the full action functional is

$$
S=S_{\mathrm{field}}+S_{\mathrm{int}}=\int_{\mathbb{R}^{4}}\left(-\frac{1}{4} F_{\mu \nu} F^{\mu \nu}-A_{\mu} j^{\mu}\right) d^{4} x
$$

(In the classical notation, $-\frac{1}{4} F_{\mu \nu} F^{\mu \nu}=\frac{1}{2}\left(|\mathbf{E}|^{2}-|\mathbf{B}|^{2}\right)$.)\\
To confirm the validity of (2.33), let us derive the field equations. If we add a perturbation $\delta A^{\mu}$ that vanishes at (spatial and temporal) infinity to $A^{\mu}$, the change in $S$ to first order is

$$
\begin{aligned}
\delta S & =\int\left(-\frac{1}{2} F^{\mu \nu} \delta F_{\mu \nu}-j^{\nu} \delta A_{\nu}\right) d^{4} x \\
& =\int\left(-\frac{1}{2} F^{\mu \nu}\left[\partial_{\mu}\left(\delta A_{\nu}\right)-\partial_{\nu}\left(\delta A_{\mu}\right)\right]-j^{\nu} \delta A_{\nu}\right) d^{4} x \\
& =\int\left(-F^{\mu \nu} \partial_{\mu}\left(\delta A_{\nu}\right)-j^{\nu} \delta A_{\nu}\right) d^{4} x \\
& =\int\left(\partial_{\mu} F^{\mu \nu}-j^{\nu}\right) \delta A_{\nu} d^{4} x
\end{aligned}
$$

In the first line used the fact that $F_{\mu \nu} \delta F^{\mu \nu}=F^{\mu \nu} \delta F_{\mu \nu}$, and in passing from the second to the third we used the fact that $-F^{\mu \nu} \partial_{\nu} \delta A_{\mu}=-F^{\nu \mu} \partial_{\mu} \delta A_{\nu}=F^{\mu \nu} \partial_{\mu} \delta A_{\nu}$ (by relabeling indices and using the skew symmetry of $F^{\mu \nu}$ ). Thus the condition $\delta S=0$ yields the Maxwell equation (2.28). (Recall that the other equation (2.27) is a consequence of the formula $F_{\mu \nu}=\partial_{\mu} A_{\nu}-\partial_{\nu} A_{\mu}$.)

We have developed a Lagrangian formulation of the evolution equations for a charged particle moving in a given electromagnetic field and for an electromagnetic field in the presence of a given charge-current distribution. An attempt to combine these into a relativistically correct Lagrangian that describes a system of charges and fields interacting with each other is problematic, largely because of the question of finding a suitable form for the pure-matter term (the analogue of $-\int m d s$ for a single particle); see Goldstein [56]. When we come to quantum electrodynamics, however, this problem will evaporate. The charge-current distribution will be derived from an electron field, and the pure-matter term will be its free-field Lagrangian.

\section*{CHAPTER 3}
\begin{displayquote}
Raffiniert ist der Herrgott, aber boshaft ist er nicht. [God is subtle, but he is not malicious.]\\
-Albert Einstein
\end{displayquote}

This chapter is a brief exposition of the fundamentals of quantum theory from a mathematical point of view. Our development of the underlying mathematical structure follows Mackey [79]; see also Jauch [68] and Varadarajan [125]. For quantum mechanics proper, there are many good texts available, so I will just mention three old classics that I have found valuable: the introductory account in the Feynman Lectures [42] and the comprehensive treatises of Messiah [83] and Landau and Lifshitz [77]; the latter is particularly good from a mathematical point of view. Reed and Simon [94] is an excellent source for the relevant functional analysis; all the results about Hilbert spaces and their operators that are quoted in this chapter without references can be found there. We also need a few facts about Lie groups, Lie algebras, and their representations; we refer the reader to Hall [62] for more background.

\subsection*{The mathematical framework}
The fundamental objects in a mathematical description of a physical system are states and observables. Observables are the physical quantities such as position, momentum, energy, etc., that one is interested in studying. In this discussion we shall consider only real-valued observables. This is not much of a restriction: for vector-valued observables one can consider their components with respect to a basis, and one can generally assign numerical labels to the values of other observables. In particular, one can think of true-false statements about the system as observables whose only values are 1 and 0 . (Mackey [79] calls such observables "questions," but we prefer to think of them as assertions.) The state of a system is supposed to be a specification of the condition of the system from which one can read off all the available information about the observables of interest; the set of all possible states is called the state space.

In Hamiltonian mechanics, the state space is a symplectic manifold $M$, the observables are (Borel measurable) real-valued functions on $M$, and the true-false observables can be identified with the (Borel) subsets of $M$ (the observable $f: M \rightarrow \{0,1\}$ corresponds to $\{x \in M: f(x)=1\}$ ). Thus the set of true-false statements on $M$ is identified with the Borel $\sigma$-algebra on $M$, and the basic logical connectives "or," "and," and "not" correspond to the Boolean operations of union, intersection, and complement.

General observables can be analyzed in terms of true-false statements. A Borel function $f: M \rightarrow \mathbb{R}$ is completely specified by its level sets $f^{-1}(\{a\}), a \in \mathbb{R}$, or\\
more generally by the sets $f^{-1}(E), E$ a Borel subset of $\mathbb{R}$. That is, an observable $f$ is specified by the sets on which the statements " $f(x)=a$ " or " $f(x) \in E$ " are true.

The striking difference between classical and quantum mechanics is that in a quantum system, the the true-false statements do not form a Boolean algebra. Specifically, what fails is the distributive law:

$$
(A \text { or } B) \text { and } C \quad \Longleftrightarrow \quad(A \text { and } C) \text { or }(B \text { and } C) .
$$

The classic counterexample is the double-slit experiment in which a beam of electrons is directed at a barrier with two slits in it and the electrons that pass through the slits arrive at a screen where they are detected. Thus, a given electron passes through one of the two slits ( $A$ or $B$ ) and arrives at the screen ( $C$ ). The resulting pattern of arrivals at the screen exhibits the sort of interference oscillations one would expect if the electrons were waves rather than particles. But if one puts detectors at the slits to see which slit an electron goes through (and thus changes the process to ( $A$ and $C$ ) or ( $B$ and $C$ )), the interference pattern disappears. See Feynman [42], Chapter 3 of vol. 3, for a detailed discussion.

Thus one must find a new model for states and observables in quantum mechanics, and what turns out to work is the following. The state space of a quantum system is a projective Hilbert space $\mathbb{P H}$, that is, the set of nonzero vectors in a Hilbert space $\mathcal{H}$ modulo the equivalence relation that $u \sim v$ if and only if $v$ is a scalar multiple of $u$. In practice, one usually thinks of $\mathcal{H}$ itself as the state space, with the understanding that two equivalent vectors define the same state. It is almost always desirable to take the vectors representing states to be unit vectors; one speaks of normalized states.

At this point we refer the reader back to §1.1 for a discussion of the physicists' notation and terminology that we will be using from now on.

The true-false statements about the quantum system correspond not to subsets of $\mathcal{H}$ but to closed linear subspaces of $\mathcal{H}$. The logical "and" still corresponds to intersection, but now "or" and "not" correspond to closed linear span and orthogonal complement, respectively. Thus the lattice of subspaces takes the place of the Boolean algebra of subsets. However, to develop the theory further, it is more convenient to identify a closed linear subspace with the orthogonal projection onto that subspace. Thus, from now on, true-false statements will be identified with orthogonal projection operators.

If the system is in the state represented by the unit vector $u$, and a true-false statement is represented by the projection $P$, the statement is definitely true if $u$ is in the range $\mathcal{R}(P)$ and definitely false if $u$ is in the nullspace $\mathcal{N}(P)$. In general, we write $u=u_{0}+u_{1}$ where $u_{0} \in \mathcal{N}(P)$ and $u_{1} \in \mathcal{R}(P)$. Then $\left\|u_{0}\right\|^{2}+\left\|u_{1}\right\|^{2}=1$, and we interpret $\left\|u_{0}\right\|^{2}$ and $\left\|u_{1}\right\|^{2}$ as the probabilities that the statement is false and true, respectively. (In experimental terms, probability is interpreted in terms of statistical frequency: If the system is prepared $N$ times in state $u$ where $N$ is large, the statement will be true about $N\left\|u_{1}\right\|^{2}$ times.) This probabilistic interpretation is an essential feature of quantum mechanics, not just the result of an incomplete description of the system.

The state $u$ can be identified with the true-false statement "The system is in the state $u$." Taking $u$ to be a unit vector, the corresponding projection is the orthogonal projection $P$ onto the line spanned by $u, P v=\langle u \mid v\rangle u$, and for any unit\\
vector $\left.v,|\langle u \mid v\rangle|^{2}=|\langle u| P| v\right\rangle\left.\right|^{2}$ is the transition probability that a system known to be in the state $u$ will also be found to be in the state $v$ (or vice versa: $|\langle u \mid v\rangle|^{2}$ is symmetric in $u$ and $v$ ). The fact that a system that is definitely in one state may also be in a different state is another peculiarity of quantum systems.

By the way, physicists refer to the inner product $\langle u \mid v\rangle$ as a probability amplitude. This usage of the word "amplitude" is in conflict with the classical one according to which the amplitude of a sinusoidal wave is the coefficient $a$ (always assumed positive) in the expression $a \sin (\omega t+c)$. This is just one of those semantic peculiarities that one has to learn to live with.

The usual device for combining simpler quantum systems into more complicated ones is the tensor product of Hilbert spaces, a concept which we now briefly review. Given two Hilbert spaces $\mathcal{H}_{1}$ and $\mathcal{H}_{2}$, one way to define their (Hilbert-space) tensor product $\mathcal{H}_{1} \otimes \mathcal{H}_{2}$ is to begin with their algebraic tensor product (the set of finite linear combinations of products $u_{1} \otimes u_{2}$ where $u_{j} \in \mathcal{H}_{j}$ ). There is a unique inner product on this space such that

$$
\left\langle u_{1} \otimes u_{2} \mid v_{1} \otimes v_{2}\right\rangle=\left\langle u_{1} \mid v_{1}\right\rangle\left\langle u_{2} \mid v_{2}\right\rangle,
$$

and $\mathcal{H}_{1} \otimes \mathcal{H}_{2}$ is its completion with respect to the associated norm. Alternatively, one can define $\mathcal{H}_{1} \otimes \mathcal{H}_{2}$ to be the space of all bounded conjugate-linear maps $A: \mathcal{H}_{1} \rightarrow \mathcal{H}_{2}$ such that $\sum\left\|A e_{j}\right\|^{2}<\infty$ for some (and hence any) orthonormal basis $\left\{e_{j}\right\}$ for $\mathcal{H}_{1}$, equipped with the inner product $\langle A \mid B\rangle=\sum\left\langle A e_{j} \mid B e_{j}\right\rangle$. (In this picture $u_{1} \otimes u_{2}$ is the map $v \mapsto\left\langle v \mid u_{1}\right\rangle u_{2}$.) Either way, if $\left\{e_{j}\right\}$ and $\left\{f_{k}\right\}$ are orthonormal bases for $\mathcal{H}_{1}$ and $\mathcal{H}_{2},\left\{e_{j} \otimes f_{k}\right\}$ is an orthonormal basis for $\mathcal{H}_{1} \otimes \mathcal{H}_{2}$. Other variations on this theme are possible; see Reed and Simon [94], §II.4, or Folland [45], §7.3 for more details.

Here are two typical uses of tensor products in quantum mechanics. First, if $\mathcal{H}_{1}$ and $\mathcal{H}_{2}$ are the quantum state spaces for two particles $p_{1}$ and $p_{2}$, the state space for the two-particle system is $\mathcal{H}_{1} \otimes \mathcal{H}_{2}$, with $u \otimes v$ representing the state in which $p_{1}$ is in state $u$ and $p_{2}$ is in state $v .^{1}$ (In this situation it is often appropriate to take $\mathcal{H}_{1}=\mathcal{H}_{2}=L^{2}\left(\mathbb{R}^{3}\right)$, in which case $\mathcal{H}_{1} \otimes \mathcal{H}_{2}$ can be identified with $L^{2}\left(\mathbb{R}^{6}\right)$ by taking $f_{1} \otimes f_{2}$ to be the function $\left(x_{1}, x_{2}\right) \mapsto f_{1}\left(x_{1}\right) f_{2}\left(x_{2}\right)$.) Second, just as the motion of a classical rigid body may be analyzed on one level by considering it as a featureless "particle" and then in more detail by considering its rotations about its center of mass, a quantum particle may have one set of "states" (i.e., one Hilbert space $\mathcal{H}_{1}$ ) to describe its motion in space and another one $\mathcal{H}_{2}$ to describe its "internal degrees of freedom" such as spin; then the full state space for the particle will be $\mathcal{H}_{1} \otimes \mathcal{H}_{2}$.

Now, how do we model observables? Given a real-valued physical quantity $O$ that we want to be an observable of our system, we follow the hint indicated earlier and analyze it in terms of the true-false statements

$$
S(E) \text { : The value of } O \text { is in } E \text {, }
$$

where $E$ a Borel subset of $\mathbb{R}$. Let $P(E)$ be the projection corresponding to $S(E)$. Obviously $S(\varnothing)$ is always false and $S(\mathbb{R})$ is always true, so

$$
P(\varnothing)=0 \quad \text { and } \quad P(\mathbb{R})=I .
$$

If $E_{1}$ and $E_{2}$ are disjoint, $S\left(E_{1}\right)$ is definitely false whenever $S\left(E_{2}\right)$ is definitely true, and vice versa, so the ranges of $P\left(E_{1}\right)$ and $P\left(E_{2}\right)$ are mutually orthogonal;

\footnotetext{${ }^{1}$ If $p_{1}$ and $p_{2}$ are identical - two electrons, for example - this picture needs to be modified; see §4.5.
}
equivalently,
$$
P\left(E_{1}\right) P\left(E_{2}\right)=0 \text { whenever } E_{1} \cap E_{2}=\varnothing .
$$

In this case, the closed linear span of the ranges is their orthogonal sum, and the projection onto the latter space is $P\left(E_{1}\right)+P\left(E_{2}\right)$. But that projection corresponds to the statement " $O$ is in $E_{1}$ or $E_{2}$," which is the same as $S\left(E_{1} \cup E_{2}\right)$ and hence corresponds to the projection $P\left(E_{1} \cup E_{2}\right)$. It then follows by induction that $P\left(\bigcup_{1}^{n} E_{j}\right)=\sum_{1}^{n} P\left(E_{j}\right)$ whenever $E_{1}, \ldots, E_{n}$ are disjoint, and it is eminently reasonable to assume that the same relation holds for countable unions:

$$
P\left(\bigcup_{1}^{\infty} E_{j}\right)=\sum_{1}^{\infty} P\left(E_{j}\right) \text { whenever } E_{i} \cap E_{j}=\varnothing \text { for } i \neq j .
$$

(The convergence of the sum on the right is in the strong operator topology, i.e., the topology of pointwise convergence in norm.) A map $E \mapsto P(E)$ from the Borel sets in $\mathbb{R}$ to the orthogonal projections on $\mathcal{H}$ that satisfies (3.1)-(3.3) is called a projection-valued measure on $\mathbb{R}$.

It is worth observing that projection-valued measures behave in the natural way with respect to complements and intersections. Namely, since $P(\mathbb{R} \backslash E)+P(E)= P(\mathbb{R})=I, P(\mathbb{R} \backslash E)$ is the projection whose range and nullspace are the nullspace and range of $P(E)$, respectively. Moreover, for any $E_{1}$ and $E_{2}$, the sets $E_{1} \backslash E_{2}$, $E_{1} \cap E_{2}$, and $E_{2} \backslash E_{1}$ are disjoint, so by (3.2) and (3.3),

$$
\begin{aligned}
P\left(E_{1}\right) P\left(E_{2}\right) & =\left[P\left(E_{1} \backslash E_{2}\right)+P\left(E_{1} \cap E_{2}\right)\right]\left[P\left(E_{1} \cap E_{2}\right)+P\left(E_{2} \backslash E_{1}\right)\right] \\
& =P\left(E_{1} \cap E_{2}\right)^{2}=P\left(E_{1} \cap E_{2}\right)
\end{aligned}
$$

If $P$ is any projection-valued measure on $\mathbb{R}$ and $u$ is a unit vector, the map $P_{u}: E \mapsto\langle u| P(E)|u\rangle$ is an ordinary probability measure on $\mathbb{R}$. It is interpreted as the probability distribution of the observable $P$ in the state $u$. This is consistent with the probabilistic interpretation of single projections mentioned earlier. The states (if any) in which the observable $P$ has a definite value $a$ are those in the range of $P(\{a\})$.

Now comes the leap into analysis: By the spectral theorem, projection-valued measures are in one-to-one correspondence with self-adjoint operators. To wit, if $A$ is a self-adjoint operator, the spectral functional calculus yields a self-adjoint operator $f(A)$ for any Borel function $f: \mathbb{R} \rightarrow \mathbb{R}$, and the operators $P(E)=\chi_{E}(A)$ ( $\chi=$ characteristic function) form a projection-valued measure. On the other hand, if $P$ is a projection-valued measure, there is a unique self-adjoint operator $A$, whose domain $\mathcal{D}(A)$ is the set of all $u \in \mathcal{H}$ with $\int_{\mathbb{R}} t^{2} d P_{u}(t)<\infty$, such that $\langle u| A|u\rangle= \int_{\mathbb{R}} t d P_{u}(t)$ for all $u \in \mathcal{D}(A)$. (The diagonal matrix elements $\langle u| A|u\rangle$ determine all other matrix elements $\langle u| A|v\rangle$, and hence $A$ itself, by polarization.) Moreover, these correspondences $A \mapsto P$ and $P \mapsto A$ are mutually inverse. In short, observables may be identified with self-adjoint operators, and we shall do so henceforth.

The diagonal matrix elements $\langle u| A|u\rangle$ have an immediate physical interpretation: $\langle u| A|u\rangle=\int_{\mathbb{R}} t d P_{u}(t)$ is the expected value of the probability distribution $P_{u}$, and it therefore represents the expected value - or, as physicists usually say, the expectation value - of the observable $A$ in the state $u$. Moreover, the states in which the observable $A$ has a definite value $a$ are simply the eigenvectors of $A$ with eigenvalue $a$.

Incidentally, this is the appropriate point to mention a bit of physicists' argot that the reader may encounter occasionally. In the early days of quantum mechanics when the idea that observables are represented by noncommuting operators was still new and strange, some people spoke of quantum observables as "quantities whose values are $q$-numbers" - the notion of " $q$-number" being meant to suggest noncommutativity - as opposed to "quantities whose values are $c$-numbers," i.e., ordinary complex-valued quantities whose algebra is commutative. One still finds the terms $q$-number (rarely) and $c$-number (more frequently) in the physics literature; in particular, to say that an operator is a $c$-number is to say that it is a scalar multiple of the identity. (E.g., "The commutator of $A$ and $B$ is just a $c$-number.") That usage is echoed in this book by our tendency to write just $\lambda$ instead of $\lambda I$ for a scalar multiple of the identity.

The reader will notice that we have entered the perilous world of unbounded operators. Bounded self-adjoint operators correspond to projection-valued measures $P$ such that $P(E)=I$ for some bounded set $E$. But many important observables can potentially have values anywhere on the real line, so the corresponding operators will be unbounded. In this situation the word "self-adjoint" has to be understood in its precise technical sense. Although these technicalities play a remarkably small role in quantum field theory, we had better take a little time to set them out clearly.

To recall the precise definitions: let $A$ be an operator defined on the domain $\mathcal{D}(A) \subset \mathcal{H}$. If $\mathcal{D}(A)$ is dense in $\mathcal{H}$, the adjoint of $A$ is the operator $A^{\dagger}$ on the domain of all $u \in \mathcal{H}$ such that the map $v \mapsto\langle u \mid A v\rangle(v \in \mathcal{D}(A))$ extends to a bounded linear functional on all of $\mathcal{H}$ (uniquely, since $\mathcal{D}(A)$ is dense); $A^{\dagger} u$ is the element of $\mathcal{H}$ corresponding to this functional, so that $\left\langle A^{\dagger} u \mid v\right\rangle=\langle u \mid A v\rangle . A$ is called symmetric if $\langle A u \mid v\rangle=\langle u \mid A v\rangle$ for all $u, v \in \mathcal{D}(A)$. This means that $A^{\dagger}$ is an extension of $A$, i.e. $\mathcal{D}\left(A^{\dagger}\right) \supset \mathcal{D}(A)$ and $A^{\dagger} \mid \mathfrak{D}(A)=A$. $A$ is called self-adjoint if $A^{\dagger}=A$. If $A$ is symmetric, it may be possible to extend $A$ to a larger domain to make it self-adjoint, perhaps in many different ways. $A$ is called essentially selfadjoint if it has a unique self-adjoint extension, or equivalently, if the graph of $A^{\dagger}$ in $\mathcal{H} \times \mathcal{H}$ is the closure of the graph of $A$. The principal reason for insisting upon these distinctions is that among the symmetric operators, the self-adjoint ones are the only ones to which the spectral theorem and its consequences apply.

In the mathematics literature the word Hermitian is usually a synonym for symmetric. In the physics literature, too, if one looks for a definition of "Hermitian" one will usually find it stated that $A$ is Hermitian if $\langle u| A|v\rangle=\langle v| A|u\rangle^{*}$. However, when physicists speak of Hermitian operators, they often want self-adjoint ones, for those are the ones with a good spectral theory and the ones that really correspond to observables. Generally this is not a serious problem; the diligent mathematician can usually figure out what the appropriate domain is with a certain amount of work. In particular, when the Hilbert space is something like $L^{2}\left(\mathbb{R}^{n}\right)$, it often happens that a symmetric operator is essentially self-adjoint when considered on an obvious domain of "nice" functions, and then there is no problem. (The standard position and momentum observables defined by (3.12) below are good examples. It takes some thought to come up with a dense subspace of their domains on which they are not essentially self-adjoint.) On the other hand, when the operator in question is a differential operator on a region in $\mathbb{R}^{n}$, the question of specifying an appropriate domain usually amounts to choosing appropriate boundary conditions, and these conditions are usually dictated by physical considerations.

In any case, we shall be careful to use the term "self-adjoint" when it is really an issue, and we shall generally follow the physicists in using the term "Hermitian" when it is not.

For those who are not on easy terms with these ideas, it may be useful to consider an example. Let $\mathcal{H}=L^{2}([0,1])$, and let $A=d^{2} / d x^{2}$ on the domain of all $f \in C^{1}([0,1])$ such that $f^{\prime}$ is absolutely continuous and $f^{\prime \prime}$ (defined a.e.) is in $L^{2}$. The fundamental formula arises from integration by parts:

$$
\begin{aligned}
\langle A f \mid g\rangle & -\langle f \mid A g\rangle \\
& =\int_{0}^{1}\left[f^{\prime \prime}(x)^{*} g(x)-f(x)^{*} g^{\prime \prime}(x)\right] d x=\left[f^{\prime}(x)^{*} g(x)-f(x)^{*} g^{\prime}(x)\right]_{0}^{1}
\end{aligned}
$$

(Remember that * denotes complex conjugation!) Thus $A$ is not symmetric. However, the restriction $A_{0}$ of $A$ to $\left\{f \in \mathcal{D}(A): f(0)=f^{\prime}(0)=f(1)=f^{\prime}(1)=0\right\}$ is symmetric, and it is not hard to show that $A=A_{0}^{\dagger}$. Note that $\mathcal{D}\left(A_{0}\right)$ is defined by the imposition of four boundary conditions; it is too small to give a self-adjoint operator. The self-adjoint extensions of $A_{0}$ are defined by restricting $A$ to a domain defined by two boundary conditions in such a way that the boundary terms in (3.4) vanish for all $f$ and $g$ satisfying the conditions. There is, in fact, a fourparameter family of such conditions (naturally parametrized by the unitary group $U(2)$ ), and the resulting operators all have different spectral resolutions. Three cases will suffice to indicate the variety of possibilities:\\
i. $A_{1}$ is the restriction of $A$ to $\{f \in \mathcal{D}(A): f(0)=f(1)=0\}$. The eigenvalues of $A_{1}$ are $-n^{2} \pi^{2}, n \geq 1$, each with multiplicity 1 , and the corresponding eigenfunctions are $\sin n \pi x$.\\
ii. $A_{2}$ is the restriction of $A$ to $\left\{f \in \mathcal{D}(A): f(0)=f^{\prime}(1)=0\right\}$. The eigenvalues of $A_{2}$ are $-\left(n-\frac{1}{2}\right)^{2} \pi^{2}, n \geq 1$, each with multiplicity 1 , and the corresponding eigenfunctions are $\sin \left(n-\frac{1}{2}\right) \pi x$.\\
iii. $A_{3}$ is the restriction of $A$ to $\left\{f \in \mathcal{D}(A): f(0)=f(1)\right.$ and $\left.f^{\prime}(0)=f^{\prime}(1)\right\}$. The eigenvalues are $-4 n^{2} \pi^{2}, n \geq 0$; the eigenvalue 0 has multiplicity 1 and all others have multiplicity 2 , and the corresponding eigenfunctions are linear combinations of $\cos 2 n \pi x$ and $\sin 2 n \pi x$.\\
In each of these cases the operator has an orthonormal eigenbasis, and each of these bases has a nice interpretation in terms of the classical physics of one-dimensional vibrations. The first one models the normal modes of a string fixed at both ends; the second models the normal modes of a cylindrical pipe (an organ pipe or clarinet, say) that is closed at one end and open at the other; and the third models the normal modes of a vibrating circular hoop. So the physicists may not talk about domains of unbounded operators, but they have no trouble distinguishing these cases! On the other hand, to bring home the point about good spectral theory, let us add $A$ and $A_{0}$ to the list.\\
iv. Every $\lambda \in \mathbb{C}$ is an eigenvalue of $A$ with multiplicity 2 ; the eigenfunctions are linear combinations of $e^{\sqrt{-\lambda} x}$ and $e^{-\sqrt{-\lambda} x}$ for $\lambda \neq 0$ and of 1 and $x$ for $\lambda=0$. There is no canonical eigenfunction expansion associated to $A$.\\
v. The operator $A_{0}$ has no eigenvalues at all. However, its spectrum, the set of all $\lambda \in \mathbb{C}$ such that $A_{0}-\lambda$ is not invertible, is the whole complex plane because $A_{0}-\lambda$ is never surjective; in fact, its range isn't even dense. (The orthogonal complement of its range consists of the eigenfunctions of $A$ with eigenvalue $\lambda^{*}$.) Again there is no good spectral resolution.

If one replaces the condition " $f \in \mathcal{D}(A)$ " with " $f \in C^{2}([0,1])$ " or " $f \in C^{\infty}([0,1])$ " in the definition of $A_{1}, A_{2}$, or $A_{3}$, the resulting operator is essentially self-adjoint. But adding an extra boundary condition, or taking one away, really destroys the self-adjointness.

Next we consider symmetries of the quantum system. These are modeled by automorphisms of the projective Hilbert space $\mathbb{P H}$, that is, the bijections from $\mathbb{P H}$ to itself that preserve the projectively invariant part of the inner product,

$$
(u, v)=\frac{|\langle u \mid v\rangle|}{\|u\|\|v\|} \quad(u, v \neq 0)
$$

(Geometrically, $(u, v)$ is the cosine of the angle between the complex lines spanned by $u$ and $v$. Physically, $(u, v)^{2}$ is the transition probability between the states $u$ and v.) Clearly any unitary operator on $\mathcal{H}$ induces an automorphism of $\mathbb{P} \mathcal{H}$, as does any anti-unitary operator (a conjugate-linear operator $V$ such that $\langle u| V|w\rangle=\langle w \mid u\rangle$ ). A theorem of Wigner (see Bargmann [7] or Jauch [68]) assures us that these are the only automorphisms of $\mathbb{P} \mathcal{H}$. Thus the automorphism group can be identified with the group of unitary or antiunitary operators on $\mathcal{H}$ modulo scalar multiples of the identity operator (which induce the identity transformation on $\mathbb{P H}$ ). We equip it with the topology inherited from the strong operator topology; the unitary and antiunitary parts are then its connected components. (The group of unitary operators is connected because, by spectral theory, every unitary operator belongs to a continuous one-parameter group of unitary operators and so is connected to the identity by a continuous path.) We denote the component of the identity, $\mathcal{U}(\mathcal{H}) /\{c I:|c|=1\}$, by $\mathbb{P} \mathcal{U}(\mathcal{H})$.

Let us pass to continuous groups of symmetries. If $G$ is a connected topological group, a projective representation or ray representation of $G$ is a continuous homomorphism $\rho: G \rightarrow \mathbb{P u}(\mathcal{H})$. The question that immediately arises is whether such a $\rho$ can be lifted to a continuous homomorphism $\widetilde{\rho}: G \rightarrow \mathcal{U}(\mathcal{H})$, i.e., a unitary representation of $G$. There are two possible obstructions, one topological and one cohomological. The topological problem is that if $G$ is not simply connected, it may not even be possible to lift $\rho$ to a continuous map from $G$ into $\mathcal{U}(\mathcal{H})$, let alone a homomorphism. (Examples are provided by the double-valued representations of the rotation group; see §3.5.) The solution to this problem is to pass to the universal cover of $G$. Assuming then that $G$ is simply connected, a continuous lifting $\widehat{\rho}: G \rightarrow \mathcal{U}(\mathcal{H})$ always exists, and it satisfies $\widehat{\rho}(x y)=\omega(x, y) \widehat{\rho}(x) \widehat{\rho}(y)$ where $\omega(x, y)$ is a numerical factor of absolute value 1. Moreover, the associative law for $G$ implies that $\omega$ satisfies the cocycle condition $\omega(x y, z) \omega(x, y)=\omega(x, y z) \omega(y, z)$. If $\omega$ is a coboundary, that is, if $\omega(x, y)=\lambda(x) \lambda(y) / \lambda(x y)$ for some continuous $\lambda: G \rightarrow \mathbf{T}^{1}$, then $\widetilde{\rho}(x)=\lambda(x) \widehat{\rho}(x)$ is the desired homomorphic lift of $\rho$. (Those familiar with group cohomology will recognize the general framework into which the words "cocycle" and "coboundary" should be placed; others need not worry about it.) If $\omega$ is not a coboundary, it determines a group $\widetilde{G}$ that fits into an exact sequence $1 \rightarrow U(1) \rightarrow \widetilde{G} \rightarrow G \rightarrow 1$ (that is, $\widetilde{G}$ is an extension of $G$ by the circle group $U(1))$, and the best one can do is to find a unitary representation $\widetilde{\rho}$ of $\widetilde{G}$ such that $\pi_{\mathcal{U}} \circ \widetilde{\rho}=\rho \circ \pi_{G}$, where $\pi_{\mathcal{U}}$ and $\pi_{G}$ are the projections from $\mathcal{U}(\mathcal{H})$ to $\mathbb{P} \mathcal{U}(\mathcal{H})$ and from $\widetilde{G}$ to $G$.

When $G$ is a simply connected Lie group, the problem in group cohomology adumbrated above can be reduced to a problem in Lie algebra cohomology, which is\\
more readily computable. The relevant cohomology group is $H^{2}(\mathfrak{g}, \mathbb{R})$, which is the quotient of the space of all skew-symmetric bilinear maps $\phi: \mathfrak{g} \times \mathfrak{g} \rightarrow \mathbb{R}$ that satisfy the cocycle condition $\phi(X,[Y, Z])+\phi(Y,[Z, X])+\phi(Z,[X, Y])=0$ by the subspace of all such $\phi$ of the form $\phi(X, Y)=f([X, Y])$ for some linear functional $f$ on $\mathfrak{g}$. The upshot is Bargmann's theorem [6] (see also Simms [112] and Varadarajan [125]):

Let $G$ be a simply connected Lie group with Lie algebra $\mathfrak{g}$. Every projective representation $\rho$ of $G$ determines an element $\beta_{\rho} \in H^{2}(\mathfrak{g}, \mathbb{R})$, and $\rho$ can be lifted to a unitary representation of $G$ if and only if $\beta_{\rho}=0$. In particular, if $H^{2}(\mathfrak{g}, \mathbb{R})=0$, every projective representation of $G$ can be lifted to a unitary representation of $G$.

The cohomology group $H^{2}(\mathfrak{g}, \mathbb{R})$ vanishes whenever $G$ is semisimple, and also when $G=\mathbb{R}$ and when $G$ is the Poincaré group. (It does not vanish when $G=\mathbb{R}^{n}$ for $n>1$, however. Examples of projective representations of $\mathbb{R}^{n}$ for even $n$ that cannot be lifted to unitary representations are provided by the representations of the canonical commutation relations that we shall discuss in §3.2.) For the time being, the essential case for us is $G=\mathbb{R}$ : Every one-parameter subgroup of $\mathbb{P} \mathcal{U}(\mathcal{H})$ comes from a one-parameter group of unitary operators on $\mathcal{H}$.

We are now in a position to apply another of the big theorems of functional analysis, the Stone representation theorem for (strongly continuous) one-parameter unitary groups. Recall that an operator $S$ is called skew-adjoint if $S^{\dagger}=-S$. If $S$ is any skew-adjoint operator, then $U(t)=e^{t S}$ (defined by the spectral functional calculus) is a strongly continuous one-parameter group of unitary operators, and Stone's theorem asserts that every strongly continuous one-parameter group of unitary operators is of this form. Moreover, skew-adjoint and self-adjoint operators are almost the same thing: multiplication by a purely imaginary scalar converts one into the other. Thus, if one fixes a nonzero real constant $b$, the map $A \mapsto i b A$ defines a bijection from the self-adjoint to the skew-adjoint operators. (Why not take $b=1$ ? There are good reasons: more about this in the next section.) Thus there is a correspondence between observables (self-adjoint operators) and oneparameter groups of symmetries of a quantum system: $A \leftrightarrow e^{i b t A}$.

Let us explore this a little further. A one-parameter group $e^{i b t A}$ can be considered as transforming the states by $u \mapsto u_{t}=e^{i b t A} u$ while the observables remain fixed, or as transforming the observables by $B \mapsto B_{t}=e^{-i b t A} B e^{i b t A}$ while the states remain fixed. These are known as the Schrödinger and Heisenberg pictures, respectively. If $P$ is the projection-valued measure associated to $B$, then $P_{t}(\cdot)=e^{-i b t A} P(\cdot) e^{i b t A}$ is the projection-valued measure associated to $B_{t}$, and the probability distribution of the observable $B$ in the state $u$, after the action of $e^{i b t A}$, is $E \mapsto\left\langle u_{t}\right| P(E)\left|u_{t}\right\rangle=\langle u| P_{t}(E)|u\rangle$.

In the Schrödinger picture, the differential equation governing the evolution of $u_{t}$ is

$$
\frac{d u_{t}}{d t}=i b A u_{t}
$$

the (abstract) Schrödinger equation. (More precisely, this equation is valid when $u$ is in the domain of $A$; this domain is invariant under the group $e^{i b t A}$.) In the Heisenberg picture, the differential equation governing the evolution of $B_{t}$ is

$$
\frac{d B_{t}}{d t}=-i b A B_{t}+i b B_{t} A=i b\left[B_{t}, A\right]=i b[B, A]_{t}
$$

(More precisely, this equation is valid as an operator identity on the domain of all $u \in D\left(B_{t}\right) \cap \mathcal{D}(A)$ such that $B_{t} u \in D(A)$ and $A u \in D\left(B_{t}\right)$ for all $t$.)

The resemblance to Hamilton's equation (2.3) is notable, and it leads us to think about the analogy between Poisson brackets and commutators. Both operations are Lie products, in the sense of being skew-symmetric and satisfying the Jacobi identity. The set of $C^{\infty}$ functions on a symplectic manifold forms a Lie algebra under Poisson brackets, and the set of bounded skew-adjoint operators on a Hilbert space forms a Lie algebra under commutators. The correspondence $A \leftrightarrow i b A$ between self-adjoint and skew-adjoint operators then induces a Lie algebra structure on the set of bounded self-adjoint operators by

$$
\left(A_{1}, A_{2}\right) \mapsto \frac{\left[i b A_{1}, i b A_{2}\right]}{i b}=i b\left[A_{1}, A_{2}\right]
$$

Alas, one sometimes needs to consider classical observables that aren't $C^{\infty}$ and quantum observables that aren't bounded, so the Lie algebra structures aren't perfect. The classical problem is not serious; one just has to remember that the Poisson bracket of two $C^{k}$ functions is only $C^{k-1}$. On the quantum side, if $A_{1}$ and $A_{2}$ are self-adjoint operators, not both bounded, $i b\left[A_{1}, A_{2}\right]$ is at least a symmetric operator on the domain of all $u \in \mathcal{D}\left(A_{1}\right) \cap \mathcal{D}\left(A_{2}\right)$ such that $A_{1} u \in \mathcal{D}\left(A_{2}\right)$ and $A_{2} u \in \mathcal{D}\left(A_{1}\right)$, but in general that domain will be too small for $i b\left[A_{1}, A_{2}\right]$ to be self-adjoint. (If one is fortunate, it will be essentially self-adjoint. There are substantial technical issues to be dealt with here, but they do not concern us now.)

Let us pause to take stock of the analogies between the formalisms of Hamiltonian and quantum mechanics.

\begin{itemize}
  \item In Hamiltonian mechanics, the state space is phase space, a symplectic manifold $M$; the observables are real-valued functions on $M$, and the symmetries are canonical transformations of $M$. Every (smooth) observable defines a Hamiltonian vector field, which generates a one-parameter group (or at least a flow) of canonical transformations, and every such flow arises from an observable in this way (at least if $M$ is simply connected). Moreover, the (smooth) observables form a Lie algebra under the Poisson bracket. If $f$ is an observable whose Hamiltonian vector field generates the flow $\Phi_{t}$, the evolution of any other observable $g$ under the flow is given by $(d / d t)\left(g \circ \Phi_{t}\right)=\{g, f\} \circ \Phi_{t}$.
  \item In quantum mechanics, the state space is a projective Hilbert space $\mathbb{P} \mathcal{H}$, the observables are self-adjoint operators on $\mathcal{H}$, and the symmetries are unitary or anti-unitary operators on $\mathcal{H}$. A choice of a nonzero real constant $b$ (to be specified shortly) sets up a bijection between self-adjoint and skew-adjoint operators, $A \leftrightarrow i b A$. Via this correspondence, observables are precisely the infinitesimal generators of one-parameter groups of symmetries: $U(t)=e^{i b t A}$. Moreover, self-adjoint operators form a Lie algebra under the Lie product $\left(A_{1}, A_{2}\right) \mapsto i b\left[A_{1}, A_{2}\right]$ (except for technical difficulties concerning the domains of unbounded operators). If $A$ is an observable, the evolution of any other observable $B$ under the one-parameter group $e^{i b t A}$ is given by $(d / d t)\left(e^{-i b t A} B e^{i b t A}\right)=e^{-i b t A}(i b[B, A]) e^{i b t A}$.\\
These analogies are a little ragged around the edges because of the differentiability assumptions, the local nature of most flows on the classical side, and the hangups with unbounded operators on the quantum side, but they are sufficiently striking to assure us that we are onto something good.
\end{itemize}

\subsection*{Quantization}
So far we have only constructed a skeleton of a quantum theory. To put some flesh on the bones, we need a way to find the quantum analogues of familiar classical observables such as position, momentum, and energy. (Once we understand the observables, the nature of the states will become clearer.) We shall take the Hamiltonian formulation of classical mechanics as a starting point so as to take advantage of the analogies outlined above, and the resulting quantum theory will therefore be nonrelativistic. Building a relativistic theory will require more work.

What we want is a way to assign to each classical observable (a real-valued function $f$ on phase space) a quantum observable $A_{f}$ (a self-adjoint operator on Hilbert space). We (optimistically) ask for this correspondence to enjoy the following properties:

\begin{itemize}
  \item Constant functions correspond to constant multiples of the identity:
\end{itemize}

$$
A_{\alpha}=\alpha I \quad(\alpha \in \mathbb{R})
$$

\begin{itemize}
  \item If $\phi: \mathbb{R} \rightarrow \mathbb{R}$ is a Borel function, and $E \subset \mathbb{R}$ is a Borel set, we have $(\phi \circ f)(x) \in E$ if and only if $f(x) \in \phi^{-1}(E)$. Let $P_{f}$ be the projectionvalued measure corresponding to the operator $A_{f}$. Recalling the way in which projection-valued measures are to be interpreted as observables, we infer that the projection-valued measure corresponding to the classical observable $\phi \circ f$ should be $E \mapsto P_{f}\left(\phi^{-1}(E)\right)$. But the self-adjoint operator corresponding to this measure is just $\phi\left(A_{f}\right)$, defined by the spectral functional calculus. In short,
\end{itemize}

$$
A_{\phi \circ f}=\phi\left(A_{f}\right)
$$

\begin{itemize}
  \item If $f$ and $g$ are classical observables, the expected value of the quantum observable $A_{f+g}$ in any state $u$ should be the sum of the expected values of $A_{f}$ and $A_{g}$, that is, $\langle u| A_{f+g}|u\rangle=\langle u| A_{f}|u\rangle+\langle u| A_{g}|u\rangle$. But the diagonal matrix elements of an operator completely determine the operator by polarization, so
\end{itemize}

$$
A_{f+g}=A_{f}+A_{g}
$$

\begin{itemize}
  \item Poisson brackets should correspond to commutators, as modified by the correspondence between self-adjoint and skew-adjoint operators:
\end{itemize}

$$
A_{\{f, g\}}=i b\left[A_{f}, A_{g}\right]
$$

(We ignore questions about domains here.)\\[0pt]
It turns out that this is too much to ask for; quantum mechanics isn't that similar to classical mechanics. (See Folland [44] for further discussion of this point.) But we are only looking for guidelines, not a functor that will magically transform classical systems into quantum ones, and we shall use this wish-list in that spirit. In particular, we expect only that (3.8) will hold in an approximate or asymptotic sense in situations where the classical and quantum pictures can be compared directly.

It is time to explore the meaning of the mysterious constant $b$. The reason for not simply taking $b=1$ is that its value depends on the units of measurement being used. Differentiation with respect to a position or momentum variable changes the dimensions by adding a factor of $\left[l^{-1}\right]$ or $\left[m^{-1} l^{-1} t\right]$, respectively. Thus, if $f$ and $g$ have dimensions $\left[d_{f}\right]$ and $\left[d_{g}\right]$, their Poisson bracket will have dimensions $\left[d_{f} d_{g} m^{-1} l^{-2} t\right]$. On the quantum side, the "dimensions" of a self-adjoint operator\\
$A$ are the dimensions of the variable $\lambda$ in its spectral resolution $A=\int \lambda d P(\lambda)$ (or equivalently, if $A$ has discrete spectrum, the dimensions of its eigenvalues). (This point may seem less than perfectly clear in the present abstract setting, but it is always obvious in the concrete models actually used in quantum theory, where the dimensions of an operator are manifest from its form and the coefficients occuring in it.) If $A$ and $B$ are self-adjoint operators with dimensions $\left[d_{A}\right]$ and $\left[d_{B}\right]$, the operator $i[A, B]=i(A B-B A)$ will have dimensions $\left[d_{A} d_{B}\right]$. But this means that if the correspondence (3.8) is to be valid, even in an approximate sense, the constant $b$ must have dimensions $\left[m^{-1} l^{-2} t\right]$. Hence, its value depends on the system of units being used, and its value in a given system must be determined empirically. Moreover, in order to make the correspondences between symmetries and observables in classical and quantum mechanics parallel to one another, it turns out to be necessary to take $b$ to be negative. (Alternatively, one can take $b$ positive but replace the $i$ in the correspondence $A \leftrightarrow i b A$ with $1 / i$.)

We are now ready to reveal the secret: the constant $-1 / b$ is called Planck's constant and is denoted by $\hbar$; in MKS units it is approximately $1.054 \times 10^{-34} \mathrm{kg} \cdot \mathrm{m}^{2} / \mathrm{s}$. (Planck's original constant was $h=2 \pi \hbar$, but now everyone uses $\hbar$. Basically, it is a question of whether one wants to measure frequency in cycles per second or radians per second.) The dimensions [ $m l^{2} t^{-1}$ ] of $\hbar$ can be interpreted as $\left[m l^{2} t^{-2} \times t\right]$ (energy times time) or as $\left[l \times m l t^{-1}\right]$ (moment arm times momentum), and they are therefore the same as the dimensions of action, the basic quantity in Lagrangian mechanics, and angular momentum. As we shall see, Planck's constant functions as a basic unit of both action and angular momentum in quantum theory.

The correspondences between observables and one-parameter symmetry groups and between Poisson brackets and commutators can now be rewritten as

$$
A \leftrightarrow e^{t A / i \hbar}, \quad A_{\{f, g\}} \leftrightarrow \frac{1}{i \hbar}\left[A_{f}, A_{g}\right] .
$$

While we are thinking of dimensions, let us point out that the dimensions of the $t$ in the one-parameter group must be related to those of $A$ so that $t A / \hbar$ is dimensionless (independent of the system of units). In particular, if $e^{-i t A / \hbar}$ is the group that describes the time evolution of the system, so that $t$ is time, then $A$ must have the dimensions of energy, in accordance with the analogy with Hamiltonian mechanics.

We now work toward the quantum description of a system of moving particles with classical position and momentum coordinates $x_{1}, \ldots, x_{n}$ and $p_{1}, \ldots, p_{n}$. (That is, the phase space is $T^{*} \mathbb{R}^{n} \cong \mathbb{R}^{2 n}$, where $n=3 k$ if there are $k$ particles moving in $\mathbb{R}^{3}$.) As a starting point, we try to quantize these basic observables in such a way that the correspondence between Poisson brackets and commutators in (3.9) is an exact equality. Since

$$
\left\{x_{j}, p_{k}\right\}=\delta_{j k}, \quad\left\{x_{j}, x_{k}\right\}=\left\{p_{j}, p_{k}\right\}=0
$$

we are looking for self-adjoint operators $X_{1}, \ldots, X_{n}$ and $P_{1}, \ldots, P_{n}$ on a Hilbert space $\mathcal{H}$ that satisfy

$$
\left[X_{j}, P_{k}\right]=i \hbar \delta_{j k} I, \quad\left[X_{j}, X_{k}\right]=\left[P_{j}, P_{k}\right]=0
$$

These are called the canonical commutation relations, and $P_{j}$ is said to be canonically conjugate to $X_{j}$ when they are satisfied.

One must immediately face the fact that the operators $X_{j}$ and $P_{j}$ in (3.10) are necessarily unbounded - not just because they are supposed to correspond to unbounded classical observables but for purely mathematical reasons. (Suppose $X$\\
and $P$ are bounded self-adjoint operators satisfying $[X, P]=i \hbar I$. An easy induction shows that $\left[X, P^{n}\right]=n i \hbar P^{n-1}$, so $n \hbar\left\|P^{n-1}\right\| \leq\left\|X P^{n}-P^{n} X\right\| \leq 2\|X\|\left\|P^{n}\right\|$. The self-adjointness of $P$ implies that $\left\|P^{k}\right\|=\|P\|^{k}$ for any $k$, so $n \hbar \leq 2\|X\|\|P\|$ for all $n$, which is absurd.) Consequently, the interpretation of the relation $\left[X_{j}, P_{k}\right]= i \hbar \delta_{j k} I$ is problematic: the operator on the right is defined on all of $\mathcal{H}$, but the operator on the left is not. To take the extreme case, there exist self-adjoint operators $X$ and $P$ whose domains contain only the zero vector in common, and then $[X, P]=i \hbar I$ is valid on $\mathcal{D}([X, P])$ for the stupid reason that $\mathcal{D}([X, P])=\{0\}$. (Less stupid pathologies can occur too, as we shall point out later.)

The best resolution of this problem is as follows. The commutation relations (3.10) imply that the operators $X_{j}, P_{j}$, and $i I$ span a Lie algebra. Let us consider this algebra in the abstract; that is, let $\mathfrak{h}_{n}$ be $\mathbb{R}^{n} \times \mathbb{R}^{n} \times \mathbb{R}$ equipped with the Lie bracket

$$
\left[(v, w, t),\left(v^{\prime}, w^{\prime}, t^{\prime}\right)\right]=\left(0,0, v \cdot w^{\prime}-w \cdot v^{\prime}\right)
$$

( $\mathfrak{h}_{n}$ is called the Heisenberg algebra of dimension $2 n+1$.) Then the standard basis vectors

$$
V_{j}=\left(e_{j}, 0,0\right), \quad W_{j}=\left(0, e_{j}, 0\right), \quad T=(0,0,1),
$$

where $\left\{e_{1}, \ldots, e_{n}\right\}$ is the standard basis for $\mathbb{R}^{n}$, satisfy

$$
\left[V_{j}, W_{k}\right]=\delta_{j k} T, \quad\left[V_{j}, V_{k}\right]=\left[W_{j}, W_{k}\right]=\left[V_{j}, T\right]=\left[W_{j}, T\right]=0
$$

so the relations (3.10) mean that the map $(v, w, t) \mapsto \sum\left(v_{j} X_{j}+w_{j} P_{j}\right)+i \hbar t I$ is a Lie algebra homomorphism. The "good" interpretation of (3.10) is that this homomorphism should come from a unitary representation of the corresponding Lie group.

That is, let $H_{n}$ (the Heisenberg group) be the simply connected Lie group with Lie algebra $\mathfrak{h}_{n}$. $H_{n}$ can be identified with $\mathbb{R}^{n} \times \mathbb{R}^{n} \times \mathbb{R}$ with the group law

$$
(v, w, t)\left(v^{\prime}, w^{\prime}, t^{\prime}\right)=\left(v+v^{\prime}, w+w^{\prime}, t+t^{\prime}+\frac{1}{2}\left(v \cdot w^{\prime}-w \cdot v^{\prime}\right)\right)
$$

(see Folland [44]). If $\rho$ is a unitary representation of $H_{n}$ on $\mathcal{H}$ and $\Xi \in \mathfrak{h}_{n}$, let $d \rho(\Xi)$ be the infinitesimal generator of the one-parameter unitary group $t \mapsto \rho(\exp (t \Xi))$ (a skew-adjoint operator on $\mathcal{H}$ ). Then the domains of the operators $d \rho(\Xi), \Xi \in \mathfrak{h}_{n}$, all contain the space $C^{\infty}(\rho)$ of $C^{\infty}$ vectors for $\rho$ and map $C^{\infty}(\rho)$ into itself; they are all essentially skew-adjoint on $C^{\infty}(\rho)$, and the relation

$$
\left[d \rho\left(\Xi_{1}\right), d \rho\left(\Xi_{2}\right)\right]=d \rho\left(\left[\Xi_{1}, \Xi_{2}\right]\right)
$$

holds when both sides are interpreted as operators on $C^{\infty}(\rho)$. Let $X_{j}=i \hbar d \rho\left(V_{j}\right)$ and $P_{j}=i \hbar d \rho\left(W_{j}\right)$; then the relations (3.10) hold on $C^{\infty}(\rho)$ provided that $d \rho(T)= I / i \hbar$, which means that $\rho(0,0, t)=e^{t / i \hbar} I$.

Thus, we are looking for unitary representations $\rho$ of $H_{n}$ such that $\rho(0,0, t)= e^{t / i \hbar} I$. One such representation is the Schrödinger representation $\sigma$ on $L^{2}\left(\mathbb{R}^{n}\right)$, defined by

$$
[\sigma(v, w, t) f](x)=\exp \left(\frac{v \cdot x+t-\frac{1}{2} v \cdot w}{i \hbar}\right) f(x-w)
$$

which yields

$$
\left[X_{j} f\right](x)=i \hbar d \sigma\left(V_{j}\right) f(x)=x_{j} f(x), \quad P_{j} f=i \hbar d \sigma\left(W_{j}\right) f=\frac{\hbar}{i} \frac{\partial f}{\partial x_{j}}
$$

(The domains of $X_{j}$ and $P_{j}$ are simply the set of all $f \in L^{2}$ such that $x_{j} f$ and the distribution derivative $\partial_{j} f$ are in $L^{2}$, respectively; the space $C^{\infty}(\sigma)$ is the Schwartz class $\mathcal{S}\left(\mathbb{R}^{n}\right)$.)

The representation $\sigma$ of $H_{n}$ defined by (3.11) is easily seen to be irreducible. Indeed, suppose $\mathcal{M}$ is a nonzero closed invariant subspace of $L^{2}$, and pick $f \neq 0 \in \mathcal{M}$. If $g \in \mathcal{M}^{\perp}$, then $g \perp \sigma(v, w, 0) f$ for all $v, w \in \mathbb{R}^{n}$, that is,

$$
0=\int e^{(v \cdot x-(v \cdot w) / 2) / i \hbar} f(x-w) g(x)^{*} d x=\int e^{v \cdot x / i \hbar} f\left(x-\frac{1}{2} w\right) g\left(x+\frac{1}{2} w\right)^{*} d x
$$

By Fourier uniqueness, $f\left(x-\frac{1}{2} w\right) g\left(x+\frac{1}{2} w\right)^{*}=0$ for a.e. $x$ and $w$. Since the linear map $(x, w) \mapsto\left(x-\frac{1}{2} w, x+\frac{1}{2} w\right)$ is invertible, we have $f(x) g(y)^{*}=0$ for a.e. $x$ and $y$, and so either $f=0$ or $g=0$ as elements of $L^{2}$. But $f \neq 0$ by assumption; hence $\mathcal{M}^{\perp}=\{0\}$.

Now, what about other representations $\rho$ of $H_{n}$ such that $\rho(0,0, t)=e^{t / i \hbar} I$ ? The answer is provided by the Stone-von Neumann theorem:

Suppose $\rho$ is a unitary representation of $H_{n}$ on a Hilbert space $\mathcal{H}$ such that $\rho(0,0, t)=e^{t / i \hbar} I$. Then $\mathcal{H}$ is a direct sum of mutually orthogonal subspaces $\mathcal{H}_{\alpha}$ that are invariant under $\rho$, such that the subrepresentation of $\rho$ on $\mathcal{H}_{\alpha}$ is unitarily equivalent to $\sigma$ for each $\alpha$. In particular, if $\rho$ is irreducible, then $\rho$ is unitarily equivalent to $\sigma$.

There are two good ways to prove this theorem: by von Neumann's original argument and as a corollary of the Mackey imprimitivity theorem. See Folland [44], [45].

A few mathematical remarks are in order at this point.\\
i. Given a unitary representation $\rho$ of $H_{n}$, let $\rho_{1}(v)=\rho(v, 0,0)$ and $\rho_{2}(w)= \rho(0, w, 0)$. Then $\rho_{1}$ and $\rho_{2}$ are unitary representations of $\mathbb{R}^{n}$. If $\rho$ satisfies $\rho(0,0, t)=e^{t / i \hbar} I$, then $\rho_{1}$ and $\rho_{2}$ satisfy the integrated form of the canonical commutation relations,

$$
\rho_{2}(w) \rho_{1}(v)=e^{-v \cdot w / i \hbar} \rho_{1}(v) \rho_{2}(w) .
$$

Conversely if $\rho_{1}$ and $\rho_{2}$ are unitary representations of $\mathbb{R}^{n}$ that satisfy (3.13), then $\rho(v, w, t)=e^{(i t-(v \cdot w / 2)) / i \hbar} \rho_{1}(v) \rho_{2}(w)$ is a representation of $H_{n}$ that satisfies $\rho(0,0, t)=e^{t / i \hbar} I$. The Stone-von Neumann theorem is often stated in terms of $\rho_{1}$ and $\rho_{2}$ rather than $\rho$.\\
ii. The map $(v, w) \mapsto \rho(v, w, 0)$ is a projective representation of $\mathbb{R}^{2 n}$ that does not arise from a unitary representation of $\mathbb{R}^{2 n}$.\\
iii. To appreciate the delicacy of passing from the commutation relations (3.10) to a representation of $H_{n}$, consider the following example. Let $\mathcal{H}=L^{2}([0,1])$, $[X f](x)=x f(x)$, and $P f=(\hbar / i) f^{\prime}$ on the domain $\mathcal{D}(P)=\left\{f \in A C: f^{\prime} \in\right. L^{2}$ and $\left.f(0)=f(1)\right\}$ ( $A C$ is the space of absolutely continuous functions on $[0,1]$.) Then $X$ and $P$ are both self-adjoint operators on $\mathcal{H}$, and $[X, P]=i \hbar I$ on the domain of $[X, P]$, which is $\left\{f \in A C: f^{\prime} \in L^{2}\right.$ and $\left.f(0)=f(1)=0\right\}$. The representation $(u, v, t) \mapsto u X+v P+(t / i \hbar) I$ of the Heisenberg algebra $\mathfrak{h}_{1}$ does not come from a unitary representation of $H_{1}$. Indeed, the reader may verify that the one-parameter groups generated by $X / i \hbar$ and $P / i \hbar$ (namely, $\left[e^{s X / i \hbar} f\right](x)=e^{s x / i \hbar} f(x)$ and $\left[e^{s P / i \hbar} f\right](x)=f(x-s)$, where $x-s$ is taken modulo 1 ) do not satisfy (3.13).\\
iv. The Schrödinger representation $\sigma$ depends on the parameter $\hbar$. That is, for each nonzero $\hbar \in \mathbb{R}$ there is a Schrödinger representation $\sigma_{\hbar}$, and these representations are all unitarily inequivalent. (Indeed, they are already inequivalent on the center of $H_{n}$, namely, $Z=\{(0,0, t): t \in \mathbb{R}\}$.) On the other hand, if $\rho$ is any irreducible unitary representation of $H_{n}$, Schur's lemma implies that $\rho(0,0, t)=e^{i a t} I$ for some $a \in \mathbb{R}$. If $a \neq 0$, then $\rho$ is unitarily equivalent to $\sigma_{-1 / a}$ by the Stone-von Neumann theorem. If $a=0, \rho$ factors through the quotient group $H_{n} / Z \cong \mathbb{R}^{2 n}$, and hence $\rho$ acts on a one-dimensional space by $\rho(v, w, t)=e^{i \alpha \cdot v+i \beta \cdot w}$ for some $\alpha, \beta \in \mathbb{R}^{n}$. Thus the Stone-von Neumann theorem yields a complete classification of the irreducible unitary representations of $H_{n}$.





Returning to physics, the upshot of all this is that we have a good candidate for the quantum description of a (nonrelativistic) particle moving in $\mathbb{R}^{n}$ : the quantum state space is $L^{2}\left(\mathbb{R}^{n}\right)$, and the position and momentum observables are given by (3.12). It is clear that these operators have the right dimensions: $X_{j}=$ multiplication by $x_{j}$, which carries dimensions of length; and $P_{j}=(\hbar / i) \partial / \partial x_{j}$, whose dimensions are $\left[m l^{2} t^{-1}\right]$ (from the $\hbar$ ) times $\left[l^{-1}\right]$ (from the $\partial / \partial x_{j}$ ), or $\left[m l t^{-1}\right]$, which is right for momentum.

The irreducibility of the Schrödinger representation is a desirable feature. By the Stone-von Neumann theorem, any other representation of the canonical commutation relations (suitably interpreted via the Heisenberg group) is equivalent to a direct sum of copies of the Schrödinger representation, and taking two or more copies would be like taking the union of two or more copies of of $\mathbb{R}^{2 n}$ for the classical phase space.

The projection-valued measures $\mathcal{P}_{j}$ (so denoted in order to distinguish them from the momentum operators $P_{j}$ ) corresponding to the position observables $X_{j} =$ multiplication by $x_{j}$ are simply given by $\mathcal{P}_{j}(E)=$ multiplication by the characteristic function of $\left\{x: x_{j} \in E\right\}$. In fact, these operators have a common spectral resolution, namely, the projection-valued measure $\mathcal{P}$ on $\mathbb{R}^{n}$ defined by $\mathcal{P}(E)=$ multiplication by $\chi_{E}$ for $E$ a Borel set in $\mathbb{R}^{n}$; the measures $\mathcal{P}_{j}$ on $\mathbb{R}$ are the pushforwards of $\mathcal{P}$ by the coordinate functions. This means that if $f$ is a unit vector in $L^{2}\left(\mathbb{R}^{n}\right)$ and $E \subset \mathbb{R}^{n}$ is a Borel set, $\langle f| \mathcal{P}(E)|f\rangle=\int_{E}|f|^{2}$ is the probability that the particle will be located in $E$ when the system is in the state $f$. Thus, in this model the normalized state vectors have the physical significance that the probability density of the position of the particle in $\mathbb{R}^{n}$ in the state $f$ is $|f|^{2}$. (Incidentally, if $f \in L^{2}\left(\mathbb{R}^{n}\right)$ is a unit vector representing a state, the function $f(x)$ carries dimensions $\left[l^{-n / 2}\right]$, so that $|f(x)|^{2}$ has the dimensions $\left[l^{-n}\right]$ appropriate to a probability per unit volume.)

A similar picture for the momentum variables is obtained by applying the Fourier transform. For this purpose we employ the following version of the Fourier transform that incorporates Planck's constant: if $f \in L^{2}\left(\mathbb{R}^{n}\right)$,

$$
\widehat{f}(p)=\int e^{-i p \cdot x / \hbar} f(x) d^{n} x
$$

The inversion and Parseval formulas are

$$
f(x)=\int e^{i p \cdot x / \hbar} \widehat{f}(p) \frac{d^{n} p}{(2 \pi \hbar)^{n}}, \quad \int|f(x)|^{2} d^{n} x=\int|\widehat{f}(p)|^{2} \frac{d^{n} p}{(2 \pi \hbar)^{n}}
$$

By integration by parts,

$$
\left(P_{j} f\right)^{\wedge}(p)=\frac{\hbar}{i} \int e^{-i p \cdot x / \hbar} \frac{\partial f}{\partial x_{j}}(x) d x=\int p_{j} e^{-i p \cdot x / \hbar} f(x) d x=p_{j} \widehat{f}(p)
$$

Hence we arrive at the following picture: We think of the copy of $\mathbb{R}^{n}$ on which $\widehat{f}$ lives as "momentum space," just as the original $\mathbb{R}^{n}$ is "position space." (Thus the coordinates $p_{j}$ have the dimensions of momentum, so that $p \cdot x / \hbar$ is dimensionless, as it should be.) In momentum space, $P_{j}$ is just multiplication by the $j$ th coordinate function, so the same analysis as in the preceding paragraph shows that $|\widehat{f}|^{2}$ is the probability density of the momentum of the particle in the state $f$ with respect to the normalized Lebesgue measure $d^{n} p /(2 \pi \hbar)^{n}$.






It should be noted that there are no states in which the particle has a definite position or momentum. The "generalized state" (or, as the physicists say, "nonnormalizable state") in which the particle is definitely at the position $x_{0} \in \mathbb{R}^{n}$ is represented by the delta-function $\delta\left(x-x_{0}\right)$, and the "generalized state" in which the particle definitely has momentum $p_{0} \in \mathbb{R}^{n}$ is represented by the inverse Fourier transform of the normalized delta-function $(2 \pi)^{n} \delta\left(p-p_{0}\right)$, i.e., the plane wave $e^{i p_{0} \cdot x / \hbar}$. It is frequently convenient to allow such idealized states into one's universe of discourse. This can, of course, be done with complete rigor by invoking the theory of distributions (or, on a more abstract level, rigged Hilbert spaces; see Gelfand and Vilenkin [52]), but we shall keep the discussion on an informal level as the physicists do and take the opportunity to introduce and explain some notation and formulas that may seem peculiar to mathematicians but are common in the physics literature.

The generalized state $\delta\left(x-x_{0}\right)$ where the particle is located at $x_{0}$ is the simultaneous eigenvector of the position operators $X_{j}$ with eigenvalues $x_{0 j}$, so it can conveniently be labeled by these eigenvalues and denoted $\left|x_{0}\right\rangle$. The corresponding bra $\left\langle x_{0}\right|$ is then nothing but the evaluation functional at $x_{0}:\left\langle x_{0} \mid f\right\rangle=f\left(x_{0}\right)$. The states $|x\rangle, x \in \mathbb{R}^{n}$, form a "generalized eigenbasis" for the position operators: just as $f=\sum\left\langle e_{k} \mid f\right\rangle e_{k}$ when we have a genuine orthonormal basis $\left\{e_{k}\right\}$, here we can write

$$
|f\rangle=\int|x\rangle\langle x \mid f\rangle d^{n} x, \quad \text { or } \quad \int|x\rangle\langle x| d^{n} x=I
$$

which is just a restatement of the fact that $\int f(x) \delta(y-x) d x=f(y)$. Likewise, the generalized state where the particle has momentum $p$ may be denoted by $|p\rangle$. (Here, of course, it must be understood that, contrary to common mathematical practice, the letters $x$ and $p$ are not just arbitrarily chosen symbols but carry some semantic content: $x$ means position and $p$ means momentum.) Here the scalar product $\langle p \mid f\rangle$ is $\widehat{f}(p)$, and the fact that the states $|p\rangle$ are again a "generalized orthonormal basis," that is,

$$
|f\rangle=\int|p\rangle\langle p \mid f\rangle \frac{d^{n} p}{(2 \pi \hbar)^{n}}, \quad \text { or } \quad \int|p\rangle\langle p| \frac{d^{n} p}{(2 \pi \hbar)^{n}}=I
$$

is a restatement of the Fourier inversion formula (3.15); cf. (1.5).\\
This is a good place to bring up an intriguing difference in the way physicists and mathematicians regard mathematical objects. Mathematicians are trained to think that all mathematical objects should be realized in some specific way as sets. When we do calculus on the real line we may not wish to think of real numbers as Dedekind cuts or equivalence classes of Cauchy sequences of rationals, but it\\
reassures us to know that we can do so. Physicists' thought processes, on the other hand, are anchored in the physical world rather than in set theory, and they see no need to tie themselves down to specific set-theoretic models. A physicist may speak of the state space $\mathcal{H}$ of a quantum system $Q$; if a mathematician asks "What Hilbert space is $\mathcal{H}$ ?", the response is "I just told you, it's the state space of $Q$." If one is studying a particular observable, one may wish to represent $\mathcal{H}$ in a way that makes the analysis more transparent. For example, if one is studying an operator with discrete spectrum - eigenvalues $\lambda_{n}$ and eigenvectors $\left|\lambda_{n}\right\rangle$ - one may wish to represent $\mathcal{H}$ as $l^{2}$ by identifying the state $|f\rangle$ with the sequence $\left\{\left\langle\lambda_{n} \mid f\right\rangle\right\}$; if one is studying the position of a particle, one will represent $\mathcal{H}$ as $L^{2}\left(\mathbb{R}^{n}\right)$ in a particular way by using the generalized states $|x\rangle$ as a "basis" to form the wave function $f(x)=\langle x \mid f\rangle$; if one is studying momentum, one will represent $\mathcal{H}$ as $L^{2}\left(\mathbb{R}^{n}\right)$ in a different way by using the generalized states $|p\rangle$ instead; and so forth. But, from this point of view, choosing a particular representation is akin to choosing a particular coordinate system for doing analysis on a manifold, and it is not unreasonable to regard committing oneself to such a choice at the outset as a mistake.\\
(Similarly, it is clear from the way physicists talk about Lie algebras that they regard a Lie algebra as an algebraic structure unattached to any particular set, rather than as a set endowed with an algebraic structure. If you ask "What is $\mathfrak{s o}(3)$ ?", a mathematician will answer "the set of skew-symmetric $3 \times 3$ real matrices, with the commutator as Lie bracket," but a physicist will answer something like "three generators $X, Y$, and $Z$ satisfying $[X, Y]=Z,[Y, Z]=X$, and $[Z, X]=Y$," and he will say that $X, Y$, and $Z$ "satisfy the Lie algebra $\mathfrak{s o}(3)$.")

Nonetheless, as this book is being written by a mathematician for mathematicians, we shall continue to take state spaces for particular quantum systems to be particular concrete Hilbert spaces.






It remains to quantize more general functions of position and momentum. For functions of position alone or momentum alone, there is an obvious way to do this: If $\phi$ is a function on $\mathbb{R}^{n}$, the quantization of the classical observable $\phi(x)$ is multiplication by $\phi$ in position space, and the quantization of the classical observable $\phi(p)$ is multiplication by $\phi$ in momentum space:

$$
[\phi(X) f](x)=\phi(x) f(x), \quad[\phi(P) f]^{\curlyvee}(p)=\phi(p) \widehat{f}(p)
$$

However, one runs into difficulties in going further. For functions of the form $\phi(x)+\psi(p)$, the symmetric operator $\phi(X)+\psi(P)$ may not be self-adjoint on the domain $\mathcal{D}(\phi(X)) \cap \mathcal{D}(\psi(P))$, and when one considers functions involving products of position and momentum variables it is impossible to satisfy (3.6) or (3.8) even on the formal level: see Folland [44], p. 17 and pp. 197-199. These no-go theorems simply reflect that fact that quantum mechanics is too different from classical mechanics for a quantization procedure to exist that has all the properties in the wish list (3.5)-(3.8). After all, classical mechanics is a limiting case of quantum mechanics, not the other way around, and there is no reason to expect that one can fully recover the quantum theory from the limiting case. From the physicists' point of view, quantization procedures are to be taken only as possible ways of arriving at a good quantum theory. Once one has constructed such a theory, the test of its correctness comes from working out its consequences and comparing them with experiment, not from any a priori considerations.

There is, however, some purely mathematical interest in quantization procedures. One reason is the desire for a useful functional calculus for the noncommuting operators $X_{j}$ and $P_{j}$, motivated by problems in the theory of partial differential equations. Indeed, the calculus of pseudodifferential operators, which sets up a correspondence between "symbols" $f(x, p)$ and operators $f(X, P)$ (usually with $\hbar=1$ or $\hbar=1 / 2 \pi$ ), is such a procedure. We refer to Folland [44] for an extensive discussion of it from a point of view similar to the one here.

The other mathematical motivation for studying quantization procedures comes from group representation theory. Suppose one wishes to find all the irreducible unitary representations of a connected Lie group $G$. In physical language, this can be considered as finding all quantum systems on which $G$ acts irreducibly as a group of symmetries. The corresponding "classical" problem is to find all symplectic manifolds on which $G$ acts transitively as a group of canonical transformations, i.e., all symplectic homogeneous $G$-spaces. The latter problem is actually quite explicitly solvable: The orbits of the co-adjoint action of $G$ on $\mathfrak{g}^{*}$ are symplectic homogeneous $G$-spaces, as are the orbits of the coadjoint action of central extensions of $G$ by $\mathbb{R}$ (essentially the same family of groups that intervenes in the issue of projective vs. unitary representations). Moreover, all symplectic homogeneous $G$ spaces are more or less of this form (the phrase "more or less" has to do with questions about covering spaces). Hence, if one has a good quantization procedure for such spaces, one has a hope of understanding the unitary representations of $G$. This idea works beautifully when $G$ is nilpotent or when $G$ is compact; the Kirillov theory for nilpotent groups and the Borel-Weil theory for compact groups can both be understood as special cases of quantization of coadjoint orbits. It also provides much illumination when $G$ is solvable or semisimple, although these cases present additional technical problems. (See Kirillov [75] together with Vogan's review [128], and also Vogan [127].)

Back to physics: The time evolution of a quantum system is given by a oneparameter group $U(t)$ of unitary maps, and the observable corresponding to its infinitesimal generator is the quantum Hamiltonian $H$; that is, $U(t)=\exp (t H / i \hbar)$. We expect this operator to be a quantized version of the Hamiltonian in classical mechanics. Here we consider the case of particles moving in a potential $V$, for which the classical Hamiltonian is $\frac{1}{2} m^{-1} p \cdot p+V(x)$. (Recall that $m$ is a diagonal matrix whose diagonal entries are the masses of the particles.) The obvious quantum analogue is the differential operator on $L^{2}\left(\mathbb{R}^{n}\right)$

$$
H=-\frac{1}{2} \hbar^{2} m^{-1} \nabla \cdot \nabla+V
$$

the Schrödinger operator with potential $V$, where $V$ is considered as the operator $f \mapsto V f$. The relation $U(t)=e^{t H / i \hbar}$ can be expressed as a differential equation for the evolution of a state vector $f$ :

$$
i \hbar \frac{\partial f}{\partial t}=-\frac{\hbar^{2}}{2} m^{-1} \nabla \cdot \nabla f+V f \quad(f(\cdot, t)=U(t) f)
$$

This is the classic Schrödinger equation with potential $V$.\\
The missing ingredient here is a specification of $H$ as a self-adjoint operator on $L^{2}$. (In the study of particles confined to a bounded region, this will involve imposition of boundary conditions, but we consider only particles moving in $\mathbb{R}^{n}$ here.) The operator $m^{-1} \nabla \cdot \nabla$, a variant of the Laplacian, is self-adjoint on the natural domain of all $f \in L^{2}$ such that $m^{-1} \nabla \cdot \nabla \in L^{2}$ (i.e., the Sobolev space of\\
order 2), and $V$ is self-adjoint on the domain of all $f \in L^{2}$ such that $V f \in L^{2}$. Hence, $H$ is at least a symmetric operator on the intersection of these two domains, and in most important cases it turns out to be essentially self-adjoint. There is a considerable body of theorems relating to this question. Here are a couple of them, with inessential constants omitted, in which $\nabla^{2}$ denotes the Laplacian on $\mathbb{R}^{n}$ :\\
i. If $V \in L_{\text {loc }}^{2}$ and $V$ is bounded below, then $-\nabla^{2}+V$ is essentially self-adjoint on $C_{c}^{\infty}$.\\
ii. Suppose $V$ is the sum of an $L^{p}$ function and an $L^{\infty}$ function, where $p=2$ if $n \leq 3, p>2$ if $n=4$, and $p=\frac{1}{2} n$ if $n \geq 5$. Then $-\nabla^{2}+V$ is essentially self-adjoint on $C_{c}^{\infty}$.\\
iii. (Kato) Suppose $V=\sum_{j=1}^{k} V_{j} \circ \pi_{j}$, where each $\pi_{j}$ is a partial isometry of $\mathbb{R}^{n}$ onto $\mathbb{R}^{3}$ and each $V_{j}$ is the sum of an $L^{2}$ function and an $L^{\infty}$ function on $\mathbb{R}^{3}$. Then $-\nabla^{2}+V$ is essentially self-adjoint on $C_{c}^{\infty}$.\\
In (iii), a "partial isometry" is a linear map that is an isometry on the orthogonal complement of its nullspace. In (ii) and (iii), $-\nabla^{2}+V$ is actually self-adjoint on $\mathcal{D}\left(\nabla^{2}\right)$, the Sobolev space of order 2. In (ii), the essential point is that by the Sobolev imbedding theorems, $\mathcal{D}\left(\nabla^{2}\right)$ is contained in $C_{0}$ (continuous functions vanishing at infinity) when $n \leq 3$, in $\bigcap_{2 \leq q<\infty} L^{q}$ when $n=4$, and in $L^{2 n /(n-4)}$ when $n \geq 5$, so that the product of a function in $\mathcal{D}\left(\nabla^{2}\right)$ and an $L^{p}$ function (with $p$ as in (ii)) is in $L^{2}$. (iii) is a consequence of (ii) and a little Fourier analysis. Proofs, as well as other related theorems, can be found in Reed and Simon [95], Sections X.2-X.4.

The result (iii) applies, in particular, to the Hamiltonian for an atom with atomic number $Z$ when the nucleus is considered as fixed at the origin and the moving particles are taken to be the electrons:

$$
H=-\frac{\hbar^{2}}{2 m} \nabla^{2}-\sum_{j=1}^{Z} \frac{Z e^{2}}{4 \pi\left|\mathbf{x}_{j}\right|}+\sum_{1 \leq j<k \leq Z} \frac{e^{2}}{4 \pi\left|\mathbf{x}_{j}-\mathbf{x}_{k}\right|} .
$$

(Here $m$ and $-e$ are the mass and charge of an electron.) More generally, it applies to any system of $N$ particles attracting or repelling each other by inverse-square-law forces; here $n=3 N$, each $\pi_{j}$ is of the form $\pi_{j}\left(\mathbf{x}_{1}, \ldots, \mathbf{x}_{N}\right)=\left(\mathbf{x}_{i_{1}}-\mathbf{x}_{i_{2}}\right) / \sqrt{2}$ for $\mathbf{x}_{1}, \ldots, \mathbf{x}_{N} \in \mathbb{R}^{3}$, and $V_{j}(\mathbf{x})=c_{j}|\mathbf{x}|^{-1}$ for $\mathbf{x} \in \mathbb{R}^{3}$. The analysis of such Hamiltonians is complicated when $N$ is large, and it is remarkable that interesting things can be deduced from them about assemblages of particles of macroscopic or even astronomical size. The most famous result along these lines is the "stability of matter" theorem, proved originally by Dyson and Lenard with subsequent improvements and extensions by various other people. Lieb [78] gives a very interesting account of this and related results.

One can also imagine physically unrealistic but conceptually reasonable situations in which essential self-adjointness fails. For example, consider a particle of mass 1 moving on a line subject to a potential $V(x)=-|x|^{a}$, with $a>1$. (Thus the particle is subject to a force $(\operatorname{sgn} x) a|x|^{a-1}$ that repels it from the origin, and the force increases in strength along with the distance from the origin.) The Hamiltonian $H=-\frac{1}{2}(d / d x)^{2}+V(x)$ is esentially self-adjoint on $C_{c}^{\infty}$ when $a \leq 2$, but not when $a>2$. (See Dunford and Schwartz [24], Corollary XIII.6.22.) When $a>2, H$ has a four-parameter family of self-adjoint extensions, which correspond to imposing "boundary conditions at $\pm \infty$." This phenomenon is actually reflected in the classical physics, where the position of the particle evolves by the differential\\
equation $x^{\prime \prime}(t)=(\operatorname{sgn} x) a|x|^{a-1}$. This differential equation can be solved more or less explicitly, and the upshot is that the position of the particle (starting at a position $x_{0}>0$ with velocity $v_{0}>0$, say) behaves roughly like $t^{2 /(2-a)}$ if $a<2$ and like $e^{t}$ if $a=2$, but it reaches infinity in a finite time if $a>2$. Classically, that's the end of the story; quantum mechanically, one can impose a boundary condition to tell the particle what to do when it gets to infinity, and the motion can then continue to further times.

At any rate, once the Hamiltonian has been properly defined as a self-adjoint operator, calculating the time-evolution group $e^{t H / i \hbar}$ amounts to solving the eigenvalue equation $H \phi=\lambda \phi, \lambda \in \mathbb{R}$, and determining how an arbitrary $f \in L^{2}$ can be expanded in terms of the solutions. When the spectrum of $H$ is purely discrete, this just means finding an orthonormal basis of eigenvectors, but when $H$ has continuous spectrum, the expansion of $f$ will involve an integral as well as (perhaps) a sum. In §3.4 and § 3.6 we shall work out two examples: the harmonic oscillator and the Coulomb potential.

\subsection*{Uncertainty inequalities}
Again we consider the Schrödinger model for position and momentum with the state space $L^{2}\left(\mathbb{R}^{n}\right)$. Either the position or the momentum of a particle may be localized to as small a region as we please (i.e., for any nonempty open set $U \subset \mathbb{R}^{n}$ there are nonzero $L^{2}$ functions $f$ such that either $f$ or $\widehat{f}$ vanishes outside $U$ ), but there are limits on the extent to which position and momentum can be localized simultaneously; this is one aspect of the uncertainty principle.

If $\mu$ is a probability measure on $\mathbb{R}$, a natural measure of the "uncertainty" in $\mu$ is the standard deviation

$$
\sigma(\mu)=\left[\inf _{a \in \mathbb{R}} \int(s-a)^{2} d \mu(s)\right]^{1 / 2}
$$

When $\mu$ is the probability distribution of an observable $A$ in the state $u$ ( $A$ a self-adjoint operator, $u$ a unit vector), i.e., $\mu(E)=\langle u| P(E)|u\rangle$ where $P$ is the projection-valued measure associated to $A$, we have\\
$\sigma(\mu)=\left[\inf _{a \in \mathbb{R}} \int(s-a)^{2}\langle u| d P(s)|u\rangle\right]^{1 / 2}=\left[\inf _{a \in \mathbb{R}}\langle u|(A-a)^{2}|u\rangle\right]^{1 / 2}=\inf _{a \in \mathbb{R}}\|(A-a) u\|$.\\
(When $\sigma(\mu)<\infty$, the infimum is achieved when $a$ is the expectation value $\langle A\rangle= \langle u| A|u\rangle$.) This notion easily yields an uncertainty inequality for any two noncommuting observables (self-adjoint operators) $A$ and $B$. Indeed, if $u$ is in the domain of $[A, B]$, i.e., the set of all $u \in \mathfrak{D}(A) \cap \mathfrak{D}(B)$ such that $A u \in \mathfrak{D}(B)$ and $B u \in \mathfrak{D}(A)$, we have

$$
|\langle u|[A, B]| u\rangle|=|\langle A u \mid B u\rangle-\langle B u \mid A u\rangle|=2| \operatorname{Im}\langle A u \mid B u\rangle \mid \leq 2\|A u\|\|B u\| .
$$

Replacing $A$ by $A-a$ and $B$ by $B-b$ for arbitrary $a, b \in \mathbb{R}$ (which does not affect $[A, B]$ ), we deduce that

$$
\left.\|(A-a) u\|\|(B-b) u\| \geq \frac{1}{2}|\langle u|[A, B]| u\right\rangle \mid \text { for all } u \in \mathcal{D}([A, B]) .
$$

This result is not as strong as one might wish, because the domain of $[A, B]$ may be "too small." Indeed, if $C$ is the closure of $[A, B]$, it is in general false that $\|A u\|\|B u\| \geq \frac{1}{2}|\langle u, C u\rangle|$ for all $u \in \mathcal{D}(A) \cap \mathcal{D}(B) \cap \mathcal{D}(C)$. (Counterexample: Let $X$ and $P$ be as in the third remark following the Stone-von Neumann theorem, and let\\
$u$ be the constant function 1.) For the position and momentum operators defined by (3.12), however, we can do better.

We state the result in the one-dimensional case for simplicity. For $u \in L^{2}(\mathbb{R})$ with $\|u\|_{2}=1$, the standard deviations of $X$ and $P$ in the state $u$ are

$$
\Delta x=\left[\inf _{a \in \mathbb{R}} \int(x-a)^{2}|u(x)|^{2} d x\right]^{1 / 2}, \quad \Delta p=\left[\inf _{b \in \mathbb{R}} \int(p-b)^{2}|\widehat{u}(p)|^{2} \frac{d p}{2 \pi \hbar}\right]^{1 / 2}
$$

which make sense for all $u \in L^{2}$ with the understanding that they may equal $+\infty$, and the uncertainty inequality

$$
\Delta x \Delta p \geq \frac{1}{2} \hbar
$$

is valid with no restriction on $u$. Moreover, equality holds if and only if $u$ is a Gaussian: $u(x)=(c / \pi)^{1 / 4} e^{i b x} e^{-c(x-a)^{2}}$ for some $a, b \in \mathbb{R}$ (namely, the points where the infimum is achieved in the formulas for $\Delta x$ and $\Delta p$ ) and some $c>0$. We refer to Folland and Sitaram [50] (where $\hbar$ is taken to be $1 / 2 \pi$ ) for the proof, the generalizations to $n$ dimensions, and a large assortment of related results.

There is another important uncertainty inequality that is of a somewhat different nature: that relating to time and energy, which is usually stated in the simple form $\Delta t \Delta E \geq \frac{1}{2} \hbar$. That there should be such a relation is strongly suggested by the classical theory of relativity, in which energy is the time component of a 4 vector of which the momenta are the space components. However, unlike position and momentum, time and energy are not observables that satisfy the canonical commutation relations. Indeed, there is no observable $T$ that is canonically conjugate to the Hamiltonian (energy) $H$, for if there were, by the Stone-von Neumann theorem the spectrum of $T$ and $H$ would have to be the whole real line; but in any reasonable quantum system the energy must be bounded below.

To make sense of the time-energy uncertainty inequality, one can proceed as follows. Let $H$ be the Hamiltonian of the system under consideration, and let $A=A(t)$ be a Heisenberg-picture observable; thus, $d A / d t=(1 / i \hbar)[A, H]$. For any state $|u\rangle$ of the system, let $\langle A\rangle=\langle u| A|u\rangle$ and $\Delta A=\|(A-\langle A\rangle) u\|$ be the expectation and standard deviation of $A$ (assumed finite) in the state $|u\rangle$, and let $\Delta E$ be the standard deviation of $H$ in $|u\rangle$. By (3.18), we have

$$
(\Delta A)(\Delta E) \geq \frac{1}{2}|\langle[A, H]\rangle|=\frac{1}{2} \hbar|\langle d A / d t\rangle|=\frac{1}{2} \hbar|d\langle A\rangle / d t| .
$$

Now, the quantity

$$
\Delta t_{A}=\left|\frac{\Delta A}{d\langle A\rangle / d t}\right|
$$

represents an amount of time: it is the time required for the expectation of $A$ to change by an amount equal to the standard deviation, and hence the time required for the statistics of $A$ to change appreciably. The preceding inequality thus says that

$$
(\Delta t)(\Delta E) \geq \frac{1}{2} \hbar, \quad \text { where } \quad \Delta t=\inf _{A} \Delta t_{A} .
$$

This is the appropriate interpretation of the time-energy uncertainty inequality.

\subsection*{The harmonic oscillator}
Harmonic oscillators are interesting and important in their own right, and as we shall see, they play a fundamental role in the modeling of quantum fields. The $n$-dimensional harmonic oscillator Hamiltonian, with a single scalar mass $m$, is

$$
H=-\left(\hbar^{2} / 2 m\right) \nabla^{2}+V
$$

where $V$ is a positive definite quadratic form on $\mathbb{R}^{n}$. By a suitable rotation of coordinates, one can assume that $V$ is diagonal, in which case $H$ decouples into a sum of one-dimensional operators. It therefore suffices to consider the one-dimensional case, as the $n$-dimensional eigenfunctions are just products of one-dimensional ones. For a particle of mass $m$ in the potential $V(x)=\frac{1}{2} \kappa x^{2}$, the Hamiltonian is

$$
H=-\frac{\hbar^{2}}{2 m} \frac{d^{2}}{d x^{2}}+\frac{\kappa}{2} x^{2}
$$

We can get rid of the constants by rescaling: setting $S f(x)=f\left(\left(\hbar^{2} / m \kappa\right)^{1 / 4} x\right)$, a simple calculation shows that

$$
S H S^{-1}=\frac{\hbar}{2 \sqrt{m / \kappa}}\left(-\frac{d^{2}}{d x^{2}}+x^{2}\right)
$$

The factor of $\hbar / \sqrt{m / \kappa}$ just comes along for the ride, but it is convenient to keep the factor of $\frac{1}{2}$ explicitly. Thus, we consider the rescaled Hamiltonian

$$
H_{0}=\frac{1}{2}\left(-\frac{d^{2}}{d x^{2}}+x^{2}\right)
$$

To analyze this operator, we introduce the operators

$$
A=\frac{1}{\sqrt{2}}\left(x+\frac{d}{d x}\right), \quad A^{\dagger}=\frac{1}{\sqrt{2}}\left(x-\frac{d}{d x}\right)
$$

(There is no need to worry about domains right now; all these operators can be considered as acting on the Schwartz space $\mathcal{S}(\mathbb{R})$. Observe that the notation is formally correct, as $\langle f| A^{\dagger}|g\rangle=\langle g| A|f\rangle^{*}$ for $f, g \in \mathcal{S}$.) Simple calculations yield the following formulas:

$$
\begin{aligned}
H_{0}=A A^{\dagger}-\frac{1}{2} I & =A^{\dagger} A+\frac{1}{2} I \\
{\left[A, A^{\dagger}\right] } & =I
\end{aligned}
$$

and more generally, by (3.25) and induction on $k$,

$$
\left[A,\left(A^{\dagger}\right)^{k}\right]=k\left(A^{\dagger}\right)^{k-1}
$$

Let

$$
\phi_{0}(x)=\pi^{-1 / 4} e^{-x^{2} / 2}, \quad \phi_{k}=\frac{1}{\sqrt{k!}}\left(A^{\dagger}\right)^{k} \phi_{0} \text { for } k=1,2,3, \ldots
$$

It is easy to verify that $\phi_{k}(x)=P_{k}(x) e^{-x^{2} / 2}$ where $P_{k}$ is a polynomial of degree $k$, odd or even according as $k$ is odd or even. It is called the $k$ th normalized Hermite polynomial, and $\phi_{k}$ is the $k$ th normalized Hermite function. (The usual Hermite polynomials are $H_{k}=\pi^{1 / 4} \sqrt{2^{k} k!} P_{k}$.) We have

$$
A^{\dagger} \phi_{k}=\sqrt{k+1} \phi_{k+1}
$$

and by (3.26) and the fact that $A \phi_{0}=0$,

$$
A \phi_{k}=\frac{1}{\sqrt{k!}}\left(\left[A,\left(A^{\dagger}\right)^{k}\right]+\left(A^{\dagger}\right)^{k} A\right) \phi_{0}=\frac{k}{\sqrt{k!}}\left(A^{\dagger}\right)^{k-1} \phi_{0}=\sqrt{k} \phi_{k-1}
$$

It follows that

$$
A^{\dagger} A \phi_{k}=k \phi_{k}, \quad A A^{\dagger} \phi_{k}=(k+1) \phi_{k},
$$

and hence, by (3.24),

$$
H_{0} \phi_{k}=\left(k+\frac{1}{2}\right) \phi_{k}
$$

We claim that $\left\{\phi_{k}\right\}_{0}^{\infty}$ is an orthonormal basis for $L^{2}$. We first observe that $\left\|\phi_{0}\right\|_{2}=1$ (that's why the $\pi^{-1 / 4}$ is there) and that $A \phi_{0}=0$. For $k>0$, we have

$$
\left\langle\phi_{0} \mid \phi_{k}\right\rangle=\frac{1}{\sqrt{k!}}\left\langle\phi_{0} \mid\left(A^{\dagger}\right)^{k} \phi_{0}\right\rangle=\frac{1}{\sqrt{k!}}\left\langle A \phi_{0} \mid\left(A^{\dagger}\right)^{k-1} \phi_{0}\right\rangle=0 .
$$

Moreover, if $k \geq l>0$, by (3.30) we have

$$
\left\langle\phi_{l} \mid \phi_{k}\right\rangle=\frac{1}{\sqrt{k l}}\left\langle A^{\dagger} \phi_{l-1} \mid A^{\dagger} \phi_{k-1}\right\rangle=\frac{1}{\sqrt{k l}}\left\langle\phi_{l-1} \mid A A^{\dagger} \phi_{k-1}\right\rangle=\sqrt{\frac{k}{l}}\left\langle\phi_{l-1} \mid \phi_{k-1}\right\rangle .
$$

By induction and the preceding calculation, it follows that

$$
\left\langle\phi_{l} \mid \phi_{k}\right\rangle=\sqrt{\frac{k!}{l!}}\left\langle\phi_{0} \mid \phi_{k-l}\right\rangle=\sqrt{\frac{k!}{l!}} \delta_{k-l, 0}=\delta_{k l},
$$

and orthonormality is proved.\\
Finally, suppose $\left\langle\phi_{k} \mid f\right\rangle=0$ for all $k$. It follows that $f$ is orthogonal to $p(x) e^{-x^{2} / 2}$ for every polynomial $p$, and hence

$$
\int f(x) e^{i \xi x} e^{-x^{2} / 2} d x=\sum_{k=0}^{\infty} \int f(x) \frac{(i \xi x)^{k}}{k!} e^{-x^{2} / 2} d x=0
$$

for every $\xi$. (The interchange of summation and integration is justified since both $f$ and $e^{|\xi x|} e^{-x^{2} / 2}$ are in $L^{2}$, so their product is in $L^{1}$.) By Fourier uniqueness we conclude that $f(x) e^{-x^{2} / 2}=0$ for a.e. $x$, and hence $f=0$ in $L^{2}$.

In short, we have shown that $\left\{\phi_{k}\right\}_{0}^{\infty}$ is an orthonormal eigenbasis for $H_{0}$. If we put the constants $\hbar, m$, and $\kappa$ back in, we see that the Hamiltonian $H$ given by (3.21) has eigenfunctions $\phi_{k}\left(\left(m \kappa / \hbar^{2}\right)^{1 / 4} x\right)$, with eigenvalues $\hbar\left(k+\frac{1}{2}\right) / \sqrt{m / \kappa}$.

Let us pause to see what this means for a macroscopic harmonic oscillator such as a weight hanging from a perfectly elastic spring, where $m$ and $\kappa$ are of the order of magnitude of unity in MKS units. The classical equation of motion is $m x^{\prime \prime}=-\kappa x$, whose solutions are trig functions of $\sqrt{\kappa / m} t$. To be specific, let us assume the weight starts at rest at position $x=a$ at time $t=0$; then the classical solution is $x(t)=a \cos (\sqrt{\kappa / m} t)$, and its energy and momentum are $\frac{1}{2} \kappa a^{2}$ and $m x^{\prime}(t)=-a \sqrt{m \kappa} \sin (\sqrt{\kappa / m} t)$. Now, $\hbar \approx 10^{-34}$ in MKS units, so the quantum energy levels are extremely closely spaced, and there is no trouble finding one that is equal to $\frac{1}{2} \kappa a^{2}$ to reasonable accuracy. However, the quantum states of definite energy are stationary - if $H \phi=\lambda \phi$ then $e^{t H / i \hbar} \phi=e^{t \lambda / i \hbar} \phi$, which represents the same state as $\phi$ - and there is no sign of the classical oscillatory motion. What is going on here?

The trouble is that the classical initial value problem, in which the position and momentum at time $t=0$ are specified exactly, makes no sense quantum mechanically because of the uncertainty principle. The best one can do is to specify an initial state in which the position and momentum are as well localized as the uncertainty principle permits; this will be a superposition of energy eigenstates, and the classical motion effectively results from interference between the oscillations $e^{t \lambda / i \hbar}$ for the different eigenvalues $\lambda$. In fact, if the initial state is a suitable uncertainty-minimizing Gaussian as described following (3.19), the quantum evolution turns out to be exactly the classical periodic motion. This little miracle is sufficiently entertaining that it is worth taking some space to work out in detail. As above, we work with the rescaled Hamiltonian $H_{0}$ and its eigenfunctions.

The essential ingredient in the calculation is a version of a classical generating function identity for Hermite polynomials. We begin by observing that

$$
A^{\dagger} f(x)=\frac{1}{\sqrt{2}}\left[x f(x)-f^{\prime}(x)\right]=-\frac{1}{\sqrt{2}} e^{x^{2} / 2} \frac{d}{d x}\left[e^{-x^{2} / 2} f(x)\right]
$$

so that

$$
\left(A^{\dagger}\right)^{k} f(x)=\frac{(-1)^{k}}{2^{k / 2}} e^{x^{2} / 2} \frac{d^{k}}{d x^{k}}\left[e^{-x^{2} / 2} f(x)\right]
$$

and hence

$$
\begin{aligned}
\phi_{k}(x)=\frac{1}{\pi^{1 / 4} \sqrt{k!}}\left(A^{\dagger}\right)^{k} e^{-x^{2} / 2} & =\frac{(-1)^{k}}{\pi^{1 / 4} \sqrt{2^{k} k!}} e^{x^{2} / 2} \frac{d^{k}}{d x^{k}} e^{-x^{2}} \\
& =\left.\frac{1}{\pi^{1 / 4} \sqrt{2^{k} k!}} e^{x^{2} / 2} \frac{d^{k}}{d w^{k}} e^{-(x-w)^{2}}\right|_{w=0}
\end{aligned}
$$

Therefore, by Taylor's theorem, for any $w \in \mathbb{C}$ we have

$$
\pi^{-1 / 4} e^{-(x-w)^{2}}=e^{-x^{2} / 2} \sum_{0}^{\infty} \sqrt{\frac{2^{k}}{k!}} \phi_{k}(x) w^{k}
$$

Setting $w=\frac{1}{2} z$, so that $e^{-(x-w)^{2}}=e^{-x^{2} / 2} e^{-(x-z)^{2} / 2} e^{z^{2} / 4}$, we obtain the result we want: For any $z \in \mathbb{C}$,

$$
\pi^{-1 / 4} e^{-(x-z)^{2} / 2}=e^{-z^{2} / 4} \sum_{0}^{\infty} \frac{1}{\sqrt{2^{k} k!}} \phi_{k}(x) z^{k}
$$

Now, for $a, b \in \mathbb{R}$, let

$$
\psi_{a, b}(x)=\pi^{-1 / 4} e^{-(x-a)^{2} / 2} e^{i b x}
$$

Then the expected values of position and momentum in the state $\psi_{a, b}$ are $a$ and $b$, respectively. We take $\psi_{a, 0}$ as the initial state for the quantum oscillator corresponding to the classical initial values $x_{0}=a$ and $p_{0}=0$. The expansion of $\psi_{a, 0}$ in terms of the energy eigenfunctions is obtained by taking $z=a$ in (3.32). It is now\\
a simple matter to compute its time evolution, using (3.31) and (3.32):

$$
\begin{aligned}
e^{-i t H_{0}} \psi_{a, 0}(x) & =e^{-a^{2} / 4} \sum_{0}^{\infty} \frac{e^{-i t\left(k+\frac{1}{2}\right)}}{\sqrt{2^{k} k!}} \phi_{k}(x) a^{k} \\
& =e^{-a^{2} / 4} e^{-i t / 2} \sum_{0}^{\infty} \frac{1}{\sqrt{2^{k} k!}} \phi_{k}(x)\left(a e^{-i t}\right)^{k} \\
& =e^{-a^{2} / 4} e^{\left(a e^{-i t}\right)^{2} / 4} e^{-i t / 2} \pi^{-1 / 4} e^{-\left(x-a e^{-i t}\right)^{2} / 2}
\end{aligned}
$$

which simplifies to

$$
e^{i\left(a^{2} \sin t \cos t-t\right) / 2} \pi^{-1 / 4} e^{-(x-a \cos t)^{2} / 2} e^{-i x a \sin t}=e^{i\left(a^{2} \sin t \cos t-t\right) / 2} \psi_{a \cos t,-a \sin t}(x) .
$$

The factor $e^{i\left(a^{2} \sin t \cos t-t\right) / 2}$ is just a scalar of modulus 1 , so $e^{-i t H_{0}} \psi_{a, 0}$ is a Gaussian wave packet with expected position $a \cos t$ and momentum $-a \sin t$ : the classical periodic motion! (Putting the constants $\hbar, m$, and $\kappa$ back in changes these to $a \cos \sqrt{\kappa / m} t$ and $-a \sqrt{m \kappa} \sin \sqrt{\kappa / m} t$ as it should.) The states $\psi_{a, 0}$, or more generally $\psi_{a, b}$, are called coherent states because of the way they maintain their shape under the time evolution.

Of course this exact quantum-classical correspondence is a highly unusual phenomenon, but one can show that the evolution of Gaussian wave packets under much more general potentials exhibits similar behavior in an approximate sense when $\hbar$ is taken sufficiently small; see Hagedorn [61].

To complete the picture, let us calculate the expected energy of the state $\psi_{a, 0}$. Using the expansion of $\psi_{a, 0}$ in terms of energy eigenfunctions given (3.32), we have

$$
\left\langle\psi_{a, 0}\right| H_{0}\left|\psi_{a, 0}\right\rangle=e^{-a^{2} / 2} \sum_{0}^{\infty}\left(k+\frac{1}{2}\right) \frac{a^{2 k}}{2^{k} k!}=\frac{e^{-a^{2} / 2}}{2} \frac{d}{d a}\left[a e^{a^{2} / 2}\right]=\frac{1}{2}\left(a^{2}+1\right) .
$$

If the constants $\hbar, m$, and $\kappa$ are reinserted appropriately, the expected energy turns out to be $\frac{1}{2}\left(\kappa a^{2}+\hbar \sqrt{\kappa / m}\right)$. The term $\frac{1}{2} \kappa a^{2}$ is the classical energy; the extra $\frac{1}{2} \hbar \sqrt{\kappa / m}$ is the ground state energy. (In the classical ground state, $x(t) \equiv 0$, the particle is absolutely at rest, but quantum mechanically this is not quite possible.)

\subsection*{Angular momentum and spin}
Classically, the angular momentum (about the origin) of a particle in $\mathbb{R}^{3}$ with position $\mathbf{x}$ and momentum $\mathbf{p}$ is $\mathbf{l}=\mathbf{x} \times \mathbf{p}$ (cross product). The Hamiltonian vector fields corresponding to the three components of this vector observable are the infinitesimal generators of rotations about the three coordinate axes. In more detail, the first component of $\mathbf{l}$ is $l_{1}=x_{2} p_{3}-x_{3} p_{2}$, and its Hamiltonian vector field is

$$
X_{l_{1}}=x_{2} \frac{\partial}{\partial x_{3}}-x_{3} \frac{\partial}{\partial x_{2}}+p_{2} \frac{\partial}{\partial p_{3}}-p_{3} \frac{\partial}{\partial p_{2}},
$$

which generates the flow in phase space given by a rotation about the $x_{1}$-axis in position space and the same rotation about the $p_{1}$-axis in momentum space. Similarly for the other two components.

Quantization of $\mathbf{l}$ in the obvious way leads to the vector-valued quantum observable $\mathbf{L}=\mathbf{X} \times \mathbf{P}$, that is,

$$
L_{1}=X_{2} P_{3}-X_{3} P_{2}=\frac{\hbar}{i}\left(x_{2} \frac{\partial}{\partial x_{3}}-x_{3} \frac{\partial}{\partial x_{2}}\right), \text { etc. }
$$

acting on $L^{2}\left(\mathbb{R}^{3}\right)$. (Note that there is no problem about ordering the factors here, because $X_{i}$ and $P_{j}$ commute when $i \neq j$.) The one-parameter groups generated by $L_{1}, L_{2}$, and $L_{3}$ are again the rotations about the coordinate axes. That is,

$$
e^{t L_{1} / i \hbar} f(\mathbf{x})=f\left(R_{t, 1}^{-1} \mathbf{x}\right), \text { where } R_{t, 1}=\left(\begin{array}{ccc}
1 & 0 & 0 \\
0 & \cos t & \sin t \\
0 & -\sin t & \cos t
\end{array}\right)
$$

and likewise, if $\mathbf{u}$ is any unit vector in $\mathbb{R}^{3}, e^{t \mathbf{u} \cdot \mathbf{L} / i \hbar} f(\mathbf{x})=f\left(R_{t, \mathbf{u}}^{-1} \mathbf{x}\right)$, where $R_{t, \mathbf{u}}$ is rotation through the angle $t$ about the $\mathbf{u}$ axis.

In group-theoretic terms: The rotation group $S O(3)$ acts on $L^{2}\left(\mathbb{R}^{3}\right)$ by

$$
[\rho(R) f](\mathbf{x})=f\left(R^{-1} \mathbf{x}\right)
$$

Its Lie algebra $\mathfrak{s o}(3)$ can be identified with $\mathbb{R}^{3}$ with Lie bracket $=$ cross product, the unit vector $\mathbf{u} \in \mathbb{R}^{3}$ corresponding to the generator of rotations about the $\mathbf{u}$ axis. The angular momentum operators are the range of the infinitesimal representation $d \rho$ of $\mathbb{R}^{3} \cong \mathfrak{s o}(3)$, that is, the angular momentum about the $\mathbf{u}$ axis is the operator

$$
u_{1} L_{1}+u_{2} L_{2}+u_{3} L_{3}=i \hbar d \rho(\mathbf{u})=\left.\frac{d}{d t} \rho\left(R_{t, \mathbf{u}}\right)\right|_{t=0}
$$

The observables $L_{j}$ are important in quantum mechanics, but they do not tell the whole story about angular momentum: most quantum particles ${ }^{2}$ have an additional intrinsic angular momentum that has nothing to do with their state of motion, known as spin. Spin is a purely quantum phenomenon with no classical analogue, so there is no point in trying to motivate the way it works; we simply state the result. This discussion pertains to nonrelativistic particles of mass $m>0$; we shall fit it into the relativistic framework in the next chapter, where we shall also discuss massless particles.

Until now we have been taking the state space for a quantum particle in $\mathbb{R}^{3}$ to be $L^{2}\left(\mathbb{R}^{3}\right)$. However, this usually does not provide a complete description of the state of the particle; there are additional observables that do not fit into this picture. Instead, the state space must be taken as $L^{2}\left(\mathbb{R}^{3}\right) \otimes \mathcal{H}$, where $\mathcal{H}$ is another Hilbert space that describes the "internal degrees of freedom" of the particle. We assume given a unitary (possibly projective) representation of the rotation group $S O(3)$ on this state space that describes how the state of the particle changes under spatial rotations. The (total) angular momentum operators will then be taken to be the image of $\mathfrak{s o}(3)$ under the corresponding infinitesimal representation.

Since we are interested only in angular momentum at this point, and not in other possible internal features of the particles, we take $\mathcal{H}$ to be the smallest Hilbert space that will encompass this phenomenon. Namely, $\mathcal{H}$ will be a finite-dimensional space equipped with an irreducible (possibly projective) representation $\pi$ of $S O(3)$, and the action of $S O(3)$ on $L^{2}\left(\mathbb{R}^{3}\right) \otimes \mathcal{H}$ will be $\rho \otimes \pi$ where $\rho$ is the natural representation on $L^{2}\left(\mathbb{R}^{3}\right)$ discussed above. We therefore need some facts about representations of $S O(3)$ and its universal cover $S U(2)$. We refer the reader back to §1.3 for the details of the relation between $S O(3)$ and $S U(2)$, and in particular for the calculation of the covering map $\kappa$.

Up to unitary equivalence, $S U(2)$ has exactly one irreducible unitary representation on $\mathbb{C}^{n}$ for each positive integer $n$; it may be realized as the action of $S U(2)$ on the homogeneous (holomorphic) polynomials of degree $n-1$ on $\mathbb{C}^{2}$. (See

\footnotetext{${ }^{2}$ Particles can be elementary or composite here.
}Folland [45] or Hall [62] for a proof.) For physical purposes it is customary to set $n=2 k+1$, where $k$ is an integer or half-integer; the irreducible representation on $\mathbb{C}^{2 k+1}$ is denoted by $\pi_{k} . \pi_{0}$ is the trivial representation; $\pi_{1 / 2}$ is the defining representation of $\operatorname{SU}(2)$ on $\mathbb{C}^{2}$, and $\pi_{1}$ is (equivalent to) the adjoint representation on the complexification of $\mathfrak{s u}(2)$.

We have $\pi_{k}(-I)=(-1)^{2 k} I$, so $\pi_{k}$ descends to a representation of $S O(3) \cong S U(2) /\{ \pm I\}$ precisely when $k$ is an integer, and in any case it defines a projective representation of $S O(3)$. We shall denote these representations also by $\pi_{k}$. Thus, if $R \in S O(3)$, by $\pi_{k}(R)$ we really mean $\pi_{k}\left(\kappa^{-1}(R)\right)$ where $\kappa$ is the covering map defined in §1.3. This is defined only up to a factor of $\pm I$ when $k$ is a half-integer; we say that $\pi_{k}$ is a double-valued representation. The representations $\pi_{k}$ exhaust the irreducible projective representations of $S O(3)$, and those with $k$ integral exhaust the irreducible unitary ones, up to equivalence.

Let $Y_{1}, Y_{2}, Y_{3} \in \mathfrak{s o}(3)$ be the infinitesimal generators of the rotations about the three coordinate axes; thus $Y_{j}=\kappa^{\prime}\left(\frac{1}{2} i \sigma_{j}\right)$ in the notation of §1.3. Let $\pi$ be any unitary representation of $S U(2)$ and $d \pi$ the corresponding representation of $\mathfrak{s o}(3) \cong \mathfrak{s u}(2)$. The operators

$$
J_{j}=\frac{\hbar}{i} d \pi\left(Y_{j}\right)=\frac{1}{2} \hbar d \pi\left(\sigma_{j}\right) \quad j=1,2,3
$$

are the three components of the angular momentum $\mathbf{J}$ in the representation $\pi$. Further, let $\mathbf{Y}^{2}=Y_{1}^{2}+Y_{2}^{2}+Y_{3}^{2}$, an element of the universal enveloping algebra of $\mathfrak{s o}(3)$. The operator

$$
\mathbf{J}^{2}=J_{1}^{2}+J_{2}^{2}+J_{3}^{2}=-\hbar^{2} d \pi\left(\mathbf{Y}^{2}\right)
$$

the squared total angular momentum, plays an important role. As is well known, $\mathbf{Y}^{2}$ is the so-called Casimir element of the universal enveloping algebra of $\mathfrak{s o}(3)$, and it belongs to (in fact, generates) the center of that algebra. It follows that $\mathbf{J}^{2}$ commutes with the representation $\pi$; therefore, by Schur's lemma, it acts as a scalar on each irreducible subspace of the representation space of $\pi$.

Let us consider the angular momentum operators for the irreducible representations $\pi_{k}$, with units of measurement chosen so that $\hbar=1$. Each operator $J_{j}$ has the $2 k+1$ distinct eigenvalues $-k,-k+1, \ldots, k-1, k$. It is enough to verify this for $J_{3}$, as the $J_{j}$ are all unitarily equivalent (rotations about all axes look the same). If we realize $\pi_{k}$ as the action of $S U(2)$ on the homogeneous polynomials of degree $2 k$ on $\mathbb{C}^{2}$, we have

$$
\pi_{k}\left(e^{i t \sigma_{3} / 2}\right)\left(z_{1}^{m} z_{2}^{2 k-m}\right)=\left(e^{i t / 2} z_{1}\right)^{m}\left(e^{-i t / 2} z_{2}\right)^{2 k-m}=e^{i(m-k) t} z_{1}^{m} z_{2}^{2 k-m}
$$

and so

$$
J_{3}\left(z_{1}^{m} z_{2}^{2 k-m}\right)=\left.\frac{1}{i} \frac{d}{d t} \pi_{k}\left(e^{i t \sigma_{3} / 2}\right)\left(z_{1}^{m} z_{2}^{2 k-m}\right)\right|_{t=0}=(m-k) z_{1}^{m} z_{2}^{2 k-m}
$$

Thus the eigenvalues of $J_{3}$ are $\{k-m: m=0, \ldots, 2 k\}$, as advertised. Furthermore, since $\pi_{k}$ is irreducible, $\mathbf{J}^{2}$ is a scalar multiple of the identity. To find the scalar, it suffices to compute the action of $\mathbf{J}^{2}$ on one vector, for example, the monomial $z_{1}^{2 k}$ :\\
$\mathbf{J}^{2}\left(z_{1}^{2 k}\right)=-\frac{d^{2}}{d t^{2}}\left[\left(e^{i t / 2} z_{1}\right)^{2 k}+\left(z_{1} \cos \frac{1}{2} t+i z_{2} \sin \frac{1}{2} t\right)^{2 k}+\left(z_{1} \cos \frac{1}{2} t-z_{2} \sin \frac{1}{2} t\right)^{2 k}\right]_{t=0}$, which the reader may verify to be $k(k+1) z_{1}^{2 k}$; thus $\mathbf{J}^{2}=k(k+1) I$.

Now, a quantum particle whose state space is $L^{2}\left(\mathbb{R}^{3}\right) \otimes \mathbb{C}^{2 k+1}$, equipped with the representation $\rho \otimes \pi_{k}$ where $\rho$ is the natural representation of $S O(3)$ on $L^{2}\left(\mathbb{R}^{3}\right)$,\\
is said to have spin $k$. We can think of $L^{2}\left(\mathbb{R}^{3}\right) \otimes \mathbb{C}^{2 k+1}$ as $L^{2}\left(\mathbb{R}^{3}, \mathbb{C}^{2 k+1}\right)$, the space of square-integrable $\mathbb{C}^{2 k+1}$-valued functions on $\mathbb{R}^{3}$, and the representation $\rho \otimes \pi_{k}$ is then given by

$$
\left[\left(\rho \otimes \pi_{k}\right)(R) f\right](\mathbf{x})=\pi_{k}(R)\left[f\left(R^{-1} \mathbf{x}\right)\right]
$$

The infinitesimal representation of $\rho \otimes \pi_{k}$ is $d \rho \otimes I+I \otimes d \pi_{k}$, so the angular momentum operators are

$$
\mathbf{J}=\mathbf{L}+\mathbf{S}, \text { where } \mathbf{L}=\frac{\hbar}{i} d \rho(\mathbf{Y}) \text { and } \mathbf{S}=\frac{\hbar}{i} d \pi_{k}(\mathbf{Y})
$$

$\mathbf{L}$ is just the angular momentum discussed at the beginning of this section, now called the orbital angular momentum, acting componentwise on the vector-valued functions in $L^{2}\left(\mathbb{R}^{3}, \mathbb{C}^{2 k+1}\right)$. $\mathbf{S}$ is the spin angular momentum, acting on the components of the functions.

We briefly describe the decomposition of the representation $\rho$ into irreducible components; see Folland [46] for more details. The basic ingredients are the spaces of spherical harmonics $\delta \mathcal{H}_{l}, l=0,1,2, \ldots$, whose elements are functions of the form $Y(\mathbf{x})=P(\mathbf{x})|\mathbf{x}|^{-l}$ where $P$ is a homogeneous harmonic polynomial of degree $l$. (Note that such functions are homogeneous of degree 0 , i.e., are effectively functions of $\mathbf{x} /|\mathbf{x}|$.) The spaces $\mathcal{S} \mathcal{H}_{l}$ are invariant under the natural action of the rotation group, and they are irreducible; the representation of $S O(3)$ on $\mathcal{S} \mathcal{H}_{l}$ is equivalent to $\pi_{l}$. For each $l$, let

$$
\mathcal{H}_{l}=\left\{f \in L^{2}\left(\mathbb{R}^{3}\right): f(\mathbf{x})=g(|\mathbf{x}|) Y(\mathbf{x}): Y \in \mathcal{S} \mathcal{H}_{l}\right\} .
$$

The condition that $f$ be in $L^{2}\left(\mathbb{R}^{3}\right)$ is equivalent to the condition that $g$ be in $L^{2}\left((0, \infty), r^{2} d r\right)$. The rotation group acts trivially on the factor $g$ and by the representation $\pi_{l}$ on the factor $Y$, so a choice of orthonormal basis for $L^{2}\left((0, \infty), r^{2} d r\right)$ exhibits $\mathcal{H}_{l}$ as a direct sum of spaces isomorphic to $\mathcal{S} \mathcal{H}_{l}$, and the representation $\rho \mid \mathcal{H}_{l}$ as a direct sum of copies of $\pi_{l}$. We have

$$
L^{2}\left(\mathbb{R}^{3}\right)=\bigoplus_{0}^{\infty} \mathcal{H}_{l}
$$

yielding the decomposition of $\rho$ into $S O(3)$-isotypic components. Observe that this can also be construed as the spectral decomposition of the squared orbital angular momentum operator $\mathbf{L}^{2}$ : the space $\mathcal{H}_{l}$ is the eigenspace of $\mathbf{L}^{2}$ with eigenvalue $l(l+1)$.

From this one can obtain the irreducible decomposition of $\rho \otimes \pi_{k}$ for any $k$. The preceding analysis yields

$$
L^{2}\left(\mathbb{R}^{3}, \mathbb{C}^{2 k+1}\right)=\bigoplus_{0}^{\infty} \mathcal{H}_{l} \otimes \mathbb{C}^{2 k+1}
$$

and the representation of $S O(3)$ on $\mathcal{H}_{l} \otimes \mathbb{C}^{2 k+1}$ is a direct sum of copies of $\pi_{l} \otimes \pi_{k}$. What remains, therefore, is to decompose $\pi_{l} \otimes \pi_{k}$ into its irreducible components. This problem is well understood, but the details of its solution will not concern us here. (See Landau and Lifshitz $[\mathbf{7 7}]$; the key-words to look for are "Clebsch-Gordan coefficients.") For us it will suffice to describe the most important case, $k=\frac{1}{2}$.

Of course $\pi_{0}$ is the trivial one-dimensional representation, so $\pi_{0} \otimes \pi_{1 / 2} \cong \pi_{1 / 2}$. For $l>0$, consider an irreducible subspace of $\mathcal{H}_{l}$, say

$$
\mathcal{X}=\left\{f: f(\mathbf{x})=g(|\mathbf{x}|) Y(\mathbf{x}), Y \in S \mathcal{H}_{l}\right\}
$$

for a fixed $g$, and consider the action of the $z$ component of the total angular momentum, $J_{3}=L_{3}+S_{3}$, on $\mathfrak{X} \otimes \mathbb{C}^{2}$. The operators $L_{3}$ and $S_{3}$, considered as acting on $X$ and $\mathbb{C}^{2}$ respectively, have orthonormal eigenbases $\left\{\phi_{m}\right\}_{m=0}^{2 l}$ and $\left\{v_{+}, v_{-}\right\}$with $L_{3} \phi_{m}=(l-m) \phi_{m}, S_{3} v_{ \pm}= \pm \frac{1}{2}$. Then $\left\{\phi_{m} \otimes v_{a}: m=0, \ldots, 2 l ; a= \pm\right\}$ is an orthonormal eigenbasis for $J_{3}$ with $J_{3}\left(\phi_{m} \otimes v_{ \pm}\right)=l-m \pm \frac{1}{2}$. The eigenvalues $\pm\left(l+\frac{1}{2}\right)$ occur with multiplicity one and the others occur with multiplicity two. The largest of them is $l+\frac{1}{2}$, which occurs only for the eigenvector $\phi_{0} \otimes v_{+}$; hence the representation of $S O(3)$ on the invariant subspace generated by this vector is a copy of $\pi_{l+(1 / 2)}$. This accounts for one copy of each of the eigenvalues $l-m \pm \frac{1}{2}$. Its orthogonal complement in $X \otimes \mathbb{C}^{2}$ is an $S O(3)$-invariant space on which $J_{3}$ has eigenvalues $l-\frac{1}{2}, l-\frac{3}{2}, \ldots,-l+\frac{1}{2}$, each with multiplicity one; the representation on it is therefore a copy of $\pi_{l-(1 / 2)}$. In short,

$$
\pi_{l} \otimes \pi_{1 / 2} \cong \pi_{l+(1 / 2)} \oplus \pi_{l-(1 / 2)} \text { for } l>0, \quad \pi_{0} \otimes \pi_{1 / 2} \cong \pi_{1 / 2}
$$

\subsection*{The Coulomb potential}
We now perform a spectral analysis of the Hamiltonian

$$
H=-\frac{\hbar^{2}}{2 m} \nabla^{2}-\frac{a}{|\mathbf{x}|}
$$

for a particle of mass $m$ moving in the potential $V(\mathbf{x})=-a /|\mathbf{x}|(a>0)$ corresponding to an attractive inverse-square-law force. The operator $H$ acts on $L^{2}\left(\mathbb{R}^{3}, \mathbb{C}^{2 k+1}\right)$ if the particle has spin $k$, but the spin variables just come along for the ride here, so there is no loss of generality in taking $k=0$. Our main task is to find and analyze the solutions of the differential equation $H f=E f$ for $E \in \mathbb{R}$ ( $E$ for energy). We sketch the main points but omit many of the details of the calculations. A more complete treatment can be found in Landau and Lifshitz [ $\mathbf{7 7 ]}$.

The basic example is an electron with mass $m_{e}$ and charge $-e$ moving in the Coulomb potential generated by an atomic nucleus with mass $m_{N}$ and charge $Z e$ where $Z$ is the number of protons in the nucleus, so that $a=Z e^{2} / 4 \pi$ in HeavisideLorentz units. It is a pretty good approximation to pretend that $m_{N}=\infty$, so that the nucleus does not move; however, this two-body problem, just like the classical one discussed at the end of §2.1, is mathematically exactly equivalent to a single particle moving in the fixed potential $-a /|\mathbf{x}|$ if one takes its mass $m$ to be the reduced mass $m_{e} m_{N} /\left(m_{e}+m_{N}\right)$.

First, we get rid of the constants by rescaling. The constant $a$ has dimensions $\left[m l^{3} t^{-2}\right]$ since $a /|\mathbf{x}|$ represents an energy, so $\hbar^{2} / m a$ has dimensions of length and $m a^{2} / \hbar^{2}$ has dimensions of energy. Setting $\mathbf{x}=\left(\hbar^{2} / m a\right) \widetilde{\mathbf{x}}$ (so $\widetilde{\mathbf{x}}$ is dimensionless) makes

$$
H=\frac{m a^{2}}{2 \hbar^{2}}\left[-\nabla_{\widetilde{\mathbf{x}}}^{2}-\frac{2}{|\widetilde{\mathbf{x}}|}\right]
$$

We now drop the tildes and choose the unit of energy to be $m a^{2} / 2 \hbar^{2}$, so

$$
H=-\nabla^{2}-\frac{2}{|\mathbf{x}|}
$$

(For a hydrogen atom consisting of one electron and one proton, with $m$ taken to be the reduced mass $m_{e} m_{p} /\left(m_{e}+m_{p}\right)$, the resulting units of length and energy are $\hbar^{2} / m a=5.29 \times 10^{-9} \mathrm{~cm}$ and $m a^{2} / 2 \hbar^{2}=13.6 \mathrm{eV}$. The latter is known as the

Rydberg energy ; as we shall see, it is the absolute value of the energy of the electron in the ground state.)

Now, $H$ is invariant under rotations, so it respects the decomposition $L^{2}\left(\mathbb{R}^{3}\right)= \bigoplus_{0}^{\infty} \mathcal{H}_{l}$, where $\mathcal{H}_{l}$ is as in (3.33). Moreover, the Laplacian is given in spherical coordinates by

$$
\nabla^{2}=\frac{d^{2}}{d r^{2}}+\frac{2}{r} \frac{d}{d r}-\frac{1}{r^{2}} \mathbf{L}^{2} \quad(r=|\mathbf{x}|)
$$

where $\mathbf{L}^{2}$ is our friend the squared orbital angular momentum operator. On $\mathscr{H}_{l}$ we have $\mathbf{L}^{2}=l(l+1) I$, so the restriction of $H$ to $\mathcal{H}_{l}$ is given by

$$
H\left(g(r) Y_{l}(\mathbf{x})\right)=\left[-g^{\prime \prime}(r)-\frac{2}{r} g^{\prime}(r)+\frac{l(l+1)}{r^{2}} g(r)-\frac{2}{r} g(r)\right] Y_{l}(\mathbf{x})
$$

Hence the partial differential equation $H f=E f$ boils down to the ordinary differential equations

$$
g^{\prime \prime}(r)+\frac{2}{r} g^{\prime}(r)+\left(\frac{2}{r}-\frac{l(l+1)}{r^{2}}+E\right) g(r)=0 \quad(l=0,1,2, \ldots) .
$$

As far as the discrete spectrum of $H$ (i.e., the stationary states of the quantum system) goes, we are interested in solutions $g$ of (3.35) such that the function $f(\mathbf{x})=g(|\mathbf{x}|) Y_{l}(\mathbf{x})$ belongs to $L^{2}\left(\mathbb{R}^{3}\right)$; this requires that $g$ be well-behaved at 0 and $\infty$. Moreover, we expect that the stationary states will occur when the particle is "trapped in the potential well," i.e., for energies $E<0$. Thus we take the condition $E<0$ as a working assumption for the time being.

The equation (3.35) has a regular singular point at the origin, where its characteristic exponents are $l$ and $-l-1$; thus it has two solutions that are asymptotic to $r^{l}$ and $r^{-l-1}$ as $r \rightarrow 0$. The second solution must be discarded as the resulting function on $\mathbb{R}^{3}$ is not square-integrable. ${ }^{3}$ Moreover, for large $r$ the equation (3.35) is approximately $g^{\prime \prime}(r)+E g(r)=0$, so we expect to find solutions that are asymptotic to $e^{ \pm \sqrt{-E} r}$ as $r \rightarrow \infty$, and we are only interested in the one with the negative exponent. (Remember that we are assuming $E<0$.)

Motivated by these considerations, we make some changes of variables. First, to standardize the behavior at infinity (and, frankly, with some foreknowledge of the final result), we set

$$
\nu=\frac{1}{\sqrt{-E}}, \quad s=\frac{2 r}{\nu}, \quad \widetilde{g}(s)=g(r)=g(\nu s / 2)
$$

This turns (3.35) into

$$
\widetilde{g}^{\prime \prime}(s)+\frac{2}{s} \widetilde{g}^{\prime}(s)+\left(\frac{\nu}{s}-\frac{l(l+1)}{s^{2}}-\frac{1}{4}\right) \widetilde{g}(s)=0
$$

The solutions are now expected to be asymptotic to $e^{ \pm s / 2}$ as $s \rightarrow \infty$ and to $s^{l}$ as $s \rightarrow 0$. This suggests the further substitution

$$
\widetilde{g}(s)=s^{l} e^{-s / 2} h(s)
$$

which turns (3.36) into

$$
s h^{\prime \prime}(s)+(2 l+2-s) h^{\prime}(s)+(\nu-l-1) h(s)=0 .
$$

\footnotetext{${ }^{3}$ When $l=0$ this argument is insufficient, as the function $f(\mathbf{x})=|\mathbf{x}|^{-1}$ is in $L^{2}\left(\mathbb{R}^{3}\right)$. However, it is not in the domain of $H$, viz., the Sobolev space of order 2. Indeed, $\nabla^{2}|\mathbf{x}|^{-1}=-4 \pi \delta(\mathbf{x})$.
}This, finally, is something in standard form: it is the confluent hypergeometric equation with parameters $a=l+1-\nu$ and $c=2 l+2$, and the solutions that are regular at the origin are constant multiples of

$$
F(l+1-\nu, 2 l+2 ; s)=1+\sum_{k=1}^{\infty} \frac{(l+1-\nu)(l+2-\nu) \cdots(l+k-\nu)}{(2 l+2)(2 l+3) \cdots(2 l+k+1)} \frac{s^{k}}{k!}
$$

(The other solutions are asymptotic to multiples of $s^{-2 l-1}$ as $s \rightarrow 0$.) A wealth of information about these confluent hypergeometric functions is known; in particular, they grow like $e^{s}$ as $s \rightarrow \infty$ except when $l+1-\nu$ is a nonpositive integer, in which case the series terminates and gives a polynomial of degree $\nu-l-1$. In this case $\nu$ is itself an integer $\geq l+1$, which (following standard practice) we denote by $n$, and the resulting polynomial is one of the (generalized) Laguerre polynomials:

$$
F(l+1-n, 2 l+2 ; s)=\frac{(2 l+1)!(n-l-1)!}{(n+l)!} L_{n-l-1}^{2 l+1}(s) .
$$

Undoing the changes of variables, we see that the numbers $E_{n}=-1 / n^{2}$ are eigenvalues of $H \mid \mathcal{H}_{l}$ for every integer $n \geq l+1$, and the corresponding (unnormalized) eigenfunctions are

$$
f_{n, Y}(\mathbf{x})=r^{l} e^{-r / n} L_{n-l-1}^{2 l+1}(2 r / n) Y(\mathbf{x}), \quad Y \in \mathcal{S} \mathcal{H}_{l} \quad(r=|\mathbf{x}|) .
$$

For $E>0$ one can make the same changes of variables to find that the solutions of (3.35) that are regular at the origin are multiples of $r^{l} e^{-r / \nu} F(l+1-\nu, 2 l+ 2 ; 2 r / \nu)$, but here $\nu=1 / \sqrt{-E}$ is pure imaginary. One can use contour integral techniques to figure out the asymptotic behavior of this function as $r \rightarrow \infty$; it turns out to be asymptotic to a constant times $r^{-1} \sin \left(\beta r+\beta^{-1} \log 2 \beta r+C_{l}\right)$ with $\beta=\sqrt{E}$, which is not in $L^{2}\left((0, \infty), r^{2} d r\right)$. For $E=0$, the equation (3.35) can be transformed into a Bessel equation, and the solution turns out to be $\sqrt{1 / r} J_{2 l+1}(\sqrt{8 r})$, which also is not in $L^{2}$. Hence there is no discrete spectrum for $E \geq 0$. However, the corresponding eigenfunctions of the differential operator $H$,

$$
f_{\beta, Y}(\mathbf{x})=r^{l} e^{-i \beta r} F(l+1+(i / \beta), 2 l+2 ; 2 i \beta r) Y(\mathbf{x}), \quad Y \in \mathcal{S} \mathcal{H}_{l}
$$

are close enough to being in $L^{2}$ that they contribute infinitesimally to the spectral resolution of the Hilbert space operator $H$, just as $e^{i \omega x}$ contributes infinitesimally to the spectral resolution of $i d / d x$ on $L^{2}(\mathbb{R})$ (viz., the Fourier transform).

In short, the discrete spectrum of $H \mid \mathcal{H}_{l}$ is $\left\{-1 / n^{2}: n \geq l+1\right\}$ and the continuous spectrum is $[0, \infty)$. Let us choose a basis $Y_{1}, \ldots, Y_{2 l+1}$ for $\mathcal{S} \mathcal{H}_{l}$ (the physicists' standard choice is the eigenbasis of the $z$-angular momentum $L_{3}$ ). Then any $\phi \in \mathcal{H}_{l}$ can be written as $\phi(\mathbf{x})=\sum_{1}^{2 l+1} \phi_{m}(r) Y_{m}(\mathbf{x})$, and the expansion of $\phi$ in terms of eigenfunctions of $H$ has the form

$$
\phi(\mathbf{x})=\sum_{m=1}^{2 l+1}\left[\sum_{n=l+1}^{\infty} C_{m n} f_{n, Y_{m}}(\mathbf{x})+\int_{0}^{\infty} \widetilde{\phi}_{m}(\beta) f_{\beta, Y_{m}}(\mathbf{x}) d \beta\right]
$$

The coefficients $C_{m n}$ and $\widetilde{\phi}_{m}(\beta)$ are easily determined once the $f_{n, Y}$ 's and $f_{\beta, Y}$ 's are properly normalized.

Putting together all the $\mathcal{H}_{l}$ 's, we see that the discrete spectrum of $H$ is $\left\{-1 / n^{2}\right.$ : $n \geq 1\}$. The eigenvalue $-1 / n^{2}$ occurs in $\mathcal{H}_{l}$ when $l \leq n-1$, and its multiplicity there is $\operatorname{dim}\left(\mathcal{S} \mathcal{H}_{l}\right)=2 l+1$, so its total multiplicity is $\sum_{0}^{n-1}(2 l+1)=n^{2}$. The continuous spectrum is $[0, \infty)$, and it occurs in each $\mathcal{H}_{l}$, so its total multiplicity is\\
infinite. Finally, if the particle in the Coulomb potential has spin $k$, the Hamiltonian is the direct sum of $2 k+1$ copies of the scalar Hamiltonian, so the multiplicities of all the eigenvalues must be multiplied by $2 k+1$.

In particular, for an electron, we have $k=\frac{1}{2}$, and the multiplicity of the eigenvalue $-1 / n^{2}$ is $2 n^{2}$. In this case it is customary to label the joint eigenstates of the energy and the angular momentum by the integers $n$ (the "principal quantum number"), $l$ (the "orbital angular momentum quantum number"), and $j$ (the "total angular momentum quantum number"). Here $n$ and $l$ are as in the preceding discussion, and $j$ is the half-integer such that the eigenvalue of the total squared angular momentum $\mathbf{J}^{2}$ is $j(j+1)$. In view of (3.34), $j$ is either $l+\frac{1}{2}$ or $l-\frac{1}{2}$, with $j=\frac{1}{2}$ being the only possibility when $l=0$. Moreover, in deference to the old spectroscopists' terminology, for $l=0,1,2,3$ it is customary to denote the states with quantum numbers $(n, l, j)$ by $n L_{j}$ where $L=S, P, D, F$ according as $l=0,1,2,3$. Thus, for example, the states with quantum numbers $\left(1,0, \frac{1}{2}\right),\left(2,1, \frac{1}{2}\right)$, and $\left(3,2, \frac{5}{2}\right)$ are denoted by $1 S_{1 / 2}, 2 P_{1 / 2}$, and $3 D_{5 / 2}$.

This description of an electron in a Coulomb potential, which accounts for the "electron shells" of the Bohr model as the eigenspaces of the Hamiltonian, was one of the early triumphs of the quantum theory. Its prediction of the electron energy levels in a hydrogen atom agrees very well with the spectroscopic data that were available when quantum mechanics was developed. Moreover, it provides a somewhat cruder but still very informative model for larger atoms when one neglects the electrical interaction of the electrons with each other except to note that the inner electrons tend to shield the outer electrons from the charge of the nucleus. The only further necessary ingredient is the Pauli exclusion principle, which we shall discuss in §5.4, §5.5, and §6.5; in the present context it says that the state vectors of the various electrons of the atom must all be orthogonal to each other. Thus there is room for 2 electrons in the lowest energy level ( $n=1$ ), 8 electrons in the next level ( $n=2$ ), and so forth. From this one can see the periodic table of elements taking shape before one's eyes.

The agreement with experiment is not perfect, however: the actual energy levels of electrons in the hydrogen atom depend to a small extent on the angular momentum quantum numbers too. The biggest of these effects, the so-called "fine splitting" between the states with the same $n$ but different $j$, is explained by the Dirac model, which we shall discuss in the next chapter.

\section*{CHAPTER 4}
\section*{Relativistic Quantum Mechanics}
\begin{displayquote}
When an equation is as successful as Dirac's, it is never simply a mistake. It may not be valid for the reason supposed by the author, it may break down in new contexts, and it may not even mean what its author thought it meant. . . . But the great equations of modern physics are a permanent part of scientific knowledge, which may outlast even the beautiful cathedrals of earlier ages.
\end{displayquote}

-Steven Weinberg ([31], p. 257)

The historical development of relativistic quantum theory is too tangled a tale to recount in detail here. (An interesting brief account can be found in Weinberg [131], and Schweber [107] has produced a masterful comprehensive history. See also Schweber [108] and the references given there.) Initially the idea was to develop a relativistic theory for the motion of a single particle, or a fixed finite collection of particles, subject to electromagnetic forces. Eventually it was realized that this is the wrong goal, as high-energy interactions necessarily involve creation, annihilation, and transmutation of particles. (One should keep in mind that in the late 1920s, when the initial attempts at a relativistic theory were made, the only known subatomic particles were electrons, protons, and atomic nuclei, and the composition of the latter was still a matter of conjecture. The experimental discovery of neutrons and positrons, not to mention other species of particles, was still in the future.) The early attempts at a one-particle theory yielded differential equations, now known by the names of Klein, Gordon, and Dirac, that were meant to be relativistic analogues of the Schrödinger equation. ${ }^{1}$ Interpreted as wave equations for single quantum particles, the Klein-Gordon equation was largely unsuccessful, and the Dirac equation was hugely successful in accounting for low-energy phenomena but presented serious problems in the high-energy regime. But you can't keep a good differential equation down: these equations turn out to be of fundamental importance for quantum field theory - but with a different interpretation that we shall describe briefly at the end of §4.3 and exploit in the following chapters. Before we get to that, however, it will be well to study the Klein-Gordon and Dirac equations as the early quantum physicists did.

From now on we shall adopt "natural" units in which Planck's constant $\hbar$ and the speed of light $c$ are both equal to 1 . If one keeps track of the dimensions of the quantities one is working with, reinserting the factors of $\hbar$ and $c$ is always just a matter of getting equations and formulas to be dimensionally correct.

\footnotetext{${ }^{1}$ The Klein-Gordon equation was proposed independently by Schrödinger, so it could equally well be named the relativistic Schrödinger equation, and it is so called by some people.
}\subsection*{The Klein-Gordon and Dirac equations}
The Schrödinger equation $i \partial \psi / \partial t=-(1 / 2 m) \nabla^{2} \psi+V \psi$ is obviously unsatisfactory from the relativistic point of view, as it treats the time and space variables on a drastically different footing. But recall how we derived it: (i) The (classical, nonrelativistic) energy of a particle of mass $m$ in a potential $V$ is $H=|\mathbf{p}|^{2} / 2 m+V(\mathbf{x})$. (ii) $H$ is the generator of time translations according to the rules of Hamiltonian mechanics. (iii) In quantum mechanics, $\mathbf{p}$ becomes $(1 / i) \nabla$, and $V(\mathbf{x})$ acts as a multiplication operator. We might therefore try to obtain a relativistic equation by modifying (i) so as to obtain a relativistically correct formula for the energy and then proceeding as in (ii) and (iii). For this purpose we need to drop the potential $V$, which is not relativistically meaningful as it stands, and replace it with something else. The only Lorentz-covariant interaction we have at hand is the electromagnetic one, so we shall consider only the cases of a free particle and a particle in an electromagnetic field; and we shall not add the electromagnetic field to the recipe until the case of a free particle is well in hand.

Recall that in relativistic mechanics, the energy $E$ and the momentum $\mathbf{p}$ of a free particle combine to give a 4 -vector $p^{\mu}=(E, \mathbf{p})$ (in units such that $c=1$ ) whose Lorentz inner product with itself is the rest mass of the particle. The quantum analogue of the momentum $\mathbf{p}$ is the operator $-i \nabla$, and the quantum analogue of the energy is the Hamiltonian operator $H$, which generates the time-translation group $e^{t H / i}$ and hence corresponds to $i \partial_{t}$. In short, in relativistic notation the recipe for quantizing energy and momentum is simply

$$
p^{\mu} \rightarrow i \partial^{\mu}
$$

where $\partial^{\mu}$ is given by (1.3). (Note that this makes the signs come out right!) Applying (4.1) to the free-particle energy-momentum relation $p^{2}=m^{2}$ yields our first relativistic quantum wave equation, the Klein-Gordon equation:

$$
-\partial^{2} \phi=m^{2} \phi, \quad \text { or } \quad \partial_{t}^{2} \phi=\nabla^{2} \phi-m^{2} \phi .
$$

The first thing that catches the eye is that (4.2) is second-order in $t$, so the properly posed initial value problem requires not only $\psi$ but also $\partial_{t} \psi$ to be given at $t=0$. What to do about this? Perhaps one should take the state vector to be not $\psi$ but the pair $\left(\psi_{1}, \psi_{2}\right)=\left(\psi, \partial_{t} \psi\right)$. One can then rewrite (4.2) as a first-order system:

$$
\frac{\partial}{\partial t}\binom{\psi_{1}}{\psi_{2}}=\binom{\psi_{2}}{\nabla^{2} \psi_{1}-m^{2} \psi_{1}}=\left(\begin{array}{cc}
0 & 1 \\
\nabla^{2}-m^{2} & 0
\end{array}\right)\binom{\psi_{1}}{\psi_{2}} .
$$

The trouble with (4.3) is that the operator on $L^{2} \times L^{2}$ represented by the matrix on the right is not skew-adjoint, so it does not generate a one-parameter unitary group. One might also think of replacing (4.2) by the equation $i \partial_{t} \psi=A \psi$ where $A$ is a self-adjoint operator whose square is the positive operator $-\nabla^{2}+m^{2}$. This idea is not completely without merit, and it will resurface at the end of §4.3. But of course $i \partial_{t} \psi=A \psi$ is no longer a differential equation, and it was not considered seriously by the pioneers of quantum theory.

Another way of looking at the problem is to think of the unitarity of time evolution as a conservation law. That is, if $\psi(\mathbf{x}, t)=e^{-i t H} \psi_{0}(\mathbf{x})$ is a solution of the ordinary Schrödinger equation $i \partial_{t} \psi=H \psi$, where $H$ is (a self-adjoint form of) $-\nabla^{2} / 2 m+V$, then $\|\psi\|^{2}=\int|\psi(\mathbf{x}, t)|^{2} d \mathbf{x}$ is a conserved quantity. We can express\\
this conservation law as a continuity equation as we did for conservation of charge in §2.4. Indeed,

$$
\begin{aligned}
\frac{\partial}{\partial t}|\psi|^{2} & =\frac{\partial \psi}{\partial t} \psi^{*}+\psi \frac{\partial \psi^{*}}{\partial t}=\frac{i}{2 m}\left[\left(\nabla^{2} \psi\right) \psi^{*}-\psi\left(\nabla^{2} \psi^{*}\right)\right] \\
& =\frac{i}{2 m} \operatorname{div}\left[\psi^{*} \nabla \psi-\psi \nabla \psi^{*}\right]
\end{aligned}
$$

Thus, taking $\|\psi\|=1$ so that $\rho=|\psi|^{2}$ represents a probability density, if we define the "probability current density" to be $\mathbf{J}=-(i / 2 m)\left(\psi^{*} \nabla \psi-\psi \nabla \psi^{*}\right)= (1 / m) \operatorname{Im}\left(\psi^{*} \nabla \psi\right)$, we have

$$
\frac{\partial \rho}{\partial t}+\operatorname{div} \mathbf{J}=0
$$

There is a similar continuity equation associated to (4.2). Namely, if $\psi$ satisfies (4.2) and we set $j^{\mu}=(1 / m) \operatorname{Im}\left(\psi^{*} \partial^{\mu} \psi\right)$, then $\partial_{\mu} j^{\mu}=0$, i.e., $\partial_{t} \rho+\operatorname{div} \mathbf{j}=0$ where $j^{\mu}=(\rho, \mathbf{j})$. Thus the conserved quantity for (4.2) is $\rho \rho=(1 / m) \int \operatorname{Im} \psi^{*} \partial_{t} \psi$. But $\rho=(1 / m) \operatorname{Im} \psi^{*} \partial_{t} \psi$ (or its negative) cannot be interpreted as a probability density because, in general, it is not positive. (If $\psi$ is real, as it might well be since (4.2) has real coefficients, $\rho$ vanishes identically!) For this reason, among others, (4.2) was found to be unsatisfactory as a relativistic wave equation for a free particle. (See Weinberg [131] and Schweber [107] for more extensive discussions of this matter.)

Dirac had the insight that what is impossible with scalar-valued functions might be feasible with vector-valued functions. What is desired is a Lorentz-covariant Schrödinger-type equation $i \partial_{t} \psi=H \psi$ for a free particle of mass $m>0$, where the wave function $\psi$ takes values in $\mathbb{C}^{n}$ for some $n$ to be determined. For this purpose the space and time variables must be on an equal footing, so $H$ must be a first-order differential operator in the space variables with constant coefficients:

$$
H=\frac{1}{i}\left(\alpha^{1} \partial_{1}+\alpha^{2} \partial_{2}+\alpha^{3} \partial_{3}\right)+m \beta
$$

Here $\alpha^{1}, \alpha^{2}, \alpha^{3}$ and $\beta$ are $n \times n$ complex matrices that we require to be Hermitian so that $H$ will be self-adjoint, and the factor of $m$ is inserted for later convenience. Moreover, $H$ represents the total energy of the particle, which in classical relativity satisfies $E^{2}=|\mathbf{p}|^{2}+m^{2}$, so we want $H^{2}=-\nabla^{2}+m^{2}$. But

$$
H^{2}=-\sum_{1}^{3}\left(\alpha^{j}\right)^{2} \partial_{j}^{2}-\sum_{1 \leq j<k \leq 3}\left(\alpha^{j} \alpha^{k}+\alpha^{k} \alpha^{j}\right) \partial_{j} \partial_{k}+\frac{m}{i} \sum_{1}^{3}\left(\alpha^{j} \beta+\beta \alpha^{j}\right) \partial_{j}+m^{2} \beta^{2},
$$

so to make $H^{2}=-\nabla^{2}+m^{2}$ we need

$$
\alpha^{j} \alpha^{k}+\alpha^{k} \alpha^{j}=2 \delta_{j k} I, \quad \alpha^{j} \beta+\beta \alpha^{j}=0, \quad \beta^{2}=I
$$

We are now on familiar mathematical ground: the conditions (4.4) say that $\alpha^{1}, \alpha^{2}$, $\alpha^{3}$, and $\beta$ are the generators of a Clifford algebra. The smallest $n$ for which one can find $n \times n$ matrices satisfying (4.4) is $n=4$, and the reader may verify that we may take

$$
\alpha^{j}=\left(\begin{array}{cc}
0 & \sigma_{j} \\
\sigma_{j} & 0
\end{array}\right), \quad \beta=\left(\begin{array}{cc}
I & 0 \\
0 & -I
\end{array}\right),
$$

(written in $2 \times 2$ blocks) where the $\sigma_{j}$ are the Pauli matrices (1.12). ${ }^{2}$ Note that these matrices are all Hermitian, as desired. We have now arrived at the Dirac

\footnotetext{${ }^{2}$ If $m=0$ we need only the $\alpha$ 's, not $\beta$, and in this case $n=2$ suffices: we can take $\alpha^{j}=\sigma_{j}$.
}
equation for a free particle of mass $m$ :
$$
i \frac{\partial \psi}{\partial t}=\frac{1}{i} \sum_{1}^{3} \alpha^{j} \frac{\partial \psi}{\partial x^{j}}+m \beta \psi
$$

There is still a slight asymmetry between the space and time variables that may be removed as follows. Define the Dirac matrices $\gamma^{\mu}$ by

$$
\gamma^{0}=\beta=\left(\begin{array}{cc}
I & 0 \\
0 & -I
\end{array}\right), \quad \gamma^{j}=\beta \alpha^{j}=\left(\begin{array}{cc}
0 & \sigma_{j} \\
-\sigma_{j} & 0
\end{array}\right) \text { for } j=1,2,3
$$

Since $\beta^{2}=I$, multiplying (4.6) on the left by $\beta$ yields

$$
i \gamma^{\mu} \partial_{\mu} \psi=m \psi
$$

the covariant form of the Dirac equation. The $\gamma^{\mu}$ are easily seen to satisfy the relations

$$
\begin{aligned}
\left\{\gamma^{\mu}, \gamma^{\nu}\right\} \equiv \gamma^{\mu} \gamma^{\nu}+\gamma^{\nu} \gamma^{\mu}=2 g^{\mu \nu} I & (\mu, \nu=0,1,2,3) \\
\left(\gamma^{\mu}\right)^{\dagger}=\left(\gamma^{\mu}\right)^{-1}=\gamma_{\mu}\left(=g_{\mu \nu} \gamma^{\nu}\right) & (\mu=0,1,2,3)
\end{aligned}
$$

where $g^{\mu \nu}=g_{\mu \nu}$ is the Lorentz metric. (Here and in the sequel, we employ the relativistic tensor notation introducted in §1.1.)

We refer to the copy of $\mathbb{C}^{4}$ on which the Dirac matrices act as spinor space, and its elements are called (Dirac) spinors. (On the other hand, Pauli spinors are elements of the copy of $\mathbb{C}^{2}$ on which the Pauli matrices act.) To avoid confusion, it is important to remember that spinor space has nothing to do with space-time. The fact that the dimensions of these two spaces are both 4 is a coincidence.

It is not hard to verify that the 16 matrices $I, \gamma^{\mu}, \gamma^{\mu} \gamma^{\nu}(\mu<\nu), \gamma^{\mu} \gamma^{\nu} \gamma^{\rho}$ ( $\mu<\nu<\rho$ ), and $\gamma^{0} \gamma^{1} \gamma^{2} \gamma^{3}$ are a basis for the space of all $4 \times 4$ complex matrices, so the Clifford algebra generated by the $\gamma^{\mu}$ is the full $4 \times 4$ matrix algebra. By the way, the last of these matrices occurs frequently enough to have its own name; multiplied by $i$ so as to make it real, it is conventionally called $\gamma^{5}$ :

$$
\gamma^{5}=i \gamma^{0} \gamma^{1} \gamma^{2} \gamma^{3}=\left(\begin{array}{ll}
0 & I \\
I & 0
\end{array}\right)
$$

(Obviously this notation was introduced by someone who took the Lorentz index $\mu$ to run from 1 to 4 rather than 0 to 3 . The convention changed but the name stuck.)

There is nothing sacred about the choice (4.7) for the $\gamma^{\mu}$; any four matrices $\widetilde{\gamma}^{\mu}$ satisfying (4.9) and (4.10) would do just as well. One can obtain such sets of matrices by taking $\widetilde{\gamma}^{\mu}=U \gamma^{\mu} U^{-1}$ where $\gamma^{\mu}$ is given by (4.7) and $U$ is any unitary $4 \times 4$ matrix; conversely, every such set $\widetilde{\gamma}^{\mu}$ is obtained in this way. (This follows from the well-known fact that the full matrix algebra has a unique irreducible representation; see Messiah [83], §XX.10, for a direct proof.) The choice (4.7) is called the Dirac representation. The other choice that we shall find useful is the Weyl or chiral representation of the Dirac matrices,

$$
\gamma^{0}=\left(\begin{array}{cc}
0 & I \\
I & 0
\end{array}\right), \quad \gamma^{j}=\left(\begin{array}{cc}
0 & \sigma_{j} \\
-\sigma_{j} & 0
\end{array}\right) \quad(j=1,2,3), \quad \gamma^{5}=\left(\begin{array}{cc}
-I & 0 \\
0 & I
\end{array}\right)
$$

which is related to the Dirac representation by

$$
\gamma_{\mathrm{Weyl}}^{\mu}=U \gamma_{\mathrm{Dirac}}^{\mu} U^{-1}, \quad U=\frac{1}{\sqrt{2}}\left(\begin{array}{cc}
I & -I \\
I & I
\end{array}\right)
$$

Many calculations with Dirac matrices are representation-independent. When a particular representation is needed, for most purposes the Weyl representation is the convenient one, particularly when questions of parity (left- or right-handedness) arise. (We shall explain this at the beginning of §4.3.) In fact, in this book we shall use the Dirac representation only in the discussion of the nonrelativistic approximation and the Coulomb potential in §4.3 and in a brief reference to the latter in §7.11; elsewhere the reader should keep the Weyl representation in mind.

Unlike the Klein-Gordon equation, the Dirac equation yields a conservation-of-probability law; this was one of the main things that persuaded Dirac that his theory was a good one. To explain it, and for other purposes later on, we need to introduce a modified adjoint for Dirac wave functions. We recall that the matrices $\alpha^{j}$ and $\beta$ in (4.6) are Hermitian. Thus $\gamma^{0}=\beta$ is Hermitian, but $\gamma^{j}=\gamma^{0} \alpha^{j} (j=1,2,3)$ is not; rather, $\left(\gamma^{j}\right)^{\dagger}=\gamma^{0} \gamma^{j} \gamma^{0}$, and this requires an extra twist in setting up certain quantities.

Suppose $\psi$ is a solution of the Dirac equation $i \gamma^{\mu} \partial_{\mu} \psi-m \psi=0$; here $\psi$ is a $4 \times 1$ column vector of functions. Taking Hermitian adjoints yields $-i \partial_{\mu} \psi^{\dagger}\left(\gamma^{\mu}\right)^{\dagger}-m \psi^{\dagger}=$ 0 , $\psi^{\dagger}$ being the $1 \times 4$ row vector whose components are the conjugates of the components of $\psi$. To turn this into an equation involving $\gamma^{\mu}$ rather than $\left(\gamma^{\mu}\right)^{\dagger}$, we multiply it on the right by $\gamma^{0}$ and set

$$
\bar{\psi}=\psi^{\dagger} \gamma^{0}
$$

obtaining

$$
-i \partial_{\mu} \bar{\psi} \gamma^{\mu}-m \bar{\psi}=0
$$

the adjoint Dirac equation.\\
$\bar{\psi}=\psi^{\dagger} \gamma^{0}$ is called the (Dirac) adjoint spinor of $\psi$; this notation will be employed throughout the sequel.

Now, if we multiply the original Dirac equation on the left by $\bar{\psi}$ and the adjoint equation on the right by $\psi$ and subtract, we obtain

$$
i \bar{\psi} \gamma^{\mu} \partial_{\mu} \psi+i \partial_{\mu} \bar{\psi} \gamma^{\mu} \psi=0
$$

in other words,

$$
\partial_{\mu} j^{\mu}=0, \quad \text { where } \quad j^{\mu}=\bar{\psi} \gamma^{\mu} \psi
$$

This is the promised conservation law, for

$$
j^{0}=\bar{\psi} \gamma^{0} \psi=\psi^{\dagger} \gamma^{0} \gamma^{0} \psi=\psi^{\dagger} \psi
$$

which is the probability density for the position of the Dirac particle; the spatial components $\left(j^{1}, j^{2}, j^{3}\right)$ therefore define the corresponding current density.

A related fact is that the Dirac equation can be derived from the Lagrangian density

$$
\mathcal{L}=i \bar{\psi} \gamma^{\mu} \partial_{\mu} \psi-m \bar{\psi} \psi
$$

in which $\bar{\psi}$ and $\psi$ are to be regarded as independent dynamical variables. Variation of $\int \mathcal{L} d^{4} x$ with respect to $\bar{\psi}$ gives (4.8), whereas variation with respect to $\psi$ gives the adjoint equation (4.15). This Lagrangian is invariant under the transformations $\psi \mapsto e^{i \theta} \psi$ for $\theta \in \mathbb{R}$, and $j^{0}$ is the corresponding conserved quantity according to Noether's theorem.

This gives the clue for incorporating an (unquantized, external) electromagnetic field into the Dirac equation. As in the classical situation that we discussed in §2.4, the rule is simply to replace the energy-momentum vector $p^{\mu}$ by $p^{\mu}-q A^{\mu}$ where $A^{\mu}$\\
is the electromagnetic potential and $q$ is the charge of the particle. At this point, however, we make a shift in notation to avoid confusion in later sections where the letter $q$ is needed to denote momenta: we denote the charge of the particle by $e$. (The letter $e$ is conventionally used to denote the charge of an electron or its absolute value, and in most applications it will indeed have that meaning since most charged subatomic particles (except quarks) have charge $\pm e$. But for the present discussion, $e$ can be anything.)

As noted at the beginning of this section, the quantized energy-momentum vector $p^{\mu}$ is $i \partial^{\mu}$, so the prescription for including electromagnetism in the Dirac equation is

$$
i \partial_{\mu} \rightarrow i \partial_{\mu}-e A_{\mu}
$$

The Lagrangian density is now

$$
\mathcal{L}=\bar{\psi} \gamma^{\mu}\left(i \partial_{\mu}-e A_{\mu}\right) \psi-m \bar{\psi} \psi
$$

which yields the Dirac equation and adjoint Dirac equation

$$
\begin{aligned}
\gamma^{\mu}\left(i \partial_{\mu}-e A_{\mu}\right) \psi-m \psi & =0 \\
\left(-i \partial_{\mu}-e A_{\mu}\right) \bar{\psi} \gamma^{\mu}-m \bar{\psi} & =0
\end{aligned}
$$

(In taking the adjoint, the $i$ acquires a minus sign, but the $A_{\mu}$ does not, because it is real.) The probability current $j^{\mu}=\bar{\psi} \gamma^{\mu} \psi$ still works the same way, because when one multiplies (4.17) by $\bar{\psi}$ and (4.18) by $\psi$ and subtracts, the $A_{\mu}$ 's cancel out.

The covariant formulation of the Dirac equation may have obscured the initial idea of an evolution equation for vectors in a Hilbert space. Considered simply as a differential operator, the Dirac operator on the left of (4.17) acts on more-or-less arbitrary $\mathbb{C}^{4}$-valued functions on $\mathbb{R}^{4}$. But as an evolution equation for quantum states, its domain should be taken to be the space of differentiable functions on $\mathbb{R}$ (the time variable) with values in $L^{2}\left(\mathbb{R}^{3}, \mathbb{C}^{4}\right)$, and it should be rewritten in the form (4.6):

$$
i \partial_{t} \psi=\left[\sum_{1}^{3} \alpha^{j}\left(-i \partial_{j}-e A_{j}\right)+e A_{0}+m \beta\right] \psi
$$

The operator in brackets is the Dirac Hamiltonian, which acts on $L^{2}\left(\mathbb{R}^{3}, \mathbb{C}^{4}\right)$. Of course, this formulation spoils the relativistic symmetry between space and time. To some extent this is inevitable when one singles out the time variable, but we shall see in §4.4 that it is possible to give a "Lorentz-friendly" description of the situation without losing sight of the Hilbert space.

If our derivation of (4.17)-(4.18) seems a bit breezy, one should keep in mind that quantization rules are only guidelines. The real justification for the Dirac equation (4.17) is the extent to which it works as a description of quantum phenomena. We shall present the fundamental evidence on this score in §4.3. First, though, we shall examine the behavior of the Dirac equation under relativistic changes of reference frame.

\subsection*{Invariance and covariance properties of the Dirac equation}
The Dirac equation was proposed as a relativistically correct evolution equation for quantum particles. If it is to fulfill that requirement, it had better maintain its form under the Poincaré group of space-time symmetries. The key idea is that when\\
applying Lorentz transformations, the 4-tuple of matrices $\left(\gamma^{0}, \gamma^{1}, \gamma^{2}, \gamma^{3}\right)$ should be treated like an ordinary 4-vector. (In particular, we can write $\gamma_{\mu}=g_{\mu \nu} \gamma^{\nu}$; we then have, for example, $\gamma_{\mu} \gamma^{\mu}=4 I$.)

To see how this works, suppose that $L \in O(1,3)$; we identify $L$ with the matrix $\left(L_{\nu}^{\mu}\right)$ so that $(L x)^{\mu}=L_{\nu}^{\mu} x^{\nu}$. If we set

$$
[L \gamma]^{\mu}=L_{\nu}^{\mu} \gamma^{\nu}
$$

the matrices $\widetilde{\gamma}^{\mu}=[L \gamma]^{\mu}$ still satisfy the anticommutation relations (4.9), for

$$
\widetilde{\gamma}^{\mu} \widetilde{\gamma}^{\nu}+\widetilde{\gamma}^{\nu} \widetilde{\gamma}^{\mu}=L_{\rho}^{\mu} L_{\sigma}^{\nu}\left(\gamma^{\rho} \gamma^{\sigma}+\gamma^{\sigma} \gamma^{\rho}\right)=2 L_{\rho}^{\mu} g^{\rho \sigma} L_{\sigma}^{\nu}=2 g^{\mu \nu}
$$

by the defining property of $O(1,3)$. It follows on general grounds that there must be a matrix $B_{L}$ such that $\widetilde{\gamma}^{\mu}=B_{L} \gamma^{\mu} B_{L}^{-1}$. Rather than presenting the general argument, we shall identify $B_{L}$ explicitly.

For this purpose we employ the Weyl representation (4.12) of the Dirac matrices, which enables us to use the double covering map $\kappa: S L(2, \mathbb{C}) \rightarrow S O^{\uparrow}(1,3)$ defined by (1.17) in an efficient way. In fact, if $x \in \mathbb{R}^{4}$, by (4.12) we have

$$
\gamma \cdot x=\left(\begin{array}{cc}
0 & M(x) \\
M(P x) & 0
\end{array}\right),
$$

where $\gamma \cdot x=\sum_{1}^{4} \gamma^{\mu} x^{\mu}$ is the Euclidean scalar product of $\gamma$ and $x, M(x)$ is defined by (1.15), and $P$ is the spatial inversion or parity operator,

$$
P\left(x^{0}, \mathbf{x}\right)=\left(x^{0},-\mathbf{x}\right) .
$$

Now, if

$$
L=\kappa(A) \in S O^{\uparrow}(1,3)
$$

we have $M(L x)=A M(x) A^{\dagger}$. Moreover, it is easily verified that if $M(x)$ is invertible then $M(P x)=(\operatorname{det} M(x)) M(x)^{-1}$, so that

$$
\begin{aligned}
M(P L x)=(\operatorname{det} M(L x)) M(L x)^{-1} & =(\operatorname{det} M(x)) A^{\dagger-1} M(x)^{-1} A^{-1} \\
& =A^{\dagger-1} M(P x) A^{-1}
\end{aligned}
$$

As the invertible matrices are dense in $M\left(\mathbb{R}^{4}\right)$ (the space of Hermitian $2 \times 2$ matrices), we have $M(P L x)=A^{\dagger-1} M(P x) A^{-1}$ in general. Therefore, since $L^{\dagger}= \kappa(A)^{\dagger}=\kappa\left(A^{\dagger}\right)$,

$$
\begin{aligned}
L_{\gamma} \cdot x=\gamma \cdot L^{\dagger} x & =\left(\begin{array}{cc}
0 & M\left(L^{\dagger} x\right) \\
M\left(P L^{\dagger} x\right) & 0
\end{array}\right)=\left(\begin{array}{cc}
0 & A^{\dagger} M(x) A \\
A^{-1} M(P x) A^{\dagger-1} & 0
\end{array}\right) \\
& =\left(\begin{array}{cc}
A^{\dagger} & 0 \\
0 & A^{-1}
\end{array}\right)\left(\begin{array}{cc}
0 & M(x) \\
M(P x) & 0
\end{array}\right)\left(\begin{array}{cc}
A^{\dagger-1} & 0 \\
0 & A
\end{array}\right)
\end{aligned}
$$

In short, if we define the homomorphism $\Phi: S L(2, \mathbb{C}) \rightarrow S L(4, \mathbb{C})$ by

$$
\Phi(A)=\left(\begin{array}{cc}
A^{\dagger-1} & 0 \\
0 & A
\end{array}\right)
$$

we have

$$
L \gamma \cdot x=\Phi\left(A^{-1}\right)(\gamma \cdot x) \Phi(A)
$$

whence, by taking $x$ to be one of the standard basis vectors,

$$
L_{\nu}^{\mu} \gamma^{\nu}=\Phi\left(A^{-1}\right) \gamma^{\mu} \Phi(A) \quad(L=\kappa(A))
$$

At this point it will be convenient to sweep the double covering map under the rug and regard the correspondence $L \mapsto \Phi(A)$ as a double-valued representation of $S O^{\uparrow}(1,3)$. That is, we shall write

$$
\Phi(L)=\Phi\left(\kappa^{-1}(L)\right)
$$

with the understanding that $\Phi(L)$ is defined only up to a factor $\pm 1$ (an ambiguity that will never cause any difficulty as long as the the same sign is used for both terms in (4.22)). Since $\kappa$ preserves adjoints, (4.22) can then be rewritten as

$$
L_{\nu}^{\mu} \gamma^{\nu}=\Phi\left(L^{-1}\right) \gamma^{\mu} \Phi(L)
$$

The representation $\Phi$ can be extended to the full Lorentz group $O(1,3)$. In fact, since the substitution $x \mapsto P x$ simply interchanges the two $M$ 's in (4.20), it is clear that for the parity operator $P$ we can take

$$
\Phi(P)=( \pm) \gamma^{0}=( \pm)\left(\begin{array}{cc}
0 & I \\
I & 0
\end{array}\right)
$$

(A direct verification of (4.23) is also easy: $\left(\gamma^{0}\right)^{3}=\gamma^{0}$ and $\gamma^{0} \gamma^{j} \gamma^{0}=-\gamma^{j}$ for $j>0$.) For the time inversion $T\left(x^{0}, \mathbf{x}\right)=\left(-x^{0}, \mathbf{x}\right)=-P\left(x^{0}, \mathbf{x}\right)$ it is then readily verified that

$$
\Phi(T)=( \pm) \gamma^{1} \gamma^{2} \gamma^{3}=( \pm)\left(\begin{array}{cc}
0 & -I \\
I & 0
\end{array}\right)
$$

does the job. ${ }^{3}$\\
We are now in a position to analyze the Lorentz covariance of the Dirac equation. First suppose that $L$ is an orthochronous Lorentz transformation. The electromagnetic potential $A$ transforms under coordinate changes as a covariant vector field, i.e., a differential 1-form. In other words, if two reference frames are related by the transformation $L$, the components $A_{\mu}$ with respect to the first frame are related to the components $A_{\mu}^{L}$ with respect to the second one by

$$
A_{\mu}^{L}(x)=L_{\mu}^{\nu} A_{\nu}(L x)
$$

Now, if $\psi$ is a $\mathbb{C}^{4}$-valued function on $\mathbb{R}^{4}$, let us set ${ }^{4}$

$$
\psi^{L}(x)=\Phi\left(L^{-1}\right) \psi(L x)
$$

Then since $\partial_{\mu}(f \circ L)=L_{\mu}^{\nu}\left(\partial_{\nu} f\right) \circ L$ (the chain rule), we have

$$
\begin{aligned}
{\left[\gamma ^ { \mu } \left(i \partial_{\mu}\right.\right.} & \left.\left.-e A_{\mu}^{L}(x)\right)-m\right] \psi^{L}(x) \\
& =\gamma^{\mu} L_{\mu}^{\nu} \Phi\left(L^{-1}\right)\left[i\left(\partial_{\nu} \psi\right)(L x)-e A_{\nu}(L x) \psi(L x)\right]-m \Phi\left(L^{-1}\right) \psi(L x) \\
& =\left(\Phi\left(L^{-1}\right) \gamma^{\nu} \Phi(L) \Phi\left(L^{-1}\right)\left[i\left(\partial_{\nu} \psi\right)(L x)-e A_{\nu}(L x) \psi(L x)\right]-m \Phi\left(L^{-1}\right) \psi(L x)\right. \\
& =\Phi\left(L^{-1}\right)\left[\gamma^{\nu}\left(i \partial_{\nu} \psi-e A_{\nu} \psi\right)-m \psi\right](L x)
\end{aligned}
$$

Thus, if $\psi$ satisfies the Dirac equation with potential $A_{\mu}$, then $\psi^{L}$ satisfies the Dirac equation with potential $A_{\mu}^{L}$.

\footnotetext{${ }^{3}$ The restriction of $\Phi$ to the orthochronous Lorentz group comes from a genuine representation of its double cover. However, the ambiguity of sign when one adds in time inversion is unavoidable, because $P$ and $T$ commute whereas $\Phi(P)$ and $\Phi(T)$ anticommute.\\
${ }^{4}$ If one wants the map taking $L$ to the transformation $\psi \mapsto \psi^{L}$ to be a group homomorphism, one should replace $L$ by $L^{-1}$ in (4.24).
}For the time inversion $T x=\left(-x^{0}, \mathbf{x}\right)$ the transformation law for $A_{\mu}$ has an extra minus sign in it:

$$
A_{\mu}^{T}(x)=-T_{\mu}^{\nu} A_{\nu}(T x)
$$

that is,

$$
A_{0}^{T}\left(x^{0}, \mathbf{x}\right)=A_{0}\left(-x_{0}, \mathbf{x}\right) \text { and } A_{j}^{T}\left(x^{0}, \mathbf{x}\right)=-A_{j}\left(-x_{0}, \mathbf{x}\right) .
$$

The reason is that $\nabla^{2} A_{0}$ is the charge density, which transforms as a scalar function under time inversion, but $\nabla^{2} \mathbf{A}$ is the current density, which acquires a minus sign because the motion of the charges is reversed. This necessitates an extra twist in the transformation law for $\psi$ that makes it antiunitary. To wit, since the complex conjugate matrices $\gamma^{\mu *}$ satisfy the same anticommutation relations as $\gamma^{\mu}$ and are still unitary, there is a unitary matrix $B$, determined up to a phase factor, such that

$$
\gamma^{\mu} B=B \gamma^{\mu *}
$$

(In the Dirac and Weyl representations, $\gamma^{0}, \gamma^{1}$, and $\gamma^{3}$ are real while $\gamma^{2}$ is imaginary, and it follows easily that $B=\gamma^{0} \gamma^{1} \gamma^{3}$ works.) If we define

$$
\psi^{T}(x)=\gamma^{0} B \psi(T x)^{*}
$$

we have

$$
\begin{aligned}
& {\left[\gamma^{\mu}\left(i \partial_{\mu}-e A_{\mu}^{T}(x)\right)-m\right] \psi^{T}(x)} \\
& \quad=\gamma^{0} B\left[\gamma^{0 *}\left(i \partial_{0}-e A_{0}(T x)\right)-\sum_{j=1}^{3} \gamma^{j *}\left(i \partial_{j}+e A_{j}(T x)\right)-m\right](\psi \circ T)(x)^{*} \\
& \quad=\gamma^{0} B\left\{\left[\gamma^{0}\left(-i \partial_{0}-e A_{0}(T x)\right)-\sum_{j=1}^{3} \gamma^{j}\left(-i \partial_{j}+e A_{j}(T x)\right)-m\right](\psi \circ T)(x)\right\}^{*} \\
& \quad=\gamma^{0} B\left[\gamma^{\mu}\left(i \partial_{\mu}-e A_{\mu}\right) \psi-m \psi\right]^{*}(T x)
\end{aligned}
$$

(Here's what has happened: $\gamma^{0}$ commutes with itself and anticommutes with the $\gamma^{j}$, and $\gamma^{\mu} B=B \gamma^{\mu *}$, which gives the first equality. Putting the conjugation on the whole expression to the right of $\gamma^{0} B$ changes the $i$ 's to $-i$ 's, and finally $\partial_{0}(\psi \circ T)=-\left(\partial_{0} \psi\right) \circ T$.) In short, if $\psi$ satisfies the Dirac equation with potential $A_{\mu}$, then $\psi^{T}$ satisfies the Dirac equation with potential $A_{\mu}^{T}$.

The time-reversal transform $\psi \mapsto \psi^{T}$ is closely related to another basic antiunitary symmetry of the system, charge conjugation, which is defined by

$$
\psi^{C}(x)=\gamma^{5} B \psi(x)^{*}
$$

where $B$ is as above and $\gamma^{5}$ is defined by (4.11). A calculation much like the preceding one shows that if $\psi$ satisfies the Dirac equation $\gamma^{\mu}\left(i \partial_{\mu}-e A_{\mu}\right) \psi=m \psi$, then $\psi^{C}$ satisfies $\gamma^{\mu}\left(i \partial_{\mu}+e A_{\mu}\right) \psi^{C}=m \psi^{C}$. The point here is that $\gamma^{5}$ anticommutes with all the $\gamma^{\mu}$; combining this with complex conjugation, which changes $i$ to $-i$, the net effect is to change $e$ to $-e$.

The transformation laws $A \mapsto A^{L}, \psi \mapsto \psi^{L}$ under Lorentz transformations are complemented by the simple transformation law $A_{\mu}^{a}(x)=A_{\mu}(x-a), \psi^{a}(x)= \psi(x-a)$ under space-time translations. Thus the Dirac equation is covariant under the full Poincaré group.

There is one other crucial invariance property that must be noted. The electromagnetic potential $A_{\mu}$ is well-defined only modulo the gauge transformations $A_{\mu} \mapsto A_{\mu}-\partial_{\mu} \chi$, where $\chi$ is an arbitrary (smooth) real-valued function. The product rule readily shows that the compensating transformation of the wave function is $\psi \mapsto \exp (i e \chi) \psi$; that is, if $\psi$ and $A_{\mu}$ satisfy $\gamma^{\mu}\left(i \partial_{\mu}-e A_{\mu}\right) \psi=m \psi$, then\\
$\psi^{\prime}=\exp (i e \chi) \psi$ and $A_{\mu}^{\prime}=A_{\mu}-\partial_{\mu} \chi$ satsify $\gamma^{\mu}\left(i \partial_{\mu}-e A_{\mu}^{\prime}\right) \psi^{\prime}=m \psi^{\prime}$. We shall place this result into a broader context in Chapter 9.

There is much more that can be said about the mathematics of the Dirac equation; a comprehensive account can be found in Thaller [118].

\subsection*{Consequences of the Dirac equation}
Spin and parity. Since $A^{\dagger-1}=A$ for $A \in S U(2)$, the restriction of the representation $\Phi$ defined by (4.21) to the rotation group $S O(3)$ is just the direct sum of two copies of the irreducible representation $\pi_{1 / 2}$ described in §3.5. It follows that the solutions of the Dirac equation describe particles of spin $\frac{1}{2}$. For Dirac the fact that spin is built into his equation in this way was an "unexpected bonus."

However, the fact that there are two copies of $\pi_{1 / 2}$ - that is, the shift from $\mathbb{C}^{2}$ valued wave functions in §3.5 to $\mathbb{C}^{4}$-valued wave functions here to describe particles of spin $\frac{1}{2}$-requires further comment. The point is that the change from two to four components does not represent an additional two degrees of freedom. The easiest way to see this is to write out the Dirac equation using the Weyl representation (4.12) of the Dirac matrices: setting

$$
\psi_{L}=\binom{\psi_{1}}{\psi_{2}}, \quad \psi_{R}=\binom{\psi_{3}}{\psi_{4}}
$$

( $L$ and $R$ stand, conventionally, for "left" and "right"), the Dirac equation (4.17) becomes

$$
\begin{aligned}
& m \psi_{L}=\left(i \partial_{t}-e A_{0}\right) \psi_{R}-\sum_{1}^{3}\left(i \partial_{j}-e A_{j}\right) \sigma_{j} \psi_{R} \\
& m \psi_{R}=\left(i \partial_{t}-e A_{0}\right) \psi_{L}+\sum_{1}^{3}\left(i \partial_{j}-e A_{j}\right) \sigma_{j} \psi_{L}
\end{aligned}
$$

Thus, either of the two-component functions $\psi_{L}$ and $\psi_{R}$ is completely determined once the other one is given. The enlargement from two to four components is necessary in order to make room for the Dirac matrices, but the spin variables still represent only two degrees of freedom for the wave functions.

The reason for associating the labels "left" and "right" to $\psi_{L}$ and $\psi_{R}$ is that they are interchanged by the parity operator $P$. More precisely, since $\Phi(P)=\left(\begin{array}{cc}0 & I \\ I & 0\end{array}\right)$, by (4.24) we have $\left(\psi^{P}\right)_{L}(x)=\psi_{R}(P x)$ and $\left(\psi^{P}\right)_{R}(x)=\psi_{L}(P x)$. Note also that the left- and right-handed wave functions are $\mp 1$-eigenvectors of $\gamma^{5}=\left(\begin{array}{cl}-I & 0 \\ 0 & I\end{array}\right)$ : if $\psi_{R}=0$ then $\gamma^{5} \psi=-\psi$, and if $\psi_{L}=0$ then $\gamma^{5} \psi=\psi$. This will be important in §9.4.

The nonrelativistic approximation. The Dirac model for a spin- $\frac{1}{2}$ particle, involving a first-order differential equation for a $\mathbb{C}^{4}$-valued wave function, does not seem to bear much resemblance to the nonrelativistic model of Chapter 3, involving a second-order differential equation for a $\mathbb{C}^{2}$-valued wave function. Nonetheless, the Dirac theory does reproduce the old theory in the low-energy limit. We now give a quick sketch of the way this works.

For this purpose we write the Dirac equation in the form with the time variable singled out:

$$
\left(i \partial_{t}-e A_{0}\right) \psi=\sum_{1}^{3} \alpha^{j}\left(-i \partial_{j}-e A_{j}\right) \psi+m \beta \psi
$$

where $\alpha^{j}$ and $\beta$ are given by (4.5). We also set

$$
\psi_{l}=\binom{\psi_{1}}{\psi_{2}}, \quad \psi_{s}=\binom{\psi_{3}}{\psi_{4}}
$$

where $l$ and $s$ now stand for "large" and "small," for reasons shortly to be explained. (These are not the $\psi_{L}$ and $\psi_{R}$ of (4.27), because there we were employing a basis for $\mathbb{C}^{4}$ that gives the Weyl representation (4.12) of the Dirac matrices, whereas here we are employing a basis that gives the Dirac representation (4.7). By (4.13), the two are related by $\psi_{l}=\left(\psi_{L}+\psi_{R}\right) / \sqrt{2}, \psi_{s}=\left(\psi_{L}-\psi_{R}\right) / \sqrt{2}$.) With the abbreviation

$$
D=\sum_{1}^{3} \sigma_{j}\left(-i \partial_{j}-e A_{j}\right),
$$

the Dirac equation then becomes the pair of coupled equations

$$
i \partial_{t} \psi_{l}=\left(m+e A_{0}\right) \psi_{l}+D \psi_{s}, \quad i \partial_{t} \psi_{s}=\left(-m+e A_{0}\right) \psi_{s}+D \psi_{l}
$$

We shall give a quick and easy account of the nonrelativistic limit of (4.29) under the assumption that the electromagnetic field $A$ is constant in time, so that we are in the usual setting of an equation in the form $i \partial_{t} \psi=H \psi$ where $H$ is a timeindependent operator. We proceed in the informal style of physicists by assuming at first that $\psi$ is an eigenfunction of the Hamiltonian with eigenvalue $E$, so that $i \partial_{t} \psi=E \psi$. (Honest $L^{2}$ wave functions are continuous superpositions of these eigenfunctions; to rework the argument in terms of them would be more laborious but not more enlightening. For a more sophisticated and rigorous mathematical discussion, see Thaller [118].) The second equation in (4.29) then says that

$$
\psi_{s}=\frac{1}{E+m-e A_{0}} D \psi_{l}
$$

We now make the nonrelativistic approximation, namely, that (1) the relativistic energy $E$ consists almost entirely of the rest mass $m$, i.e., $E-m \ll m$, (2) the electromagnetic field is weak so that $\left|e A_{\mu}(\mathbf{x})\right| \ll m$, at least in the region where $\psi(\mathbf{x})$ is nonnegligible, (3) $\|\nabla \psi\|$, the root-mean-square momentum in the state $\psi$, is small in comparison with $m\|\psi\|$. (With the factors of $c$ and $\hbar$ reinserted, these approximations are $E-m c^{2} \ll m c^{2},\left|e A_{\mu}\right| \ll m c^{2}$ and $\hbar\|\nabla \psi\| \leq m c\|\psi\|$.) Then $\left\|D \psi_{l}\right\| \ll m\|\psi\|$, whereas $E+m-e A_{0} \approx 2 m$, so $\psi_{s}$ is small in comparison with $\psi$ :

$$
\left\|\psi_{s}\right\| \ll\|\psi\| .
$$

In the nonrelativistic limit, $\psi_{s}$ becomes negligible, and we are left with the twocomponent wave function $\psi_{l}$. The equation it satisfies is obtained by substituting (4.30) into the first equation in (4.29), again replacing $i \partial_{t}$ by $E$ :

$$
(E-m) \psi_{l}=e A_{0} \psi_{l}+D \frac{1}{E+m-e A_{0}} D \psi_{l}
$$

Again making the approximation $E+m-e A_{0} \approx 2 m$, we obtain

$$
(E-m) \psi_{l}=e A_{0} \psi_{l}+\frac{1}{2 m} D^{2} \psi_{l}
$$

A straightforward calculation shows that

$$
D^{2}=(-i \nabla-e \mathbf{A})^{2}-e \sum_{1}^{3} B_{j} \sigma_{j}=(-i \nabla-e \mathbf{A})^{2}-2 e \mathbf{B} \cdot \mathbf{S} \text {, }
$$

where $\mathbf{B}=\operatorname{curl} \mathbf{A}$ is the magnetic field and $\mathbf{S}=\frac{1}{2} \boldsymbol{\sigma}$ is the spin operator for the 2 -component wave function $\psi_{l}$. Finally, we drop the assumption that $\psi$ is an energy eigenstate and turn $E$ back into $i \partial_{t}$. Setting $\psi_{0}=e^{i m t} \psi_{l}$ to compensate for the shift from $i \partial_{t}$ to $i \partial_{t}-m$ (which does not affect the state represented by $\psi_{l}$ at time $t$ ), we obtain

$$
i \partial_{t} \psi_{0}=e A_{0} \psi_{0}+\frac{1}{2 m}(-i \nabla-e \mathbf{A})^{2} \psi_{0}-\frac{e}{m} \mathbf{B} \cdot \mathbf{S} \psi_{0}
$$

This equation, with the final term omitted, is the classical Schrödinger equation for a particle in the electromagnetic potential $A_{\mu}$, obtained from the free-particle equation by the canonical substitution $p_{\mu} \mapsto p_{\mu}-e A_{\mu}$ without taking spin into account. The final term represents the interaction of the spin with the magnetic field.

Let us consider the special case of a spin- $\frac{1}{2}$ charged particle moving in a constant magnetic field $\mathbf{B}$ (together with, perhaps, an electrostatic field such as a Coulomb field). We can then take $\mathbf{A}=\frac{1}{2} \mathbf{B} \times \mathbf{x}$, which gives, after a little algebraic manipulation,

$$
(-i \nabla-e \mathbf{A})^{2}=-\nabla^{2}-e \mathbf{B} \cdot \mathbf{L}+\frac{e^{2}}{4}|\mathbf{B} \times \mathbf{x}|^{2}
$$

The Hamiltonian for $\psi_{0}$ is therefore

$$
e A_{0}-\frac{1}{2 m} \nabla^{2}-\frac{e}{2 m} \mathbf{B} \cdot(\mathbf{L}+2 \mathbf{S})+\frac{e^{2}}{8 m}|\mathbf{B} \times \mathbf{x}|^{2}
$$

(When the magnetic field is weak, the last term is generally negligible.) The striking feature here is the term $\mathbf{L}+2 \mathbf{S}$ : the spin interacts with the magnetic field twice as strongly as the orbital angular momentum. This phenomenon, for atomic electrons, was one of the outstanding puzzles of atomic physics in the 1920s, and Dirac's explanation of it was one of the triumphs of his theory.

Classically, a magnet in a constant magnetic field $\mathbf{B}$ experiences a force $\nabla(\boldsymbol{\mu} \cdot \mathbf{B})$ where $\boldsymbol{\mu}$ is the magnetic moment of the magnet, and the associated potential energy is $-\boldsymbol{\mu} \cdot \mathbf{B}$. The intrinsic magnetic moment of a spin- $\frac{1}{2}$ particle (i.e., the moment attributable to spin rather than the motion of the particle) has the form

$$
\boldsymbol{\mu}=\frac{g e}{2 m} \mathbf{S}=\frac{g e}{4 m} \boldsymbol{\sigma},
$$

where $g$, the Landé $g$-factor, is determined experimentally and depends on the particle in question. Thus, the preceding calculations show that the Dirac theory predicts $g=2$. This is quite accurate for electrons, but the true value of $g_{\text {electron }}$ is a bit bigger than 2 . The discrepancy is known as the anomalous magnetic moment, and the theoretical calculation of it is one of the triumphs of quantum electrodynamics; we shall discuss it in §7.11.

On the other hand, the $g$-factor of a proton is about 5.59 , which is nowhere near 2 , and the magnetic moment of a neutron, which ought to be 0 since the neutron has no charge, is actually about $-2 / 3$ times that of a proton. These figures, which were a mystery for a long time, are now understood to be reflections of the fact that protons and neutrons are not truly elementary but are made up of quarks.

An approximate calculation of these magnetic moments in the quark model can be found in Griffiths [59], §5.10.

The Coulomb potential revisited. The Dirac equation yields a model for the bound states of electrons in atoms similar to the one we derived from the Schrödinger equation in §3.6, but more refined. To wit, it takes the wave function for an electron in an atom of atomic number $Z$ (i.e., with $Z$ protons in the nucleus) to be a solution of the Dirac equation with $e=$ the charge of the electron, $A_{0}=Z|e| / 4 \pi r$, and $\mathbf{A}=\mathbf{0}$, where $r$ is the distance to the origin. The assumption that the nucleus is infinitely heavy results in errors of the same order of magnitude as the fine-structure effects that the Dirac equation captures, so one should use a reference frame in which the center of mass of the electron-nucleus system is fixed and take the mass $m$ in the Dirac equation to be the reduced mass $m_{e} m_{N} /\left(m_{e}+m_{N}\right)$, as in the classical two-body problem discussed at the end of §2.1. The assumption that the charge of the nucleus is all located precisely at the origin is also an idealization, but it gives a very good approximation for realistic values of $Z$. Since we are now taking $\hbar=c=1$, the constant $e^{2} / 4 \pi$ in the coefficient of the potential $e A_{0}=-Z e^{2} / 4 \pi r$ is just the fine structure constant $\alpha$ (see §1.2):

$$
\alpha=\frac{e^{2}}{4 \pi} \approx \frac{1}{137}
$$

In short, the mathematical problem to be solved is the spectral theory of the Dirac Hamiltonian

$$
H=\frac{1}{i} \sum_{1}^{3} \alpha^{j} \partial_{j}+m \beta-\frac{Z \alpha}{r} \quad\left(Z \in \mathbb{Z}_{+}, \alpha \approx \frac{1}{137}\right)
$$

particularly the eigenfunctions and eigenvalues of $H$ that represent stationary states of the electron. We shall concentrate on the structural features rather than going through all the calculations in detail. Fuller discussions can be found in Thaller [118], Sakurai [103], or Messiah [83].

The first point to be made is that unlike the Schrödinger Hamiltonian with a $-Z \alpha / r$ potential, this $H$ is scale-invariant, because both $\partial_{j}$ and $1 / r$ are homogeneous of degree -1 under dilations, and so is $m$ in the natural units where $[m]=\left[l^{-1}\right]$. Thus, whereas in the Schrödinger theory we could make the coefficient of $1 / r$ equal to 2 by choosing the units appropriately, this is not possible for the Dirac theory, and the size of the coefficient $Z \alpha$ is significant.

The spherical symmetry of the problem points to a consideration of the decomposition of the state space $L^{2}\left(\mathbb{R}^{3}, \mathbb{C}^{4}\right)$ into irreducible pieces under the action of the rotation group. We recall from (4.24) that the rotation $R \in S O(3)$ acts on the state vector $\psi$ by

$$
\psi^{R}(t, \mathbf{x})=\Phi(R) \psi\left(t, R^{-1} \mathbf{x}\right)
$$

This action is the tensor product of the natural representation $\rho(R) \phi(\mathbf{x})=\phi\left(R^{-1} \mathbf{x}\right)$ of $S O(3)$ on $L^{2}\left(\mathbb{R}^{3}\right)$ and the representation $\Phi$ on $\mathbb{C}^{4}$. The Schrödinger Hamiltonian of §3.6 commutes with each of these actions separately, but the Dirac Hamiltonian does not, so we need to take spin into account more carefully here than we did in §3.6.

As we showed in §3.6, $\rho$ decomposes via spherical harmonics as $\rho \cong \bigoplus_{0}^{\infty} \pi_{l}$, where $\pi_{l}$ is as in §3.5. Moreover, by the remarks on spin at the beginning of this section, $\Phi \mid S O(3)$ is the direct sum of two copies of $\pi_{1 / 2}$; but when $\rho \otimes \Phi$ is restricted\\
to solutions of the Dirac equation, where the first two components of $\psi$ in the Weyl representation determine the last two, there is only one copy. Hence, by (3.34), the representation of $S O(3)$ given by (4.34) decomposes as

$$
\bigoplus_{l=0,1,2, \ldots} \pi_{l} \otimes \pi_{1 / 2} \cong \bigoplus_{j=\frac{1}{2}, \frac{3}{2}, \frac{5}{2}, \ldots} 2 \pi_{j}
$$

The 2 indicates multiplicity; one copy of $\pi_{j}$ comes from $l=j-\frac{1}{2}$, the other from $l=j+\frac{1}{2}$.

One can think of this in terms of angular momentum. The angular momentum operators $\mathbf{J}=\left(J^{1}, J^{2}, J^{3}\right)$ are the infinitesimal generators of the action (4.34) of the rotations about the coordinate axes. Corresponding to the two occurrences of $R$ on the right side of (4.34), we have

$$
\mathbf{J}=\mathbf{L}+\mathbf{S}
$$

where $\mathbf{L}$ is the orbital angular momentum and $\mathbf{S}$ is the spin,

$$
\mathbf{S}=\frac{1}{2}\left(\begin{array}{ll}
\boldsymbol{\sigma} & 0 \\
0 & \boldsymbol{\sigma}
\end{array}\right)
$$

The irreducible representation spaces of type $\pi_{j}$ are eigenspaces of the total squared angular momentum $\mathbf{J}^{2}$ with eigenvalue $j(j+1)$, just as in §3.5. The Hamiltonian (4.33) is invariant under the action (4.34) of the rotations and hence commutes with $\mathbf{J}^{2}$, but it does not commute with $\mathbf{L}^{2}$.

We can resolve the multiplicity-two representations $2 \pi_{j}$ by passing to the full orthogonal group $O(3)$, that is, by considering the parity or spatial inversion operator $P(t, \mathbf{x})=(t,-\mathbf{x})$. Its action on the state vectors commutes with the action of $S O(3)$ and with the Dirac Hamitonian, and the space of state vectors on which $2 \pi_{j}$ acts decomposes into the $\pm 1$ eigenspaces of the action of $P$, that is, into the states of even and odd parity. Identifying "even" with 0 and "odd" with 1 , we denote the parity of such an eigenstate by $\epsilon$ :

$$
\psi^{P}=(-1)^{\epsilon} \psi
$$

At this point we need to be a little more concrete. We work in the Dirac representation, where the coefficients $\alpha^{j}$ and $\beta$ in the Hamiltonian (4.33) are given by (4.5), and we write $\psi_{l}$ and $\psi_{s}$ for the first and second pair of components of $\psi$ as in (4.28). We recall from §4.2 that $\Phi(P)=\gamma^{0}=\beta$, which in the Dirac representation is $\left(\begin{array}{cc}I & 0 \\ 0 & -I\end{array}\right)$, so that

$$
\psi^{P}(t, \mathbf{x})=\binom{\psi_{l}(t,-\mathbf{x})}{-\psi_{s}(t,-\mathbf{x})}
$$

Thus if $\psi$ is a state of either even or odd parity, the components $\psi_{l}$ and $\psi_{s}$ must have opposite spatial parity: if $\psi^{P}=\psi$, then $\psi_{l}$ is even and $\psi_{s}$ is odd, and vice versa if $\psi^{P}=-\psi$.

When one works out the implications of these remarks in more detail, one finds that the eigenstates of $\mathbf{J}^{2}$ with eigenvalue $j(j+1)$ and definite parity $\epsilon$ are of the form

$$
\psi=\binom{\psi_{l}}{\psi_{s}}=\frac{1}{r}\binom{f(r) Y_{l}(\theta, \phi)}{i g(r) Y_{s}(\theta, \phi)}
$$

where $Y_{l}$ and $Y_{s}$ are $\mathbb{C}^{2}$-valued functions whose components are spherical harmonics of degree $j \pm \frac{1}{2}$. More precisely, if $j+\frac{1}{2} \equiv \epsilon(\bmod 2)$, then $Y_{l}$ is of degree $j+\frac{1}{2}$ and\\
$Y_{s}$ is of degree $j-\frac{1}{2}$; vice versa if $j-\frac{1}{2} \equiv \epsilon(\bmod 2)$. The requirement that $\psi$ be an eigenfunction of the Hamiltonian $H$ with eigenvalue $E$ is then equivalent to the following pair of differential equations for $f$ and $g$ :

$$
\begin{aligned}
f^{\prime}(r)+\frac{\kappa}{r} f(r) & =\left(E+m+\frac{Z \alpha}{r}\right) g(r), \\
-g^{\prime}(r)+\frac{\kappa}{r} g(r) & =\left(E-m+\frac{Z \alpha}{r}\right) f(r),
\end{aligned}
$$

where

$$
\kappa=(-1)^{\epsilon+j+(1 / 2)}\left(j+\frac{1}{2}\right)
$$

The analysis up to this point is valid with $-Z \alpha / r$ replaced by any central potential $V(r)$.

The equations (4.36) are analytic with a regular singular point at the origin, so they yield readily to standard power series techniques. Beginning with the Ansatz

$$
f(r)=a_{0} r^{\lambda}+\text { higher order }, \quad g(r)=b_{0} r^{\mu}+\text { higher order },
$$

one sees first that $\mu$ must equal $\lambda$. (The left side of the equation involving $g^{\prime}$ is $O\left(r^{\mu-1}\right)$, whereas the right side is $\geq C r^{\lambda-1}$ for $r$ small; hence $\lambda \geq \mu$. Similarly, $\mu \leq \lambda$.) With $\mu=\lambda$, then, examination of the terms of order $\lambda-1$ in (4.36) shows that

$$
(\lambda+\kappa) a_{0}=Z \alpha b_{0}, \quad(-\lambda+\kappa) b_{0}=Z \alpha a_{0}
$$

This pair of equations for $a_{0}$ and $b_{0}$ has nonzero solutions only when $\kappa^{2}-\lambda^{2}= (Z \alpha)^{2}$, that is, when $\lambda= \pm \sqrt{\kappa^{2}-(Z \alpha)^{2}}$.

At this point we must address the square-integrability of the solutions at the origin. The factor of $1 / r$ in (4.35) compensates for the factor of $r^{2}$ in the measure in spherical coordinates, so the $\psi$ in (4.35) will be $L^{2}$ near the origin with respect to volume measure if and only if $f$ and $g$ are in $L^{2}$ near the origin with respect to linear measure; this happens precisely when $\operatorname{Re} \lambda>-\frac{1}{2}$. Thus, if $\kappa^{2}-(Z \alpha)^{2}>\frac{1}{4}$ we take the solution with $\lambda=\sqrt{\kappa^{2}-(Z \alpha)^{2}}$ and discard the one with $\lambda=-\sqrt{\kappa^{2}-(Z \alpha)^{2}}$. If $\kappa^{2}-(Z \alpha)^{2}<\frac{1}{4}$, however, both solutions are $L^{2}$ near the origin, and this is a disaster. It means that the differential operators in (4.36), and hence the Dirac Hamiltonian itself, do not determine a unique self-adjoint operator on $L^{2}$ and so do not have a well-defined spectral theory. One could specify a self-adjoint form by imposing "boundary conditions" at the origin, but their physical significance is unclear.

Fortunately, this is not a serious cause for worry. Recall that $|\kappa|=j+\frac{1}{2}$ is a positive integer, so we are always safe when $(Z \alpha)^{2}<\frac{3}{4}$, i.e., $Z<\sqrt{3} / 2 \alpha \approx 118 .^{5}$ All atomic nuclei with a more-than-ephemeral existence satisfy this condition with room to spare. Even for large $Z$, one avoids the mathematical catastrophe by recognizing that the charge of the nucleus is not really located at a single point. Taking the nucleus to be a solid ball of charge of radius $\sim 10^{-13} \mathrm{~cm}$ gives a more realistic model in which the potential has no singularity at the origin. (However, the resulting model for an atom is still physically bizarre for $Z>\sqrt{3} / 2 \alpha$, as the

\footnotetext{${ }^{5}$ Some physics books say that the critical condition is $(Z / \alpha)^{2}<1$, i.e., $Z<137$, because for $Z / \alpha>1$ the solutions are all bounded (and wildly oscillatory) near the origin. But the real problem occurs earlier.
}
effective diameters of the electron orbits for small $|\kappa|$ are of the same order of magnitude as that of the nucleus.)

We henceforth assume that $Z<118$. Then next step is to solve the equations (4.36), with

$$
\lambda=\sqrt{\kappa^{2}-(Z \alpha)^{2}} .
$$

For this purpose it is convenient to make a change of variables: keeping in mind that for bound states, the total energy $E$ will be less than the rest mass $m$, we set $\rho=r \sqrt{m^{2}-E^{2}}$ and take

$$
f(r)=e^{-\rho} \rho^{\lambda} \sum_{0}^{\infty} a_{n} \rho^{n}, \quad g(r)=e^{-\rho} \rho^{\lambda} \sum_{0}^{\infty} b_{n} \rho^{n} .
$$

Plugging these formulas into (4.36) gives simple recursion formulas for $a_{n}$ and $b_{n}$ that can be solved explicitly; the resulting power series can be expressed in terms of confluent hypergeometric functions. From this one finds that $f$ and $g$ grow exponentially at infinity, and hence are not in $L^{2}$, except in those cases where the power series actually terminate; moreover, the series terminate with the $\rho^{N}$ term $(N=0,1,2, \ldots)$ precisely when $\sqrt{m^{2}-E^{2}}(\lambda+N)=E Z \alpha$. Solving this equation for $E$ gives

$$
E=m\left[1+\left(\frac{Z \alpha}{\lambda+N}\right)^{2}\right]^{-1 / 2}=m\left[1+\frac{(Z \alpha)^{2}}{\left(N+\sqrt{\kappa^{2}-(Z \alpha)^{2}}\right)^{2}}\right]^{-1 / 2}
$$

Recall from (4.37) that $\kappa= \pm\left(j+\frac{1}{2}\right)$. For $N>0$ there is a solution of the differential equation for each of these values of $\kappa$, but for $N=0$ there is a solution only for $\kappa<0$. (In brief, the reason is that when $N=0$, the recursion formula together with the fact that $a_{1}=b_{1}=0$ implies that $a_{0}$ is a negative multiple of $b_{0}$; this is incompatible with the first equation in (4.38) if $\kappa>0$.)

To put this formula into a more recognizable form, we assume that $Z \alpha \ll 1$ and expand the right side in powers of $Z \alpha$. Using the fact that $\kappa^{2}=\left(j+\frac{1}{2}\right)^{2}$ and setting

$$
n=N+j+\frac{1}{2}
$$

after some manipulation of Taylor polynomials we obtain

$$
E=E(n, j)=m\left[1-\frac{(Z \alpha)^{2}}{2 n^{2}}+(Z \alpha)^{4}\left(\frac{3}{8 n^{4}}-\frac{1}{n^{3}(2 j+1)}\right)+O\left((Z \alpha)^{6}\right)\right] .
$$

Here $n=1,2,3, \ldots$ and $j=\frac{1}{2}, \frac{3}{2}, \frac{5}{2}, \ldots$, subject to the relation $j<n$ (since $n-j-\frac{1}{2}=N \geq 0$ ). For each such $n$ and $j$ there are four independent eigenfunctions with eigenvalue $E(n, j)$ except when $n=j+\frac{1}{2}$, in which case there are two: more precisely, there are two for each allowable value of $\kappa$, which can be taken to be the states with spin up or down along some specified axis.

This is now amenable to comparision with the results of the Schrödinger equation. The first term on the right of (4.39) is, of course, just the rest mass. The next term, $-m(Z \alpha)^{2} / 2 n^{2}$, gives the energy levels predicted by the Schrödinger equation, $n$ being the "principal quantum number" that tells which "shell" the electron lives in. (In the discussion in §3.6 we took the unit of energy to be $m(Z \alpha)^{2} / 2$; here this factor appears explicitly.) The $(Z \alpha)^{4}$ term gives Dirac's corrections to the energy levels. The interesting feature here is the dependence on $j$, which is not present in the Schrödinger theory; this is the "fine splitting" of energy levels whose agreement\\
with experiment is one of the major successes of the Dirac theory. In particular, the difference in the energies of the $2 P_{1 / 2}$ and $2 P_{3 / 2}$ states ${ }^{6}$ for a hydrogen atom (i.e., $Z=1, n=2, j=\frac{1}{2}, \frac{3}{2}$ ) is about $m \alpha^{4} / 32 \approx 4.5 \times 10^{-5} \mathrm{eV}$, which is about $10^{-5}$ times smaller than the main term $-m \alpha^{2} / 8 \approx-3.4 \mathrm{eV}$ for their energies.

But as with the magnetic moment of the electron, this is not the end of the story. The Dirac theory gives the same energy to the $2 S_{1 / 2}$ and $2 P_{1 / 2}$ states, but experimentally there is a difference of about $4.4 \times 10^{-6} \mathrm{eV}$ (the "Lamb shift"). Explaining this is another job for quantum electrodynamics; we shall say more about it in §6.2 and §7.11.

There is one respect in which the predictions of the Dirac and Schrödinger equations for the electron in a hydrogen atom are both completely wrong. According to both of them, if the electron is in an energy eigenstate, it will remain there for all time unless some external force comes along to knock it out. But in fact, if the electron is in an eigenstate other than the ground state (an "excited state"), it eventually gets tired of being there and drops down to a lower energy level with the emission of a photon. ${ }^{7}$ The problem is that the Schrödinger and Dirac models as presented here recognize only the Coulomb force generated by the nucleus and not the whole electromagnetic field associated to the nucleus and the electron together. The full analysis of this situation requires quantum field theory; we shall present a simplified model of it in §6.2.

Negative-energy states. A problem with the Dirac equation that was recognized from the beginning is the existence of states of negative energy - that is, the fact that the spectrum of the free Dirac Hamiltonian

$$
H=-i \sum \alpha^{j} \partial_{j}+m \beta
$$

is not bounded below. Indeed, on applying the Fourier transform on $\mathbb{R}^{3}$ and using the Dirac matrices (4.5), we see that

$$
\widehat{H \psi}(\mathbf{p})=(\mathbf{p} \cdot \boldsymbol{\alpha}+m \beta) \widehat{\psi}(\mathbf{p})=\left(\begin{array}{cc}
m I & \mathbf{p} \cdot \boldsymbol{\sigma} \\
\mathbf{p} \cdot \boldsymbol{\sigma} & -m I
\end{array}\right) \widehat{\psi}(\mathbf{p}) .
$$

A straightforward calculation shows that for each $\mathbf{p}$, the $4 \times 4$ Hermitian matrix multiplying $\widehat{\psi}(\mathbf{p})$ has eigenvalues $\pm \sqrt{p^{2}+m^{2}}(p=|\mathbf{p}|)$, each with multiplicity 2 , and that the unitary matrix

$$
U(\mathbf{p})=\frac{\left(m+\sqrt{p^{2}+m^{2}}\right) I+\beta \mathbf{p} \cdot \boldsymbol{\alpha}}{\left[2 m \sqrt{p^{2}+m^{2}}+2\left(p^{2}+m^{2}\right)\right]^{1 / 2}}
$$

diagonalizes it. Let $P_{+}\left(z_{1}, z_{2}, z_{3}, z_{4}\right)=\left(z_{1}, z_{2}\right)$ and $P_{-}\left(z_{1}, z_{2}, z_{3}, z_{4}\right)=\left(z_{3}, z_{4}\right)$; then the transformation $V$ defined by

$$
V \psi(\mathbf{p})=\left(P_{+} U(\mathbf{p}) \widehat{\psi}(\mathbf{p}), P_{-} U(\mathbf{p}) \widehat{\psi}(\mathbf{p})\right)
$$

is a unitary map from $L^{2}\left(\mathbb{R}^{3}, \mathbb{C}^{4}\right)$ to $L^{2}\left(\mathbb{R}^{3}, \mathbb{C}^{2}\right) \oplus L^{2}\left(\mathbb{R}^{3}, \mathbb{C}^{2}\right)$ that intertwines $H$ with the operator $M$ defined by

$$
M\left(f_{+}, f_{-}\right)(\mathbf{p})=\left(\sqrt{p^{2}+m^{2}} f_{+}(\mathbf{p}),-\sqrt{p^{2}+m^{2}} f_{-}(\mathbf{p})\right) .
$$

\footnotetext{${ }^{6}$ Recall that the $P$ in " $2 P_{1 / 2}$ " referred to the orbital angular momentum quantum number in the nonrelativistic theory. Here it is the orbital angular momentum quantum number of the large component $\psi_{l}$.

7 "Eventually" is typically around $10^{-9}$ second.
}Thus, if we denote the inverse images of the two summands in $L^{2}\left(\mathbb{R}^{3}, \mathbb{C}^{2}\right) \oplus L^{2}\left(\mathbb{R}^{3}, \mathbb{C}^{2}\right)$ under $V$ by $\mathcal{H}_{ \pm}$, we see that the restrictions of $H$ to $\mathcal{H}_{+}$and $\mathcal{H}_{-}$ are operators on these spaces with spectra $[m, \infty)$ and $(-\infty,-m]$, respectively.

The existence of states with arbitrarily large negative energy is physically very unsatisfactory. One might try to do away with them by assuming that the physical Hilbert space consists only of the subspace $\mathcal{H}_{+}$, but this assumption is not tenable, for various reasons. For example, $\mathcal{H}_{+}$does not contain states that are highly localized in space, nor does it contain the ground state wave function for the electron in a hydrogen atom that we discussed earlier. Therefore, there must be some mechanism for forbidding a particle to release an infinite amount of energy by falling to lower and lower energy levels.

Dirac's remarkable (not to say fantastic) solution to this problem was to invoke the Pauli exclusion principle and to posit that almost all the negative-energy states are already occupied by a completely invisible "sea" of electrons. The only observable things are the electrons with positive energy and the unoccupied negativeenergy states or "holes" in the negative-energy sea, which are perceived as particles with positive charge. A positive-energy electron could fall into one of these unoccupied states by emitting a pair of photons; this would be perceived as mutual annihilation of the electron and the positively charged particle. Dirac originally proposed that the positively charged particles should be protons, but it quickly became clear that this is incompatible both with the discrepancy in mass between electrons and protons and the complete lack of observational evidence for electronproton annihilation. Dirac was then forced to accept the possible existence of what we now call positrons, and he did so in a paper in 1931. The discovery of the positron by Anderson in 1932 was yet another triumph of the Dirac theory.

The "hole" theory, however, presents problems of its own. Aside from the epistemological difficulties associated to the "negative energy sea," it implies that the Dirac equation cannot really work as it was originally intended, as an evolution equation for a single electron, for the possibility of electron-positron creation or annihilation is always present. The interpretation of antiparticles as holes in a negative energy sea is also untenable for particles of integer spin, to which the exclusion principle does not apply. The idea that negative-energy states have something to do with antiparticles, however, is supported by the fact that the charge conjugation operator $\psi \mapsto \psi^{C}$ that we introduced in (4.26) intertwines $H$ with $-H$ and hence interchanges $\mathcal{H}_{+}$and $\mathcal{H}_{-}$. (The calculation that $H\left(\psi^{C}\right)=-(H \psi)^{C}$ is easy and left to the reader.)

We refer the reader to physics books for a fuller discussion of these matters, which took a fair amount of time for the physicists to sort out. The final upshot, however, was that a much more radical reinterpretation of the Dirac equation is necessary to produce a satisfactory theory, and that a similar reinterpretation of the Klein-Gordon equation gives it a new life as a fundamental theoretical tool. To wit, the Dirac and Klein-Gordon equations must be considered as wave equations not for wave functions of a quantum particle but for classical fields, even though these classical fields do not represent anything in classical physics. One must then quantize these classical fields to produce quantum fields whose quanta are the particles one wishes to study. (This reinterpretation is sometimes given the rather misleading name of "second quantization." There is only one quantization; it just takes place at a different level than one was originally expecting.)

\subsection*{Single-particle state spaces}
So far we have taken the state space for a quantum particle to be $L^{2}\left(\mathbb{R}^{3}, V\right)$ where $V$ is a suitable finite-dimensional vector space. The $\mathbb{R}^{3}$ here has to do with the position of the particle in space; the time variable appears only as a separate parameter. To do things on a more relativistically invariant footing, this picture needs to be revised. More precisely, the state of a particle is determined by various observations made on the particle from a particular inertial reference frame. Two observers in two different reference frames will assign different state vectors to the particle, but the underlying physics must be invariant - at least when the two frames are related by a transformation $T$ in the connected component of the identity $\mathcal{P}_{0}$ of the Poincaré group $\mathcal{P}$. (Space and time inversion may present some problems that need to be considered separately.) The correspondence between the state vectors in the two reference frames must be given by a unitary map $U_{T}$ of the state space, determined up to a phase factor, and clearly $U_{T_{1} T_{2}}$ must be $U_{T_{1}} U_{T_{2}}$ up to a phase factor. (Antiunitary maps are a possibility in principle, but since $\mathcal{P}_{0}$ is connected, continuity considerations force all the $U_{T}$ to be unitary.)

In short, relativistic invariance implies that the state space $\mathcal{H}$ must come equipped with a projective unitary representation of the group $\mathcal{P}_{0}$. Moreover, the theorem of Bargmann [6] quoted in §3.1 guarantees that every projective unitary representation of $\mathcal{P}_{0}$ comes from a genuine unitary representation of its universal cover $\mathbb{R}^{4} \ltimes S L(2, \mathbb{C})$. Finally, we shall assume that this representation is irreducible, which means that the particle is in some sense "elementary." (The particle may, in fact, have some additional structure, as long as that structure is irrelevant for the physical problems under consideration. Thus, for some purposes a proton, or even a whole atomic nucleus, can be considered "elementary." The situation is analogous to the problems in classical celestial mechanics in which the sun and the planets can be assumed to be point masses.)

We are therefore faced with the problem of describing the irreducible unitary representations of the group

$$
G=\mathbb{R}^{4} \ltimes S L(2, \mathbb{C})
$$

This was first accomplished by Wigner [134] and is now usually treated as a classic application of the "Mackey machine." We refer to Folland [45] and Varadarajan [125] for a detailed account of the mathematics and its physical implications; here we just present the results succinctly.

In this discussion we will be dealing with two copies of $\mathbb{R}^{4}$, "position space" and "momentum space," which should be considered to be in duality with each other via the Lorentz inner product $(p, x) \mapsto p_{\mu} x^{\mu}$. Since this notation with indices will become very cumbersome in places, we shall also write the Lorentz inner product as $\langle p, x\rangle$ :

$$
\langle p, x\rangle=p_{\mu} x^{\mu}
$$

The group $S L(2, \mathbb{C})$ acts as Lorentz transformations on position space, and hence also on momentum space by the dual or contragredient action. We denote the action of $A \in S L(2, \mathbb{C})$ on $x \in \mathbb{R}^{4}$ simply by $A x$, and the contragredient action of $A$ on $p \in\left(\mathbb{R}^{4}\right)^{*}$ by $A^{\dagger-1} p$. Thus, strictly speaking, $A x=\kappa(A) x$ where $\kappa: S L(2, \mathbb{C}) \rightarrow S O^{\uparrow}(1,3)$ is the covering map defined by (1.17), and $A^{\dagger-1} p= \kappa(A)^{\dagger-1} p=\kappa\left(A^{\dagger-1}\right) p$.

Given $p \in \mathbb{R}^{4}$, let $H_{p}$ (the "little group" of $p$ ) be the subgroup of $\operatorname{SL}(2, \mathbb{C})$ that fixes $p$, let $G_{p}=\mathbb{R}^{4} \ltimes H_{p}$, and let $\sigma$ be an irreducible unitary representation of $H_{p}$ on a Hilbert space $\mathcal{H}_{\sigma}$. Then one can define an irreducible unitary representation $\rho_{p, \sigma}$ of $G_{p}$ on $\mathcal{H}_{\sigma}$ by

$$
\rho_{p, \sigma}(a, A)=\exp \left(-i p_{\mu} a^{\mu}\right) \sigma(A)
$$

Finally, let $\pi_{p, \sigma}$ be the unitary representation of $G$ induced from $\rho_{p, \sigma}$ (to be described in more detail later):

$$
\pi_{p, \sigma}=\operatorname{ind}_{G_{p}}^{G}\left(\rho_{p, \sigma}\right)
$$

Then $\pi_{p, \sigma}$ is irreducible, and every irreducible unitary representation of $G$ is equivalent to some $\pi_{p, \sigma}$. Moreover, $\pi_{p, \sigma}$ and $\pi_{p^{\prime}, \sigma^{\prime}}$ are equivalent if and only if $p$ and $p^{\prime}$ belong to the same $S L(2, \mathbb{C})$-orbit, say $p^{\prime}=B p$, and $A \mapsto \sigma(A)$ and $A \mapsto \sigma^{\prime}\left(B A B^{-1}\right)$ are equivalent representations of $H_{p}$.

We recall from §1.3 that the $S L(2, \mathbb{C})$-orbits in $\mathbb{R}^{4}$ are as follows:

$$
X_{m}^{+}=\left\{p: p^{2}=m^{2}, p_{0}>0\right\}, \text { and } X_{m}^{-}=\left\{p: p^{2}=m^{2}, p_{0}<0\right\}
$$

for $m \in[0, \infty)$,

$$
Y_{m}=\left\{p: p^{2}=-m^{2}\right\}
$$

for $m \in(0, \infty)$, and

$$
\{0\} .
$$

For $m>0$, we may take the representative point in the orbit $X_{m}^{ \pm}$to be $p_{m}^{ \pm}= ( \pm m, 0,0,0)$ and the representative point in $Y_{m}$ to be $q_{m}=(0,0, m, 0)$; the corresponding little groups are $S U(2)$ (the double cover of $S O(3)$, acting in the spatial variables) and $S L(2, \mathbb{R})$ (the double cover of $S O^{\uparrow}(1,2)$, acting in $p_{0} p_{1} p_{3}$-space). For $m=0$, we may take the representative point in $X_{0}^{ \pm}$to be $p_{0}^{ \pm}=( \pm 1,0,0, \pm 1)$; the little group is

$$
H_{p_{0}^{ \pm}}=\left\{M_{\theta, b}=\left(\begin{array}{cc}
e^{i \theta} & b \\
0 & e^{-i \theta}
\end{array}\right): \theta \in \mathbb{R}, b \in \mathbb{C}\right\}
$$

The reader may verify that this is a double cover of the group of rigid motions (translations and rotations) in the complex plane. Finally, the little group for the orbit $\{0\}$ is $S L(2, \mathbb{C})$ itself.

To proceed further we need to describe $\pi_{p, \sigma}$ more explicitly. In general, if $G$ is a locally compact group and $H$ is a closed subgroup such that $G / H$ has a $G$-invariant measure $\mu$, and $\sigma$ is a unitary representation of $H$ on $\mathcal{H}$, the induced representation $\operatorname{ind}_{H}^{G}(\sigma)$ may be described as follows. For $g \in G$, let $\bar{g}$ be the image of $g$ in $G / H$, and let $\mathcal{F}_{0}$ be the space of continuous $\mathcal{H}$-valued functions $f$ on $G$ that satisfy

$$
f(g h)=\sigma(h)^{-1} f(g) \quad(g \in G, h \in H)
$$

such that $\{\bar{g}: g \in \operatorname{supp} f\}$ is compact in $G / H$. Since $\sigma$ is unitary, for $f \in \mathcal{F}_{0}$ we have $\|f(g h)\|_{\mathcal{H}}=\|f(g)\|_{\mathcal{H}}$, so $\|f(g)\|_{\mathcal{H}}$ depends only on $\bar{g}$. We define $\mathcal{F}$ to be the completion of $\mathcal{F}_{0}$ with respect to the Hilbert norm

$$
\|f\|_{\mathcal{F}}=\left[\int_{G / H}\|f(g)\|_{\mathcal{H}}^{2} d \mu(\bar{g})\right]^{1 / 2}
$$

Then $\operatorname{ind}_{H}^{G}(\sigma)$ is the representation of $G$ on $\mathcal{F}$ by left translation:

$$
\left[\operatorname{ind}_{H}^{G}(\sigma)(g)\right] f\left(g^{\prime}\right)=f\left(g^{-1} g^{\prime}\right)
$$

An alternative, more geometrically appealing description of $\operatorname{ind}_{H}^{G}(\sigma)$ is as the representation of $G$ by left translations on the $L^{2}$ sections of the homogeneous vector bundle $B$ on $G / H$ determined by $\sigma . B$ is the quotient of $G \times \mathcal{H}$ by the equivalence relation

$$
(g, v) \sim\left(g h, \sigma(h)^{-1} v\right) \quad(g \in G, h \in H, v \in \mathcal{H}),
$$

and $f \in \mathcal{F}$ is identified with the section whose value at $\bar{g}$ is the image of ( $g, f(g)$ ) under the equivalence relation.

In our case, the representation $\pi_{p, \sigma}$ of $G=\mathbb{R}^{4} \ltimes S L(2, \mathbb{C})$ acts on $\mathcal{H}_{\sigma}$-valued functions $f$ on $G$ that satisfy

$$
f((a, A)(b, B))=e^{i\langle p, b\rangle} \sigma(B)^{-1} f(a, A) \quad\left(a, b \in \mathbb{R}^{4}, A \in S L(2, \mathbb{C}), B \in H_{p}\right) .
$$

Since

$$
(b, I)(a, A)=(a+b, A)=(a, A)\left(A^{-1} b, I\right)
$$

the action of the space-time translation group is given by

$$
\begin{aligned}
{\left[\pi_{p, \sigma}(b, I) f\right](a, A) } & =f((-b, I)(a, A))=f\left((a, A)\left(-A^{-1} b, I\right)\right)=e^{-i\left\langle p, A^{-1} b\right\rangle} f(a, A) \\
& =e^{-i\left\langle A^{\dagger-1} p, b\right\rangle} f(a, A)
\end{aligned}
$$

Hence, the energy and momentum operators - the infinitesimal generators of the time and space translations - are multiplication by the functions $(a, A) \mapsto \left(A^{\dagger-1} p\right)_{\mu}$. In particular, the operator corresponding to $E^{2}-|\mathbf{p}|^{2}$, which is the square of the rest mass of a particle, is the scalar $m^{2}$ on $X_{m}^{ \pm}$and $-m^{2}$ on $Y_{m}$, and the energy operator is positive on $X_{m}^{+}$and negative on $X_{m}^{-}$. It follows that the representations associated to the orbits $X_{m}^{+}$describe particles of mass $m$. The orbits $X_{m}^{-}$correspond to the negative-energy solutions of the relativistic wave equations whose explanation requires quantum fields and antiparticles as discussed in the following chapter. The orbits $Y_{m}$ would correspond to particles of imaginary mass, and the orbit $\{0\}$ would correspond to particles whose state is invariant under all space-time translations, neither of which exist in reality. ${ }^{8}$





For $m>0$, the little group of the point $p_{m}^{+}=(m, 0,0,0)$ (the energy and momentum in the rest frame of the particle) is $\operatorname{SU}(2)$, and the action of $\operatorname{SU}(2)$, the double cover of $S O(3)$, is what determines the spin of the particle. More precisely, if we take $\sigma=\sigma_{s}$ to be the unique irreducible representation of $\operatorname{SU}(2)$ of dimension $2 s+1\left(s=0, \frac{1}{2}, 1, \ldots\right)$, the representation $\pi_{p_{m}^{+}, \sigma_{s}}$ describes a particle of mass $m$ and spin $s$. To nail this down completely, one should compute $\pi_{p_{m}^{+}, \sigma_{s}}(0, A)$ for $A \in S U(2)$ and show that its infinitesimal generators (the angular momentum operators) are the sum of the appropriate orbital angular momentum and spin operators. For this purpose it is better to give an alternate description of $\pi_{p_{m}^{+}, \sigma_{s}}$ as acting on a space of $\mathcal{H}_{\sigma}$-valued functions on the orbit $X_{m}^{+}$rather than a space of functions on $G$. We shall work this out for $s=0$ and $s=\frac{1}{2}$ below; see Varadarajan [125] or Simms [112] for the general case.

Henceforth we shall simplify the notation by writing

$$
\pi_{m, s}=\pi_{p_{m}^{+}, \sigma_{s}}
$$

In the simplest case $s=0$, where $\sigma_{s}$ is the trivial one-dimensional representation of $\operatorname{SU}(2)$, the space on which $\pi_{m, 0}$ acts consists of functions on $G$ that are

\footnotetext{${ }^{8}$ However, the trivial one-dimensional representation, which is associated to the orbit $\{0\}$, may be said to represent the vacuum.
}
constant on the cosets of $G_{p_{m}^{+}}$and hence can be considered as functions on the orbit $X_{m}^{+} \cong G / G_{p_{m}^{+}}$. With this identification, the Hilbert space on which $\pi_{m, 0}$ acts is $L^{2}\left(X_{m}^{+}, d^{3} \mathbf{p} /(2 \pi)^{3} \omega_{\mathbf{p}}\right)$, where
$$
\omega_{\mathbf{p}}=\sqrt{|\mathbf{p}|^{2}+m^{2}}
$$
and $d^{3} \mathbf{p} /(2 \pi)^{3} \omega_{\mathbf{p}}$ is the invariant measure on $X_{m}^{+}$. (See $\oint 1.3 ;$ we have chosen a different normalization of the invariant measure here.) For each fixed $t$, the inverse Fourier transform
$$
\phi(t, \mathbf{x})=\int \exp \left(-i \omega_{\mathbf{p}} t+\mathbf{p} \cdot \mathbf{x}\right) \frac{f\left(\omega_{\mathbf{p}}, \mathbf{p}\right)}{\sqrt{\omega_{\mathbf{p}}}} \frac{d^{3} \mathbf{p}}{(2 \pi)^{3}}
$$
defines a unitary map $f \mapsto \phi(t, \cdot)$ from $L^{2}\left(X_{m}^{+}, d^{3} \mathbf{p} /(2 \pi)^{3} \omega_{\mathbf{p}}\right)$ onto $L^{2}\left(\mathbb{R}^{3}\right)$. The maps $f\left(\omega_{\mathbf{p}}, \mathbf{p}\right) \mapsto e^{-i \omega_{\mathbf{p}} t} f\left(\omega_{\mathbf{p}}, \mathbf{p}\right)$ form a one-parameter unitary group on the former space whose infinitesimal generator is multiplication by $\omega_{\mathbf{p}}=\sqrt{|\mathbf{p}|^{2}+m^{2}}$, and the inverse Fourier transform intertwines it with the time-translation group $\phi(s, \mathbf{x}) \mapsto \phi(s+t, \mathbf{x})$, whose infinitesimal generator is therefore $\sqrt{-\nabla^{2}+m^{2}}$. Hence $\phi$ satisfies the Klein-Gordon equation, and we have recovered the description of the state space as position-space wave functions. A similar description is available for higher spin.




For $m=0$, the little group of the point $p_{0}^{+}=(1,0,0,1)$ is given by (4.40). It has two sets of irreducible representations: the one-dimensional representations $\tau_{n} (n \in \mathbb{Z})$ given by $\tau_{n}\left(M_{\theta, b}\right)=e^{i n \theta}$, and a family of infinite-dimensional ones induced from the one-dimensional representations of the subgroup $\left\{M_{0, b}: b \in \mathbb{C}\right\}$. The latter do not correspond to any known physical particles, but the representation $\pi_{p_{0}^{+}, \tau_{n}}$ describes a massless particle of spin $\frac{1}{2}|n|$. Again we shall abbreviate: for $s=0, \frac{1}{2}, 1, \frac{3}{2}, \ldots$,

$$
\pi_{0, \pm s}=\pi_{p_{0}^{+}, \tau_{2 s}}
$$

This requires some more explanation. For massive particles, the spin is the particle's angular momentum when the particle is at rest. It is a vector-valued operator, and for any unit vector $\mathbf{u} \in \mathbb{R}^{3}$ one can consider its component in the direction $\mathbf{u}$, i.e., the spin about the $\mathbf{u}$-axis. When one says that the particle has spin $s$, one means that $s$ is the largest eigenvalue of one (and hence all) of these components. Massless particles, however, are never at rest. They always travel with the speed of light, and the only component of spin that makes sense is the one about the axis along which they are traveling. We are therefore dealing with irreducible representations not of $S O(3)$ but of $S O(2)$, and these are parametrized by integers. In particular, the momentum vector $p_{0}^{+}$corresponds to travel along the $x^{3}$-axis, and rotation through an angle $\theta$ about this axis is the image of $M_{\theta / 2,0}$ in $S O(3)$. Thus, in the representation induced from $\tau_{n}$, the $x^{3}$-spin has the value $\frac{1}{2} n$. To distinguish this from the situation pertaining to massive particles, one often says that the particle has helicity $\frac{1}{2} n$.

Now, $n$ can be either positive or negative, so one must ask what is the physical difference beteween $\pi_{0, s}$ and $\pi_{0,-s}$. To understand this, one needs to consider the effect of the parity transformation $(t, \mathbf{x}) \mapsto(t,-\mathbf{x})$, that is, to enlarge the symmetry group to the double cover $\widetilde{G}$ of the orthochronous Poincaré group. The representations $\pi_{m, s}$ with $m>0$ all extend to representations of this larger group, but $\pi_{0, \pm s}$ does not. Rather, the direct sum $\pi_{0, s} \oplus \pi_{0,-s}$ extends to a representation of $\widetilde{G}$ in which the parity operator interchanges the two summands. Since parity is a symmetry of electromagnetism, gravity, and quantum chromodynamics, it is\\
this direct sum that properly describes the associated massless particles: photons, gluons, and gravitons, which have spins 1,1 , and 2 , respectively. However, parity is emphatically not a symmetry of the weak interaction, so in theories in which neutrinos are treated as massless, the two representations $\pi_{0, \pm 1 / 2}$ describe two different species of particle: $\pi_{0,-1 / 2}$ gives neutrinos and $\pi_{0,1 / 2}$ gives antineutrinos. ${ }^{9}$






How does the Dirac model for massive spin- $\frac{1}{2}$ particles, which uses functions with values in $\mathbb{C}^{4}$ rather than $\mathbb{C}^{2}$, fit into this picture? In terms of the representation $\pi_{m, 1 / 2}$, the answer is quite simple. The direct sum of two copies of $\pi_{m, 1 / 2}$ acts on certain functions on $G$ with values in $\mathbb{C}^{4}=\mathbb{C}^{2} \times \mathbb{C}^{2}$, and the representation $\pi_{D}$ that corresponds to the positive-energy solutions of the Dirac equation is the subrepresentation of $\pi_{m, 1 / 2} \oplus \pi_{m, 1 / 2}$ on the subspace of functions $f=\left(f_{1}, f_{2}, f_{3}, f_{4}\right)$ such that $\left(f_{1}, f_{2}\right)=\left(f_{3}, f_{4}\right)$. Here the little group is $\operatorname{SU}(2)$, and $\sigma_{1 / 2}$ is simply the identity representation of $\operatorname{SU}(2)$ on $\mathbb{C}^{2}$.

To see how this works, recall the representation $\Phi$ of $S L(2, \mathbb{C})$ defined by (4.21). Since $A^{\dagger-1}=A$ for $A \in S U(2)$, we can state the covariance condition (4.41) for the representation $\pi_{D}$ in terms of $\Phi$ as $f((a, A)(b, B))=e^{i\left\langle p_{m}^{+}, b\right\rangle} \Phi\left(B^{-1}\right) f(a, A)$. That is, the Hilbert space for $\pi_{D}$ is

$$
\begin{aligned}
\mathcal{H}_{\pi_{D}}= & \left\{f: G \rightarrow \mathbb{C}^{4}: f((a, A)(b, B))=e^{i\left\langle p_{m}^{+}, b\right\rangle} \Phi\left(B^{-1}\right) f(a, A) \text { for } B \in S U(2),\right. \\
& \left.\left(f_{1}, f_{2}\right)=\left(f_{3}, f_{4}\right), \quad \int_{X_{m}^{+}}|f(a, A)|^{2} \frac{d^{3} \mathbf{p}}{(2 \pi)^{3} \omega_{\mathbf{p}}}<\infty \quad\left(\left(\omega_{\mathbf{p}}, \mathbf{p}\right)=A^{\dagger-1} p_{m}^{+}\right)\right\} .
\end{aligned}
$$

The connection with the Dirac equation is as follows. Let $V$ be the subbundle of the product bundle $X_{m}^{+} \times \mathbb{C}^{4}$ over $X_{m}^{+}$whose fiber at $p$ is $V_{p}=\left\{v: p_{\mu} \gamma^{\mu} v=m v\right\}$, where the $\gamma^{\mu}$ are the Weyl form (4.12) of the Dirac matrices. Thus, the sections of $V$ are the $\mathbb{C}^{4}$-valued functions $f$ on $X_{m}^{+}$that satisfy the Fourier-transformed Dirac equation

$$
p_{\mu} \gamma^{\mu} f(p)=m f(p)
$$

We define a Hilbert norm on the fiber $V_{p}$ by

$$
\|v\|_{p}^{2}=\frac{m|v|^{2}}{p_{0}},
$$

where $|v|$ is the Euclidean norm on $\mathbb{C}^{4}$. The crucial point is that if $A \in S L(2, \mathbb{C})$, the operator $\Phi(A)$ is an isometry from $V_{p}$ to $V_{A^{\dagger-1} p}$. Indeed, if $v \in V_{p}$, by (4.22) - recalling that $A \gamma$ is really $\kappa(A) \gamma$ - we have

$$
\left(A^{\dagger-1} p\right)_{\mu} \gamma^{\mu} \Phi(A) v=p_{\mu}\left[\kappa(A)^{-1}\right]_{\nu}^{\mu} \gamma^{\nu} \Phi(A) v=p_{\mu} \Phi(A) \gamma^{\mu} v=m \Phi(A) v,
$$

so $\Phi(A) v \in V_{A^{\dagger-1} p}$. Moreover, since $\gamma^{j}$ is skew-adjoint for $j>0$, for $v \in V_{p}$ we have $m|v|^{2}=p_{\mu}\left\langle\gamma^{\mu} v \mid v\right\rangle=p_{0}\left\langle\gamma^{0} v \mid v\right\rangle$, so $\|v\|_{p}^{2}=\left\langle\gamma^{0} v \mid v\right\rangle$. (Here $\langle\cdot \mid \cdot\rangle$ is the Euclidean inner product on $\mathbb{C}^{4}$. The condition $v \in V_{p}$ is essential for the positivity of $\left\langle\gamma^{0} v \mid v\right\rangle$ !) The fact that $\Phi(A)$ is an isometry now follows since

$$
\Phi(A)^{\dagger} \gamma^{0} \Phi(A)=\left(\begin{array}{cc}
A^{-1} & 0 \\
0 & A^{\dagger}
\end{array}\right)\left(\begin{array}{cc}
0 & I \\
I & 0
\end{array}\right)\left(\begin{array}{cc}
A^{\dagger-1} & 0 \\
0 & A
\end{array}\right)=\left(\begin{array}{cc}
0 & I \\
I & 0
\end{array}\right)=\gamma^{0} .
$$

\footnotetext{${ }^{9}$ Photons, gravitons, and some gluons are their own antiparticles; other gluons are distinguished from their antiparticles by their color charge.
}We now transfer the representation $\pi_{D}$ to a space of sections of $V$. Let $\mathcal{H}_{\Pi}$ be the Hilbert space of sections $f$ of $V$ that are square-integrable with respect to the pointwise norm $\|\cdot\|_{p}$ and the invariant measure $d^{3} \mathbf{p} /(2 \pi)^{3} \omega_{\mathbf{p}}$ on $X_{m}^{+}$, that is,

$$
\int\left\|f\left(\omega_{\mathbf{p}}, \mathbf{p}\right)\right\|_{p}^{2} \frac{d^{3} \mathbf{p}}{(2 \pi)^{3} \omega_{\mathbf{p}}}=\int \frac{m\left|f\left(\omega_{\mathbf{p}}, \mathbf{p}\right)\right|^{2}}{\omega_{\mathbf{p}}^{2}} \frac{d^{3} \mathbf{p}}{(2 \pi)^{3}}<\infty
$$

In view of the results of the preceding paragraph, the following formula defines a unitary representation $\Pi$ of $\mathbb{R}^{4} \ltimes S L(2, \mathbb{C})$ on $\mathcal{H}_{\Pi}$ :

$$
\Pi(b, B) f(p)=e^{-i\langle p, b\rangle} \Phi(B) f\left(B^{\dagger} p\right)
$$

We claim that the map $T$ defined by

$$
T f(a, A)=\exp \left[i\left\langle p_{m}^{+}, A^{-1} a\right\rangle\right] \Phi\left(A^{-1}\right) f\left(A^{\dagger-1} p_{m}^{+}\right)
$$

gives a unitary equivalence of $\Pi$ and $\pi_{D}$.





To prove this, first observe that $T f$ has the right $S U(2)$-covariance property: if $B \in S U(2)$,
$$
\begin{aligned}
T f((a, A)(b, B)) & =T f(a+A b, A B) \\
& =\exp \left[i\left\langle p_{m}^{+},(A B)^{-1}(a+A b)\right\rangle\right] \Phi\left((A B)^{-1}\right) f\left((A B)^{\dagger-1} p_{m}^{+}\right) \\
& =\exp \left[i\left\langle B^{-1 \dagger} p_{m}^{+}, A^{-1} a+b\right\rangle\right] \Phi\left(B^{-1}\right) \Phi\left(A^{-1}\right) f\left(A^{-1 \dagger} B^{-1 \dagger} p_{m}^{+}\right) \\
& =\exp \left[i\left\langle p_{m}^{+}, b\right\rangle\right] \Phi\left(B^{-1}\right) T f(a, A)
\end{aligned}
$$

since $B^{-1 \dagger} p_{m}^{+}=p_{m}^{+}$. Next, observe that since $p_{m}^{+}=(m, 0,0,0)$ and $\gamma^{0}=\left(\begin{array}{ll}0 & I \\ I & 0\end{array}\right)$, the fiber $V_{p_{m}^{+}}$is just $\left\{v: \gamma^{0} v=v\right\}=\left\{v:\left(v_{1}, v_{2}\right)=\left(v_{3}, v_{4}\right)\right\}$, and $\|v\|_{p_{m}^{+}}=|v|$. Since $\Phi\left(A^{-1}\right)$ is an isometry from $V_{A^{\dagger-1} p_{m}^{+}}$to $V_{p_{m}^{+}}$, it follows that $T$ is an isometry from $\mathcal{H}_{\Pi}$ into $\mathcal{H}_{\pi_{D}}$. Since $\pi_{D}$ is irreducible, it remains only to check that $T$ is an intertwining operator from $\Pi$ to $\pi_{D}$ :

$$
\begin{aligned}
T[\Pi(b, B) f](a, A) & =\exp \left[i\left\langle p_{m}^{+}, A^{-1} a\right\rangle S\left(A^{-1}\right)[\Pi(b, B) f]\left(A^{\dagger-1} p_{m}^{+}\right)\right. \\
& =\exp \left[i\left\langle p_{m}^{+}, A^{-1} a\right\rangle-i\left\langle A^{\dagger-1} p_{m}^{+}, b\right\rangle\right] \Phi\left(A^{-1}\right) \Phi(B) f\left(B^{\dagger} A^{\dagger-1} p_{m}^{+}\right) \\
& =\exp \left[i\left\langle p_{m}^{+}, A^{-1}(a-b)\right\rangle\right] \Phi\left(\left(B^{-1} A\right)^{-1}\right) f\left(\left(B^{-1} A\right)^{\dagger-1} p_{m}^{+}\right) \\
& =T f\left(B^{-1}(a-b), B^{-1} A\right)=T f\left((b, B)^{-1}(a, A)\right) \\
& =\left[\pi_{D}(b, B) T f\right](a, A)
\end{aligned}
$$

The Hilbert space $\mathcal{H}_{\Pi}$ consists of momentum-space wave functions, and one can return to the position-space wave functions via the Fourier transform. That is, if $f \in \mathcal{H}_{\Pi}$, let

$$
\psi(t, \mathbf{x})=\int \exp \left(-i \omega_{\mathbf{p}} t+\mathbf{p} \cdot \mathbf{x}\right) \frac{\sqrt{m} f\left(\omega_{\mathbf{p}}, \mathbf{p}\right)}{\omega_{\mathbf{p}}} \frac{d^{3} \mathbf{p}}{(2 \pi)^{3}}
$$

The map $f \mapsto \psi(t, \cdot)$ is an isometry from $\mathcal{H}_{\Pi}$ into $L^{2}\left(\mathbb{R}^{3}, \mathbb{C}^{4}\right)$ for each $t$, and the fact that $p_{\mu} \gamma^{\mu} f(p)=m f(p)$ implies that $\psi$ satisfies the free Dirac equation $i \gamma^{\mu} \partial_{\mu} \psi=m \psi$. The space of $\psi$ 's thus obtained is the space of positive-energy solutions to the Dirac equation; that is, for each $t, \psi(t, \cdot)$ belongs to the subspace of $L^{2}\left(\mathbb{R}^{3}, \mathbb{C}^{4}\right)$ on which the Dirac Hamiltonian is a positive operator (called $\mathcal{H}_{+}$in §4.3).

One obtains the negative-energy solutions by playing the same game with the representation $\pi_{m, 1 / 2}^{-}=\pi_{p_{m}^{-}, \sigma_{1 / 2}}$ and the orbit $X_{m}^{-}$. Since the equation $p_{\mu} \gamma^{\mu} v=m v$ becomes $-\gamma^{0} v=v$ when $p=p_{m}^{-}$, the analogue of $\pi_{D}$ here is the subrepresentation\\
of $\pi_{m, 1 / 2}^{-} \oplus \pi_{m, 1 / 2}^{-}$on the space of functions $f$ such that $\left(f_{3}, f_{4}\right)=-\left(f_{1}, f_{2}\right)$. The resolution of the negative-energy paradox by a reinterpretation of the Dirac equation in terms of fields will be explained in the next chapter.

\subsection*{Multiparticle state spaces}
Having discussed state spaces for single particles, we now turn to the question of constructing state spaces for systems of particles. As we mentioned in §3.1, if we have $k$ particles with state spaces $\mathcal{H}_{1}, \ldots, \mathcal{H}_{k}$, we can describe states of the $k$-particle system using the tensor product $\mathcal{H}_{1} \otimes \cdots \otimes \mathcal{H}_{k}$. This is the completion of the algebraic tensor product of the $\mathcal{H}_{j}$ 's with respect to the inner product

$$
\left\langle u_{1} \otimes u_{2} \cdots \otimes u_{k}, v_{1} \otimes v_{2} \cdots \otimes v_{k}\right\rangle=\left\langle u_{1} \mid v_{1}\right\rangle\left\langle u_{2} \mid v_{2}\right\rangle \cdots\left\langle u_{k} \mid v_{k}\right\rangle .
$$

To put it another way, if $\left\{e_{n}^{j}\right\}_{n=1}^{\infty}$ is an orthonormal basis for $\mathscr{H}_{j}$, then

$$
\left\{e_{n_{1}}^{1} \otimes \cdots \otimes e_{n_{k}}^{k}: 1 \leq n_{1}, \ldots, n_{k}<\infty\right\}
$$

is an orthonormal basis for $\mathcal{H}_{1} \otimes \cdots \otimes \mathcal{H}_{k}$. In this setting, $u_{1} \otimes \cdots \otimes u_{k}$ describes the compound state in which the $j$ th particle is in state $u_{j}$ for each $j$. Superpositions of states that are not reducible to single products also occur; for example, if $u_{ \pm}$ are states in which particle A has spin "up" or "down" along some axis, and $v_{ \pm}$ similarly for particle B , then $\left(u_{+} \otimes v_{+}+u_{-} \otimes v_{-}\right) / \sqrt{2}$ is a state in which the spins of the particles are aligned but the direction is unspecified. On the macroscopic level, the famous Schrödinger cat paradox [68] involves states of the form

$$
(\text { undecayed atom }) \otimes(\text { live cat })+(\text { decayed atom }) \otimes(\text { dead cat }) .
$$

If we are describing a system of $k$ particles of the same species $(k$ electrons, for example), we can take the Hilbert spaces $\mathcal{H}_{j}$ all to be the same space $\mathcal{H}$. However, in this case the putative $k$-particle state space $\bigotimes^{k} \mathcal{H}$ is too big. It is a fundamental fact that particles of the same species are truly indistinguishable; there is no way to attach labels to two electrons so as to tell which one is which. Mathematically, the permutation group on $k$ letters, $S_{k}$, acts on $\bigotimes^{k} \mathcal{H}$ in a canonical way,

$$
\sigma\left(u_{1} \otimes \cdots \otimes u_{k}\right)=u_{\sigma(1)} \otimes \cdots \otimes u_{\sigma(k)}
$$

and the only unit vectors that can represent sets of $k$ identical particles are those that are mapped into scalar multiples of themselves by the permutation group: $\sigma(\mathbf{v})=R(\sigma) \mathbf{v}$, where $R(\sigma)$ is a scalar of absolute value 1 . Clearly $R$ is a onedimensional representation of $S_{k}$, of which there are only two, the trivial representation and the representation $R(\sigma)=\operatorname{sgn} \sigma$, which correspond to symmetric and antisymmetric tensors respectively. In short, we are led to consider the subspaces

$$
\begin{aligned}
\left(\mathbb{S}^{k} \mathcal{H}\right. & =\left\{\mathbf{v} \in \bigotimes^{k} \mathcal{H}: \sigma(\mathbf{v})=\mathbf{v}, \forall \sigma \in S_{k}\right\} \\
\bigwedge^{k} \mathcal{H} & =\left\{\mathbf{v} \in \bigotimes^{k} \mathcal{H}: \sigma(\mathbf{v})=(\operatorname{sgn} \sigma) \mathbf{v}, \forall \sigma \in S_{k}\right\} .
\end{aligned}
$$

Which of these is appropriate depends on the species of particle in question. Particles whose multiparticle states are symmetric are called Bosons; those whose multiparticle states are antisymmetric are called Fermions. It turns out that Bosons are precisely those particles whose spin is an integer; this is the spin-statistics theorem, and we shall say more about it in §5.3 and §6.5.

Let us take a closer look at the Boson space (S) ${ }^{k} \mathcal{H}$. There is a canonical projection $P_{s}$ ( $s$ for symmetric) from $\bigotimes^{k} \mathcal{H}$ onto $\mathrm{S}^{k} \mathcal{H}$ :

$$
P_{s}\left(u_{1} \otimes \cdots \otimes u_{k}\right)=\frac{1}{k!} \sum_{\sigma \in S_{k}} u_{\sigma(1)} \otimes \cdots \otimes u_{\sigma(k)} .
$$

To analyze $P_{s}$ it is convenient to consider an orthonormal basis $\left\{e_{j}\right\}$ for $\mathcal{H}$. Suppose $u_{1}, \ldots, u_{k}$ are all elements of this basis, say $u_{j}=e_{i_{1}}$ for $n_{1}$ values of $j, \ldots, u_{j}=e_{i_{m}}$ for $n_{m}$ values of $j$, where $i_{1}, \ldots, i_{m}$ are distinct and $n_{1}+\cdots+n_{m}=k$, and let $\mathbf{u}=u_{1} \otimes \cdots \otimes u_{k}$. The sum defining $P_{s}(\mathbf{u})$ breaks up into groups of identical terms, corresponding to permutations that only permute the factors of $e_{i}$ among themselves. Each group has $n_{1}!\cdots n_{m}!$ terms, there are $k!/ n_{1}!\cdots n_{m}!$ groups, and terms in different groups are orthogonal to each other. Hence the sum is a sum of $k!/ n_{1}!\cdots n_{m}!$ distinct terms, all orthogonal to each other, each of norm $n_{1}!\cdots n_{m}!$, and so

$$
\left\|P_{s}(\mathbf{u})\right\|=\frac{n_{1}!\cdots n_{m}!}{k!} \sqrt{\frac{k!}{n_{1}!\cdots n_{m}!}}=\sqrt{\frac{n_{1}!\cdots n_{m}!}{k!}} .
$$

Moreover, if $v_{1}, \ldots, v_{k}$ are also elements of $\left\{e_{j}\right\}$ and $\mathbf{v}=v_{1} \otimes \cdots \otimes v_{k}$, we have

$$
\left\langle P_{s}(\mathbf{u}) \mid \mathbf{v}\right\rangle=\left\{\begin{array}{ll}
n_{1}!\cdots n_{m}!/ k! & \text { if } \mathbf{v}=\sigma(\mathbf{u}) \text { for some } \sigma \in S_{k} \\
0 & \text { otherwise }
\end{array}\right\}=\left\langle\mathbf{u} \mid P_{s}(\mathbf{v})\right\rangle
$$

Thus $P_{s}$ is an orthogonal projection.\\
For any sequence $\left\{n_{j}\right\}$ of nonnegative integers with $\sum_{1}^{\infty} n_{j}=k$, let

$$
\left|n_{1}, n_{2}, \ldots\right\rangle=\sqrt{\frac{k!}{n_{1}!n_{2}!\cdots}} P_{s}\left(\left(\otimes^{n_{1}} e_{1}\right) \otimes\left(\otimes^{n_{2}} e_{2}\right) \otimes \cdots\right),
$$

where $\otimes^{n} e_{j}=e_{j} \otimes \cdots \otimes e_{j}$ ( $n$ factors). (Here we are using the convenient Dirac convention of labeling a vector simply by the indices that specify it.) Then

$$
\left\{\left|n_{1}, n_{2}, \ldots\right\rangle: \sum n_{j}=k\right\}
$$

is an orthonormal basis for $\left(\mathrm{S}^{k} \mathcal{H}\right.$.

\pagebreak

The Boson Fock space. 

In order to form a state space that can accomodate any number of identical Bosons, we put all the spaces $\mathrm{S}^{k} \mathcal{H}$ together to form the complete symmetric tensor algebra over $\mathcal{H}$, also known as the Boson Fock space over $\mathcal{H}$ :

$$
\mathcal{F}_{s}(\mathcal{H})=\bigoplus_{k=0}^{\infty} \mathrm{S}^{k} \mathcal{H}
$$

Here (S) ${ }^{0} \mathcal{H}=\mathbb{C}$, and the object on the right is the orthogonal direct sum of Hilbert spaces. We combine the orthonormal bases of the preceding paragraph for the spaces $\left(\mathrm{S}^{k} \mathcal{H}\right.$, together with

$$
|0,0, \ldots\rangle=1 \in \mathrm{~S}^{0} \mathcal{H}
$$

to obtain the orthonormal basis

$$
\left\{\left|n_{1}, n_{2}, \ldots\right\rangle: n_{j} \geq 0, \sum n_{j}<\infty\right\}
$$

for $\mathcal{F}_{s}(\mathcal{H}) . \mathcal{F}_{s}(\mathcal{H})$ is a state space that describes arbitrary numbers of identical Bosons, including none at all. The basis vector $\left|n_{1}, n_{2}, \ldots\right\rangle$ specifies the state in which there are $n_{j}$ particles in state $e_{j}$ for each $j$.

There are no states in $\mathcal{F}_{s}(\mathcal{H})$ that contain infinitely many particles, but there are states in which there is a positive probability of finding arbitrarily large numbers of particles, namely, infinite linear combinations of basis vectors $\left|n_{1}, n_{2}, \ldots\right\rangle$ with $\sum n_{j}$ arbitrarily large. It is convenient also to consider the finite-particle subspace of states in which the total number of particles is bounded above, that is,

$$
\mathcal{F}_{s}^{0}(\mathcal{H})=\text { the algebraic direct sum of the spaces }\left(\mathrm{S}^{k} \mathcal{H} .\right.
$$

On this space we define the number operator $N$ by

$$
N\left(\sum a_{n_{1} n_{2} \ldots}\left|n_{1}, n_{2}, \ldots\right\rangle\right)=\sum\left(n_{1}+n_{2}+\cdots\right) a_{n_{1} n_{2} \ldots}\left|n_{1}, n_{2}, \ldots\right\rangle .
$$

If $\mathbf{u} \in\left(S^{k} \mathcal{H}\right.$ we have $N \mathbf{u}=k \mathbf{u}$, so $N$ is the observable that tells how many particles are in a given state. ( $N$ has a unique extension to an unbounded selfadjoint operator on $\mathcal{F}_{s}(\mathcal{H})$.)

We now introduce some operators that raise or lower the number of particles. For this purpose we retreat temporarily to the full tensor algebra $\mathcal{F}(\mathcal{H})= \oplus_{0}^{\infty} \otimes^{k} \mathcal{H}$ and its dense subspace $\mathcal{F}^{0}(\mathcal{H})$, the algebraic direct sum of the tensor spaces $\bigotimes^{k} \mathcal{H}$. The number operator $N$ is defined on $\mathcal{F}^{0}(\mathcal{H})$ just as above:

$$
N=k I \text { on } \bigotimes^{k} \mathcal{H} .
$$

For $v \in \mathcal{H}$, the operators $B(v)$ and $B(v)^{\dagger}$ on $\mathcal{F}^{0}(\mathcal{H})$ are defined by

$$
\begin{aligned}
B(v)\left(u_{1} \otimes \cdots \otimes u_{k}\right) & =\left\langle v \mid u_{1}\right\rangle u_{2} \otimes \cdots \otimes u_{k} \\
B(v)^{\dagger}\left(u_{1} \otimes \cdots \otimes u_{k}\right) & =v \otimes u_{1} \otimes \cdots \otimes u_{k}
\end{aligned}
$$

(When $k=0, B(v) 1=0$ and $B(v)^{\dagger} 1=v$, and when $k=1, B(v) u=\langle v, u\rangle 1$, where 1 is the unit element in $\mathbb{C}=\bigotimes^{0} \mathcal{H}$.) It is easily verified that $\langle B(v) \mathbf{u} \mid \mathbf{w}\rangle=\left\langle\mathbf{u} \mid B(v)^{\dagger} \mathbf{w}\right\rangle$ for $\mathbf{u}, \mathbf{w} \in \mathcal{F}^{0}(\mathcal{H})$, so the notation is consistent. Note that $B(v)$ depends antilinearly on $v$, whereas $B(v)^{\dagger}$ depends linearly on $v$.

The operator $B(v)^{\dagger}$ does not preserve the symmetric subspace $\mathcal{F}_{s}^{0}(\mathcal{H})$, as the factor $v$ is inserted only on the left, but $B(v)$ does preserve this subspace:

$$
\begin{aligned}
B(v) P_{s}\left(u_{1} \otimes \cdots \otimes u_{k}\right) & =\frac{1}{k!} \sum_{j=1}^{k}\left\langle v, u_{j}\right\rangle \sum_{\sigma(1)=j} u_{\sigma(2)} \otimes \cdots \otimes u_{\sigma(k)} \\
& =\frac{1}{k} \sum_{j=1}^{k}\left\langle v, u_{j}\right\rangle P_{s}\left(u_{1} \otimes \cdots \widehat{u}_{j} \cdots \otimes u_{k}\right),
\end{aligned}
$$

where the hat indicates that the term is omitted. If we take the $u_{j}$ and $v$ to be elements of an orthonormal basis $\left\{e_{l}\right\}$ for $\mathcal{H}$, say $v=e_{l}$ and $u_{j}=e_{l}$ for $n$ values of $j$, the sum on the right consists of $n$ identical nonzero terms and $k-n$ zero terms. Taking into account the normalizations in (4.45), we see that the action of $B\left(e_{j}\right)$ on the basis $\left\{\left|n_{1}, n_{2}, \ldots\right\rangle\right\}$ for $\mathcal{F}_{s}^{0}(\mathcal{H})$ is

$$
B\left(e_{j}\right)\left|n_{1}, \ldots, n_{j}, \ldots\right\rangle=\sqrt{\frac{n_{j}}{k}}\left|n_{1}, \ldots, n_{j}-1, \ldots\right\rangle \quad\left(k=n_{1}+n_{2}+\cdots\right)
$$

For reasons that will become clearer shortly, it is convenient to get rid of the factor of $1 / \sqrt{k}$ in this formula by defining the operator $A(v)$ on $\mathcal{F}_{s}^{0}(\mathcal{H})$ as

$$
A(v)=B(v) \sqrt{N}=\sqrt{N+I} B(v) \quad \text { on } \mathcal{F}_{s}^{0}(\mathcal{H})
$$

that is,

$$
A(v) \mathbf{w}=\sqrt{k} B(v) \mathbf{w} \text { for } \mathbf{w} \in\left(S^{k} \mathcal{H}\right.
$$

so that

$$
A\left(e_{j}\right)\left|n_{1}, \ldots, n_{j}, \ldots\right\rangle=\sqrt{n_{j}}\left|n_{1}, \ldots, n_{j}-1, \ldots\right\rangle .
$$

The adjoint of $A(v)$ as an operator on $\mathcal{F}_{s}^{0}(\mathcal{H})$ (rather than $\mathcal{F}^{0}(\mathcal{H})$ ) is given by

$$
A(v)^{\dagger}=P_{s} B(v)^{\dagger} \sqrt{N+I}=P_{s} \sqrt{N} B(v)^{\dagger} \quad \text { on } \mathcal{F}_{s}^{0}(\mathcal{H}) .
$$

(An extra factor $P_{s}$ could be inserted on the right of the formula defining $A(v)$; it has no effect since the domain of $A(v)$ is contained in the range of $P_{s}$.) Informally speaking, $A(v)^{\dagger}$ creates a particle in the state $v$, whereas $A(v)$ destroys a particle in the state $v$ and annihilates any multiparticle state in which no particle has any probability of being in the state $v$. For this reason, $A(v)^{\dagger}$ and $A(v)$ are called creation and annihilation operators.

The most important consequence of introducing the factor of $\sqrt{N}$ into $A(v)$ is that the operators $A(v)$ satisfy the following variant of the canonical commutation relations:

$$
\left[A(v), A(w)^{\dagger}\right]=\langle v \mid w\rangle I, \quad[A(v), A(w)]=\left[A(v)^{\dagger}, A(w)^{\dagger}\right]=0
$$

for any $v, w \in \mathcal{H}$. To see this, we compute the action of $A(v) A(w)^{\dagger}$ and $A(w)^{\dagger} A(v)$ on $\mathbf{u}=P_{s}\left(u_{1} \otimes \cdots \otimes u_{k}\right) \in\left(\mathrm{S}^{k} \mathcal{H}\right.$ :

$$
\begin{aligned}
A(v) A(w)^{\dagger} \mathbf{u}= & (k+1) B(v) P_{s} B(w)^{\dagger} P_{s}\left(u_{1} \otimes \cdots \otimes u_{k}\right) \\
= & (k+1) B(v) P_{s}\left[\frac{1}{k!} \sum_{\sigma \in S_{k}} w \otimes u_{\sigma(1)} \otimes \cdots \otimes u_{\sigma(k)}\right] \\
= & \frac{1}{k!} B(v) \sum_{j=0}^{k} \sum_{\sigma \in S_{k}} u_{\sigma(1)} \otimes \cdots \otimes w \otimes \cdots \otimes u_{\sigma(k)} \quad(w \text { in } j \text { th place }) \\
= & \frac{1}{k!} \sum_{\sigma \in S_{k}}\langle v \mid w\rangle u_{\sigma(1)} \otimes \cdots \otimes u_{\sigma(k)} \\
& +\frac{1}{k!} \sum_{j=1}^{k} \sum_{\sigma \in S_{k}}\left\langle v \mid u_{\sigma(1)}\right\rangle u_{\sigma(2)} \otimes \cdots \otimes w \otimes \cdots \otimes u_{\sigma(k)} \\
= & \langle v, w\rangle P_{s}\left(u_{1} \otimes \cdots \otimes u_{k}\right)+\sum_{j=1}^{k}\left\langle v \mid u_{j}\right\rangle P_{s}\left(w \otimes u_{1} \otimes \cdots \widehat{u}_{j} \cdots \otimes u_{k}\right),
\end{aligned}
$$

where $\widehat{u}_{j}$ means that $u_{j}$ is omitted. On the other hand,

$$
\begin{aligned}
A(w)^{\dagger} A(v) \mathbf{u} & =k P_{s} B(w)^{\dagger} B(v) P_{s}\left(u_{1} \otimes \cdots \otimes u_{k}\right) \\
& =\frac{k}{k!} P_{s} B(w)^{\dagger} \sum_{\sigma \in S_{k}}\left\langle v, u_{\sigma(1)}\right\rangle u_{\sigma(2)} \otimes \cdots \otimes u_{\sigma(k)} \\
& =P_{s} B(w)^{\dagger} \sum_{j=1}^{k}\left\langle v \mid u_{j}\right\rangle P_{s}\left(u_{1} \otimes \cdots \widehat{u}_{j} \cdots \otimes u_{k}\right) \\
& =\sum_{j=1}^{k}\left\langle v \mid u_{j}\right\rangle P_{s}\left(w \otimes u_{1} \otimes \cdots \widehat{u}_{j} \cdots \otimes u_{k}\right) .
\end{aligned}
$$

Thus $A(v) A(w)^{\dagger} \mathbf{u}-A(w)^{\dagger} A(v) \mathbf{u}=\langle v \mid w\rangle \mathbf{u}$.

This proves the first equality in (4.48). Moreover, an easy calculation similar to the preceding ones shows that the restrictions of $B(v)$ and $B(w)$ to $\mathcal{F}_{s}^{0}$ commute, and hence

$$
A(v) A(w)=\sqrt{N+I} B(v) B(w) \sqrt{N}=\sqrt{N+I} B(w) B(v) \sqrt{N}=A(w) A(v) .
$$

By taking adjoints it follows that $A(v)^{\dagger} A(w)^{\dagger}=A(w)^{\dagger} A(v)^{\dagger}$, so (4.48) is established.

Rather than considering $A(v)$ for arbitrary $v \in \mathcal{H}$, it is often convenient to take $v$ to belong to some useful orthonormal basis. Thus, given an orthonormal basis $\left\{e_{j}\right\}$ for $\mathcal{H}$, we set

$$
A_{j}=A\left(e_{j}\right), \quad A_{j}^{\dagger}=A\left(e_{j}\right)^{\dagger} .
$$

The action of these operators on the basis (4.45) is simple:

$$
\begin{aligned}
A_{j}\left|n_{1}, \ldots, n_{j}, \ldots\right\rangle & =\sqrt{n_{j}}\left|n_{1}, \ldots, n_{j}-1, \ldots\right\rangle, \\
A_{j}^{\dagger}\left|n_{1}, \ldots, n_{j}, \ldots\right\rangle & =\sqrt{n_{j}+1}\left|n_{1}, \ldots, n_{j}+1, \ldots\right\rangle .
\end{aligned}
$$

It follows that

$$
A_{j}^{\dagger} A_{j}\left|n_{1}, \ldots, n_{j}, \ldots\right\rangle=n_{j}\left|n_{1}, \ldots, n_{j}, \ldots\right\rangle,
$$

and hence

$$
\sum A_{j}^{\dagger} A_{j}=N .
$$

Moreover, by (4.48),

$$
\left[A_{j}, A_{k}^{\dagger}\right]=\delta_{j k} I, \quad\left[A_{j}, A_{k}\right]=\left[A_{j}^{\dagger}, A_{k}^{\dagger}\right]=0 .
$$

There is one other important construction on Fock space that needs to be discussed here. Suppose $U$ is a unitary operator on $\mathcal{H}$. Then $U$ induces a unitary operator on all the tensor powers of $\mathcal{H}$ and hence a unitary operator $\mathcal{F}(U)$ on the full tensor algebra $\mathcal{F}(\mathcal{H})=\bigoplus_{0}^{\infty} \bigotimes^{k} \mathcal{H}$, by the formula

$$
\mathcal{F}(U)\left(u_{1} \otimes \cdots \otimes u_{k}\right)=U u_{1} \otimes \cdots \otimes U u_{k} .
$$

$\mathcal{F}(U)$ clearly commutes with the number operator $N$ and with the projections onto the symmetric and antisymmetric subspaces. Moreover, for any $v \in \mathcal{H}$,

$$
\begin{aligned}
\mathcal{F}(U) B(v) \mathcal{F}(U)^{-1}\left(u_{1} \otimes \cdots \otimes u_{k}\right) & =\mathcal{F}(U)\left\langle v \mid U^{-1} u_{1}\right\rangle U^{-1} u_{2} \otimes \cdots \otimes U^{-1} u_{k} \\
& =\left\langle U v \mid u_{1}\right\rangle u_{2} \otimes \cdots \otimes u_{k}=B(U v)\left(u_{1} \otimes \cdots \otimes u_{k}\right),
\end{aligned}
$$

so that $\mathcal{F}(U) B(v) \mathcal{F}(U)^{-1}=B(U v)$. It follows that $\mathcal{F}(U) B(v)^{\dagger} \mathcal{F}(U)^{-1}=B(U v)^{\dagger}$ also. Combining these facts, we obtain

$$
\mathcal{F}(U) A(v) \mathcal{F}(U)^{-1}=A(U v), \quad \mathcal{F}(U) A(v)^{\dagger} \mathcal{F}(U)^{-1}=A(U v)^{\dagger} .
$$


\pagebreak


The Fermion Fock space. 

Just as for Bosons, we can define a state space for arbitrary numbers of identical Fermions: the Fermion Fock space

$$
\mathcal{F}_{a}(\mathcal{H})=\bigwedge \mathcal{H}=\bigoplus_{0}^{\infty} \bigwedge^{k} \mathcal{H} \quad\left(\bigwedge^{0} \mathcal{H}=\mathbb{C}\right) .
$$

As before, the direct sum on the right is an orthogonal sum of Hilbert spaces, and the corresponding algebraic direct sum is denoted by $\mathcal{F}_{a}^{0}$.

All of the preceding notions have analogues here. First, we have the natural orthogonal projection $P_{a}: \bigotimes^{k} \mathcal{H} \rightarrow \bigwedge^{k} \mathcal{H}$ :

$$
P_{a}\left(u_{1} \otimes \cdots \otimes u_{k}\right)=u_{1} \wedge \cdots \wedge u_{k}=\frac{1}{k!} \sum_{\sigma \in S_{k}}(\operatorname{sgn} \sigma) u_{\sigma(1)} \otimes \cdots \otimes u_{\sigma(k)} .
$$

The number operator $N$ is again defined by $N=k I$ on $\bigwedge^{k} \mathcal{H}$, and for $v \in \mathcal{H}$, the annihilation and creation operators $A(v)$ and $A(v)^{*}$ are defined on $\mathcal{F}_{a}^{0}(\mathcal{H})$ by

$$
A(v)=B(v) \sqrt{N}, \quad A(v)^{\dagger}=P_{a} \sqrt{N} B(v)^{\dagger}
$$

The exterior product notation leads to simple formulas for $A(v)$ and $A(v)^{*}$ :

$$
\begin{aligned}
A(v)\left(u_{1} \wedge \cdots \wedge u_{k}\right) & =\frac{1}{\sqrt{k}} \sum_{j=1}^{k}(-1)^{j}\left\langle v \mid u_{j}\right\rangle u_{1} \wedge \cdots \widehat{u}_{j} \cdots \wedge u_{k} \\
A(v)^{\dagger}\left(u_{1} \wedge \cdots \wedge u_{k}\right) & =\sqrt{k+1} v \wedge u_{1} \wedge \cdots \wedge u_{k}
\end{aligned}
$$

$A(v)$ destroys a particle in the state $v$ and gives 0 if there is no such particle; $A(v)^{\dagger}$ creates a particle in the state $v$ but gives 0 if there is already a particle in that state.

The big difference here is that the analogue of (4.48) involves anticommutation relations. To wit,

$$
\left\{A(v), A(w)^{\dagger}\right\}=\langle v, w\rangle I, \quad\{A(v), A(w)\}=\left\{A(v)^{\dagger}, A(w)^{\dagger}\right\}=0
$$

for any $v, w \in \mathcal{H}$, where

$$
\{A, B\}=A B+B A
$$

To see this, observe that since $v \wedge w+w \wedge v=0$, we have $\left\{A(v)^{\dagger}, A(w)^{\dagger}\right\}=0$ and hence also $\{A(v), A(w)\}=0$. Next,

$$
\begin{aligned}
A(v) A(w)^{\dagger}\left(u_{1}\right. & \left.\wedge \cdots \wedge u_{k}\right)=\sqrt{k+1} A(v) w \wedge u_{1} \wedge \cdots \wedge u_{k} \\
& =\langle v, w\rangle u_{1} \wedge \cdots \wedge u_{k}+\sum_{1}^{k}(-1)^{j}\left\langle v, u_{j}\right\rangle w \wedge u_{1} \wedge \cdots \widehat{u}_{j} \cdots \wedge u_{k}
\end{aligned}
$$

whereas

$$
A(w)^{\dagger} A(v)\left(u_{1} \wedge \cdots \wedge u_{k}\right)=\sum_{1}^{k}(-1)^{j-1}\left\langle v \mid u_{j}\right\rangle w \wedge u_{1} \wedge \cdots \widehat{u}_{j} \cdots \wedge u_{k}
$$

(The $\sqrt{k}$ 's or $\sqrt{k+1}$ 's cancel out in both cases.) Hence $\left\{A(v), A(w)^{\dagger}\right\}=\langle v, w\rangle I$.\\
By the way, in marked contrast to the Boson case, the operators $A(v)$ and $A(v)^{\dagger}$ extend to bounded operators on the whole Fermion Fock space. Indeed, for any $\mathbf{w} \in \mathcal{F}_{a}^{0}(\mathcal{H})$, (4.56) implies that

$$
\|A(v) \mathbf{w}\|^{2}+\left\|A(v)^{\dagger} \mathbf{w}\right\|^{2}=\left\langle\mathbf{w} \mid A(v)^{\dagger} A(v) \mathbf{w}\right\rangle+\left\langle\mathbf{w} \mid A(v) A(v)^{\dagger} \mathbf{w}\right\rangle=\|v\|^{2}\|\mathbf{w}\|^{2}
$$

so that

$$
\|A(v)\| \leq\|v\|, \quad\left\|A(v)^{\dagger}\right\| \leq\|v\|
$$

The operators $\mathcal{F}(U)$ defined by (4.53) can be considered as unitary operators on $\mathcal{F}_{a}(\mathcal{H})$, and the formulas (4.54) remain valid in this setting (with the same proof).

As in the Boson case, it is often convenient to express things in terms of an orthonormal basis for $\mathcal{H}$. First, if $\left\{e_{j}\right\}$ is such a basis,

$$
\left\{e_{i_{1}} \wedge \cdots \wedge e_{i_{k}}: i_{1}<\cdots<i_{k}\right\}
$$

is an orthogonal basis for $\bigwedge^{k} \mathcal{H}$. Moreover, since $e_{i_{1}} \wedge \cdots \wedge e_{i_{k}}$ is $1 / k!$ times a sum of $k!$ orthogonal terms $(\operatorname{sgn} \sigma) e_{i_{\sigma(1)}} \otimes \cdots \otimes e_{i_{\sigma(k)}}$, each of which has norm 1 , we have $\left\|e_{i_{1}} \wedge \cdots \wedge e_{i_{k}}\right\|=1 / \sqrt{k!}$. The corresponding orthonormal basis for $\mathcal{F}_{a}(\mathcal{H})$ is therefore

$$
\left\{\left|n_{1}, n_{2}, \ldots\right\rangle: n_{j}=0 \text { or } 1, \sum n_{j}<\infty\right\}
$$

given by

$$
\left|n_{1}, n_{2}, \ldots\right\rangle=\sqrt{k!} e_{i_{1}} \wedge \cdots \wedge e_{i_{k}}
$$

where $i_{1}<\cdots<i_{k}$ are the indices $i$ for which $n_{i}=1$. The fact that each $n_{j}$ must be 0 or 1 is a precise statement of the Pauli exclusion principle, which says that "no two identical Fermions can occupy the same state."

As before, we set

$$
A_{j}=A\left(e_{j}\right), \quad A_{j}^{\dagger}=A\left(e_{j}\right)^{\dagger}
$$

By (4.56), we then have

$$
\left\{A_{j}, A_{k}^{\dagger}\right\}=\delta_{j k} I, \quad\left\{A_{j}, A_{k}\right\}=\left\{A_{j}^{\dagger}, A_{k}^{\dagger}\right\}=0
$$

It is also easily verified that

$$
\begin{aligned}
& A_{j}\left|n_{1}, \ldots, n_{j}, \ldots\right\rangle= \begin{cases}(-1)^{m}\left|n_{1}, \ldots, n_{j}-1, \ldots\right\rangle & \text { if } n_{j}=1, \\
0 & \text { if } n_{j}=0,\end{cases} \\
& A_{j}^{\dagger}\left|n_{1}, \ldots, n_{j}, \ldots\right\rangle= \begin{cases}(-1)^{m}\left|n_{1}, \ldots, n_{j}+1, \ldots\right\rangle & \text { if } n_{j}=0, \\
0 & \text { if } n_{j}=1,\end{cases}
\end{aligned}
$$

where $m$ is the integer such that $n_{i}=1$ for $m$ values of $i$ with $i<j$. It follows that

$$
A_{j}^{\dagger} A_{j}\left|n_{1}, n_{2}, \ldots\right\rangle=n_{j}\left|n_{1}, n_{2}, \ldots\right\rangle,
$$

and hence, as in (4.51),

$$
\sum A_{j}^{\dagger} A_{j}=N
$$


\pagebreak


Finally, we consider the state spaces for arbitrary numbers of particles of several different species. More precisely, we consider systems of free particles. The state space for a rigorous model of an interacting system might turn out to be rather different. However, our approach to interactions will be to apply perturbation theory to noninteracting systems, so we still need the free-particle state space to get started.

There is not much new to be said here. If one has $K$ species of particles, one takes an appropriate single-particle state space $\mathcal{H}_{k}$ for each species, $k=1, \ldots, K$, and forms the appropriate (Boson or Fermion) Fock space $\mathcal{F}_{k}$ over $\mathcal{H}_{k}$ for each $k$. Then the state space for the whole system is

$$
\mathcal{F}=\mathcal{F}_{1} \otimes \cdots \otimes \mathcal{F}_{K} .
$$

For each $k$ one has annihilation and creation operators $A_{k}(v), A_{k}^{\dagger}(v)\left(v \in \mathcal{H}_{k}\right)$ on $\mathcal{F}_{k}$, and these are taken to act on $\mathcal{F}$ by the identification

$$
A_{k}(v) \leftrightarrow I \otimes \cdots \otimes I \otimes A_{k}(v) \otimes I \otimes \cdots \otimes I .
$$

With these definitions, creation and/or annihilation operators pertaining to two different species of particle always commute with each other. However, it simplifies some things to postulate that such operators pertaining to Fermions should anticommute, and this convention is commonly adopted in the physics literature. This can be achieved by a small modification in the definitions of the Fermionic annihilation and creation operators. To wit, it suffices to list the Fermion Fock spaces in some definite order, say $\mathcal{F}_{1}, \ldots, \mathcal{F}_{J}$, and then to replace the operators $A_{j}(v)$ on $\mathcal{F}_{j}$ by $A_{j}(v)(-1)^{N_{1}+\cdots+N_{j-1}}$ where $N_{i}$ is the number operator on $\mathcal{F}_{i}$, and likewise for $A_{j}(v)^{\dagger}$. That is, the operators $A_{j}(v)$ and $A_{j}^{\dagger}(v)$ acquire an extra minus sign when applied to states where there is an odd number of particles of the previously listed Fermion species. This does not affect the states defined by these operators in a significant way, as $\mathbf{v}$ and $-\mathbf{v}$ define the same physical state; for the same reason, the choice of ordering is immaterial. Only the (anti)commutation relations among the operators are important.



\pagebreak


A footnote for fans of category theory. There are some abstract structural aspects of the constructions in this section that may be worth pointing out. To begin with, let $\mathbf{H}$ be the category whose objects are Hilbert spaces and whose morphisms are (linear) contractions, i.e., linear maps of norm $\leq 1$. (It might seem more natural to take the morphisms to be the linear isometries, and one could do so. But if one wants the set of morphisms to be closed under adjoints, and orthogonal direct sums to be products in the sense of category theory, one needs to include partial isometries; and every contraction can be approximated in norm by compositions of partial isometries. Either way, the isomorphisms are the unitary maps.) To each $\mathcal{H} \in \mathbf{H}$ one can associate the full Fock space $\mathcal{F}(\mathcal{H})=\bigoplus_{0}^{\infty} \bigotimes^{k} \mathcal{H}$ as well as the Boson and Fermion Fock spaces $\mathcal{F}_{s}(\mathcal{H})$ and $\mathcal{F}_{a}(\mathcal{H})$, and to each morphism $A: \mathcal{H}_{1} \rightarrow \mathcal{H}_{2}$ one can associate the morphism $\mathcal{F}(A): \mathcal{F}\left(\mathcal{H}_{1}\right) \rightarrow \mathcal{F}\left(\mathcal{H}_{2}\right)$, defined just as in (4.53), as well as its restrictions to $\mathcal{F}_{s}\left(\mathcal{H}_{1}\right)$ and $\mathcal{F}_{a}\left(\mathcal{H}_{1}\right)$. (It is obvious that $\mathcal{F}(A)$ is a contraction whenever $A$ is. Note, however, that if one tries to define $\mathcal{F}(A)$ on $\mathcal{F}\left(\mathcal{H}_{1}\right)$ by (4.53) when $\|A\|>1$, the result is generally an unbounded operator on $\mathcal{F}\left(\mathcal{H}_{1}\right)$.) In short, $\mathcal{F}, \mathcal{F}_{s}$, and $\mathcal{F}_{a}$ are functors from $\mathbf{H}$ to itself. They are, in fact, "exponential functors" in the sense that they convert direct sums into tensor products: $\mathcal{F}\left(\mathcal{H}_{1} \oplus \mathcal{H}_{2}\right) \cong \mathcal{F}\left(\mathcal{H}_{1}\right) \otimes \mathcal{F}\left(\mathcal{H}_{2}\right)$, and likewise for $\mathcal{F}_{s}$ and $\mathcal{F}_{a}$. (We leave the verification of this as an exercise for the interested reader.) See Nelson [87] for a related functorial construction that is relevant to quantum field theory.



\pagebreak

\section*{CHAPTER 5}
\section*{Free Quantum Fields}
\begin{displayquote}
He had bought a large map representing the sea, Without the least vestige of land:\\
And the crew were much pleased when they found it to be A map they could all understand.\\
-Lewis Carroll, The Hunting of the Snark (Fit the Second)
\end{displayquote}

In this chapter we take our first step into quantum field theory by discussing free fields, i.e., fields without interactions. The good thing about free fields is that they can be constructed with complete mathematical rigor, although that task requires more sophistication than one might imagine at the outset; the bad thing is that they don't exhibit any interesting physics. But the time spent on them is not wasted, because they are an essential ingredient in the interacting field theories that we shall consider in the next chapter.

\subsection*{Scalar fields}
We begin by constructing the free quantum scalar fields of mass $m>0$. What this means is the following: we think of solutions $\phi$ of the Klein-Gordon equation $\left(\partial^{2}+m^{2}\right) \phi=0$ as representing a classical field, and we are going to construct the corresponding quantum field. The quanta of these fields, as we shall see, are particles of spin zero and mass $m$.

There is a little conceptual barrier here in that there is no actual classical physical field that is described by the Klein-Gordon equation, and spin-zero elementary particles are also rather scarce. ${ }^{1}$ (There are some; the pions are the least exotic examples, but since the charged pions have a mean lifetime of about $10^{-8}$ seconds and the neutral ones have a mean lifetime of about $10^{-16}$ seconds, they aren't often seen by anyone but particle physicists.) The reader is best advised to consider what we are doing as constructing a toy model for the electromagnetic field. The free electromagnetic potential (not interacting with any charged matter) in Lorentz gauge satisfies the wave equation $\partial^{2} A_{\mu}=0$, which is the Klein-Gordon equation but (i) with $m=0$ and (ii) for a vector-valued rather than scalar-valued function. These conditions correspond to the fact that photons have mass 0 and spin 1. The vector nature of the field necessitates some additional algebra that does not seriously affect the underlying ideas but makes their execution a little messier, and the condition $m=0$ complicates the situation considerably. The scalar field

\footnotetext{${ }^{1}$ The ideas of quantum field theory also have interesting applications to condensed matter physics, however, and certain excitations of crystal lattices ("phonons") can be modeled by spinzero particles. We refer the reader to Mattuck [82] and Zee [138] for more information about this subject.
}
with positive mass exhibits all the essential ideas and difficulties of quantum fields in general, but in the simplest possible context.

We start by deriving the quantum field as physicists (going back to Dirac) usually do, paying more attention to intuitive clarity than mathematical rigor; then we shall show how to perform the construction rigorously. The point of following the physicists' path is twofold: first, it shows how to arrive at the quantum field starting from more familiar objects and makes clear what are the technical difficulties along the way; second, it leads to the standard notations that are used throughout the physics literature and in this book. In any case, the rigorous construction is easier to understand if one already knows what one is trying to construct.

We shall first construct a neutral scalar field, which corresponds to a real-valued classical Klein-Gordon field, and whose quanta are their own antiparticles (such as neutral pions). We shall then modify the construction to obtain a "charged" scalar field, which corresponds to a complex-valued classical Klein-Gordon field, and whose quanta have distinct antiparticles (although they need not possess electric charge).

To make the technical details a bit easier to begin with, we consider the KleinGordon equation not in all of $\mathbb{R}^{4}$ but in a box. That is, we fix a bounded region $\mathcal{B} \subset \mathbb{R}^{3}$ and consider the differential equation $\left(\partial_{t}^{2}-\nabla^{2}+m^{2}\right) \phi(t, \mathbf{x})=0$ for $t \in \mathbb{R}$ and $\mathbf{x} \in \mathcal{B}$, subject to real boundary conditions that make $-\nabla^{2}$ positive and selfadjoint. (We shall make more specific choices of $\mathcal{B}$ and the boundary conditions later.) The boundedness of $\mathcal{B}$ guarantees that $-\nabla^{2}$ has only a discrete spectrum, that is, there is an orthonormal eigenbasis $\left\{f_{j}\right\}$ for $-\nabla^{2}$ with eigenvalues $\left\{\lambda_{j}^{2}\right\}$. Since the differential equation and boundary conditions are real, the eigenfunctions $f_{j}$ may be chosen to be real, and for the time being we assume that they are.

Now, a classical real field $\phi_{\text {clas }}(t, \mathbf{x})$ can be expanded in terms of the eigenfunctions $f_{j}$ :

$$
\phi_{\mathrm{clas}}(t, \mathbf{x})=\sum q_{j}(t) f_{j}(\mathbf{x}),
$$

where the coefficients $q_{j}(t)$ are also real. $\phi_{\text {clas }}$ satisfies the Klein-Gordon equation if and only if these coefficients satisfy

$$
q_{j}^{\prime \prime}(t)+\omega_{j}^{2} q_{j}=0, \quad \omega_{j}^{2}=\lambda_{j}^{2}+m^{2}
$$

But this is the equation for a classical harmonic oscillator. To turn $\phi_{\text {clas }}$ into a quantum field, we replace these classical oscillators with quantum oscillators.

This necessitates a digression to discuss systems of quantum harmonic oscillators. First let us consider a system of $K$ one-dimensional oscillators, where $K$ is a positive integer. The state space is $L^{2}\left(\mathbb{R}^{K}\right)$ and the Hamiltonian is

$$
H=-\frac{1}{2} \nabla^{2}+\frac{1}{2} \sum \omega_{j} x_{j}^{2}=\sum \frac{\omega_{j}}{2}\left(-\frac{1}{\omega_{j}} \frac{\partial^{2}}{\partial x_{j}^{2}}+\omega_{j} x_{j}^{2}\right)
$$

where $\omega_{1}, \ldots, \omega_{K}$ are positive numbers. The substitution $y_{j}=\sqrt{\omega_{j}} x_{j}$ makes

$$
H=\sum \frac{\omega_{j}}{2}\left(-\frac{\partial^{2}}{\partial y_{j}^{2}}+y_{j}^{2}\right)
$$

a sum of standard one-dimensional oscillator Hamiltonians, to which we can apply the analysis of §3.4. Namely, let $X_{j}$ and $P_{j}$ be the position and momentum\\
operators for the $j$ th oscillator,

$$
X_{j}=\text { mult. by } x_{j}, \quad P_{j}=\frac{1}{i} \frac{\partial}{\partial x_{j}},
$$

and let $Y_{j}$ and $Q_{j}$ be the corresponding operators for the $y$-coordinates,

$$
Y_{j}=\sqrt{\omega_{j}} X_{j}=\text { mult. by } \sqrt{\omega_{j}} x_{j}, \quad Q_{j}=\frac{1}{\sqrt{\omega_{j}}} P_{j}=\frac{1}{i \sqrt{\omega_{j}}} \frac{\partial}{\partial x_{j}} .
$$

We then set

$$
A_{j}=\frac{1}{\sqrt{2}}\left(Y_{j}+i Q_{j}\right), \quad A_{j}^{\dagger}=\frac{1}{\sqrt{2}}\left(Y_{j}-i Q_{j}\right)
$$

For future reference we note that

$$
X_{j}=\frac{1}{\sqrt{2 \omega_{j}}}\left(A_{j}+A_{j}^{\dagger}\right), \quad P_{j}=\frac{1}{i} \sqrt{\frac{\omega_{j}}{2}}\left(A_{j}-A_{j}^{\dagger}\right) .
$$

These operators satisfy the canonical commutation relations

$$
\begin{gathered}
{\left[X_{j}, P_{k}\right]=i \delta_{j k} I, \quad\left[A_{j}, A_{k}^{\dagger}\right]=\delta_{j k} I} \\
{\left[X_{j}, X_{k}\right]=\left[P_{j}, P_{k}\right]=\left[A_{j}, A_{k}\right]=\left[A_{j}^{\dagger}, A_{k}^{\dagger}\right]=0}
\end{gathered}
$$

and the Hamiltonian is

$$
H=\sum \omega_{j}\left(A_{j}^{\dagger} A_{j}+\frac{1}{2}\right)
$$

The eigenfunctions of $H$ are simply the products of the one-dimensional eigenfunctions of the one-dimensional oscillator given by (3.27):

$$
\phi_{n_{1} n_{2} \cdots n_{K}}(x)=\frac{\left(A_{1}^{\dagger}\right)^{n_{1}} \cdots\left(A_{K}^{\dagger}\right)^{n_{K}}}{\sqrt{n_{1}!\cdots n_{K}!}} \phi_{00 \cdots 0}(x),
$$

where $\phi_{00 \cdots 0}(x)=\left[\prod\left(\omega_{j} / \pi\right)\right]^{1 / 4} \exp \left(-\frac{1}{2} \sum \omega_{j} x_{j}^{2}\right)$. By (3.28) and (3.29) we have

$$
\begin{gathered}
A_{j} \phi_{n_{1} \cdots n_{K}}=\sqrt{n_{j}} \phi_{n_{1} \cdots n_{j}-1 \cdots n_{K}}, \quad A_{j}^{\dagger} \phi_{n_{1} \cdots n_{K}}=\sqrt{n_{j}+1} \phi_{n_{1} \cdots n_{j}+1 \cdots n_{K}}, \\
H \phi_{n_{1} \cdots n_{K}}=\left[\sum \omega_{j}\left(n_{j}+\frac{1}{2}\right)\right] \phi_{n_{1} \cdots n_{K}} .
\end{gathered}
$$

If the reader has studied §4.5, this should look very familiar. Indeed, the correspondence

$$
\phi_{n_{1} \ldots n_{K}} \mapsto\left|n_{1}, \ldots, n_{K}\right\rangle
$$

defines a unitary map from $L^{2}\left(\mathbb{R}^{K}\right)$ to the Boson Fock space $\mathcal{F}_{s}\left(\mathbb{C}^{K}\right)$ that intertwines the operators $A_{j}$ and $A_{j}^{\dagger}$ with the creation and annihilation operators on $\mathcal{F}_{s}\left(\mathbb{C}^{K}\right)$ ! But what we really need for field theory is the case $K=\infty$, an infinite collection of harmonic oscillators. The material of §4.5 gives us exactly the mathematical machinery that we need: the Boson Fock space $\mathcal{F}_{s}(\mathcal{H})$, where now $\mathcal{H}$ is a separable infinite-dimensional Hilbert space, equipped with its creation and annihilation operators $A_{j}$ and $A_{j}^{\dagger}$ relative to a fixed orthonormal basis with the same index set as the eigenfunctions $f_{j}$.

There is just one problem that has to be addressed. Formally, the Hamiltonian ought to be $H=\sum_{1}^{\infty} \omega_{j}\left(A_{j}^{\dagger} A_{j}+\frac{1}{2}\right)$, but this does not make sense when the series $\sum \omega_{j}$ diverges, as it will in the situations we need since $\omega_{j}^{2}=\lambda_{j}^{2}+m^{2} \geq m^{2}$ for all $j$. This is our first brush with the dreaded divergences of quantum field theory, but unlike most of them, this one is easy to fix. Adding a constant to the Hamiltonian\\
does not affect the dynamics, so we simply throw away the infinite $\frac{1}{2} \sum \omega_{j}$ and consider the "renormalized" Hamiltonian

$$
H=\sum \omega_{j} A_{j}^{\dagger} A_{j}
$$

instead. This is a well defined positive operator on the finite-particle space $\mathcal{F}_{s}^{0}(\mathcal{H})$, and it has a unique extension to a positive self-adjoint operator on $\mathcal{F}_{s}(\mathcal{H})$; the vectors $\left|n_{1}, n_{2}, \cdots\right\rangle$ are an eigenbasis for it with eigenvalues $\sum \omega_{j} n_{j}$ (a finite sum since only finitely many $n_{j}$ are nonzero).

We are now in a position to quantize the classical field (5.1). By (5.2), the position operator for the $j$ th oscillator is $X_{j}=\left(A_{j}+A_{j}^{\dagger}\right) / \sqrt{2 \omega_{j}}$. Substituting this into (5.1) in place of the classical position variables $q_{j}$, we obtain the quantum field, denoted by $\phi$ :

$$
\phi(\mathbf{x})=\sum \frac{1}{\sqrt{2 \omega_{j}}} f_{j}(\mathbf{x})\left(A_{j}+A_{j}^{\dagger}\right)
$$

(We shall not worry about the convergence of this series until a little later.)\\
Before proceeding further, we need to free ourselves from the assumption that the eigenfunctions $f_{j}$ are real, that is, to obtain a formula for $\phi(\mathbf{x})$ in terms of an arbitrary complex-valued orthonormal eigenbasis $\left\{g_{j}\right\}$ for $-\nabla^{2}$ on the box $\mathcal{B}$ subject to the given boundary conditions. It cannot be right simply to substitute $g_{j}$ for $f_{j}$ in (5.4) because the resulting operator will not be Hermitian, and "Hermitian" is the quantum analogue of the classical "real-valued." Rather, let us observe that $g_{j}=\sum_{k} u_{j k} f_{k}$ where $U=\left(u_{j k}\right)$ is a unitary matrix such that $u_{j k}=0$ unless the eigenvalue for $g_{j}$ is the same as the eigenvalue for $f_{k}$, and hence the corresponding $\omega_{j}$ 's are also equal. Since $f_{k}$ is real and $U^{-1}=U^{\dagger}$,

$$
\sum_{j} u_{j k}^{*} g_{j}=f_{k}=f_{k}^{*}=\sum_{j} u_{j k} g_{j}^{*}
$$

Substituting these equations into (5.4), we obtain

$$
\phi=\sum_{k} \frac{1}{\sqrt{2 \omega_{k}}} f_{k}\left(A_{k}+A_{k}^{\dagger}\right)=\sum_{j, k} \frac{1}{\sqrt{2 \omega_{j}}}\left(g_{j} u_{j k}^{*} A_{k}+g_{j}^{*} u_{j k} A_{k}^{\dagger}\right)
$$

Thus if we set

$$
A_{j}^{\prime}=\sum_{k} u_{j k}^{*} A_{k},
$$

we have

$$
\phi=\sum_{j} \frac{1}{\sqrt{2 \omega_{j}}}\left(g_{j} A_{j}^{\prime}+g_{j}^{*}\left(A_{j}^{\prime}\right)^{\dagger}\right) .
$$

Moreover, we recall from §4.5 that $A_{j}=A\left(e_{j}\right)$ where $\left\{e_{j}\right\}$ is some orthonormal basis of $\mathcal{H}$. Since the map $v \mapsto A(v)$ is antilinear, we then have $A_{j}^{\prime}=A\left(e_{j}^{\prime}\right)$ where $e_{j}^{\prime}=\sum_{k} u_{j k} e_{k}$, and $\left\{e_{j}^{\prime}\right\}$ is again an orthonormal basis of $\mathcal{H}$.

In short, we see that no matter what basis $\left\{f_{j}\right\}$ we pick, the correct formula for the field is

$$
\phi(\mathbf{x})=\sum_{j} \frac{1}{\sqrt{2 \omega_{j}}}\left(f_{j}(\mathbf{x}) A_{j}+f_{j}(\mathbf{x})^{*} A_{j}^{\dagger}\right)
$$

where the $A_{j}$ 's and $A_{j}^{\dagger}$ 's are the annihilation and creation operators for an orthonormal basis of $\mathcal{H}$. The choice of such a basis is irrelevant; only the formal structure matters.

The $t$-dependence of the classical field seems to have disappeared, but it is restored by passing from the Schrödinger-picture operators $A_{j}, A_{j}^{\dagger}$ to the timedependent Heisenberg-picture operators

$$
A_{j}(t)=e^{i t H} A_{j} e^{-i t H}, \quad A_{j}^{\dagger}(t)=e^{i t H} A_{j}^{\dagger} e^{-i t H}
$$

This can easily be made more concrete. Indeed, from (4.52) and the fact that $H=\sum \omega_{j} A_{j}^{\dagger} A_{j}$, we see that

$$
\left[A_{j}, H\right]=\omega_{j} A_{j}, \quad\left[A_{j}^{\dagger}, H\right]=-\omega_{j} A_{j}^{\dagger}
$$

In view of this, (5.6) implies that

$$
\frac{d A_{j}(t)}{d t}=\frac{1}{i}\left[A_{j}(t), H\right]=-i \omega_{j} A_{j}(t), \quad \frac{d A_{j}^{\dagger}(t)}{d t}=\frac{1}{i}\left[A_{j}^{\dagger}(t), H\right]=i \omega_{j} A_{j}^{\dagger}(t)
$$

and hence

$$
A_{j}(t)=e^{-i \omega_{j} t} A_{j}, \quad A_{j}^{\dagger}(t)=e^{i \omega_{j} t} A_{j}^{\dagger}
$$

The Heisenberg-picture field $\phi(t, \mathbf{x})$ is therefore

$$
\phi(t, \mathbf{x})=\sum \frac{1}{\sqrt{2 \omega_{j}}}\left(f_{j}(\mathbf{x}) e^{-i \omega_{j} t} A_{j}+f_{j}(\mathbf{x})^{*} e^{i \omega_{j} t} A_{j}^{\dagger}\right)
$$

Now we make some specific choices. We take the box $\mathcal{B}$ to be the cube $\left[-\frac{1}{2} L, \frac{1}{2} L\right]^{3}$ in $\mathbb{R}^{3}$, and we impose periodic boundary conditions. (Don't worry about the nature of the boundary conditions; we're going to let $L \rightarrow \infty$ presently.) Then we can take the normalized eigenfunctions to be $f_{\mathbf{p}}(\mathbf{x})=L^{-3 / 2} e^{i \mathbf{p} \cdot \mathbf{x}}$, with eigenvalues $-|\mathbf{p}|^{2}$, where $\mathbf{p}$ lies in the lattice

$$
\Lambda=\left[\frac{2 \pi}{L} \mathbb{Z}\right]^{3}
$$

The corresponding numbers $\omega_{\mathbf{p}}$ are

$$
\omega_{\mathbf{p}}=\sqrt{|\mathbf{p}|^{2}+m^{2}}
$$

and the field takes the form

$$
\phi(t, \mathbf{x})=\sum_{\Lambda} \frac{1}{\sqrt{2 \omega_{\mathbf{p}} L^{3}}}\left(e^{i \mathbf{p} \cdot \mathbf{x}-i \omega_{\mathbf{p}} t} A_{\mathbf{p}}+e^{-i \mathbf{p} \cdot \mathbf{x}+i \omega_{\mathbf{p}} t} A_{\mathbf{p}}^{\dagger}\right)
$$

This is starting to look promising. Although our calculations have been completely non-Lorentz-invariant, we can see Lorentz form re-emerging in the exponents. Moreover, it is clear that (5.9) formally satisfies the Klein-Gordon equation. However, it is time we paid a little attention to questions of convergence, and here the situation is not so happy. Suppose, for instance, we apply the operator $\phi(t, \mathbf{x})$ to the vacuum state:

$$
\phi(t, \mathbf{x})|0,0, \ldots\rangle=\sum_{\Lambda} \frac{1}{\sqrt{2 \omega_{\mathbf{p}} L^{3}}} e^{-i \mathbf{p} \cdot \mathbf{x}+i \omega_{\mathbf{p}} t}|0, \ldots, 0,1,0, \ldots\rangle \quad(1 \text { in the } \mathbf{p} \text { th slot })
$$

The square of the norm of this alleged vector is $\sum_{\Lambda} 1 / 2\left(|\mathbf{p}|^{2}+m^{2}\right)$, which is infinite (by comparison to $\int_{\mathbb{R}^{3}} d \mathbf{p} /\left(|\mathbf{p}|^{2}+m^{2}\right)=4 \pi \int_{0}^{\infty} r^{2} d r /\left(r^{2}+m^{2}\right)$ ). If we apply $\phi(t, \mathbf{x})$ to a multiparticle state, the result is even worse, because the operators $A_{j}$ and $A_{j}^{*}$ introduce factors of $\sqrt{n_{j}}$ that are generally bigger than 1 .

The way out of this is to interpret $\phi$ as an operator-valued distribution rather than an operator-valued function. That is, $\phi$ is the linear map that assigns to each\\
compactly supported $C^{\infty}$ function $\chi_{1}$ on $\mathbb{R}$ and each $C^{\infty} \Lambda$-periodic function $\chi_{2}$ on $\mathbb{R}^{3}$ the operator

$$
\begin{aligned}
\int_{B} \int_{\mathbb{R}} \phi(t, \mathbf{x}) \chi_{1}(t) \chi_{2}(\mathbf{x}) & d t d \mathbf{x} \\
= & \sum_{\Lambda} \frac{1}{\sqrt{2 \omega_{\mathbf{p}} L^{3}}}\left[\widehat{\chi}_{1}\left(\omega_{\mathbf{p}}\right) \widehat{\chi}_{2}(-\mathbf{p}) A_{\mathbf{p}}+\widehat{\chi}_{1}\left(-\omega_{\mathbf{p}}\right) \widehat{\chi}_{2}(\mathbf{p}) A_{\mathbf{p}}^{\dagger}\right]
\end{aligned}
$$

with the obvious interpretation of the Fourier coefficients $\widehat{\chi}_{1}$ and $\widehat{\chi}_{2}$. The rapid decay of these coefficients as $|\mathbf{p}| \rightarrow \infty$ guarantees that this series converges nicely as an operator on the finite-particle space $\mathcal{F}_{s}^{0}(\mathcal{H})$.

The final step in the construction of the quantum field is the removal of the box $\mathcal{B}$. That is, we think of (5.9) as a Riemann sum for an integral over $\mathbf{p}$-space and pass to the limit as $L \rightarrow \infty$. The discrete volume element implicit in (5.9) is $\Delta V=(2 \pi / L)^{3}$, the volume of a fundamental cube of the lattice $\Lambda$, so we can rewrite (5.9) as

$$
\phi(t, \mathbf{x})=\sum_{\Lambda} \frac{1}{\sqrt{2 \omega_{\mathbf{p}}}} L^{3 / 2}\left(e^{i \mathbf{p} \cdot \mathbf{x}-i \omega_{\mathbf{p}} t} A_{\mathbf{p}}+e^{-i \mathbf{p} \cdot \mathbf{x}+i \omega_{\mathbf{p}} t} A_{\mathbf{p}}^{*}\right) \frac{\Delta V}{(2 \pi)^{3}}
$$

As $L \rightarrow \infty$, the sum becomes an integral over $\mathbb{R}^{3}$, and the rescaled annihilation and creation operators $L^{3 / 2} A_{\mathbf{p}}$ and $L^{3 / 2} A_{\mathbf{p}}^{\dagger}$ turn into operator-valued distributions on $\mathbb{R}^{3}$, denoted by $a$ and $a^{\dagger}$ (actually, the action of $a$ on test functions is antilinear rather than linear). Indeed, if we take the underlying Hilbert space $\mathcal{H}$ to be $L^{2}\left(\mathbb{R}^{3}, d^{3} \mathbf{p} /(2 \pi)^{3}\right)$, the evaluation of the distribution $a$ on the test function $\chi \in C_{c}^{\infty}\left(\mathbb{R}^{3}\right)$ is the operator denoted by $A(\chi)$ in §4.5. They satisfy the distributional commutation relations

$$
\left[a(\mathbf{p}), a^{\dagger}\left(\mathbf{p}^{\prime}\right)\right]=(2 \pi)^{3} \delta\left(\mathbf{p}-\mathbf{p}^{\prime}\right) I, \quad\left[a(\mathbf{p}), a\left(\mathbf{p}^{\prime}\right)\right]=\left[a^{\dagger}(\mathbf{p}), a^{\dagger}\left(\mathbf{p}^{\prime}\right)\right]=0 .
$$

which are a restatement of (4.48).\\
Maintaining the notational fiction that distributions are functions, we can therefore write the field with the box removed as

$$
\phi(t, \mathbf{x})=\int_{\mathbb{R}^{3}} \frac{1}{\sqrt{2 \omega_{\mathbf{p}}}}\left(e^{i \mathbf{p} \cdot \mathbf{x}-i \omega_{\mathbf{p}} t} a(\mathbf{p})+e^{-i \mathbf{p} \cdot \mathbf{x}+i \omega_{\mathbf{p}} t} a^{\dagger}(\mathbf{p})\right) \frac{d^{3} \mathbf{p}}{(2 \pi)^{3}}
$$

or, with $x=(t, \mathbf{x})$ and $p=\left(\omega_{\mathbf{p}}, \mathbf{p}\right)$,

$$
\phi(x)=\int \frac{1}{\sqrt{2 \omega_{\mathbf{p}}}}\left(e^{-i p_{\mu} x^{\mu}} a(\mathbf{p})+e^{i p_{\mu} x^{\mu}} a^{\dagger}(\mathbf{p})\right) \frac{d^{3} \mathbf{p}}{(2 \pi)^{3}}
$$

Warning: It is a sad fact of life that if one examines $n$ books on quantum field theory, one will probably find $n$ variants of (5.12) that differ in the factors of $2 \pi$ and $\omega_{\mathbf{p}}$. The point is that one can redefine $a(\mathbf{p})$ by incorporating such factors into it. If one wants to think of (5.12) as an integral over the mass shell $X_{m}^{+}$with its Lorentz-invariant measure $d^{3} \mathbf{p} /(2 \pi)^{3} \omega_{\mathbf{p}}$, for example, one can replace $a(\mathbf{p})$ with $\widetilde{a}(\mathbf{p})=\omega_{\mathbf{p}}^{1 / 2} a(\mathbf{p})$; or if one wants to eliminate the $2 \pi$ 's from (5.10), one can replace $a(\mathbf{p})$ with $a(\mathbf{p}) /(2 \pi)^{3 / 2}$. In this book we shall keep the canonical relations (5.10) for the annihilation and creation operators and let the consequences be what they may.

The good news about the formula (5.12) is its obviously Lorentz-friendly form. The bad news is its analytic prickliness: the operator-valued functions in the integrand are actually distributions, and the integral itself must be interpreted in a distributional sense. We shall reassure the reader by showing how to construct $\phi$ by rigorous mathematics in the next section. However, it must be understood that formulas such as (5.12) are quite convenient for the sort of calculations one actually has to perform to extract the physics, and they are ubiquitous in the physics literature. They will appear frequently in this book, and hopefully the reader will acquire greater comfort with them by seeing how they are used in practice.

Now, the Hilbert space on which $\phi(x)$ (or rather the regularized version obtained by integrating against a test function) acts is $\mathcal{F}_{s}(\mathcal{H})$, the Boson Fock space over a Hilbert space $\mathcal{H}$. We introduced this Fock space simply because it is the place where the required operators arising from the harmonic oscillators live. But the physics of the quantum field comes from taking the Fock space seriously as a multiparticle state space built from a single-particle state space $\mathcal{H}$; the particles in question are the quanta of the field. The operators $a(\mathbf{p})$ and $a^{\dagger}(\mathbf{p})$ annihilate and create particles with momentum $\mathbf{p}$ (where again this is an idealization; $a$ and $a^{\dagger}$ are distributions, and real particles don't have completely definite momenta), and $\phi(x)$ creates and annihilates particles at the space-time point $x$. The physical meaning of these assertions may seem rather obscure, but we must ask the reader to accept them for the time being on the level of hand-waving. We are building a house, and at present we are merely laying the foundation.


\pagebreak


Let us now turn to the quantization of a complex Klein-Gordon field. We employ the informal language of distributions as in (5.12).

If $\phi_{\text {clas }}$ is a complex-valued solution of the Klein-Gordon equation, its real and imaginary parts $\phi_{1, \text { clas }}$ and $\phi_{2, \text { clas }}$ are also solutions. To quantize them, we introduce annihilation and creation operators $a_{1}(\mathbf{p}), a_{1}^{\dagger}(\mathbf{p})$ and $a_{2}(\mathbf{p}), a_{2}^{\dagger}(\mathbf{p})$, which are taken to satisfy the canonical commutation relations (5.10) among themselves and to commute with each other:

$$
\left[a_{1}(\mathbf{p}), a_{1}^{\dagger}\left(\mathbf{p}^{\prime}\right)\right]=\left[a_{2}(\mathbf{p}), a_{2}^{\dagger}\left(\mathbf{p}^{\prime}\right)\right]=(2 \pi)^{3} \delta\left(\mathbf{p}-\mathbf{p}^{\prime}\right), \quad \text { all other commutators }=0 .
$$

(The Hilbert space $\mathcal{F}$ on which these operators act is obtained by starting with two Fock spaces $\mathcal{F}_{1}$ and $\mathcal{F}_{2}$ on which the $a_{1}$ 's and $a_{2}$ 's act, respectively, and taking $\mathcal{F}=\mathcal{F}_{1} \otimes \mathcal{F}_{2}$.) We can then form the quantum fields

$$
\phi_{j}(x)=\int \frac{1}{\sqrt{2 \omega_{\mathbf{p}}}}\left(e^{-i p_{\mu} x^{\mu}} a_{j}(\mathbf{p})+e^{i p_{\mu} x^{\mu}} a_{j}^{\dagger}(\mathbf{p})\right) \frac{d^{3} \mathbf{p}}{(2 \pi)^{3}} \quad(j=1,2)
$$

We put these together into a "complex" quantum field by setting

$$
\phi(x)=\frac{\phi_{1}(x)+i \phi_{2}(x)}{\sqrt{2}}, \quad a(\mathbf{p})=\frac{a_{1}(\mathbf{p})+i a_{2}(\mathbf{p})}{\sqrt{2}}, \quad b(\mathbf{p})=\frac{a_{1}(\mathbf{p})-i a_{2}(\mathbf{p})}{\sqrt{2}},
$$

so that

$$
\phi^{\dagger}(x)=\frac{\phi_{1}(x)-i \phi_{2}(x)}{\sqrt{2}}, \quad a^{\dagger}(\mathbf{p})=\frac{a_{1}^{\dagger}(\mathbf{p})-i a_{2}^{\dagger}(\mathbf{p})}{\sqrt{2}}, \quad b^{\dagger}(\mathbf{p})=\frac{a_{1}^{\dagger}(\mathbf{p})+i a_{2}^{\dagger}(\mathbf{p})}{\sqrt{2}},
$$

The factors of $\sqrt{2}$ are there to normalize the commutation relations so that (5.14)

$$
\left[a(\mathbf{p}), a^{\dagger}\left(\mathbf{p}^{\prime}\right)\right]=\left[b(\mathbf{p}), b^{\dagger}\left(\mathbf{p}^{\prime}\right)\right]=(2 \pi)^{3} \delta\left(\mathbf{p}-\mathbf{p}^{\prime}\right), \quad \text { all other commutators }=0 .
$$

The fields $\phi$ and $\phi^{\dagger}$ are expressed in terms of the $a$ 's and $b$ 's as follows:

$$
\begin{aligned}
\phi(x) & =\int \frac{1}{\sqrt{2 \omega_{\mathbf{p}}}}\left(e^{-i p_{\mu} x^{\mu}} a(\mathbf{p})+e^{i p_{\mu} x^{\mu}} b^{\dagger}(\mathbf{p})\right) \frac{d^{3} \mathbf{p}}{(2 \pi)^{3}} \\
\phi^{\dagger}(x) & =\int \frac{1}{\sqrt{2 \omega_{\mathbf{p}}}}\left(e^{-i p_{\mu} x^{\mu}} b(\mathbf{p})+e^{i p_{\mu} x^{\mu}} a^{\dagger}(\mathbf{p})\right) \frac{d^{3} \mathbf{p}}{(2 \pi)^{3}}
\end{aligned}
$$

Let us look at this from a slightly different angle. Suppose $\phi$ is a tempered distribution on $\mathbb{R}^{4}$ that satisfies the Klein-Gordon equation $\left(\partial^{2}+m^{2}\right) \phi=0$. Then the Fourier transform of $\phi$ is a tempered distribution that satisfies $\left(-p^{2}+m^{2}\right) \widehat{\phi}(p)=$ 0 and hence is supported on the two-sheeted hyperboloid

$$
\begin{gathered}
\left\{p: p^{2}=m^{2}\right\}=X_{m}^{+} \cup X_{m}^{-} \\
X_{m}^{+}=\left\{p: p^{2}=m^{2}, p^{0}>0\right\}, \quad X_{m}^{-}=\left\{p: p^{2}=m^{2}, p^{0}<0\right\}
\end{gathered}
$$

In particular, we may consider the space of solutions of the Klein-Gordon equation whose Fourier transforms are of the form $u \lambda$ where $\lambda$ is a Lorentz-invariant measure on $X_{m}^{+} \cup X_{m}^{-}$, considered as a distribution on $\mathbb{R}^{4}$, and $u$ is a (reasonable) function. If $\phi$ is such a solution, we have

$$
\begin{aligned}
\phi(x) & =\int_{X_{m}^{+} \cup X_{m}^{-}} e^{-i p_{\mu} x^{\mu}} u(p) d \lambda(p) \\
& =\int_{X_{m}^{+}}\left(e^{-i p_{\mu} x^{\mu}} u(p)+e^{i p_{\mu} x^{\mu}} u(-p)\right) d \lambda(p)
\end{aligned}
$$

Now, if $\phi$ is real-valued, we have $\widehat{\phi}(-p)=\widehat{\phi}(p)^{*}$, so $\phi$ is completely determined by the restriction of $\widehat{\phi}$ to $X_{m}^{+}$. We obtain the corresponding quantum field by replacing the Fourier coefficients $u(p), p \in X_{m}^{+}$, by suitably normalized annihilation operators and their complex conjugates $u(-p)$ by their adjoint creation operators. However, if $\phi$ is allowed to be complex valued, the restrictions of $\widehat{\phi}$ to $X_{m}^{+}$and $X_{m}^{-}$ are independent, and both of them are needed to recover $\phi$. Hence the Fourier coefficients $u(-p)$ must be replaced by a different set of creation operators.

The physical interpretation is as follows: $a, a^{\dagger}$ and $b, b^{\dagger}$ are the annihilation and creation operators for two different species of particle with mass $m$ and spin 0 , which are antiparticles of each other. We make the conventional choice that $a, a^{\dagger}$ are associated with "particles" and $b, b^{\dagger}$ are associated with "antiparticles." Thus the field $\phi$ destroys particles and creates antiparticles; vice versa for $\phi^{\dagger}$.

This interpretation gives the solution of the problem of "negative energy states," which plagued the early attempts to use the Klein-Gordon equation just as it did the Dirac equation. When the Klein-Gordon equation is used as a single-particle wave equation, the negative energy states present a real difficulty. But when we think of the Klein-Gordon equation as describing a "classical" field and pass to the corresponding quantum field, there is a way out: the part of the solution whose Fourier transform lives on the negative energy shell contributes not particle states with negative energy but antiparticle states with positive energy. We shall see how this works in more detail in the next section.


\pagebreak

\subsection*{The rigorous construction}
We now present the rigorous construction of a neutral (real) quantum scalar field with mass parameter $m$. Afterwards we shall modify the construction to obtain a charged (complex) field.

We have already assembled all the ingredients we need. The first is the state space for a single spin-zero particle of mass $m$ (see §4.4):

$$
\mathcal{H}=L^{2}\left(X_{m}^{+}, \lambda\right),
$$

where $X_{m}^{+}$is the mass shell and $\lambda$ is the normalized Lorentz-invariant measure on it:

$$
d \lambda(p)=d \lambda\left(\omega_{\mathbf{p}}, \mathbf{p}\right)=\frac{d^{3} \mathbf{p}}{(2 \pi)^{3} \omega_{\mathbf{p}}} \quad\left(\omega_{\mathbf{p}}=p^{0}=\sqrt{|\mathbf{p}|^{2}+m^{2}}\right)
$$

Based on this, we construct the Boson Fock space $\mathcal{F}_{s}(\mathcal{H})$, on which we have annihilation and creation operators $A(v)$ and $A(v)^{\dagger}$ for $v \in \mathcal{H}$ as defined in §4.5. Finally, we define $R: \mathcal{S}\left(\mathbb{R}^{4}\right) \rightarrow \mathcal{H}$ by

$$
R f=\widehat{f} \mid X_{m}^{+}
$$

where $\widehat{f}(p)=\int e^{i p_{\mu} x^{\mu}} f(x) d^{4} x$ is the Lorentz-covariant Fourier transform of $f$.\\
The quantum field $\Phi$ (which differs from the $\phi$ of the preceding section by a change of variable, as we explain below) is defined as a real tempered distribution on $\mathbb{R}^{4}$ with values in the space of operators on the finite-particle Fock space $\mathcal{F}_{s}^{0}(\mathcal{H})$ - that is, an $\mathbb{R}$-linear map that takes a real-valued Schwartz-class function $f$ on $\mathbb{R}^{4}$ to an operator $\Phi(f)$ on $\mathcal{F}_{s}^{0}(\mathcal{H})$ - by

$$
\Phi(f)=\frac{1}{\sqrt{2}}\left[A(R f)+A(R f)^{\dagger}\right]
$$

(One can extend $\Phi$ to a $\mathbb{C}$-linear map on complex-valued Schwartz functions by setting $\Phi(f+i g)=\Phi(f)+i \Phi(g)$, of course, but (5.16) is only $\mathbb{R}$-linear as it stands because $A$ is antilinear.)

Observe that $\Phi$ is a distribution solution of the Klein-Gordon equation, that is, $\Phi\left(\left(\partial^{2}+m^{2}\right) f\right)=0$ for any $f \in \mathcal{S}$, because $\left[\left(\partial^{2}+m^{2}\right) f\right]^{\wedge}=\left(-p^{2}+m^{2}\right) \widehat{f}=0$ on $X_{m}^{+}$.

Let us see how to go from (5.16) to (5.12). First, we map $X_{m}^{+}$to $\mathbb{R}^{3}$ in the obvious way, $\left(\omega_{\mathbf{p}}, \mathbf{p}\right) \mapsto \mathbf{p}$, and correspondingly identify $L^{2}\left(X_{m}^{+}, \lambda\right)$ with $L^{2}\left(\mathbb{R}^{3}\right)$, where the measure on $\mathbb{R}^{3}$ is taken to be $d^{3} \mathbf{p} /(2 \pi)^{3}$. Explicitly, this correspondence is the unitary map $J: L^{2}\left(X_{m}^{+}, \lambda\right) \rightarrow L^{2}\left(\mathbb{R}^{3}\right)$ defined by

$$
J u(\mathbf{p})=\frac{1}{\sqrt{\omega_{\mathbf{p}}}} u\left(\omega_{\mathbf{p}}, \mathbf{p}\right)
$$

We then have the induced unitary map $\mathcal{F}(J): \mathcal{F}_{s}\left(L^{2}\left(X_{m}^{+}, \lambda\right)\right) \rightarrow \mathcal{F}_{s}\left(L^{2}\left(\mathbb{R}^{3}\right)\right)$, as explained in §4.5. We use it to transfer the annihilation, creation, and field operators to $L^{2}\left(\mathbb{R}^{3}\right)$; that is, we define

$$
a(v)=\mathcal{F}(J) A\left(J^{-1} v\right) \mathcal{F}(J)^{-1}, \quad a^{\dagger}(v)=\mathcal{F}(J) A\left(J^{-1} v\right)^{\dagger} \mathcal{F}(J)^{-1} \quad\left(v \in L^{2}\left(\mathbb{R}^{3}\right)\right)
$$

(it will be more convenient to write $a^{\dagger}(v)$ rather than $a(v)^{\dagger}$ ) and

$$
\phi(f)=\mathcal{F}(J) \Phi(f) \mathcal{F}(J)^{-1} \quad\left(f \in \mathcal{S}_{\mathbb{R}}\left(\mathbb{R}^{4}\right)\right)
$$

Now, $a^{\dagger}$ can be interpreted as an operator-valued distribution on $\mathbb{R}^{3}$ by restricting its argument to $\mathcal{S}\left(\mathbb{R}^{3}\right)$. Adopting the notational fiction that it is a function, we write

$$
a^{\dagger}(v)=\int a^{\dagger}(\mathbf{p}) v(\mathbf{p}) \frac{d^{3} \mathbf{p}}{(2 \pi)^{3}}
$$

Likewise, $a$ is an anti-linear distribution, and we write

$$
a(v)=\int a(\mathbf{p}) v(\mathbf{p})^{*} \frac{d^{3} \mathbf{p}}{(2 \pi)^{3}}
$$

By (4.48) we have

$$
\begin{aligned}
\iint\left[a(\mathbf{p}), a^{\dagger}(\mathbf{q})\right] u(\mathbf{p})^{*} v(\mathbf{q}) \frac{d^{3} \mathbf{p} d^{3} \mathbf{q}}{(2 \pi)^{6}} & =\left[a(u), a^{\dagger}(v)\right]=\langle u \mid v\rangle=\int u(\mathbf{p})^{*} v(\mathbf{p}) \frac{d^{3} \mathbf{p}}{(2 \pi)^{3}} \\
& =\iint(2 \pi)^{3} \delta(\mathbf{p}-\mathbf{q}) v(\mathbf{p})^{*} v(\mathbf{q}) \frac{d^{3} \mathbf{p} d^{3} \mathbf{q}}{(2 \pi)^{6}}
\end{aligned}
$$

for all $u$ and $v$, and similarly for $[a(\mathbf{p}), a(\mathbf{q})]$ and $\left[a^{\dagger}(\mathbf{p}), a^{\dagger}(\mathbf{q})\right]$. (The last three expressions are written as scalars but, as usual, are to be interpreted as scalar multiples of the identity.) Hence the commutation relations (4.48) expressed in colloquial distribution language are

$$
\left[a(\mathbf{p}), a^{\dagger}(\mathbf{q})\right]=(2 \pi)^{3} \delta(\mathbf{p}-\mathbf{q}), \quad[a(\mathbf{p}), a(\mathbf{q})]=\left[a^{\dagger}(\mathbf{p}), a^{\dagger}(\mathbf{q})\right]=0
$$

Next, observe that for $f \in \mathcal{S}_{\mathbb{R}}\left(\mathbb{R}^{4}\right)$ we have

$$
\phi(f)=\frac{1}{\sqrt{2}} \mathcal{F}(J)\left[A(R f)+A(R f)^{\dagger}\right] \mathcal{F}(J)^{-1}=\frac{1}{\sqrt{2}}\left[a(J R f)+a^{\dagger}(J R f)\right]
$$

and

$$
\operatorname{JRf}(\mathbf{p})=\frac{1}{\sqrt{\omega_{\mathbf{p}}}} \widehat{f}\left(\omega_{\mathbf{p}}, \mathbf{p}\right)
$$

so

$$
\begin{aligned}
\phi(f) & =\int \frac{1}{\sqrt{2 \omega_{\mathbf{p}}}}\left(\widehat{f}\left(\omega_{\mathbf{p}}, \mathbf{p}\right)^{*} a(\mathbf{p})+\widehat{f}\left(\omega_{\mathbf{p}}, \mathbf{p}\right) a^{\dagger}(\mathbf{p})\right) \frac{d^{3} \mathbf{p}}{(2 \pi)^{3}} \\
& =\iint \frac{1}{\sqrt{2 \omega_{\mathbf{p}}}}\left(e^{-i p_{\mu} x^{\mu}} a(\mathbf{p})+e^{i p_{\mu} x^{\mu}} a^{\dagger}(\mathbf{p})\right) f(x) d^{4} x \frac{d^{3} \mathbf{p}}{(2 \pi)^{3}}
\end{aligned}
$$

In other words,

$$
\phi(x)=\int \frac{1}{\sqrt{2 \omega_{\mathbf{p}}}}\left(e^{-i p_{\mu} x^{\mu}} a(\mathbf{p})+e^{i p_{\mu} x^{\mu}} a^{\dagger}(\mathbf{p})\right) \frac{d^{3} \mathbf{p}}{(2 \pi)^{3}}
$$

which is (5.12).

\pagebreak


We now modify the construction to obtain the charged (complex) field. For this purpose we need both the positive and negative energy mass shells $X_{m}^{+}$and $X_{m}^{-}$, each equipped with the invariant measure $d \lambda(p)=d^{3} \mathbf{p} /(2 \pi)^{3} \omega_{\mathbf{p}}$. We set

$$
\mathcal{H}_{+}=L^{2}\left(X_{m}^{+}, \lambda\right), \quad \mathcal{H}_{-}=L^{2}\left(X_{m}^{-}, \lambda\right), \quad \mathcal{H}=L^{2}\left(X_{m}^{+} \cup X_{m}^{-}\right)=\mathcal{H}_{+} \oplus \mathcal{H}_{-},
$$

and define the anti-unitary operator $C: \mathcal{H}_{ \pm} \rightarrow \mathcal{H}_{\mp}$ by

$$
C u(p)=u(-p)^{*}
$$

(On the inverse Fourier transform side, i.e., in position space, $C$ is just complex conjugation, $f \mapsto f^{*}$.) $\mathcal{H}_{+}$is the state space for a single particle of mass $m$ and spin 0 ; we build the Fock space $\mathcal{F}_{+}=\mathcal{F}_{s}\left(\mathcal{H}_{+}\right)$and the annihilation and creation\\
operators $A(u), A(u)^{\dagger}$ for $u \in \mathcal{H}_{+}$on it as before. Now, here comes the twist. Let $\mathcal{F}_{-}$be another copy of $\mathcal{F}_{s}\left(\mathcal{H}_{+}\right)$(not $\mathcal{F}_{s}\left(\mathcal{H}_{-}\right)$) and define annihilation and creation operators $B(v), B(v)^{\dagger}$ for $v \in \mathcal{H}_{-}$( not $\mathcal{H}_{+}$) on it by

$$
B(v)=A(C v), \quad B(v)^{\dagger}=A(C v)^{\dagger} .
$$

We think of $\mathcal{F}_{+}$and $A, A^{\dagger}$ as the Fock space and annihilation-creation operators for particles, and $\mathcal{F}_{-}$and $B, B^{\dagger}$ as the Fock space and annihilation-creation operators for antiparticles. $C$ is the "charge-conjugation" operator (although there may be no electric charge involved) that turns a particle into an antiparticle and vice versa. Observe that since $A(v)$ and $A(v)^{\dagger}$ depend antilinearly and linearly on $v$, respectively, and $C$ is antilinear, $B(v)$ and $B(v)^{\dagger}$ depend linearly and antilinearly on $v$, respectively.

We combine the two Fock spaces into a single Fock space $\mathcal{F}=\mathcal{F}_{+} \otimes \mathcal{F}_{-}$. Observe that $\mathcal{F}_{ \pm}=\bigoplus_{0}^{\infty} \mathcal{F}_{ \pm}^{k}$ where $\mathcal{F}_{ \pm}^{k}=\left(S^{k} \mathcal{H}_{ \pm}\right.$is the space of $k$-particle (for + ) or $k$ antiparticle (for -) states, and hence

$$
\mathcal{F}=\bigoplus_{j, k=0}^{\infty} \mathcal{F}_{+}^{j} \otimes \mathcal{F}_{-}^{k}
$$

where the summands on the right are the spaces of states with $j$ particles and $k$ antiparticles. We consider $A(u)$ and $B(v)$ and their adjoints as operators on $\mathcal{F}$ by letting $A(u)$ act on the first factor and $B(v)$ on the second one. Finally, we define $R_{ \pm}: \mathcal{S}\left(\mathbb{R}^{4}\right) \rightarrow \mathcal{H}_{ \pm}$by

$$
R_{ \pm} f=\hat{f} \mid X_{m}^{ \pm}
$$

and for $f \in \mathcal{S}\left(\mathbb{R}^{4}\right)$ (now complex-valued) we define the field operators $\Phi(f)$ and $\Phi^{\dagger}(f)$ on the finite-particle subspace of $\mathcal{F}$ by

$$
\begin{aligned}
\Phi(f) & =\frac{1}{\sqrt{2}}\left[A\left(R_{+} f\right)+B\left(R_{-} f\right)^{\dagger}\right], \\
\Phi^{\dagger}(f) & =\frac{1}{\sqrt{2}}\left[A\left(R_{+} f\right)^{\dagger}+B\left(R_{-} f\right)\right] .
\end{aligned}
$$

Thus $\Phi^{\dagger}$ is a distribution on $\mathbb{R}^{4}$ with values in the space of operators on the finiteparticle subspace of $\mathcal{F} ; \Phi$ is too, except that it depends antilinearly on its argument $f$. The reader may verify by a calculation similar to the one we performed for the neutral field that (5.18) is equivalent to (5.15).

Further discussion of the mathematics of free scalar fields can be found in Reed and Simon [95], §X.7.

\subsection*{Lagrangians and Hamiltonians}
In nonrelativistic quantum mechanics the dynamics of a system is described by the Schrödinger equation, which we derived by starting with the Hamiltonian formulation of the classical mechanical system and quantizing the Hamiltonian. Quantum field theory, on the other hand, is relativistic from the outset, so the Lagrangian formulation provides a better starting point. That is, one begins with a Lagrangian $L$ that is the integral of a Lorentz-invariant Lagrangian density $\mathcal{L}$, derives a Hamiltonian from it, and then quantizes the latter. The requirement of Lorentz invariance places severe constraints on the possible form of the Lagrangian and hence provides much guidance on the possible forms of quantum field theories.

Let us see how this works for free scalar fields. There is no interesting dynamics here, but it is a useful exercise to see how the process works in this simple case. It will lead to some useful formulas and provide some practice with the sort of freewheeling calculations with operator-valued distributions that physicists like to perform. We emphasize that all such calculations in this chapter can be made perfectly rigorous. The reader may wish to do so as an exercise, but here the idea is to see how the physicists operate.

We begin with the real scalar field. The Lagrangian that yields the KleinGordon equation is

$$
L=\int \mathcal{L} d^{3} \mathbf{x}, \quad \mathcal{L}=\frac{1}{2}\left[\left(\partial_{\mu} \phi\right)\left(\partial^{\mu} \phi\right)-m^{2} \phi^{2}\right]=\frac{1}{2}\left[\left(\partial_{t} \phi\right)^{2}-\left|\nabla_{\mathbf{x}} \phi\right|^{2}-m^{2} \phi^{2}\right]
$$

Here, by analogy with Lagrangian mechanics with finitely many degrees of freedom, $\phi(\cdot, \mathbf{x})$ plays the role of a "position" variable labeled by $\mathbf{x}$; the corresponding canonically conjugate "momentum" is

$$
\pi=\frac{\partial \mathcal{L}}{\partial\left(\partial_{t} \phi\right)}=\partial_{t} \phi
$$

and the Hamiltonian is

$$
H=\int \mathcal{H} d^{3} \mathbf{x}, \quad \mathcal{H}=\pi\left(\partial_{t} \phi\right)-\mathcal{L}=\frac{1}{2}\left[\left(\partial_{t} \phi\right)^{2}+\left|\nabla_{\mathbf{x}} \phi\right|^{2}+m^{2} \phi^{2}\right]
$$

We can quantize these things. Let $\phi$ now denote the quantum field (5.12). The canonically conjugate field (not the physical momentum associated with the field) is

$$
\pi(x)=\partial_{t} \phi(x)=\frac{1}{i} \int \sqrt{\frac{\omega_{\mathbf{p}}}{2}}\left(e^{i p_{\mu} x^{\mu}} a(\mathbf{p})-e^{-i p_{\mu} x^{\mu}} a^{\dagger}(\mathbf{p})\right) \frac{d^{3} \mathbf{p}}{(2 \pi)^{3}}
$$

We could also arrive at this formula in another way. At the beginning of our derivation of the quantum field $\phi$, we replaced Fourier coefficients $q_{j}(t)$ for the classical field with position operators $\left(A_{j}+A_{j}^{\dagger}\right) / \sqrt{2 \omega_{j}}$ for harmonic oscillators. The conjugate momentum operators are $\sqrt{\omega_{j} / 2}\left(A_{j}-A_{j}^{\dagger}\right)$ (see (5.2)); if we insert them in the Fourier series in place of the position operators and proceed as before, we end up with (5.19). Similarly, one can modify the rigorous construction of $\phi(f)$ for $f \in \mathcal{S}\left(\mathbb{R}^{4}\right)$ to obtain $\pi(f)$; we leave the details to the reader. (They involve replacing $R f$ by $\widetilde{R} f(p)=\omega_{\mathbf{p}} R f(p)$.)

The fields $\phi$ and $\pi$ at a given time $t$ satisfy "canonical commutation relations." The formal calculation is as follows: $\phi(t, \mathbf{x})$ and $\pi(t, \mathbf{y})$ are each the sum of two integrals, the first involving annihilation operators and the second involving creation operators. If one writes out the commutator $[\phi(t, \mathbf{x}), \pi(t, \mathbf{y})]$ as a double integral and uses the commutation relations (5.10) to simplify it, one finds that it is equal\\
to

$$
\begin{aligned}
& \frac{1}{2 i} \iint \sqrt{\frac{\omega_{\mathbf{q}}}{\omega_{\mathbf{p}}}}\left(e^{-i p_{\mu} x^{\mu}+i q_{\mu} y^{\mu}}\left[a^{\dagger}(\mathbf{p}), a(\mathbf{q})\right]-e^{i p_{\mu} x^{\mu}-i q_{\mu} y^{\mu}}\left[a(\mathbf{p}), a^{\dagger}(\mathbf{q})\right]\right) \frac{d^{3} \mathbf{p} d^{3} \mathbf{q}}{(2 \pi)^{6}} \\
= & \frac{i}{2} \iint \sqrt{\frac{\omega_{\mathbf{q}}}{\omega_{\mathbf{p}}}}\left(e^{-i p_{\mu} x^{\mu}+i q_{\mu} y^{\mu}}+e^{i p_{\mu} x^{\mu}-i q_{\mu} y^{\mu}}\right) \delta(\mathbf{p}-\mathbf{q}) \frac{d^{3} \mathbf{p} d^{3} \mathbf{q}}{(2 \pi)^{3}} \\
= & \frac{i}{2} \int\left(e^{i p_{\mu}\left(y^{\mu}-x^{\mu}\right)}+e^{i p_{\mu}\left(x^{\mu}-y^{\mu}\right)}\right) \frac{d^{3} \mathbf{p}}{(2 \pi)^{3}} \\
= & i \delta(\mathbf{x}-\mathbf{y})
\end{aligned}
$$

The last equality is (1.5), taking into account that $x^{0}-y^{0}=t-t=0$. Similarly one sees that $\phi(t, \mathbf{x})$ commutes with $\phi(t, \mathbf{y})$, and $\pi(t, \mathbf{x})$ with $\pi(t, \mathbf{y})$, for all $\mathbf{x}$ and $\mathbf{y}$. In short:

$$
[\phi(t, \mathbf{x}), \pi(t, \mathbf{y})]=i \delta(\mathbf{x}-\mathbf{y}), \quad[\phi(t, \mathbf{x}), \phi(t, \mathbf{y})]=[\pi(t, \mathbf{x}), \pi(t, \mathbf{y})]=0
$$

(One may worry about the fact that (5.20) is not Lorentz-invariant, as the notion of simultaneity at different points of space does not make sense relativistically. In fact, one may replace $(t, \mathbf{y})$ by $(s, \mathbf{y})$ in (5.20) as long as $|t-s|<|\mathbf{x - y}|$, i.e., as long as $(t, \mathbf{x})$ and $(s, \mathbf{y})$ are space-like separated. We shall return to this point later.)

The formulation of these results in mathematicians' language is a rather easy consequence of (5.16). There is one little technicality to be overcome: the relations (5.20) hold only for a fixed time $t$, whereas smearing out $\phi(t, \mathbf{x})$ into $\phi(f)= \int \phi(t, \mathbf{x}) f(t, \mathbf{x}) d t d^{3} \mathbf{x}$ involves the values of $\phi$ at different times. But in fact, there is no need to smear $\phi$ out in $t$. A glance at the definition of $\phi(f)$ shows that it makes sense not just for $f \in \mathcal{S}\left(\mathbb{R}^{4}\right)$ but for distributions of the form $f(t, \mathbf{x})=\delta\left(t-t_{0}\right) F(\mathbf{x})$ with $F \in \mathcal{S}\left(\mathbb{R}^{3}\right)$. Indeed, for such $f$ we have $\widehat{f}\left(\omega_{p}, \mathbf{p}\right)=e^{-i t_{0} \omega_{\mathbf{p}}} \widehat{F}(-\mathbf{p})$ where $\widehat{F}$ denotes the Euclidean Fourier transform on $\mathbb{R}^{3}$, and this is still a nice function on the mass shell $X_{m}^{+}$.

Now let us calculate the Hamiltonian. There are three possible starting points. First, the Hamiltonian is the infinitesimal generator of the time-translation group. Thus, on the one-particle space $L^{2}\left(X_{m}^{+}, \lambda\right)$ it is multiplication by $\omega_{\mathbf{p}}$, and on the $k$-particle space it is the sum of multiplications by $\omega_{\mathbf{p}}$ on all the factors:\\
(5.21) $H\left(P_{s}\left(u_{1} \otimes \cdots \otimes u_{k}\right)\right)=P_{s}\left(\left(\omega_{\mathbf{p}} u_{1}\right) \otimes \cdots \otimes u_{k}\right)+\cdots+P_{s}\left(u_{1} \otimes \cdots \otimes\left(\omega_{\mathbf{p}} u_{k}\right)\right)$.

Second, we can start from the renormalized Hamiltonian (5.3) for a discrete but infinite system of harmonic oscillators and pass to the continuum limit, obtaining

$$
H=\int \omega_{\mathbf{p}} a^{\dagger}(\mathbf{p}) a(\mathbf{p}) \frac{d^{3} \mathbf{p}}{(2 \pi)^{3}}
$$

It is not hard to see that (5.21) and (5.22) are equivalent. Third, we can quantize the formula for the classical Hamiltonian derived from the Lagrangian:

$$
\left.H=\left.\frac{1}{2} \int\left[\left(\partial_{t} \phi(t, \mathbf{x})\right)^{2}+\mid \nabla_{\mathbf{x}} \phi(t, \mathbf{x})\right)\right|^{2}+m^{2} \phi(t, \mathbf{x})^{2}\right] d^{3} \mathbf{x}
$$

by substituting the expression (5.12) for $\phi(t, \mathbf{x})$ in (5.23).\\
Let us see how this works. We have $\phi=\phi^{+}+\phi^{-}$, where

$$
\phi^{+}(x)=\int \frac{1}{\sqrt{2 \omega_{\mathbf{p}}}} e^{-i p_{\mu} x^{\mu}} a(\mathbf{p}) \frac{d^{3} \mathbf{p}}{(2 \pi)^{3}}, \quad \quad \phi^{-}(x)=\phi^{+}(x)^{\dagger}
$$

Now,

$$
\partial_{t} \phi^{+}(x)=\int \frac{1}{\sqrt{2 \omega_{\mathbf{p}}}}\left(-i \omega_{\mathbf{p}}\right) e^{-i p_{\mu} x^{\mu}} a(\mathbf{p}) \frac{d^{3} \mathbf{p}}{(2 \pi)^{3}},
$$

so

$$
\int\left(\partial_{t} \phi^{+}(x)\right)^{2} d^{3} \mathbf{x}=-\iiint \frac{\sqrt{\omega_{\mathbf{p}} \omega_{\mathbf{p}}^{\prime}}}{2} e^{-i p_{\mu} x^{\mu}} e^{-i p_{\mu}^{\prime} x^{\mu}} a(\mathbf{p}) a\left(\mathbf{p}^{\prime}\right) \frac{d^{3} \mathbf{p} d^{3} \mathbf{p}^{\prime}}{(2 \pi)^{6}} d^{3} \mathbf{x}
$$

By the Fourier inversion formula

$$
\int e^{i\left(\mathbf{p}+\mathbf{p}^{\prime}\right) \cdot \mathbf{x}} d^{3} \mathbf{x}=(2 \pi)^{3} \delta\left(\mathbf{p}+\mathbf{p}^{\prime}\right)
$$

this reduces to

$$
\int\left(\partial_{t} \phi^{+}(x)\right)^{2} d^{3} \mathbf{x}=-\int \frac{\omega_{\mathbf{p}}}{2} e^{-2 i \omega_{\mathbf{p}} t} a(\mathbf{p}) a(-\mathbf{p}) \frac{d^{3} \mathbf{p}}{(2 \pi)^{3}}
$$

After similar calculations to evaluate the integrals of $\left(\partial^{t} \phi^{-}\right)^{2},\left(\partial_{t} \phi^{+}\right)\left(\partial_{t} \phi^{-}\right)$, and $\left(\partial_{t} \phi^{-}\right)\left(\partial_{t} \phi^{+}\right)$, we obtain

$$
\begin{aligned}
& \int\left(\partial_{t} \phi(x)\right)^{2} d^{3} \mathbf{x} \\
= & \int \frac{\omega_{\mathbf{p}}^{2}}{2 \omega_{\mathbf{p}}}\left[-e^{-2 i \omega_{\mathbf{p}} t} a(\mathbf{p}) a(-\mathbf{p})+a(\mathbf{p}) a^{\dagger}(\mathbf{p})+a^{\dagger}(\mathbf{p}) a(\mathbf{p})-e^{2 i \omega_{\mathbf{p}} t} a^{\dagger}(\mathbf{p}) a^{\dagger}(-\mathbf{p})\right] \frac{d^{3} \mathbf{p}}{(2 \pi)^{3}}
\end{aligned}
$$

Likewise,

$$
\nabla_{\mathbf{x}} \phi^{+}(x)=\int \frac{1}{\sqrt{2 \omega_{\mathbf{p}}}}(i \mathbf{p}) e^{-i p_{\mu} x^{\mu}} a(\mathbf{p}) \frac{d^{3} \mathbf{p}}{(2 \pi)^{3}}
$$

so an analogous calculation gives

$$
\begin{aligned}
& \int\left|\nabla_{\mathbf{x}} \phi(x)\right|^{2} d^{3} \mathbf{x} \\
= & \int \frac{|\mathbf{p}|^{2}}{2 \omega_{\mathbf{p}}}\left[e^{-2 i \omega_{\mathbf{p}} t} a(\mathbf{p}) a(-\mathbf{p})+a(\mathbf{p}) a^{\dagger}(\mathbf{p})+a^{\dagger}(\mathbf{p}) a(\mathbf{p})+e^{2 i \omega_{\mathbf{p}} t} a^{\dagger}(\mathbf{p}) a^{\dagger}(-\mathbf{p})\right] \frac{d^{3} \mathbf{p}}{(2 \pi)^{3}} .
\end{aligned}
$$

The signs of the coefficients of $a(\mathbf{p}) a(-\mathbf{p})$ and $a^{\dagger}(\mathbf{p}) a^{\dagger}(-\mathbf{p})$ are different here because the $\delta\left(\mathbf{p}+\mathbf{p}^{\prime}\right)$ makes $\mathbf{p} \cdot \mathbf{p}^{\prime}$ into $-|\mathbf{p}|^{2}$ rather than $|\mathbf{p}|^{2}$. Yet another calculation of the same sort gives

$$
\begin{aligned}
& \int m^{2} \phi(x)^{2} d^{3} \mathbf{x} \\
= & \int \frac{m^{2}}{2 \omega_{\mathbf{p}}}\left[e^{-2 i \omega_{\mathbf{p}} t} a(\mathbf{p}) a(-\mathbf{p})+a(\mathbf{p}) a^{\dagger}(\mathbf{p})+a^{\dagger}(\mathbf{p}) a(\mathbf{p})+e^{2 i \omega_{\mathbf{p}} t} a^{\dagger}(\mathbf{p}) a^{\dagger}(-\mathbf{p})\right] \frac{d^{3} \mathbf{p}}{(2 \pi)^{3}}
\end{aligned}
$$

Adding these results and recalling that $\omega_{\mathbf{p}}^{2}=|\mathbf{p}|^{2}+m^{2}$, we see that the expression on the right of (5.23) is equal to

$$
\frac{1}{2} \int \omega_{\mathbf{p}}\left[a(\mathbf{p}) a^{\dagger}(\mathbf{p})+a^{\dagger}(\mathbf{p}) a(\mathbf{p})\right] \frac{d^{3} \mathbf{p}}{(2 \pi)^{3}}
$$

Finally, since $\left[a(\mathbf{p}), a^{\dagger}\left(\mathbf{p}^{\prime}\right)\right]=(2 \pi)^{3} \delta\left(\mathbf{p}-\mathbf{p}^{\prime}\right), a(\mathbf{p}) a^{\dagger}(\mathbf{p})$ is equal to $a^{\dagger}(\mathbf{p}) a(\mathbf{p})$ plus an infinite constant times the identity (an "infinite $c$-number" in physicists' parlance). After discarding the constant, we obtain (5.22). This infinite constant is exactly the same one we discarded in renormalizing the Hamiltonian for the infinite family of oscillators in §5.1.

The correction that needs to be applied to (5.23) to remove the infinite constant can be simply described as follows: replace all $a(\mathbf{p}) a^{\dagger}\left(\mathbf{p}^{\prime}\right)$ by $a^{\dagger}\left(\mathbf{p}^{\prime}\right) a(\mathbf{p})$. This procedure is sufficiently commonly encountered that it has a name. In general, a sum of products of creation and annihilation operators is said to be Wick ordered or normally ordered if all creation operators occur to the left of all annihilation operators in each product, and to Wick order such a sum of products is to replace it with the corresponding Wick ordered product. This is well defined since creation operators commute with each other, as do annihilation operators. Wick ordering is indicated by putting a colon on either side of the expression in question. Thus, for example,

$$
: a(\mathbf{p}) a^{\dagger}(\mathbf{p}):=: a^{\dagger}(\mathbf{p}) a(\mathbf{p}):=a^{\dagger}(\mathbf{p}) a(\mathbf{p})
$$

and the corrected version of (5.23) is

$$
\left.H=\frac{1}{2} \int:\left(\partial_{t} \phi(t, \mathbf{x})\right)^{2}+\mid \nabla_{\mathbf{x}} \phi(t, \mathbf{x})\right)\left.\right|^{2}+m^{2} \phi(t, \mathbf{x})^{2}: d^{3} \mathbf{x}
$$

We now turn to the complex Klein-Gordon field. The Lagrangian for a classical complex field is the same as for a real field except that squares of real numbers are replaced by absolute squares of complex numbers. Taking into account the factor of $\sqrt{2}$ that we introduced in (5.13), we find that the classical Lagrangian is

$$
L=\int \mathcal{L} d^{3} \mathbf{x}, \quad \mathcal{L}=\left(\partial_{\mu} \phi\right) \partial^{\mu} \phi^{*}-m^{2}|\phi|^{2}
$$

so the canonically conjugate field is

$$
\pi=\frac{\partial \mathcal{L}}{\partial\left(\partial_{t} \phi\right)}=\partial_{t} \phi^{*}
$$

and the Hamiltonian is

$$
H=\int \mathcal{H} d^{3} \mathbf{x}, \quad \mathcal{H}=\left|\partial_{t} \phi\right|^{2}+\left|\nabla_{\mathbf{x}} \phi\right|^{2}+m^{2}|\phi|^{2} .
$$

When we quantize, the complex conjugates turn into adjoints: The canonical conjugate of the quantum field $\phi$ of (5.15) is

$$
\pi(x)=\partial_{t} \phi^{\dagger}(x)=\int \frac{i \omega_{\mathbf{p}}}{\sqrt{2 \omega_{\mathbf{p}}}}\left[e^{i p_{\mu} x^{\mu}} a^{\dagger}(\mathbf{p})-e^{-i p_{\mu} x^{\mu}} b(\mathbf{p})\right] \frac{d^{3} \mathbf{p}}{(2 \pi)^{3}}
$$

The Hamiltonian can be expressed by formulas analogous to (5.22) and (5.24). The analogue of (5.22) is

$$
H=\int \omega_{\mathbf{p}}\left[a^{\dagger}(\mathbf{p}) a(\mathbf{p})+b^{\dagger}(\mathbf{p}) b(\mathbf{p})\right] \frac{d^{3} \mathbf{p}}{(2 \pi)^{3}}
$$

and the analogue of (5.24) is

$$
H=\int: \partial_{t} \phi^{\dagger} \partial_{t} \phi+\nabla_{\mathbf{x}} \phi^{\dagger} \cdot \nabla_{\mathbf{x}} \phi+m^{2} \phi^{\dagger} \phi: d^{3} \mathbf{x}
$$

Let us examine this a little more closely. If one starts from the classical Hamiltonian and quantizes the fields as in (5.23), without Wick ordering, one finds two possibilities for $H$ depending on whether one replaces $|\phi|^{2}$ by $\phi^{\dagger} \phi$ or $\phi \phi^{\dagger}$ :

$$
\begin{array}{r}
\int\left(\partial_{t} \phi^{\dagger} \partial_{t} \phi+\nabla_{\mathbf{x}} \phi^{\dagger} \cdot \nabla_{\mathbf{x}} \phi+m^{2} \phi^{\dagger} \phi\right) d^{3} \mathbf{x} \\
\text { and } \int\left(\partial_{t} \phi \partial_{t} \phi^{\dagger}+\nabla_{\mathbf{x}} \phi \cdot \nabla_{\mathbf{x}} \phi^{\dagger}+m^{2} \phi \phi^{\dagger}\right) d^{3} \mathbf{x}
\end{array}
$$

Calculations like the ones above that led from (5.23) to (5.22) show that these two expressions are respectively equal to

$$
\begin{array}{r}
\int \omega_{\mathbf{p}}\left[a^{\dagger}(\mathbf{p}) a(\mathbf{p})+b(\mathbf{p}) b^{\dagger}(\mathbf{p})\right] \frac{d^{3} \mathbf{p}}{(2 \pi)^{3}} \\
\text { and } \int \omega_{\mathbf{p}}\left[a(\mathbf{p}) a^{\dagger}(\mathbf{p})+b^{\dagger}(\mathbf{p}) b(\mathbf{p})\right] \frac{d^{3} \mathbf{p}}{(2 \pi)^{3}}
\end{array}
$$

Wick ordering turns both of these into (5.25); either way, there is an infinite constant to be discarded.

There is another quantity of interest here. Recall that when we introduced the Klein-Gordon equation in §4.1 we found the conserved quantity $\int \operatorname{Im} \phi^{*} \partial_{t} \phi d^{3} \mathbf{x}$. The quantum analogue of this is

$$
Q=i \int: \phi^{\dagger} \partial_{t} \phi-\partial_{t} \phi^{\dagger} \phi: d^{3} \mathbf{x}
$$

which, by a calculation like the one leading from (5.23) to (5.22), is equal to

$$
\int\left[a^{\dagger}(\mathbf{p}) a(\mathbf{p})-b^{\dagger}(\mathbf{p}) b(\mathbf{p})\right] \frac{d^{3} \mathbf{p}}{(2 \pi)^{3}}=N_{a}-N_{b}
$$

Here $N_{a}$ and $N_{b}$ are simply the number operators for particles and antiparticles (the operators with eigenvalue $k$ on the states with $k$ particles and $k$ antiparticles, respectively). It is easily verified that $N_{a}-N_{b}$ commutes with the Hamiltonian, so it represents a conserved quantity. (In fact, $N_{a}$ and $N_{b}$ commute with $H$ separately, but that is because we are dealing with a free field where nothing is happening. The net particle number $N_{a}-N_{b}$ is conserved in many situations where $N_{a}$ and $N_{b}$ are not.) If the particles and antiparticles carry unit electric charges of opposite signs, then $N_{a}-N_{b}$ is the total net charge (if the charge of the "particles" is taken as positive). In this situation, it is clear that the classical quantity $\operatorname{Im}\left[\phi^{*}\left(\partial_{t} \phi\right)\right]$ should be interpreted as the charge density of the field, and the corresponding expression with space derivatives, $\operatorname{Im}\left[\phi^{*} \nabla_{\mathbf{x}} \phi\right]$, as the current density.

One final remark. In classical mechanics, the Lagrangian and Hamiltonian are close cousins. They are used for different purposes, but they are similar in form, and it is usually easy to pass from one to the other. In quantum field theory, on the other hand, the Hamiltonian turns into an operator on some Hilbert space, while the Lagrangian remains a "classical," unquantized object. The Lagrangian is, so to speak, the classical root from which the quantum theory grows. This idea will be illustrated in the following chapters, and it will attain a deeper and more compelling significance in Chapter 8.

\subsection*{Spinor and vector fields}
Spinor fields. The construction of the Dirac spinor field whose quanta are spin- $\frac{1}{2}$ particles is similar to the construction of the complex scalar field, with only minor algebraic complications, so we shall be brief. As before, the idea is to think of a solution of the Dirac equation as a classical field, write it as a Fourier integral, and quantize it by replacing the coefficients by creation and annihilation operators.

We begin with the Fourier representation of a (reasonably) general solution of of the Dirac equation $i \gamma^{\mu} \partial_{\mu} \psi=m \psi$. That is, we take $\psi$ to be the inverse Fourier\\
transform of a $\mathbb{C}^{4}$-valued function $f$ on the extended mass shell $X_{m}^{+} \cup X_{m}^{-}$times the invariant measure $d \lambda(p)=d^{3} \mathbf{p} /(2 \pi)^{3} \omega_{\mathbf{p}}$, satisfying $p_{\mu} \gamma^{\mu} f=m f$ :\\
$\psi(x)=\int_{X_{m}^{+} \cup X_{m}^{-}} e^{-i p_{\mu} x^{\mu}} f(p) \frac{\sqrt{m} d^{3} \mathbf{p}}{(2 \pi)^{3} \omega_{\mathbf{p}}}=\int_{X_{m}^{+}}\left[e^{-i p_{\mu} x^{\mu}} f(p)+e^{i p_{\mu} x^{\mu}} f(-p)\right] \frac{\sqrt{m} d^{3} \mathbf{p}}{(2 \pi)^{3} \omega_{\mathbf{p}}}$.\\
(The $\sqrt{m}$ is as in (4.44); we shall comment on its significance later.) Our first job is to rewrite this in a way that displays the spin states explicitly and encodes the equation $p_{\mu} \gamma^{\mu} f=m f$.

Recall the discussion of the Dirac equation at the end of §4.4 - in particular, the action of $S L(2, \mathbb{C})$ on $X_{m}^{ \pm}$, which we denoted by $(A, p) \mapsto A^{\dagger-1} p$; the base points $p_{m}^{ \pm}=( \pm m, 0,0,0)$; the spin representation $\Phi$ of $S L(2, \mathbb{C})$ defined by (4.21); and the fact that if $v \in \mathbb{C}^{4}$ satisfies $p_{\mu} \gamma^{\mu} v=m v$, then $w=\Phi(A) v$ satisfies $p_{\mu}^{\prime} \gamma^{\mu} w=m w$ where $p^{\prime}=A^{\dagger-1} p$. Let

$$
u(\mathbf{0},+)=\left(\begin{array}{l}
1 \\
0 \\
1 \\
0
\end{array}\right), \quad u(\mathbf{0},-)=\left(\begin{array}{l}
0 \\
1 \\
0 \\
1
\end{array}\right), \quad v(\mathbf{0},+)=\left(\begin{array}{c}
1 \\
0 \\
-1 \\
0
\end{array}\right), \quad v(\mathbf{0},-)=\left(\begin{array}{c}
0 \\
1 \\
0 \\
-1
\end{array}\right) .
$$

Then $u(\mathbf{0}, \pm)$ and $v(\mathbf{0}, \pm)$, respectively, are bases for the subspaces of $\mathbb{C}^{4}$ defined by the equations $\gamma^{0} v=v$ and $-\gamma^{0} v=v$, i.e., $\left(p_{m}^{+}\right)_{\mu} \gamma^{\mu} v=m v$ and $\left(p_{m}^{-}\right)_{\mu} \gamma^{\mu} v=m v$, with $\gamma^{\mu}$ in the Weyl representation. Here the $\mathbf{0}$ is the space component of $p_{m}^{ \pm}$and the ± is a label for the spin state. (More precisely, $\pm \frac{1}{2}$ is the eigenvalue of the $x^{3}$-component of the spin, since the $u$ 's and $v$ 's are eigenvectors of $\left(\begin{array}{cc}\sigma & 0 \\ 0 & \sigma\end{array}\right)$ where $\sigma=\sigma_{3}$ is the 3rd Pauli matrix.) Now, for each $p \in X_{m}^{ \pm}$there is a unique positive Hermitian $B_{p} \in S L(2, \mathbb{C})$ such that $B_{p}^{\dagger-1} p_{m}^{ \pm}=p$ - or rather $\kappa\left(B_{p}^{\dagger-1}\right) p_{m}^{ \pm}=p$. (Indeed, let $A$ be any element of $S L(2, \mathbb{C})$ such that $A^{\dagger-1} p_{m}^{ \pm}=p$, and let $A=B U$ be its polar decomposition. Since $S U(2)$ is the subgroup that fixes $p_{m}^{ \pm}$, we can take $B_{p}=B$, and any other $A^{\prime}$ with $\left(A^{\prime}\right)^{\dagger-1} p_{m}^{ \pm}=p$ gives the same $B$ because $A^{\prime}=A U^{\prime}$ for some $U \in S U(2)$. The Lorentz transformation defined by $B_{p}$ is a pure boost with no rotation.) We set

$$
u(\mathbf{p}, \pm)=\Phi\left(B_{\left(\omega_{\mathbf{p}}, \mathbf{p}\right)}\right) u(\mathbf{0}, \pm), \quad v(\mathbf{p}, \pm)=\Phi\left(B_{\left(-\omega_{\mathbf{p}}, \mathbf{p}\right)}\right) v(\mathbf{0}, \pm)
$$

Then $u(\mathbf{p}, \pm)$ and $v(\mathbf{p}, \pm)$ are bases for the subspaces of $\mathbb{C}^{4}$ defined by the equation $p_{\mu} \gamma^{\mu} v=m v$ for $p \in X_{m}^{+}$and $p \in X_{m}^{-}$, respectively. Thus, any $\mathbb{C}^{4}$-valued function $f(p)$ on $X_{m}^{+}$or $X_{m}^{-}$that satisfies $p_{\mu} \gamma^{\mu} f(p)=m f(p)$ can be written as a linear combination of $u(\mathbf{p}, \pm)$ or $v(\mathbf{p}, \pm)$ with scalar coefficients that depend on $\mathbf{p}$.

We can therefore rewrite (5.28) as

$$
\psi(x)=\int \sum_{s= \pm} \sqrt{\frac{m}{2 \omega_{\mathbf{p}}}}\left[e^{-i p_{\mu} x^{\mu}} f(\mathbf{p}, s) u(\mathbf{p}, s)+e^{i p_{\mu} x^{\mu}} g(\mathbf{p}, s) v(\mathbf{p}, s)\right] \frac{d^{3} \mathbf{p}}{(2 \pi)^{3}}
$$

where $p=\left(\omega_{\mathbf{p}}, \mathbf{p}\right)$ and the coefficients $f(\mathbf{p}, s)$ and $g(\mathbf{p}, s)$ are scalar-valued functions in which a factor of $\sqrt{2 / \omega_{\mathbf{p}}}$ has been incorporated. Just as in the case of the complex scalar field, the quantum Dirac field is now obtained by replacing the coefficient $f(\mathbf{p}, s)$ by the annihilation operator $a(\mathbf{p}, s)$ for a particle with momentum $\mathbf{p}$ and spin state $s$, and the coefficient $g(\mathbf{p}, s)$ by the creation operator $b^{\dagger}(\mathbf{p}, s)$ for an antiparticle with momentum $\mathbf{p}$ and spin state $s$. The quanta of the Dirac field are going to be Fermions (a point to which we shall return shortly), so these operators\\
need to satisfy the anticommutation relations

$$
\begin{gathered}
\left\{a(\mathbf{p}, s), a^{\dagger}\left(\mathbf{p}^{\prime}, s^{\prime}\right)\right\}=\left\{b(\mathbf{p}, s), b^{\dagger}\left(\mathbf{p}^{\prime}, s^{\prime}\right)\right\}=(2 \pi)^{3} \delta\left(\mathbf{p}-\mathbf{p}^{\prime}\right) \delta_{s s^{\prime}}, \\
\left\{a(\mathbf{p}, s), a\left(\mathbf{p}^{\prime}, s^{\prime}\right)\right\}=\left\{b(\mathbf{p}, s), b\left(\mathbf{p}^{\prime}, s^{\prime}\right)\right\}=\left\{a^{\sharp}(\mathbf{p}, s), b^{b}\left(\mathbf{p}^{\prime}, s^{\prime}\right)\right\}=0,
\end{gathered}
$$

where $a^{\#}$ is either $a$ or $a^{\dagger}$ and $b^{b}$ is either $b$ or $b^{\dagger}$. Thus, the quantized Dirac field is

$$
\psi(x)=\int \sum_{s= \pm} \sqrt{\frac{m}{2 \omega_{\mathbf{p}}}}\left[e^{-i p_{\mu} x^{\mu}} u(\mathbf{p}, s) a(\mathbf{p}, s)+e^{i p_{\mu} x^{\mu}} v(\mathbf{p}, s) b^{\dagger}(\mathbf{p}, s)\right] \frac{d^{3} \mathbf{p}}{(2 \pi)^{3}}
$$

and its Dirac adjoint (recall (4.14)) is

$$
\bar{\psi}(x)=\int \sum_{s= \pm} \sqrt{\frac{m}{2 \omega_{\mathbf{p}}}}\left[e^{i p_{\mu} x^{\mu}} \bar{u}(\mathbf{p}, s) a^{\dagger}(\mathbf{p}, s)+e^{-i p_{\mu} x^{\mu}} \bar{v}(\mathbf{p}, s) b(\mathbf{p}, s)\right] \frac{d^{3} \mathbf{p}}{(2 \pi)^{3}}
$$

(5.31)-(5.33) can be expressed in rigorous mathematical language just as we did for the scalar field: the operators $a(\mathbf{p}, s), b(\mathbf{p}, s)$, and $\psi(x)$ are actually distributions with values in the operators on the finite-particle subspace of $\mathcal{F}=\mathcal{F}_{1} \otimes \mathcal{F}_{2}$, where $\mathcal{F}_{1}$ and $\mathcal{F}_{2}$ are copies of the Fermion Fock space over the state space for a single Dirac particle as discussed in §4.4, §4.5, and §5.2. These distributions are actually better behaved than their Bosonic analogues in the following respect: The smeared-out creation and annihilation operators, $a(f, s)=\int a(\mathbf{p}, s) f(\mathbf{p}) d^{3} \mathbf{p}$ and so forth, are bounded operators on $\mathcal{F}$ because of (4.57), and hence so are the spatially smearedout field operators $\psi(t, f)=\int \psi(t, \mathbf{x}) f(\mathbf{x}) d^{3} \mathbf{x}$.

At this point we insert a calculation of "spin sums" that will be used here and later. Recall that $u(\mathbf{p}, \pm)$ and $v(\mathbf{p}, \pm)$ are column vectors; their Dirac adjoints $\bar{u}(\mathbf{p}, \pm)=u(\mathbf{p}, \pm)^{\dagger} \gamma_{0}$ and $\bar{v}(\mathbf{p}, \pm)=v(\mathbf{p}, \pm)^{\dagger} \gamma^{0}$ are row vectors, and the products $u(\mathbf{p}, \pm) \bar{u}(\mathbf{p}, \pm)$ and $v(\mathbf{p}, \pm) \bar{v}(\mathbf{p}, \pm)$ are therefore $4 \times 4$ matrices. An easy calculation from the definitions (5.29) shows that

$$
\sum_{s= \pm} u(\mathbf{0}, s) \bar{u}(\mathbf{0}, s)=\gamma^{0}+I, \quad \sum_{s= \pm} v(\mathbf{0}, s) \bar{v}(\mathbf{0}, s)=\gamma^{0}-I .
$$

From this we can obtain the corresponding sums for arbitrary $\mathbf{p}$. By (5.30), with $B=B_{\left(\omega_{\mathbf{p}}, \mathbf{p}\right)}$ we have

$$
\sum_{s= \pm} u(\mathbf{p}, s) \bar{u}(\mathbf{p}, s)=\sum_{s= \pm} \Phi(B) u(\mathbf{0}, s) u(\mathbf{0}, s)^{\dagger} \Phi(B)^{\dagger} \gamma^{0}
$$

But by (4.12) and (4.15) we have $\Phi(B)^{\dagger}=\Phi\left(B^{\dagger}\right)=\gamma^{0} \Phi(B)^{-1} \gamma^{0}$, so

$$
\sum_{s= \pm} u(\mathbf{p}, s) \bar{u}(\mathbf{p}, s)=\sum_{s= \pm} \Phi(B) u(\mathbf{0}, s) \bar{u}(\mathbf{0}, s) \Phi(B)^{-1}=\Phi(B) \gamma^{0} \Phi(B)^{-1}+I .
$$

Moreover, if $L=\kappa(B)^{-1} \in S O(1,3)$, by (4.22) we have $\Phi(B) \gamma^{0} \Phi(B)^{-1}=L_{\mu}^{0} \gamma^{\mu}$, and in view of the fact that $B$ is Hermitian, the defining condition $B^{\dagger-1} p_{m}^{+}=p$ just means that $m L_{\mu}^{0}=p_{\mu}$. Therefore, after an isomorphic calculation with the $v$ 's, we conclude that

$$
\sum_{s= \pm} u(\mathbf{p}, s) \bar{u}(\mathbf{p}, s)=m^{-1} p_{\mu} \gamma^{\mu}+I, \quad \sum_{s= \pm} v(\mathbf{p}, s) \bar{v}(\mathbf{p}, s)=m^{-1} p_{\mu} \gamma^{\mu}-I
$$

with $p_{0}=\omega_{\mathbf{p}}$.\\
We recall from §4.1 that the Lagrangian density associated to the free Dirac equation is

$$
\mathcal{L}=i \bar{\psi} \gamma^{\mu} \partial_{\mu} \psi-m \bar{\psi} \psi,
$$

so the field $\pi$ canonically conjugate to $\psi$ is

$$
\pi=\frac{\partial \mathcal{L}}{\partial\left(\partial_{t} \psi\right)}=i \psi^{\dagger}
$$

The fields $\psi$ and $\pi$ satisfy canonical anticommutation relations analogous to (5.20):

$$
\begin{aligned}
\left\{\psi_{j}(t, \mathbf{x}), \pi_{k}(t, \mathbf{y})\right\}=i \delta(\mathbf{x}-\mathbf{y}) \delta_{j k} & \\
\left\{\psi_{j}(t, \mathbf{x}), \psi_{k}(t, \mathbf{y})\right\}=\left\{\pi_{j}(t, \mathbf{x}), \pi_{k}(t, \mathbf{y})\right\}=0 & (j, k=1, \cdots, 4)
\end{aligned}
$$

To prove this, observe that by (5.32), $\left\{\psi_{j}(t, \mathbf{x}), i \psi_{k}^{\dagger}(t, \mathbf{y})\right\}$ is equal to

$$
i \iint \sum_{s, s^{\prime}= \pm} \frac{m e^{-i \omega_{\mathbf{p}} t+i \mathbf{p} \cdot \mathbf{x}} e^{i \omega_{\mathbf{p}^{\prime}} t-i \mathbf{p}^{\prime} \cdot \mathbf{y}}}{2 \sqrt{\omega_{\mathbf{p}} \omega_{\mathbf{p}^{\prime}}}} u(\mathbf{p}, s)_{j} u\left(\mathbf{p}^{\prime}, s^{\prime}\right)_{k}^{\dagger}\left\{a(\mathbf{p}, s), a^{\dagger}\left(\mathbf{p}^{\prime}, s^{\prime}\right)\right\} \frac{d^{2} \mathbf{p} d^{3} \mathbf{p}^{\prime}}{(2 \pi)^{6}}
$$

plus a similar term with $u$ 's and $a$ 's replaced by $v$ 's and $b$ 's and the signs in the exponents reversed. By (5.31), (5.34), and the fact that

$$
u(\mathbf{p}, s)_{j} u(\mathbf{p}, s)_{k}^{\dagger}=u(\mathbf{p}, s)_{j}\left[\bar{u}(\mathbf{p}, s) \gamma^{0}\right]_{k}=\left[u(\mathbf{p}, s) \bar{u}(\mathbf{p}, s) \gamma^{0}\right]_{j k}
$$

integration over $\mathbf{p}^{\prime}$ and summation over $s$ yield

$$
\begin{aligned}
\left\{\psi_{j}(t, \mathbf{x}), i \psi_{k}^{\dagger}(t, \mathbf{y})\right\}= & i \int \frac{m}{2 \omega_{\mathbf{p}}} e^{i \mathbf{p} \cdot(\mathbf{x}-\mathbf{y})}\left[\left(m^{-1} p_{\mu} \gamma^{\mu}+m I\right) \gamma^{0}\right]_{j k} \frac{d^{3} \mathbf{p}}{(2 \pi)^{3}} \\
& +i \int \frac{m}{2 \omega_{\mathbf{p}}} e^{-i \mathbf{p} \cdot(\mathbf{x}-\mathbf{y})}\left[\left(m^{-1} p_{\mu} \gamma^{\mu}-m I\right) \gamma^{0}\right]_{j k} \frac{d^{3} \mathbf{p}}{(2 \pi)^{3}}
\end{aligned}
$$

The substitution $\mathbf{p} \rightarrow-\mathbf{p}$ in the second integral shows that the terms involving $m I$ and $p_{\mu} \gamma^{\mu}$ with $\mu>0$ in the two integrals cancel out. The terms involving $p_{0} \gamma^{0}$ add up, and since $\left(\gamma^{0}\right)^{2}=I$ and $p_{0}=\omega_{\mathbf{p}}$, the integrals reduce to

$$
\left\{\psi_{j}(t, \mathbf{x}), i \psi_{k}^{\dagger}(t, \mathbf{y})\right\}=i \int e^{i \mathbf{p} \cdot(\mathbf{x}-\mathbf{y})} \delta_{j k} \frac{d^{3} \mathbf{p}}{(2 \pi)^{3}}=i \delta(\mathbf{x}-\mathbf{y}) \delta_{j k}
$$

The other equations in (5.35) are obtained similarly.\\
The factor $\sqrt{m}$ in (5.32) that we carried over from (4.44) is what gives the proper normalization in (5.35). One can also see that the Dirac field (5.32) should contain a factor with dimensions $\left[m^{1 / 2}\right]$ in comparison to the scalar field (5.15) from the fact that the mass terms in their respective Lagrangians are $m \bar{\psi} \psi$ and $m^{2}|\phi|^{2}$, both of which are supposed to yield an energy density. However, there are even more variations on the formula (5.32) in the literature than there are for the scalar field, because one can use different normalizations not only for $a$ and $b$ but also for the spinors $u$ and $v$. In particular, some people incorporate the factor of $\sqrt{m}$ into $u$ and $v$.

The Hamiltonian density for the Dirac field $\psi$ is

$$
\mathcal{H}=\pi \partial_{t} \psi-\mathcal{L}=\pi \partial_{t} \psi=i \psi^{\dagger} \partial_{t} \psi
$$

since $\mathcal{L}=\mathcal{L}\left(\psi, \partial_{\mu} \psi\right)=0$ when $\psi$ is a solution of the Dirac equation. After substituting (5.32) for $\psi$ in this expression, a short calculation involving the Fourier inversion formula $\int e^{i \mathbf{p} \cdot \mathbf{x}} d^{3} \mathbf{x}=(2 \pi)^{3} \delta(\mathbf{p})$ shows that the Hamiltonian is

$$
H=\int \mathcal{H} d^{3} \mathbf{x}=\int \sum_{\sigma} \omega_{\mathbf{p}}\left[a^{\dagger}(\mathbf{p}, s) a(\mathbf{p}, s)-b(\mathbf{p}, s) b^{\dagger}(\mathbf{p}, s)\right] \frac{d^{3} \mathbf{p}}{(2 \pi)^{3}}
$$

At first sight this looks like it might not be bounded below, which would be most unfortunate. But by (5.31) we have

$$
b(\mathbf{p}, s) b^{\dagger}\left(\mathbf{p}^{\prime}, s^{\prime}\right)=-b^{\dagger}\left(\mathbf{p}^{\prime}, s^{\prime}\right) b(\mathbf{p}, s)+\delta\left(\mathbf{p}-\mathbf{p}^{\prime}\right) \delta_{s s^{\prime}}
$$

so the cost of throwing away an infinite constant, as we did for the scalar field, we can rewrite (5.36) as

$$
H=\int \sum_{\sigma} \omega_{\mathbf{p}}\left[a^{\dagger}(\mathbf{p}, s) a(\mathbf{p}, s)+b^{\dagger}(\mathbf{p}, s) b(\mathbf{p}, s)\right] \frac{d^{3} \mathbf{p}}{(2 \pi)^{3}}
$$

which is manifestly positive. (5.37) is, of course, the Wick-ordered form of (5.36) - with the understanding that where Fermionic operators are concerned, an interchange of two operators in the Wick-ordering process introduces a minus sign.

This brings up an important point. If we had tried to quantize the Dirac field as a Boson field, so that the anticommutation relations (5.31) were replaced by commutation relations, this trick would not work and we would end up with energies that are not bounded below. For exactly the same reason, if we had tried to quantize the scalar Klein-Gordon field as a Fermion field, we would be in trouble after arriving at the expressions (5.27). This is one aspect of the spin-statistics theorem, which we shall discuss further in §5.5 and §6.5.

Vector fields. We would now like to quantize the electromagnetic field. The first point that has to be appreciated is that the field to be quantized is not the electromagnetic field $F_{\mu \nu}$ itself but its potential $A_{\mu}$. For one thing, the vector field $A_{\mu}$ rather than the tensor field $F_{\mu \nu}$ is the appropriate thing to yield spin-one quanta, which photons are known by experiment to be. Moreover, the AharonovBohm effect (see Sakurai [103]) shows that the potential $A_{\mu}$ has a real effect on quantum phenomena even in regions where $F_{\mu \nu}$ vanishes. As Weinberg [131], p. 340, observes, one could construct a consistent theory "by demanding that all interactions involve only $F_{\mu \nu}(x)=\partial_{\mu} A_{\nu}(x)-\partial_{\nu} A_{\mu}(x)$ and its derivatives, not $A_{\mu}(x)$, but this is not the most general possibility, and not the one realized in nature."

But various difficulties arise from the fact that $A_{\mu}$ is defined only up to gauge transformations and from the closely related fact that the photon is massless. To get a handle on the situation, let us first consider an analogous field whose quanta have nonzero mass. (Such fields describe the vector bosons that appear in the WeinbergSalam theory of weak interactions.) On the classical level, such a field is a $\mathbb{C}^{4}$-valued function $A$ on space-time that transforms under Lorentz transformations by the identity representation of the Lorentz group: $A(x) \mapsto L A\left(L^{-1} x\right)$. Its Lagrangian is

$$
\mathcal{L}=-\frac{1}{4} F_{\mu \nu} F^{\mu \nu}+\frac{1}{2} m^{2} A_{\mu} A^{\mu} \quad\left(F^{\mu \nu}=\partial^{\mu} A^{\nu}-\partial^{\nu} A^{\mu}\right)
$$

and the resulting field equations are a massive version of Maxwell's equations known as the Proca equations:

$$
\partial_{\mu} F^{\mu \nu}+m^{2} A^{\nu}=0
$$

Since $\partial_{\mu} \partial_{\nu} F^{\mu \nu}=0$, (5.39) implies that $\partial_{\nu} A^{\nu}=0$, and since $\partial_{\mu} F^{\mu \nu}=\partial^{2} A^{\nu}- \partial^{\nu}\left(\partial_{\mu} A^{\mu}\right)$, we see that (5.39) is equivalent to the Klein-Gordon equation plus the condition of zero 4-divergence:

$$
\partial^{2} A^{\nu}+m^{2} A^{\nu}=0, \quad \partial_{\nu} A^{\nu}=0
$$

In the limit $m \rightarrow 0$, this yields Maxwell's equations with the Lorentz gauge condition. On the Fourier transform side, the Klein-Gordon equation implies that $\widehat{A}$ lives on the mass shells $X_{m}^{ \pm}$, and the equation $\partial_{\nu} A^{\nu}=0$ becomes the linear constraint $p_{\nu} \widehat{A}^{\nu}(p)=0$, which reduces the number of independent components of $A$ from 4 to 3 . With this in mind, one sees that the action of the rotation group on $A$ is given by its (complexified) identity representation on $\mathbb{C}^{3}$, i.e., the representation associated to spin one.

The quantization of the massive vector field proceeds in much the same way as that of the Dirac field. The result is

$$
A^{\nu}(x)=\int \sum_{j=1}^{3} \frac{1}{\sqrt{2 \omega_{\mathbf{p}}}}\left[e^{-i p_{\mu} x^{\mu}} \boldsymbol{\epsilon}^{\nu}(\mathbf{p}, j) a(\mathbf{p}, j)+e^{i p_{\mu} x^{\mu}} \boldsymbol{\epsilon}^{\nu}(\mathbf{p}, j)^{*} b^{\dagger}(\mathbf{p}, j)\right] \frac{d^{3} \mathbf{p}}{(2 \pi)^{3}}
$$

where $a$ and $b$ represent annihilation operators for particles and antiparticles - the possibility $b=a$ for neutral particles is allowed - and the $\boldsymbol{\epsilon}^{\nu}(\mathbf{p}, j)$ are polarization vectors. The nontrivial commutation relations for $a$ and $b$ are the usual ones,

$$
\left[a(\mathbf{p}, j), a^{\dagger}\left(\mathbf{p}^{\prime}, j^{\prime}\right)\right]=\left[b(\mathbf{p}, j), b^{\dagger}\left(\mathbf{p}^{\prime}, j^{\prime}\right)\right]=(2 \pi)^{3} \delta\left(\mathbf{p}-\mathbf{p}^{\prime}\right) \delta_{j j^{\prime}}
$$

and the polarization vectors are specified as follows. For $\mathbf{p}=\mathbf{0},\left\{\boldsymbol{\epsilon}^{\nu}(\mathbf{0}, j): j=\right. 1,2,3\}$ is a given orthonormal basis for $\{0\} \times \mathbb{C}^{3}$ (the same one for each $\nu$ ), and then $\boldsymbol{\epsilon}^{\nu}(\mathbf{p}, j)=B_{\mu}^{\nu} \boldsymbol{\epsilon}^{\mu}(\mathbf{0}, j)$ where $B \in S O^{\uparrow}(1,3)$ is the pure boost that takes $p_{m}^{+}=(m, \mathbf{0})$ to $\left(\omega_{\mathbf{p}}, \mathbf{p}\right)$ (similarly to the spinor coefficients of the Dirac field). (This automatically guarantees that $p_{\nu} \boldsymbol{\epsilon}^{\nu}(\mathbf{p}, j)=0$.)

Now, what happens if we let $m \rightarrow 0$ ? The Proca equations (5.40) turn into Maxwell's equations with the Lorentz gauge condition, which is fine. But the singleparticle states for massive vector particles have three independent components, corresponding to the three polarization vectors in (5.41), whereas the single-particle states for massless ones have only two, specified by the positive or negative helicity along the direction of motion. (On the classical level, this corresponds to the fact that the electromagnetic field admits only transverse oscillations, not longitudinal ones.) Thus if we try to quantize the electromagnetic field in the form (5.41) (with $b=a$ ), we have too many components, and some of them will represent unphysical "ghost fields" that create "ghost states" whose contributions to later computations must ultimately cancel out. It is indeed possible to proceed in this way, picking one's way among the technical obstacles carefully; the result is known as Gupta-Bleuler quantization.

Another possibility is to discard Lorentz invariance in favor of a quantization that will yield the physical fields more directly. On the classical level, the Lorentz gauge condition $\partial_{\mu} A^{\mu}=0$ does not specify $A^{\mu}$ completely; one can still add a term $\partial^{\mu} \chi$ where $\partial^{2} \chi=0$. For a free field, in the absence of charges, one can thereby arrange to make $A^{0}=0$, that is, to put $A$ into the Coulomb gauge or radiation gauge:

$$
A=(0, \mathbf{A}), \quad \nabla \cdot \mathbf{A}=0
$$

Quantization of this field yields

$$
\mathbf{A}(x)=\int \sum_{j=1}^{2} \frac{1}{\sqrt{2|\mathbf{p}|}}\left[e^{-i p_{\mu} x^{\mu}} \boldsymbol{\epsilon}(\mathbf{p}, j) a(\mathbf{p}, j)+e^{i p_{\mu} x^{\mu}} \boldsymbol{\epsilon}(\mathbf{p}, j)^{*} a^{\dagger}(\mathbf{p}, j)\right] \frac{d^{3} \mathbf{p}}{(2 \pi)^{3}}
$$

where $\boldsymbol{\epsilon}(\mathbf{p}, 1)$ and $\boldsymbol{\epsilon}(\mathbf{p}, 2)$ are an orthonormal basis for $\{\boldsymbol{\epsilon}: \boldsymbol{\epsilon} \cdot \mathbf{p}=0\}$. (One can choose these vectors to depend smoothly on $\mathbf{p} \in \mathbb{R}^{3} \backslash\{0\}$ by allowing them to be complex - that is, by allowing general elliptical polarizations rather than just linear ones. Note also that $\omega_{\mathbf{p}}=|\mathbf{p}|$ since $m=0$ here.) In effect, the condition $\boldsymbol{\epsilon} \cdot \mathbf{p}=0$ means that one is keeping only that part of the field (5.41) whose polarization is transverse to the direction of motion.

Using the Coulomb gauge, of course, destroys the manifest Lorentz invariance of the theory. Moreover, both this procedure and the Lorentz-invariant quantization with ghost fields present some technical problems in connection with constructing canonically conjugate fields for use in constructing the Hamiltonian. There is no general agreement among physics texts about the best way of handling these problems, and as mathematical tourists we are better off not getting into these muddy waters. The most essential information that must be extracted from the fields is the associated particle propagators; we shall address that problem in §6.5 and §8.3.

Concluding remarks. One can construct a free quantum field whose quanta have any mass $m>0$ and spin $s \in \frac{1}{2} \mathbb{Z}_{+}$. It acts on the Fock space over the one-particle state space associated to the representation $\pi_{m, s}$ of §4.4 (or the tensor product of two copies of the Fock space for particle-antiparticle pairs). The general idea is the same as in the cases $s=0, \frac{1}{2}, 1$ that we have described; only the algebra is different. See Streater and Wightman [116] and Weinberg [131].

Let us review what we have accomplished. We started from the idea of a classical field, that is, a function $\phi$ on (some region in) space-time whose value at a point $x$ is an observable quantity - say, a force, a velocity, a temperature - that can be determined by measurements performed at $x$. One would expect quantization to yield a function $\Phi$ on space-time whose value at $x$ is the quantum observable - i.e., self-adjoint operator - corresponding to the classical observable $\phi(x)$. Thus, when the system is in a state $v$ the expectation value $\langle v| \Phi(x)|v\rangle$ should be somehow closely related to the classical field $\phi(x)$ in the corresponding classical state.

What we seem to have ended up with, however, is rather different. The fact that we need operator-valued distributions rather than functions is a technical obstacle rather than a conceptual one, for even classically the notion of making measurements at a single point specified to infinite precision is an idealization. What is more serious is that the values of the quantum fields we have constructed are self-adjoint (or, more loosely, Hermitian) only when the quanta of the field are their own antiparticles; in other cases the values of the fields cannot directly represent observable quantities. This is perhaps just as well, for what would we be observing? The "classical fields" we have quantized to produce the scalar and Dirac quantum fields have no meaning in classical physics; even for the electromagnetic field, what we have quantized is not the directly observable field $F_{\mu \nu}$ but the more abstract potential $A_{\mu}$. Rather, these "classical fields" arise as frameworks for describing single quantum particles, and the quantum fields we have constructed turn out to be machines for creating and destroying such particles.

At this point readers may well be wondering what these quantum fields have to do with physical reality. So far, the answer is: not much. We have constructed only free fields, whereas all the real physics comes from interactions. Getting to the interesting stuff will require a little more patience and a lot more willingness to accept mathematical incompleteness. We shall take the plunge in the next chapter.

\subsection*{The Wightman axioms}
The success of quantum electrodynamics naturally led to attempts to develop the theory in a mathematically rigorous way and to generalize it to encompass other kinds of quantum fields. In the late 1950s Wightman and Gårding formulated a list of basic properties that any physically reasonable and mathematically welldefined quantum field theory should have. This list has come to be known as the Gårding-Wightman axioms or the Wightman axioms.

The ingredients for the Wightman axioms are a Hilbert space $\mathcal{F}$ equipped with\\
i. a dense subspace $\mathcal{D}$,\\
ii. a set of operator-valued distributions $\phi_{1}, \ldots, \phi_{N}$ on $\mathbb{R}^{4}$ (the meaning of this is explained more precisely in Axiom 1 below),\\
iii. a unitary representation $U$ of the double cover $\mathbb{R}^{4} \ltimes S L(2, \mathbb{C})$ of the restricted Poincaré group on $\mathcal{F}$, and\\
iv. a (nonunitary) representation $S$ of $S L(2, \mathbb{C})$ on $\mathbb{C}^{N}$.

Since the unitary operators $U(a, I)\left(a \in \mathbb{R}^{4}\right)$ all commute, they have a common spectral resolution; that is, there is a projection-valued measure $E$ on $\mathbb{R}^{4}$ such that $U(a, I)=\int e^{-i p_{\mu} a^{\mu}} d E(p)$.

The axioms are as follows. All of them should look rather familiar except perhaps the last one, which we shall discuss later.

\begin{enumerate}
  \item Each $\phi_{n}$ is a map from the Schwartz space $\mathcal{S}\left(\mathbb{R}^{4}\right)$ to the set of (perhaps unbounded) linear operators on $\mathcal{F}$ such that\\
i. for all $f \in \mathcal{S}\left(\mathbb{R}^{4}\right)$ we have\\
$\mathcal{D} \subset \operatorname{Dom}\left(\phi_{n}(f)\right) \cap \operatorname{Dom}\left(\phi_{n}(f)^{\dagger}\right), \quad \phi_{n}(f) \mathcal{D} \cup \phi_{n}(f)^{\dagger} \mathcal{D} \subset \mathcal{D} ;$\\
ii. for all $\xi, \eta \in \mathcal{D}$, the map $f \mapsto\langle\xi| \phi_{n}(f)|\eta\rangle$ is a tempered distribution on $\mathbb{R}^{4}$.
  \item There is a unit vector $\Omega \in D$ (the "vacuum state"), unique up to a phase factor, such that $U(a, A) \Omega=\Omega$ for all $(a, A) \in \mathbb{R}^{4} \ltimes S L(2, \mathbb{C})$.
  \item The linear span of the set of vectors of the form $\phi_{n_{1}}\left(f_{1}\right) \cdots \phi_{n_{k}}\left(f_{k}\right) \Omega$ with $k \geq 0, n_{j} \in\{1, \ldots, N\}$, and $f_{j} \in \mathcal{S}\left(\mathbb{R}^{4}\right)$ is dense in $\mathcal{H}$.
  \item For any $f \in \mathcal{S}\left(\mathbb{R}^{4}\right),(a, A) \in \mathbb{R}^{4} \ltimes S L(2, \mathbb{C})$, and $n \in\{1, \ldots, N\}$,
\end{enumerate}

$$
U(a, A) \phi_{n}(f) U(a, A)^{-1}=\sum_{m=1}^{N} S\left(A^{-1}\right)_{n}^{m} \phi_{m}((a, A) \cdot f)
$$

as operators on $\mathcal{D}$, where $[(a, A) \cdot f](x)=f\left(\kappa(A)^{-1}(x-a)\right)$ and $\kappa$ : $S L(2, \mathbb{C}) \rightarrow S O^{\uparrow}(1,3)$ is the covering map. If we adopt the notational pretense that $\phi_{n}$ is a function rather than a distribution, this means that

$$
U(a, A) \phi_{n}(x) U(a, A)^{-1}=\sum_{m=1}^{N} S\left(A^{-1}\right)_{n}^{m} \phi_{m}(A x+a) .
$$

\begin{enumerate}
  \setcounter{enumi}{4}
  \item The support of the projection-valued measure $E$ is contained in the region $\left\{\left(p^{0}, \mathbf{p}\right): p^{0} \geq|\mathbf{p}|\right\}$ inside the forward light cone. (Equivalently: if $P_{\mu}$ is the infinitesimal generator of the one-parameter group $t \mapsto U\left(t e_{\mu}, I\right)$, where $e_{0}, \ldots, e_{3}$ are the standard basis for $\mathbb{R}^{4}$, then $P_{0}$ and $P_{0}^{2}-\sum_{1}^{3} P_{j}^{2}$ are both positive operators.)
  \item If $f$ and $g$ are in $\mathcal{S}\left(\mathbb{R}^{4}\right)$ and the supports of $f$ and $g$ are spacelike separated (i.e., $x-y$ is spacelike whenever $f(x) g(y) \neq 0$ ), and $m, n \in\{1, \ldots, N\}$, then as operators on $\mathfrak{D}$,
\end{enumerate}

$$
\begin{aligned}
\text { either } & {\left[\phi_{m}(f), \phi_{n}(g)\right]=\left[\phi_{m}(f), \phi_{n}(g)^{\dagger}\right]=0 } \\
\text { or } & \left\{\phi_{m}(f), \phi_{n}(g)\right\}=\left\{\phi_{m}(f), \phi_{n}(g)^{\dagger}\right\}=0
\end{aligned}
$$

(Which of these two possibilities occurs may depend on $m$ and $n$.) Again, if we pretend that the $\phi_{n}$ are functions, this means that $\phi_{m}(x)$ (and $\left.\phi_{m}(x)^{\dagger}\right)$ commutes or anticommutes with $\phi_{n}(y)$ and $\phi_{n}(y)^{\dagger}$ whenever $x-y$ is spacelike.\\
This framework can accomodate any number of fields representing any number of particle types, because there is no assumption that the representation $S$ should be irreducible. If $S$ is a direct sum of irreducible representations $S_{k}$ on $\mathbb{C}^{n_{k}}$ with $k=1, \ldots, K$ and $n_{1}+\cdots n_{K}=N$, the $N$-tuple of fields $\phi_{n}$ separates out into $n_{k}$-tuples, each of which may be the components of the field representing one of the particles of the theory.

The free massive scalar, Dirac, and vector fields that we have discussed earlier in this chapter are all examples of systems that satify the Wightman axioms, as are the free fields of arbitrary mass $m>0$ and spin $s \in \frac{1}{2} \mathbb{Z}_{+}$. The Hilbert space $\mathcal{F}$ is the (Boson or Fermion) Fock space over the appropriate single-particle state space $\mathcal{H}$, and the dense subspace $\mathcal{D}$ is the finite-particle subspace. The vacuum $\Omega$ is the no-particle state, i.e., a unit vector in $\bigotimes^{0} \mathcal{H}=\mathbb{C}$, and the closed linear span of the vectors $\phi_{n_{1}}\left(f_{1}\right) \cdots \phi_{n_{k}}\left(f_{k}\right) \Omega$ with $k \leq K$ is the sum of the $k$-particle spaces for $k \leq K$. The representation $U$ is given by $U(a, A)=\mathcal{F}(\pi(a, A))$, where $\pi$ is the appropriate irreducible representation of $\mathbb{R}^{4} \ltimes S L(2, \mathbb{C})$ on $\mathcal{H}$ as described in §4.4 and $\mathcal{F}(\pi(\cdot))$ is the corresponding representation on $\mathcal{F}$ given by (4.53). The representation $S$ is the trivial representation on $\mathbb{C}$ for the scalar field, the representation $\Phi$ defined by (4.21) for the Dirac field, and the covering map $\kappa: S L(2, \mathbb{C}) \rightarrow S O^{\uparrow}(1,3)$ for the vector field. The joint spectrum of the position-momentum operators on the single-particle space $\mathcal{H}$ is the mass shell $X_{m}^{+}$, and hence the joint spectrum on the $k$-particle space is $\left\{p_{1}+\cdots+p_{k}: p_{1}, \ldots, p_{k} \in X_{m}^{+}\right\}$. All of these sets lie inside the forward light cone, and the support of the measure $E$ in Axiom 5 is the closure of their union. The verification of Axiom 6 is a calculation that we shall perform in §6.5. (See Streater and Wightman [116] or Bogolubov et al. [11], [12]; ${ }^{2}$ also Reed and Simon [95] for a thorough treatment of the free scalar field.)

Let us examine the meaning of Axiom 6, known as the microscopic causality condition, in more detail. The idea is that for $x \in \mathbb{R}^{4}$ the fields $\phi_{n}(x)$ are supposed to represent phenomena that take place at $x$; more precisely, if $f \in \mathcal{S}\left(\mathbb{R}^{4}\right)$, the $\phi_{n}(f)$ are supposed to represent phenomena that take place in the space-time support of $f$. (For example, if one unravels all the Fourier analysis in the definition of the free scalar field $\phi$, one finds that $\phi(t, \mathbf{x}) \Omega$ is the state containing one particle located at position $\mathbf{x}$ at time $t$.) If $x$ and $y$ are space-like separated, nothing that happens at $x$ can influence what happens at $y$ and vice versa, so (according to the discussion of uncertainty in §3.3) observables connected with phenomena at $x$ and $y$ must commute. This argument does not immediately imply (5.43), for as we

\footnotetext{${ }^{2}$ As with many Russian names, there are several possible spellings of "Bogolubov" in the Roman alphabet; the most common variations involve the replacement of the "u" by "iu," "ju," or "yu." We use the spelling in the works just cited.
}
have observed in the preceding section, the field operators usually do not represent observable quantities. However, as we shall see in the next chapter, quantities of direct physical significance are constructed out of field operators, and for them these (anti)commutation relations are important. ${ }^{3}$ In particular, they are needed to guarantee the Lorentz invariance of the scattering matrix; according to Weinberg [131], p. 198, this is the real reason to insist on microscopic causality.

In this setting one can prove a rigorous form of the spin-statistics theorem that particles of integer spin must be Bosons and particles of half-integer spin must be Fermions. ${ }^{4}$ The precise statement is as follows. Suppose $\phi_{1}, \ldots, \phi_{N}$ are the fields of a theory satisfying the Wightman axioms that describes one or several particle types; thus each $\phi_{n}$ is a component of an irreducible subset of fields that describes one of the particle types. If $\phi_{n}$ is a component of a field describing a particle of half-integer spin and $\left[\phi_{n}(f), \phi_{n}(g)^{\dagger}\right]=0$ whenever $f$ and $g$ have spacelike-separated supports, then $\phi_{n}(f) \Omega=0$ for all $f$; likewise if $\phi_{n}$ is a component of a field describing a particle of integer spin and $\left\{\phi_{n}(f), \phi_{n}(g)^{\dagger}\right\}=0$. In either case, if the different fields of the theory all either commute or anticommute (normally the case), it follows that $\phi_{n}=0$.

The Wightman axioms have served as the foundation for a large body of research in the mathematically rigorous theory of quantum fields. We shall not pursue this subject here, but we would be remiss not to mention an especially fundamental result, the so-called $P C T$ theorem. This states that any quantum field theory satisfying the axioms is invariant under the combined operation $P C T$ (the factors may be permuted in any order) where $P$ and $T$ are the space and time inversion operators coming from the action of the Lorentz group and $C$ is the charge-conjugation operator that interchanges particles and antiparticles. (We discussed the latter for Dirac wave functions in §4.2 and for scalar fields in §5.2; the definition in general is similar.) A half-century ago it was believed that these three operators individually should be symmetries of any reasonable physical theory, so that the PCT theorem might have seemed like much ado about the obvious. But then it was discovered that the weak interaction has a definite "handedness" and so is not $P$-invariant (see §9.4), and more recently, experimentalists found that certain meson decays are not $C P$-invariant - at which point the significance of the PCT theorem could no longer be doubted.

A different but related mathematical framework in which rigorous quantum field theory can be studied is provided by the notion of algebras of local observables. The fundamental data here are:\\
i. a $\mathrm{C}^{*}$ algebra $\mathcal{A}$ of operators on a Hilbert space $\mathcal{H}$,\\
ii. a $\mathrm{C}^{*}$ subalgebra $\mathcal{A}(O) \subset \mathcal{A}$ for each bounded open set $O \subset \mathbb{R}^{4}$, and\\
iii. a representation $\alpha$ of the restricted Poincaré group $\mathcal{P}_{0}$ as automorphisms of $\mathcal{A}$, subject to the following axioms:

\begin{enumerate}
  \item $\mathcal{A}\left(O_{1}\right) \subset \mathcal{A}\left(O_{2}\right)$ whenever $O_{1} \subset O_{2}$.
  \item $\bigcup_{O} \mathcal{A}(O)$ is dense in $\mathcal{A}$.
  \item $\alpha(g)(\mathcal{A}(O))=\mathcal{A}(g(O))$ for all $g \in \mathcal{P}_{0}$.
\end{enumerate}

\footnotetext{${ }^{3}$ These quantities almost always involve products of even numbers of Fermion field operators, and anticommutation relations for the latter yield commutation relations for the products.\\
${ }^{4}$ This result is due to Fierz and Pauli; the rigorous proof in the framework of the Wightman axioms is due to Burgoyne and to Lüders and Zumino. See Streater and Wightman [116] for references.
}
4. $\left[A_{1}, A_{2}\right]=0$ for all $A_{1} \in \mathcal{A}\left(O_{1}\right)$ and $A_{2} \in \mathcal{A}\left(O_{2}\right)$ whenever $O_{1}$ and $O_{2}$ are spacelike separated.\\
For example, if $\left\{\phi_{1}, \ldots, \phi_{N}\right\}$ is a set of fields satisfying the Wightman axioms for Bosons (i.e., with commutators in (5.43)), one can take $\mathcal{A}(O)$ to be the $\mathrm{C}^{*}$ algebra generated by the operators $\phi_{j}(f)$ with $1 \leq j \leq N$ and $f \in C_{c}^{\infty}(O)$, with $\alpha(g) A=U(g) A U(g)^{-1}$. The study of quantum fields in this framework is known as local or algebraic quantum field theory. (In spite of initial appearances, it can also accomodate Fermion fields.)

There is an extensive mathematical theory based on the Wightman or localobservable axioms. The reader who wishes to learn more about it may consult Streater and Wightman [116], Jost [70], Bogolubov et al. [11], [12], and Araki [3]. All of these contain proofs of the spin-statistics and PCT theorems. In addition, Haag [60] gives a comprehensive survey of the subject, with many arguments sketched or omitted.

The bad news is that it has turned out to be a remarkably difficult task to construct examples of field theories that satisfy the Wightman axioms and have nontrivial interactions. Attempts to produce such examples in physical space-time $\mathbb{R}^{4}$ have yet to succeed, so most of the work has dealt with field theories in $\mathbb{R}^{d}$ (equipped with the Lorentz group $O(1, d-1)$ ) for $d=2$ or 3 , where the singularities that plague quantum fields tend to be somewhat tamer. (It is easy to adapt the Wightman axioms to any number of space dimensions.) Even there, things are not easy, and for several years after the formulation of the Wightman axioms it was an open question whether they were satisfied by any systems other than the free fields and minor modifications thereof. The first example, a self-interacting scalar field in dimension $d=2$, was constructed by Glimm and Jaffe. Since then a variety of other examples have been explored: other self-interacting fields, a Dirac field and a scalar field with a Yukawa interaction, and various gauge fields - but all in dimensions 2 and 3. Much of this work has made use of the functional integral approach, and we shall say more about it in Chapter 8. There is a lot of interesting mathematics in it as well as interesting physics, including connections to other areas of physics such as statistical mechanics, but none of it does much to help with the down-to-earth calculations of quantum electrodynamics that we are aiming towards. We therefore leave it aside and refer the reader to the appendix of Streater and Wightman [116] and Rosen [101] for concise surveys and to Glimm and Jaffe [55] and Simon [113] for more extensive accounts.

\section*{CHAPTER 6}
\section*{Quantum Fields with Interactions}
\begin{displayquote}
To have to stop to formulate rigorous demonstrations would put a stop to most physico-mathematical inquiries. ... The physics will guide the physicist along somehow to useful and important results, by the constant union of physical and geometrical or analytical ideas. —Oliver Heaviside, Electromagnetic Theory (§224)
\end{displayquote}

\begin{displayquote}
Mephistopheles: Allwissend bin ich nicht; doch viel ist mir bewusst. [I am not omniscient, but much is known to me.] -J. W. Goethe, Faust (Part I, 4th scene)
\end{displayquote}

\begin{displayquote}
The last temptation is the greatest treason: To do the right deed for the wrong reason. -T. S. Eliot, Murder in the Cathedral (Part I)
\end{displayquote}

Everything we have done so far is mathematically respectable, although some of the results have been phrased in informal language. To make further progress, however, it is necessary to make a bargain with the devil. The devil offers us effective and conceptually meaningful techniques for calculating physically interesting quantities. In return, however, he requires us to compromise our mathematical souls by accepting the validity of certain approximation procedures and certain formal calculations without proof and - what is a good deal more disconcerting - by working with some putative mathematical objects that lack a rigorous definition. The situation is in some ways similar to the mathematical analysis of the eighteenth century, which developed without the support of a rigorous theory of limits and with the use of poorly defined infinitesimals.

The first and most essential tool the devil offers us is perturbation theory.

\subsection*{Perturbation theory}
Suppose we have a quantum system whose Hamiltonian $H$ is the sum of two terms:

$$
H=H_{0}+H_{I} .
$$

We take $H_{0}$ to be known and understood, and we wish to study the effect of adding in the extra term $H_{I}$. In most of our applications, $H_{0}$ will be the Hamiltonian for a free field and $H_{I}$ will be the interaction term, but the discussion below applies also to other situations where $H$ is taken to be a perturbation of $H_{0}$. There is a considerable mathematical literature dealing with this kind of situation, in which precise hypotheses are placed on $H, H_{0}$, and $H_{I}$. Unfortunately, it is largely irrelevant for our purposes, because the $H_{I}$ 's that we shall need are too singular. We simply have to proceed step by step, taking some care to recognize when we are taking something on faith.

We begin with some general considerations. For psychological comfort, the reader may wish to think of the case where $H_{0}$ is a self-adjoint operator and $H_{I}$ is a bounded self-adjoint operator, in which case $H$ is a self-adjoint operator with the same domain as $H_{0}$. However, this is far from the case we shall need in the sequel.

What we are really interested in is not $H$ and $H_{0}$ but the one-parameter unitary groups they generate, i.e., the time evolution operators for the unperturbed system and the perturbed system:

$$
U_{0}(t)=e^{-i t H_{0}}, \quad U(t)=e^{-i t H}
$$

We assume that the group $U_{0}(t)$ is "well known," and our problem is to compute $U(t)$. For this purpose it is convenient to adopt a point of view that is intermediate between the Schrödinger picture (in which the states evolve in time, $\psi(t)=U(t) \psi$, while the observables remain fixed) and the Heisenberg picture (in which the states remain fixed while the observables evolve in time, $A(t)=U(-t) A U(t)$ ), called the interaction picture. Namely, we let the observables evolve in time according to the unperturbed Hamiltonian,

$$
A(t)=U_{0}(-t) A U_{0}(t)
$$

and the states evolve in such a way as to correct $U_{0}(t)$ to $U(t)$ :

$$
\psi(t)=V(t) \psi, \quad \text { where } \quad V(t)=U_{0}(-t) U(t) .
$$

The matrix element of the observable $A$ between the states $\psi_{1}$ and $\psi_{2}$ at time $t$ is thus

$$
\left\langle\psi_{2}(t)\right| A(t)\left|\psi_{1}(t)\right\rangle=\left\langle\psi_{2}\right| V(t)^{\dagger} U_{0}(-t) A U_{0}(t) V(t)|\psi\rangle=\left\langle\psi_{2}\right| U(-t) A U(t)\left|\psi_{1}\right\rangle,
$$

which is the matrix element of $A$ between $U(t) \psi_{1}$ and $U(t) \psi_{2}$ (Schrödinger picture) or the matrix element of $U(-t) A U(t)$ between $\psi_{1}$ and $\psi_{2}$ (Heisenberg picture), as it should be.

In the interaction picture, the problem is to compute $V(t)$ (from which $U(t)= U_{0}(t) V(t)$ is easily derived since $U_{0}(t)$ is assumed known). Now, $V(t)$ satisfies the differential equation

$$
\frac{d V(t)}{d t}=i U_{0}(-t) H_{0} U(t)-i U_{0}(-t) H U(t)=\frac{1}{i} U_{0}(-t) H_{I} U(t)=\frac{1}{i} H_{I}(t) V(t)
$$

where the $t$-dependence of $H_{I}(t)$ is defined by (6.1). This is the real basis for our calculations. That is, instead of starting with the Hamiltonians $H$ and $H_{0}$, we start with a one-parameter unitary group $U_{0}(t)$ and a Hermitian operator $H_{I}$, and we look for a one-parameter family $V(t)$ of unitary operators such that $i V^{\prime}(t)=H_{I}(t) V(t)$. ( $H_{I}$ may be unbounded, in which case there are issues about domains that we ignore here.) Then $U(t)=U_{0}(t) V(t)$ will be a one-parameter unitary group, and $H_{0}$ and $H$ will merely be the infinitesimal generators of $U_{0}(t)$ and $U(t)$; we need not worry about them further.

The differential equation $i V^{\prime}(t)=H_{I}(t) V(t)$ is equivalent to the integral equation

$$
V(t)=V(0)+\int_{0}^{t} \frac{d}{d \tau} V(\tau) d \tau=I+\frac{1}{i} \int_{0}^{t} H_{I}(\tau) V(\tau) d \tau
$$

which may be iterated in the usual way:

$$
V(t)=I+\frac{1}{i} \int_{0}^{t}\left[H_{I}(\tau) d \tau+\frac{1}{i} \int_{0}^{\tau} H_{I}(\tau) H_{I}\left(\tau^{\prime}\right) V\left(\tau^{\prime}\right) d \tau^{\prime}\right] d \tau
$$

and so forth, leading to the formal expansion

$$
\begin{gathered}
V(t) \sim I+\sum_{1}^{\infty} V_{n}(t) \\
V_{n}(t)=\frac{1}{i^{n}} \int_{0}^{t} \int_{0}^{\tau_{n}} \cdots \int_{0}^{\tau_{2}} H_{I}\left(\tau_{n}\right) H_{I}\left(\tau_{n-1}\right) \cdots H_{I}\left(\tau_{1}\right) d \tau_{1} \cdots d \tau_{n-1} d \tau_{n}
\end{gathered}
$$

The series $I+\sum_{1}^{\infty} V_{n}(t)$ is called the Dyson series ${ }^{1}$ for $V(t)$.\\
What does this really mean? If $H_{I}$ is actually a bounded operator, the series $I+\sum V_{n}(t)$ converges in the norm topology since $\left\|V_{n}(t)\right\| \leq\left\|H_{I}\right\|^{n} t^{n} / n!$, and its sum is indeed the operator $V(t)$. If $H_{I}$ is unbounded, but there is a dense domain $\mathcal{D}$ mapped into itself by $H_{I}$ and the operators $U_{0}(t)$, then $V_{n}(t)$ at least makes sense as an operator on $\mathcal{D}$. If $H_{I}$ also contains a coefficient $\lambda$ (i.e., $H_{I}=\lambda \widetilde{H}_{I}$ ), it may be that the formula $V(t) \sim I+\sum V_{n}(t)$ is an asymptotic expansion in powers of $\lambda$, valid on $\mathcal{D}$. Of course, making this into a rigorous theorem requires some estimates for the error term. All this is mathematically interesting, but it does not really serve our purposes.

The gist of perturbation theory as a practical art is to take the first few partial sums of the series (6.4) as effective approximations to the operator $V(t)$. More precisely, the theory should generate numbers representing physically significant quantities that can be compared with experiment. Typically these numbers come from matrix elements of $V(t)$; in any case, the series (6.4) is used to generate numerical series $\sum_{0}^{\infty} c_{n}$ that formally represent a physical quantity $C$. The first part of our bargain with the devil is this:

\begin{displayquote}
If the first few terms $c_{n}$ decrease rapidly in magnitude as $n$ increases, the corresponding partial sums of the series $\sum c_{n}$ are to be accepted as effective approximations to the quantity $C$.
\end{displayquote}

The hypothesis of this Ansatz is a serious business. The progress of quantum field theory remained almost at a standstill for twenty years because the terms $c_{n}$, at first sight, have an unfortunate predilection for being infinite, and it takes some hard work to prune away the divergences and generate $c_{n}$ 's that are finite and meaningful. Even then, there is no guarantee that they will be suitably small. In quantum electrodynamics they are, and the theory is very successful; but for the interaction that holds atomic nuclei together they are usually not, and our understanding of nuclear phenomena at a fundamental level remains incomplete for this reason.

Let us be perfectly clear about one thing here: the validity of (6.5) has absolutely nothing to do with the convergence of the series $\sum c_{n}$ or the series $\sum V_{n}$ from which it is derived. To bring this point home, here is a simple parable. Consider the two series

$$
\sum_{0}^{\infty} \frac{(-100)^{n}}{n!}, \quad \sum_{0}^{\infty} \frac{n!}{(-100)^{n}} .
$$

The first series converges to $e^{-100}$, but its first few partial sums are $1,-99,4901$, $-161765 \frac{2}{3}, 4004901$, etc. To call these "approximations" to $e^{-100}\left(\approx 10^{-44}\right)$ is nothing but a bad joke. On the other hand, the second series diverges, but its

\footnotetext{${ }^{1}$ This notion antedates Dyson; the name is more properly applied to the perturbation series for the S-matrix that we shall introduce in §6.3. See Dyson [25], [26].
}
first few partial sums - $1,0.99,0.9902,0.990194,0.99019424$ - provide excellent approximations to the integral
$$
\int_{0}^{\infty} \frac{100 e^{-t} d t}{100+t}
$$

Indeed, the error is less in magnitude than the first neglected term of the series (to see this, write $1 /(1+(t / 100))$ as a finite geometric series plus a remainder), and these magnitudes decrease up to the 100 th term $100!/(100)^{100} \approx 10^{-42}$.

We return to the Dyson series. The formulas (6.4) can be rewritten in a convenient way by using the notion of time-ordered product. Suppose $A_{1}(t), \ldots, A_{n}(t)$ are operators depending on $t \in \mathbb{R}$. If $t_{1}, \ldots, t_{n}$ are distinct points in $\mathbb{R}$, the timeordered product $\mathcal{T}\left[A_{1}\left(t_{1}\right) \cdots A_{n}\left(t_{n}\right)\right]$ is defined by

$$
\mathcal{T}\left[A_{1}\left(t_{1}\right) \cdots A_{n}\left(t_{n}\right)\right]=A_{i_{1}}\left(t_{i_{1}}\right) \cdots A_{i_{n}}\left(t_{i_{n}}\right), \text { where } t_{i_{1}}>t_{i_{2}} \cdots>t_{i_{n}} .
$$

That is, the factors are ordered so that the operators are applied in the order of increasing time parameter. We also declare the time-ordering operation to be linear, so that the time-ordering of a sum or integral of such products is the sum or integral of the time-ordered products. If $V_{n}(t)$ is defined as in (6.4), we then have

$$
\begin{aligned}
V_{n}(t) & =\frac{1}{i^{n} n!} \int_{0}^{t} \int_{0}^{t} \cdots \int_{0}^{t} \mathcal{T}\left[H_{I}\left(\tau_{1}\right) \cdots H_{I}\left(\tau_{n}\right)\right] d \tau_{1} \cdots d \tau_{n} \\
& =\mathcal{T}\left[\frac{1}{i^{n} n!} \int_{0}^{t} \int_{0}^{t} \cdots \int_{0}^{t} H_{I}\left(\tau_{1}\right) \cdots H_{I}\left(\tau_{n}\right) d \tau_{1} \cdots d \tau_{n}\right]
\end{aligned}
$$

Indeed, the integral over the cube $(0, t)^{n}$ is the sum of the integrals over the $n!$ simplices where the coordinates have a particular ordering such as $\tau_{1}>\tau_{2}>\cdots> \tau_{n}$, and each of the latter integrals is equal to $V_{n}(t)$. The expansion of $V(t)$ then takes the form

$$
\begin{aligned}
V(t) & \sim \mathcal{T}\left[\sum_{0}^{\infty} \frac{1}{i^{n} n!} \int_{0}^{t} \int_{0}^{t} \cdots \int_{0}^{t} H_{I}\left(\tau_{1}\right) \cdots H_{I}\left(\tau_{n}\right) d \tau_{1} \cdots d \tau_{n}\right] \\
& \equiv \mathcal{T} \exp \left[\frac{1}{i} \int_{0}^{t} H_{I}(\tau) d \tau\right]
\end{aligned}
$$

This last expression is called the time-ordered exponential of $i^{-1} \int_{0}^{t} H_{I}(\tau) d \tau$, and it is, by definition, the sum of the time-orderings of the terms in the Taylor series of the exponential. This is merely a convenient and suggestive notation, not a new formula for $V$. A time-ordered exponential has just as much, or as little, meaning as the series that defines it.

Recall that what we really want to calculate is the group $U(t)$. Since $U(t)= U_{0}(t) V(t)$, the Dyson series (6.4) yields an expansion of $U(t)$ in terms of the free evolution group $U_{0}(t)$ and the interaction Hamiltonian $H_{I}$. Indeed, we have

$$
U(t) \sim U_{0}(t)+\sum_{1}^{\infty} U_{n}(t), \text { where } U_{n}(t)=U_{0}(t) V_{n}(t) \text { for } n \geq 1
$$

Recalling that $H_{I}(t)=U_{0}(-t) H_{I} U_{0}(t)$, for $t>0$ we have

$$
\begin{aligned}
U_{1}(t) & =\frac{1}{i} \int_{0}^{t} U_{0}(t-\tau) H_{I} U_{0}(\tau) d \tau \\
U_{2}(t) & =\frac{1}{i^{2}} \int_{t>\tau_{2}>\tau_{1}>0} U_{0}\left(t-\tau_{2}\right) H_{I} U\left(\tau_{2}-\tau_{1}\right) H_{I} U_{0}\left(\tau_{1}\right) d \tau_{1} d \tau_{2}
\end{aligned}
$$

and so forth. These formulas admit a very suggestive graphical interpretation, shown in Figure 6.1.

\begin{figure}[h]
\begin{center}
  \includegraphics[width=\textwidth]{2025_11_18_7960b2c161913da1e816g-138}

\caption{Figure 6.1. Graphical interpretation of the Dyson series.}
\end{center}
\end{figure}

In the formula for $U_{1}(t)$, the integrand $U_{0}(t-\tau) H_{I} U_{0}(\tau)$ represents the system propagating freely until time $\tau$, being hit with the interaction $H_{I}$, then propagating freely again until time $t$. Likewise, in $U_{2}(t)$, the system propagates freely until time $\tau_{1}$, is hit with $H_{I}$, propagates freely again until time $\tau_{2}$, is hit with $H_{I}$ again, and then propagates freely until time $t$. Thus, at least to the extent to which the Dyson series can be trusted, the picture that emerges is as follows. The system can propagate freely $\left(U_{0}(t)\right)$ over the time interval $(0, t)$, interact once at some intermediate time $\left(U_{1}(t)\right)$, interact twice $\left(U_{2}(t)\right)$, and so forth. We add up all the ways these things can happen over all possible intermediate times, and the result is the interacting propagator $U(t)$. This is the simplest version of the philosophy for computing time evolutions that was pioneered by Feynman, and Figure 6.1 consists of embryonic Feynman diagrams.

We conclude this section with an elementary but technical calculation that will be needed in the next section. Let us rewrite the interaction Hamiltonian as $g H_{I}$, where $g$ is a small parameter. Suppose that the unperturbed Hamiltonian $H_{0}$ has an eigenvalue $\lambda_{0}$ with unit eigenvector $v_{0}$, and suppose that $H=H_{0}+g H_{I}$ has an eigenvalue $\lambda_{g}$ with eigenvector $v_{g}$, normalized so that $\left\langle v_{0} \mid v_{g}\right\rangle=1$, such that $\lambda_{g}$ and $v_{g}$ depend smoothly on $g$. (Of course, there are theorems that guarantee that this will happen under suitable hypotheses - for example, if $H_{I}$ is bounded, or just bounded relative to $H_{0}$, and the eigenvalue $\lambda_{0}$ is simple; see Reed and Simon [97], §XII.2.) Thus we have

$$
\lambda_{g}=\lambda_{0}+\epsilon_{1} g+\epsilon_{2} g^{2}+O\left(g^{3}\right), \quad v_{g}=v_{0}+g w_{1}+g^{2} w_{2}+O\left(g^{3}\right),
$$

for some $\epsilon_{j} \in \mathbb{C}$ and $w_{j} \in \mathcal{H}$. We wish to calculate $\epsilon_{1}$ and $\epsilon_{2}$. (This calculation can be extended recursively to find the higher-order coefficients in the Taylor expansion of $\lambda_{g}$ too.)

We first observe that since

$$
1=\left\langle v_{0} \mid v_{g}\right\rangle=1+g\left\langle v_{0} \mid w_{1}\right\rangle+g^{2}\left\langle v_{0} \mid w_{2}\right\rangle+\cdots
$$

for all small $g$, we have

$$
\left\langle v_{0} \mid w_{1}\right\rangle=\left\langle v_{0} \mid w_{2}\right\rangle=0 .
$$

Now,

$$
\begin{aligned}
H v_{g} & =\left(H_{0}+g H_{I}\right)\left(v_{0}+g w_{1}+g^{2} w_{2}+\cdots\right) \\
& =\lambda_{0} v_{0}+g\left(H_{I} v_{0}+H_{0} w_{1}\right)+g^{2}\left(H_{I} w_{1}+H_{0} w_{2}\right)+\cdots
\end{aligned}
$$

and on the other hand,

$$
H v_{g}=\lambda_{g} v_{g}=\lambda_{0} v_{0}+g\left(\epsilon_{1} v_{0}+\lambda_{0} w_{1}\right)+g^{2}\left(\epsilon_{2} v_{0}+\epsilon_{1} w_{1}+\lambda_{0} w_{2}\right)+\cdots
$$

Hence

$$
\begin{aligned}
H_{0} w_{1}+H_{I} v_{0} & =\epsilon_{1} v_{0}+\lambda_{0} w_{1} \\
H_{0} w_{2}+H_{I} w_{1} & =\epsilon_{2} v_{0}+\epsilon_{1} w_{1}+\lambda_{0} w_{2}
\end{aligned}
$$

Since $\left\langle v_{0}\right| H_{0}\left|w_{1}\right\rangle=\lambda_{0}\left\langle v_{0} \mid w_{1}\right\rangle=0$, by taking the inner product of (6.9) with $v_{0}$ and using (6.8) we obtain

$$
\epsilon_{1}=\left\langle v_{0}\right| H_{I}\left|v_{0}\right\rangle
$$

and hence

$$
\left(H_{0}-\lambda_{0}\right) w_{1}=-\left(H_{I}-\epsilon_{1}\right) v_{0}=-\left(H_{I} v_{0}-\left\langle v_{0}\right| H_{I}\left|v_{0}\right\rangle v_{0}\right)=-(I-P) H_{I} v_{0}
$$

where $P$ is the orthogonal projection onto $\mathbb{C} v_{0}$. The restriction of $H_{0}-\lambda_{0}$ to the orthogonal complement of $v_{0}$, i.e., the range of $I-P$, is invertible on that space, so with a small abuse of notation,

$$
w_{1}=-\left(H_{0}-\lambda_{0}\right)^{-1}(I-P) H_{I} v_{0}=-(I-P)\left(H_{0}-\lambda_{0}\right)^{-1}(I-P) H_{I} v_{0}
$$

Finally, we take the inner product of (6.10) with $v_{0}$ and use (6.8) as above to obtain

$$
\epsilon_{2}=\left\langle v_{0}\right| H_{I}\left|w_{1}\right\rangle=-\left\langle v_{0}\right| H_{I}(I-P)\left(H_{0}-\lambda_{0}\right)^{-1}(I-P) H_{I}\left|v_{0}\right\rangle .
$$

\subsection*{A toy model for electrons in an atom}
In this section we analyze a nonrelativistic particle of mass $M$ moving in a potential $V$ and interacting with a neutral scalar field $\phi$ of mass $\mu$, as a toy model for an electron in an atom interacting with the electromagnetic radiation field. In more detail, in the latter situation we take the atomic nucleus to be infinitely heavy and therefore stationary. The potential $V$ is the Coulomb potential generated by the nucleus (perhaps modified by the presence of other electrons), the particle is an electron moving nonrelativistically (according to the Schrödinger equation) in the potential $V$, and the field $\phi$ is the rest of the electromagnetic field - which classically is the radiation field produced by the motion of the electron, but here is treated as a quantum field. However, in our toy model we neglect the complications due to the spins of electrons and photons, and hence take the wave function for the particle to be scalar-valued and the field to be a scalar field. We also assign a positive mass $\mu$ to the field quanta, but this should be envisioned as being very small in comparision with all other masses and energies under consideration. (An analogous discussion for genuine electromagnetism, without the simplifications, can be found in Sakurai [103], Chapter 2.) In what follows, it will be tempting to refer to the nonrelativistic particle as an electron and the particles associated to the field as photons, but in order not to create false impressions, we shall refer to them as "the particle" and "the (field) quanta," respectively.

We are going to perform two calculations with this model: the transition rate for the particle to emit or absorb field quanta, and the shift in the energy levels $E_{n}$ due to the presence of the field. This is not fundamental physics, but it will provide\\
a useful illustration of the workings of a quantum field and the art of approximation as well as some insight into real phenomena of atomic physics.

Here is the setup. To minimize technicalities, we take the particle and the field to live in a large cubical box $\mathcal{B}$ of side length $L$, with periodic boundary conditions, as in $\oint 5.1$ - that is, on a 3 -torus rather than $\mathbb{R}^{3}$. The state space is $\mathcal{H}=\mathcal{H}_{\text {part }} \otimes \mathcal{H}_{\text {field }}$, where $\mathcal{H}_{\text {part }}=L^{2}(\mathcal{B})$ and $\mathcal{H}_{\text {field }}$ is the Boson Fock space over $L^{2}(\mathcal{B})$. The Hamiltonian is $H=H_{\text {part }}+H_{\text {field }}+H_{I}$, where $H_{\text {part }}=-\nabla^{2} / 2 M+V$ acting on $\mathcal{H}_{\text {part }}, H_{\text {field }}=\sum_{\mathbf{p} \in \Lambda} \omega_{\mathbf{p}} a^{\dagger}(\mathbf{p}) a(\mathbf{p})$ as in (5.3) with $\Lambda=[(2 \pi / L) \mathbb{Z}]^{3}$ being the lattice of allowable momenta, and $H_{I}$ is the interaction term. Our first choice for $H_{I}$ is the simplest possible thing: $H_{I}=g \phi$, where $g$ is a scalar called the coupling constant. It corresponds to the charge of the electron, and the form $H_{I}=g \phi$ of the interaction is motivated by classical electrodynamics, where the interaction of a charged body with the electromagnetic field is proportional to the charge and to the field strength.

More precisely, if $\phi(\mathbf{x})$ is the Schrödinger-picture field,

$$
\phi(\mathbf{x})=\frac{1}{L^{3 / 2}} \sum_{\mathbf{p} \in \Lambda} \frac{1}{\sqrt{2 \omega_{\mathbf{p}}}}\left(e^{i \mathbf{p} \cdot \mathbf{x}} a(\mathbf{p})+e^{-i \mathbf{p} \cdot \mathbf{x}} a^{\dagger}(\mathbf{p})\right)
$$

we would like to take

$$
H_{I}(\psi \otimes \mathbf{v})(\mathbf{x})=g \psi(\mathbf{x}) \phi(\mathbf{x}) \mathbf{v} \quad\left(\psi \in L^{2}(\mathcal{B})=\mathcal{H}_{\text {part }}, \mathbf{v} \in \mathcal{H}_{\text {field }}\right)
$$

To put it another way, we can think of $\mathcal{H}$ as $L^{2}\left(\mathcal{B}, \mathcal{H}_{\text {field }}\right)$, and then $H_{I}$ should be simply pointwise multiplication by the operator-valued function $g \phi(\cdot)$.

But already there is a difficulty: $\phi(\cdot)$ is not a genuine function, as the series (6.13) converges only in the sense of distributions, as we pointed out in §5.1. (If we replaced the box by $\mathbb{R}^{3}$, the situation would be even worse.) When we deal with interacting fields in a relativistically correct manner, we shall have to bite the bullet and proceed somehow in the face of such singularities. But since we are presently treating the particle as nonrelativistic, there cannot be much harm in discarding the high-frequency (i.e., high-energy) components of the field and replacing (6.17) by a finite sum over the $\mathbf{p}$ such that $|\mathbf{p}| \leq K$, for some large $K$. Another way to look at this is as follows: since we are taking the the particle's velocity to be much less than 1 (the speed of light) in magnitude, its momentum must be much less than its mass $M$ in magnitude; in particular, the uncertainty in momentum is much less than $M$. But then by (3.19), the uncertainty in position must be larger than (roughly) $\hbar / M=1 / M$. It is therefore reasonable to smooth the field $\phi$ out by convolving it with a smooth approximation $\chi$ to the delta-function on $\mathbb{R}^{3}$ that is negligibly small outside the set $|\mathbf{x}|<1 / M$ :

$$
\phi * \chi(\mathbf{x})=\frac{1}{L^{3 / 2}} \sum_{\mathbf{p}} \frac{\widehat{\chi}(\mathbf{p})}{\sqrt{2 \omega_{\mathbf{p}}}}\left(e^{i \mathbf{p} \cdot \mathbf{x}} a(\mathbf{p})+e^{-i \mathbf{p} \cdot \mathbf{x}} a^{\dagger}(\mathbf{p})\right)
$$

We can specify such a $\chi$ by requiring that its Fourier transform $\widehat{\chi}(\mathbf{p})$ be equal to 1 on a ball $|\mathbf{p}| \leq K$ and supported in a slightly larger ball, where $K \geq M$, and then $\phi * \chi$ is essentially the finite sum suggested above. The details of the smoothing will be of no importance to us, so we shall simply replace $\phi$ by this finite sum and take

$$
H_{I}=\frac{g}{L^{3 / 2}} \sum_{\mathbf{p} \in \Lambda,|\mathbf{p}| \leq K} \frac{1}{\sqrt{2 \omega_{\mathbf{p}}}}\left(e^{i \mathbf{p} \cdot \mathbf{x}} a(\mathbf{p})+e^{-i \mathbf{p} \cdot \mathbf{x}} a^{\dagger}(\mathbf{p})\right)
$$

Since the system is in a box, the spectra of $H_{\text {part }}$ and $H_{\text {field }}$ are discrete. Let $E_{0} \leq E_{1} \leq E_{2} \leq \cdots$ be the eigenvalues of $H_{\text {part }}$, with eigenvectors $\psi_{0}, \psi_{1}, \psi_{2}, \ldots$. (The eigenvalues of interest are the negative ones, corresponding to the bound states, and the nonrelativistic condition entails their being much less than the mass $M$ in absolute value.) Also, let $\Omega$ be the vacuum state in the Fock space $\mathcal{H}_{\text {field }}$. Then a basis for the state space $\mathcal{H}$ is obtained by taking the tensor products of the $\psi_{n}$ 's with the $k$-quantum states $a^{\dagger}\left(\mathbf{p}_{1}\right) \cdots a^{\dagger}\left(\mathbf{p}_{k}\right) \Omega(k=0,1,2, \ldots)$ as the $\mathbf{p}_{j}$ 's range over the lattice $\Lambda$ of allowable momenta. We employ Dirac-style shorthand notation for these basis vectors:

$$
\left|n ; \mathbf{p}_{1}, \ldots, \mathbf{p}_{k}\right\rangle=c_{\mathbf{p}_{1}, \ldots, \mathbf{p}_{k}} \psi_{n} \otimes a^{\dagger}\left(\mathbf{p}_{1}\right) \cdots a^{\dagger}\left(\mathbf{p}_{k}\right) \Omega, \quad|n\rangle=\psi_{n} \otimes \Omega
$$

( $c_{\mathbf{p}_{1}, \ldots, \mathbf{p}_{k}}$ is a normalization constant, equal to $\left[\prod_{\mathbf{p} \in \Lambda} n_{\mathbf{p}}!\right]^{-1 / 2}$ where $n_{\mathbf{p}}$ is the number of $\mathbf{p}_{j}$ equal to $\mathbf{p}$.) These vectors are eigenvectors for the free Hamiltonian $H_{0}=H_{\text {part }}+H_{\text {field }}$ :

$$
\left(H_{\mathrm{part}}+H_{\mathrm{field}}\right)\left|n ; \mathbf{p}_{1}, \ldots, \mathbf{p}_{k}\right\rangle=\left(E_{n}+\omega_{\mathbf{p}_{1}}+\cdots+\omega_{\mathbf{p}_{k}}\right)\left|n ; \mathbf{p}_{1}, \ldots, \mathbf{p}_{k}\right\rangle,
$$

where $\omega_{\mathbf{p}}=\sqrt{|\mathbf{p}|^{2}+\mu^{2}}$. On the other hand, the interaction Hamiltonian is given by

$$
\begin{aligned}
& H_{I}\left|n ; \mathbf{p}_{1}, \ldots, \mathbf{p}_{k}\right\rangle(\mathbf{x}) \\
& \quad=\frac{g}{L^{3 / 2}} \sum_{|\mathbf{p}| \leq K} \frac{1}{\sqrt{2 \omega_{\mathbf{p}}}}\left[e^{i \mathbf{p} \cdot \mathbf{x}} \sqrt{m_{\mathbf{p}}}\left|n ; \mathbf{p}_{1}^{\prime}, \ldots, \mathbf{p}_{k-1}^{\prime}\right\rangle\right. \\
& \left.\quad+e^{-i \mathbf{p} \cdot \mathbf{x}} \sqrt{m_{\mathbf{p}}+1}\left|n ; \mathbf{p}, \mathbf{p}_{1}, \ldots, \mathbf{p}_{k}\right\rangle\right]
\end{aligned}
$$

where $m_{\mathbf{p}}$ is the number of $\mathbf{p}_{j}$ that are equal to $\mathbf{p}$ and, when $m_{\mathbf{p}}>0,\left(\mathbf{p}_{1}^{\prime}, \ldots, \mathbf{p}_{k-1}^{\prime}\right)$ is obtained from $\left(\mathbf{p}_{1}, \ldots, \mathbf{p}_{k}\right)$ by omitting one of the $\mathbf{p}_{j}$ 's that is equal to $\mathbf{p}$.

Emission and absorption of quanta. Let us find the transition probability for emission of a field quantum when there are no quanta present initially. That is, we assume that at time 0 the state is $|n\rangle$ (the particle has energy $E_{n}$, and there are no quanta), and ask for the probability that at time $t>0$ the state is $|m, \mathbf{p}\rangle$ (the particle has energy $E_{m}$ and there is one quantum with momentum $\mathbf{p}$ ), namely, $|\langle m, \mathbf{p}| U(t)| n\rangle\left.\right|^{2}$. (We assume that $\left|E_{n}\right|,\left|E_{m}\right|$, and $|\mathbf{p}|$ are less than the cutoff energy $K$.) According to (6.4), the first-order approximation to $\langle m, \mathbf{p}| U(t)|n\rangle$ is

$$
\langle m, \mathbf{p}| U(t)|n\rangle \approx\langle m, \mathbf{p}| U_{0}(t)|n\rangle+\frac{1}{i} \int_{0}^{t}\langle m, \mathbf{p}| U_{0}(t-\tau) H_{I} U_{0}(\tau)|n\rangle d \tau
$$

The term $\langle m, \mathbf{p}| U_{0}(t)|n\rangle$ vanishes because the vacuum state of the field is orthogonal to the 1 -quantum states. The action of $U_{0}(t)=e^{-i t H_{0}}$ and $H_{I}$ in the other term are given by (6.15) and (6.16):

$$
\begin{aligned}
& \frac{1}{i} \int_{0}^{t}\langle m, \mathbf{p}| U_{0}(t-\tau) H_{I} U_{0}(\tau)|n\rangle d \tau \\
& \quad=\frac{1}{i} e^{-i\left(E_{m}+\omega_{\mathbf{p}}\right) t} \int_{0}^{t} e^{i\left(E_{m}+\omega_{\mathbf{p}}-E_{n}\right) \tau}\langle m, \mathbf{p}| H_{I}|n\rangle d \tau \\
& \quad=\frac{g}{L^{3 / 2} i} e^{-i\left(E_{m}+\omega_{\mathbf{p}}\right) t} \frac{e^{i\left(E_{m}-E_{n}+\omega_{\mathbf{p}}\right) t}-1}{i\left(E_{m}-E_{n}+\omega_{\mathbf{p}}\right)} \int_{B} \psi_{m}(\mathbf{x})^{*} \frac{e^{-i \mathbf{p} \cdot \mathbf{x}}}{\sqrt{2 \omega_{\mathbf{p}}}} \psi_{n}(\mathbf{x}) d^{3} \mathbf{x}
\end{aligned}
$$

Thus, if we set

$$
\alpha=E_{m}-E_{n}+\omega_{\mathbf{p}}, \quad C(m, n, \mathbf{p})=\int_{B} \psi_{m}(\mathbf{x})^{*} e^{-i \mathbf{p} \cdot \mathbf{x}} \psi_{n}(\mathbf{x}) d^{3} \mathbf{x}
$$

the transition probability to first order is

$$
|\langle m, \mathbf{p}| U(t)| n\rangle\left.\right|^{2} \approx \frac{g^{2}}{L^{3}} \frac{|C(m, n, \mathbf{p})|^{2}}{2 \omega_{\mathbf{p}}}\left|\frac{e^{i \alpha t}-1}{i \alpha}\right|^{2}=\frac{g^{2}}{L^{3}} \frac{|C(m, n, \mathbf{p})|^{2}}{2 \omega_{\mathbf{p}}} \frac{\sin ^{2}(\alpha t / 2)}{(\alpha / 2)^{2}} .
$$

Let us examine the meaning of this. First, when $t$ is reasonably large the function

$$
f_{t}(\alpha)=\frac{\sin ^{2}(\alpha t / 2)}{(\alpha / 2)^{2}}
$$

has a spike at $\alpha=0$ of height $t^{2}$ and width $\approx 1 / t$, and is small elsewhere; moreover, by a standard contour integral, $\int_{-\infty}^{\infty} f_{t}(\alpha) d \alpha=2 \pi t$. Therefore, $f_{t}(\alpha) \approx 2 \pi t \delta(\alpha)$. This expresses the fact that for the transition to take place we must have (approximately) $\alpha=0$, that is, the energy $\omega_{\mathbf{p}}$ of the quantum must equal the difference $E_{n}-E_{m}$ between the initial and final energies of the particle. The fact that this relation becomes blurred for $t$ small is a reflection of the time-energy uncertainty principle: for very short times there is an unremovable uncertainty in the energy. In real life, this phenomenon manifests itself in the fact that the photons emitted when an atomic electron makes a transition from one energy level to another one are not all of exactly the same frequency; the spectroscopists speak of the line width in the resulting spectral plot.

The approximation $f_{t}(\alpha) \approx 2 \pi t \delta(\alpha)$ is useful provided $t$ is large enough so that almost all the mass of $f_{t}(\alpha)$ is concentrated in an interval where $\alpha$ (the deviation from perfect energy conservation) is small in comparison with the energy difference $E_{n}-E_{m}$ of the of the initial and final states of the particle, that is, when $t \gg 1 /\left(E_{n}-E_{m}\right)$. In this situation, (6.17) says that the transition probability per unit time is approximately

$$
\frac{\pi g^{2}}{L^{3}} \frac{|C(m, n, \mathbf{p})|^{2}}{\omega_{\mathbf{p}}} \delta\left(E_{m}-E_{n}+\omega_{\mathbf{p}}\right) .
$$

Of course this makes sense only when $t$ is small enough so that the probability at time $t$ is substantially less than one. Thus we are in a slightly uncomfortable situation where $t$ must be neither too small nor too large; without some more specific numbers it is impossible to tell whether we have accomplished anything at all. We shall consider this question in more detail after carrying out the calculation one step further for the limiting situation in which we dispense with the box, i.e., let $L \rightarrow \infty$.

We now have the emission rate for the particular quantum whose momentum is $\mathbf{p}$. However, if we are just interested in the decay of an excited state, a quantum with momentum $\mathbf{p}^{\prime}$ is just as good as a quantum with momentum $\mathbf{p}$ provided that $\omega_{\mathbf{p}}=\omega_{\mathbf{p}^{\prime}}=E_{n}-E_{m}$. The number of quantum states per unit volume in momentum space is $(L / 2 \pi)^{3}$ (for $L$ very large), so the number of quantum states whose momentum points into a small solid angle $d \Omega$ and whose energy is in a small band $[\omega, \omega+d \omega]$ (that is, $\mathbf{p} /|\mathbf{p}| \in d \Omega, \omega_{\mathbf{p}} \in[\omega, \omega+d \omega]$ ) is

$$
\left(\frac{L}{2 \pi}\right)^{3}|\mathbf{p}|^{2} \frac{d|\mathbf{p}|}{d \omega} d \omega d \Omega=\left(\frac{L}{2 \pi}\right)^{3} \omega|\mathbf{p}| d \omega d \Omega
$$

since $|\mathbf{p}|=\sqrt{\omega_{\mathbf{p}}^{2}-\mu^{2}}, \mu$ being the rest mass of a field quantum. The rate of emission of any quantum in this region of momentum space is the quantity (6.18) multiplied by the number of states in this region, that is,

$$
\frac{g^{2}}{8 \pi^{2}}|C(m, n, p)|^{2}|\mathbf{p}| \delta\left(E_{n}-E_{m}+\omega\right) d \omega d \Omega
$$

Observe that at this point the box has disappeared, and the total emission rate for a quantum in any direction, written out explicitly, is

$$
\Gamma=\frac{g^{2}}{8 \pi^{2}} \int_{|\mathbf{p}|^{2}=\left(E_{n}-E_{m}\right)^{2}-\mu^{2}}|\mathbf{p}|\left|\int \psi_{m}(\mathbf{x})^{*} e^{-i \mathbf{p} \cdot \mathbf{x}} \psi_{n}(\mathbf{x}) d^{3} \mathbf{x}\right|^{2} d \Omega(\mathbf{p})
$$

We emphasize that $d \Omega(\mathbf{p})$ is solid angle measure, not surface area measure; i.e., $\int_{|\mathbf{p}|=r} d \Omega(\mathbf{p})=4 \pi$ no matter what $r$ is.

As we said above, $\Gamma$ is a "transition probability per unit time," so that for reasonably short times $\Delta t$ the probability of a transition in time $\Delta t$ is approximately $\Gamma \Delta t$ - assuming that the restriction $\Delta t \gg 1 /\left(E_{n}-E_{m}\right)$ doesn't interfere with the notion of "reasonably short." Under this assumption, the interpretation of $\Gamma$ for longer time intervals is easily obtained. Suppose we start with a large number $N$ identical systems ( $N$ atoms, if you will) with particles in the state $\psi_{n}$. After a short time interval $\Delta t, N \Gamma \Delta t$ of the particles will have undergone the transition to state $\psi_{m}$. This leads to the differential equation $d N / d t=-\Gamma N$, so that $N(t)=N_{0} e^{-\Gamma t}$. That is, $\Gamma$ is the transition rate or decay rate of the excited state.

Now let us plug in numbers for the parameters in this simplified model that correspond to the quantities in the real physical situation that inspired it: an electron in the Coulomb potential of an atomic nucleus interacting with the electromagnetic field to make a transition from an excited state with energy $E_{n}$ to the ground state with energy $E_{0}$. That is, we take $g^{2} / 4 \pi$ to be the fine structure constant $1 / 137$ (so $\left.g^{2} \approx 10^{-1}\right), M$ to be the mass of the electron ( $M=5 \times 10^{5} \mathrm{eV}$ ), and $\mu$ to be the mass of the photon ( $\mu=0$ ). (Yes, our model requires $\mu \neq 0$, but we can assume $\mu \ll E_{n}-E_{0}$ and replace it by 0 as an approximation here.) The difference in energy levels $E_{n}-E_{0}$ is on the order of a few eV , or $10^{-5} M$. The magnitude of the integral $\int \psi_{0}(\mathbf{x})^{*} e^{i \mathbf{p} \cdot \mathbf{x}} \psi_{n}(\mathbf{x}) d^{3} \mathbf{x}$ depends strongly on the shape of the excited state wave function $\psi_{n}$. The integral always vanishes at $\mathbf{p}=\mathbf{0}$, so when $\psi_{n}$ is concentrated near the origin (i.e., when $n$ is not too large), it can be estimated by replacing $e^{i \mathbf{p} \cdot \mathbf{x}}$ by $i \mathbf{p} \cdot \mathbf{x}$ :

$$
\left|\int \psi_{0}(\mathbf{x})^{*}(\mathbf{p} \cdot \mathbf{x}) \psi_{n}(\mathbf{x}) d^{3} \mathbf{x}\right| \leq|\mathbf{p}|\langle r\rangle
$$

where $\langle r\rangle$ is the root-mean-square distance of the particle from the origin in the state $\psi_{n}$. (If the integral vanishes to higher order in $\mathbf{p}$, it is even smaller.) We may take $\langle r\rangle$ to be at most a few times the Bohr radius of the hydrogen atom, $r_{0}=\hbar^{2} / M g^{2} \approx 10 / M$ (with $M$ and $g$ as above), so with $|\mathbf{p}|=E_{n}-E_{0} \approx 10^{-5} M$, we have $|\mathbf{p}|\langle r\rangle \approx 10^{-3}$ or less. Thus the integral (6.19) is estimated by

$$
\Gamma \lesssim \frac{g^{2}}{2 \pi}\left(10^{-5} M\right)\left(10^{-3}\right)^{2} \sim 10^{-11} M
$$

This is some $10^{6}$ times smaller than $E_{n}-E_{0}$. In other words, the decay time is some $10^{6}$ times longer than the time $1 /\left(E_{n}-E_{0}\right)$ needed to dispose of the energy uncertainty, so the assumption that there are times that are neither too small nor\\
too large is tenable. The mass-time conversion (1.6) shows that the decay time $\left(10^{-11} M\right)^{-1}$ is on the order of $10^{-10}$ second, which is in the right ball park for emission of photons in real atoms.

To do a thorough job of analyzing this situation, one should go further. By a more careful analysis of the operator $P e^{-i t H} P$, where $P$ is the orthogonal projection onto the state $|n\rangle$, one can show that $\langle n| e^{-i t H}|n\rangle \approx e^{-(\Gamma+i \beta) t}$ where $\Gamma$ is the decay rate (6.19) found above and $\beta=E_{n}+\Delta E_{n}, \Delta E_{n}$ being a correction to the energy level $E_{n}$ that we shall discuss below. See Messiah [83], $\S$ XXI.13.

We took the initial state to be one with no field quanta present. If we take the initial state to contain field quanta, the calculations are very similar, but there are some points to be noted. First, if the initial state has $k$ quanta, there is a nonzero transition probability to states with either $k+1$ or $k-1$ quanta, that is, a new quantum can be emitted or an existing quantum can be absorbed. If all the existing quanta are in different states initially, the calculations go through with no essential change. But if some of them - say, $m$ of them - are in the same state, there is an extra factor of $\sqrt{m+1}$ or $\sqrt{m}$ in (6.16) that shows up as an extra factor of $m+1$ or $m$ in the transition rate. It is no surprise that a quantum in a particular state is more likely to be absorbed when there are lots of such quanta around; the interesting thing is that the presence of quanta in a particular state speeds up the rate at which such quanta are emitted too. This is the principle behind the laser.

What happens if we go to higher orders? The second-order term in the Dyson series (6.4) contains the interaction Hamiltonian - that is, the field operators twice, so it has nonzero matrix elements between states with $k$ quanta and states with $k+2, k$, or $k-2$ quanta. It therefore has no effect on the transition rate for emission or absorption of a single quantum, but it introduces new processes: emission of two quanta, absorption of two quanta, and emission and reabsorption (or absorption and reemission) of a single quantum. If one works out the calculations, in the first two cases one finds formulas for the emission or absorption rate of the same nature as (6.18); they involve delta-functions (or approximate delta-functions) that express conservation of energy. But in the third case, the uncertainty principle allows the particle and/or quantum to be temporarily "off mass shell" between the times of emission and absorption or vice versa. In particular, when a quantum is emitted and reabsorbed, it can have any energy $E$ as long as it is reabsorbed in a time $\approx 1 / E$. We say that it is a virtual quantum. This is the paradigm for all "virtual particles," which we shall discuss in more detail in §6.6.

Rather than considering absorption and emission processes further, we shall do another calculation to show how the presence of the field affects the energy levels of the particle. This will provide an introduction to the idea of renormalization.

Energy level shifts. We consider an eigenvalue $E_{n}$ of $H_{\text {part }}$, which is also an eigenvalue of $H_{0}=H_{\text {part }}+H_{\text {field }}$ with eigenvector $|n\rangle$, the state where the particle has energy $E_{n}$ and there are no field quanta present. We assume as before that $\left|E_{n}\right| \ll M$ so that the nonrelativistic approximation is valid, and we wish to determine how the presence of the field affects $E_{n}$, to second order in the coupling constant $g$. We indicated at the end of §6.1 how to perform this calculation, provided that the perturbed eigenvalue and eigenvector depend smoothly on $g$. Here we proceed on the assumption that this condition is valid. ${ }^{2}$

\footnotetext{${ }^{2}$ No physicist would waste a moment worrying about this, but here are a few remarks for the mathematically fastidious. The calculations that follow involve only states with at most two
}The results at the end of §6.1, with $\lambda_{0}=E_{n}$ and $v_{0}=|n\rangle$, show that the perturbed eigenvalue is of the form $E_{n}+\epsilon_{1} g+\epsilon_{2} g^{2}$ to second order in $g$, where $\epsilon_{1}$ and $\epsilon_{2}$ involve $H_{I}|n\rangle$. We have

$$
H_{I}|n\rangle=L^{-3 / 2} \sum_{|\mathbf{p}| \leq K}\left(2 \omega_{\mathbf{p}}\right)^{-1 / 2} e^{-i \mathbf{p} \cdot \mathbf{x}}|n ; \mathbf{p}\rangle,
$$

where $e^{-i \mathbf{p} \cdot \mathbf{x}}$ denotes the operation of pointwise multiplication by the function $e^{-i \mathbf{p} \cdot(\cdot)}$ on $L^{2}\left(B, \mathcal{H}_{\text {field }}\right)$. The 1 -quantum state $H_{I}|n\rangle$ is orthogonal to the 0 quantum state $|n\rangle$, so $\epsilon_{1}=0$ by (6.11). Hence, by (6.12), second-order correction is\\
$\Delta E_{n}=\epsilon_{2} g^{2}=-\frac{g^{2}}{2 L^{3}} \sum_{\mathbf{p}, \mathbf{p}^{\prime}} \frac{1}{\sqrt{\omega_{\mathbf{p}^{\prime}} \omega_{\mathbf{p}}}}\left\langle n ; \mathbf{p}^{\prime}\right| e^{i \mathbf{p}^{\prime} \cdot \mathbf{x}}(I-P)\left(H_{0}-E_{n}\right)^{-1}(I-P) e^{-i \mathbf{p} \cdot \mathbf{x}}|n ; \mathbf{p}\rangle$,\\
where $P$ is the orthogonal projection onto $|n\rangle$. The operator $e^{-i \mathbf{p} \cdot \mathbf{x}}$ takes $|n ; \mathbf{p}\rangle$ to a linear combination of the vectors $|m ; \mathbf{p}\rangle$ with different $m$ 's but the same $\mathbf{p}$. These are all orthogonal to $|n\rangle$, so the factors $I-P$ can be omitted; they are also preserved up to scalar multiples by $H_{0}$, so the relation $\left\langle\mathbf{p}^{\prime}, \mathbf{p}\right\rangle=\delta_{\mathbf{p}^{\prime} \mathbf{p}}$ implies that only the terms with $\mathbf{p}^{\prime}=\mathbf{p}$ survive, and we have

$$
\Delta E_{n}=-\frac{g^{2}}{2 L^{3}} \sum_{|\mathbf{p}| \leq K} \frac{1}{\omega_{\mathbf{p}}}\langle n ; \mathbf{p}| e^{i \mathbf{p} \cdot \mathbf{x}}\left(H_{0}-E_{n}\right)^{-1} e^{-i \mathbf{p} \cdot \mathbf{x}}|n ; \mathbf{p}\rangle .
$$

Moreover, the operator $e^{i \mathbf{p} \cdot \mathbf{x}}$ commutes with $H_{\text {field }}$ and with multiplication by the potential $V(\mathbf{x})$, and $e^{i \mathbf{p} \cdot \mathbf{x}} \circ \nabla \circ e^{-i \mathbf{p} \cdot \mathbf{x}}=\nabla-i \mathbf{p}$, so

$$
\begin{aligned}
e^{i \mathbf{p} \cdot \mathbf{x}}\left(H_{0}-E_{n}\right)^{-1} e^{-i \mathbf{p} \cdot \mathbf{x}} & =\left(-\frac{(\nabla-i \mathbf{p})^{2}}{2 M}+V(\mathbf{x})+H_{\mathrm{field}}-E_{n}\right)^{-1} \\
& =\left(H_{0}-E_{n}-\frac{\mathbf{p} \cdot \nabla}{i M}+\frac{|\mathbf{p}|^{2}}{2 M}\right)^{-1}
\end{aligned}
$$

Therefore,

$$
\Delta E_{n}=-\frac{g^{2}}{2 L^{3}} \sum_{|\mathbf{p}| \leq K} \frac{1}{\omega_{\mathbf{p}}}\langle n ; \mathbf{p}|\left(\omega_{\mathbf{p}}+\frac{|\mathbf{p}|^{2}}{2 M}-\frac{\mathbf{p} \cdot \nabla}{i M}\right)^{-1}|n ; \mathbf{p}\rangle
$$

Now, $\nabla / i M$ is the velocity operator for the particle, and the nonrelativistic approximation entails the velocity being small. Thus, as a crude estimate of $\Delta E_{n}$, we drop the term $\mathbf{p} \cdot \nabla / i M$ :

$$
\Delta E_{n} \sim-\frac{g^{2}}{2 L^{3}} \sum_{|\mathbf{p}| \leq K} \frac{1}{\omega_{\mathbf{p}}\left(\omega_{\mathbf{p}}+\left(|\mathbf{p}|^{2} / 2 M\right)\right)}
$$

field quanta, so they are unaffected if we multiply the interaction Hamiltonian (6.14) fore and aft by the orthogonal projection onto the space of these states. Since we have already cut off the high frequencies, the result is a bounded operator. Hence standard results of perturbation theory quoted in §6.1 apply provided that $E_{n}$ is a simple eigenvalue of $H_{0}$, and this can be artificially arranged by adding a small generic perturbation to the potential $V$. Then the perturbed eigenvalue depends smoothly on $g$ for sufficiently small $g$, but whether the physical value $g=\sqrt{4 \pi / 137}$ is "sufficiently small" is another matter.

At this point we might as well pass to the limit $L \rightarrow \infty$, so that the sum becomes an integral and $(2 \pi / L)^{3}$ becomes the volume element $d^{3} \mathbf{p}$ :

$$
\Delta E_{n} \sim-\frac{g^{2}}{16 \pi^{3}} \int_{|\mathbf{p}| \leq K} \frac{d^{3} \mathbf{p}}{\omega_{\mathbf{p}}\left(\omega_{\mathbf{p}}+\left(|\mathbf{p}|^{2} / 2 M\right)\right.}
$$

We can calculate this exactly in the limit of massless field quanta, for which $\omega_{\mathbf{p}}= |\mathbf{p}|$ :

$$
\begin{aligned}
-\frac{g^{2}}{16 \pi^{3}} \int_{|\mathbf{p}| \leq K} \frac{d^{3} \mathbf{p}}{|\mathbf{p}|^{2}(1+(|\mathbf{p}| / 2 M)} & =-\frac{g^{2}}{4 \pi^{2}} \int_{0}^{K} \frac{d \rho}{1+(\rho / 2 M)} \\
& =-\frac{g^{2}}{2 \pi^{2}} M \log \left(1+\frac{K}{2 M}\right)
\end{aligned}
$$

There is clearly something fishy about this result, because it depends on the artifical cutoff parameter $K$ (although it is quite insensitive to the precise value of $K)$. In fact, if we remove the cutoff by letting $K \rightarrow \infty$, it diverges. Moreover, if we take $K=M$ and use values for the other parameters appropriate to atomic electrons ( $g^{2} / 4 \pi=1 / 137, M=5 \times 10^{5} \mathrm{eV}, E_{n} \sim-10 \mathrm{eV}$ ) we get $\Delta E_{n} \sim-2 \times 10^{-3} M= -10^{3} \mathrm{eV}$, which is around 100 times as big as $E_{n}$ itself. Thus the validity of (6.22) as a "correction" to $E_{n}$ seems highly dubious.

On the other hand, the integral (6.22) doesn't depend on the particular energy eigenstate $|n\rangle$, or even on the potential $V$. In fact, it is exactly the second-order correction (in the limit $L \rightarrow \infty$ ) to the ground-state energy level of a free particle. Indeed, in this case the eigenfunctions for $H_{\text {part }}=-\nabla^{2} / 2 M$ are indexed by the same lattice in momentum space as the one-quantum states of the field; the normalized eigenfunction with momentum $\mathbf{p}$ is $e_{\mathbf{p}}=L^{-3 / 2} e^{-i \mathbf{p} \cdot \mathbf{x}}$, with eigenvalue $E_{\mathbf{p}}=\left|\mathbf{p}^{2}\right| / 2 M$. Noting that $e_{\mathbf{p}}$ is the result of applying the operator $e^{-i \mathbf{p} \cdot \mathbf{x}}$ to the ground state $|0\rangle=e_{\mathbf{0}}$, we see that (6.20) for the ground state of a free particle (i.e., $V=0, n=0$, and $E_{n}=0$ ) becomes

$$
\begin{aligned}
\Delta E_{0} & =\frac{g^{2}}{2 L^{3}} \sum_{\mathbf{p}} \frac{1}{\omega_{\mathbf{p}}}\left\langle e_{\mathbf{p}} ; \mathbf{p}\right|\left(H_{\mathrm{part}}+H_{\mathrm{field}}\right)^{-1}\left|e_{\mathbf{p}} ; \mathbf{p}\right\rangle \\
& =-\frac{g^{2}}{2 L^{3}} \sum_{|\mathbf{p}| \leq K} \frac{1}{\omega_{\mathbf{p}}\left(\omega_{\mathbf{p}}+\left(|\mathbf{p}|^{2} / 2 M\right)\right)}
\end{aligned}
$$

which is (6.22) in the limit $L \rightarrow \infty$.\\
What can this shift in the ground state energy of a free particle possibly mean? The only energy such a particle has is its mass, so $\Delta E_{0}$ must be interpreted as a change in the rest mass of the particle due to its interaction with the field:

$$
\Delta E_{0}=\Delta M
$$

Now, thinking in terms of the real world of atomic electrons of which this is a simplified model, we can move an electron into or out of an atomic potential well and measure what happens when we do so, but we cannot decouple it from the electromagnetic field, as its electric charge is always present. Hence the physical mass that is measured in the laboratory is not the $M$ in our Hamiltonian but the corrected - or renormalized - mass $M+\Delta M$. (The same issue arose in classical electrodynamics, where it was an unresolved puzzle to figure out how much of an electron's mass should be attributed to the energy of the electrostatic field that it generates. We shall return to this point in §7.10.)

With this in mind, it is clear that the mass shift $\Delta M$ accounts for a large fraction of all the energy level shifts for a particle in an arbitrary potential. Indeed, the approximation we made by neglecting the term $\mathbf{p} \cdot \nabla / M$ in (6.21) amounts precisely to replacing $\Delta E_{n}$ by $\Delta M$, and the physically interesting part of the shift $\Delta E_{n}$ is the remainder $\Delta E_{n}-\Delta M$. Actually this is not quite right either. The numbers $E_{n}$ are the eigenvalues of the Hamiltonian $H_{\text {part }}=-\left(\nabla^{2} / 2 M\right)+V$ containing the original "bare" mass $M$, but the physically meaningful energy levels are the eigenvalues $E_{n}^{\prime}$ of the Hamiltonian obtained by replacing $M$ by the physical mass $M+\Delta M$. The "true" energy level shifts are the differences

$$
\Delta^{\prime} E_{n}=\left(\Delta E_{n}-\Delta M\right)+\left(E_{n}^{\prime}-E_{n}\right)
$$

and it is their differences $\Delta^{\prime} E_{n}-\Delta^{\prime} E_{m}$, as corrections to the differences $E_{n}^{\prime}-E_{m}^{\prime}$, that are observed spectroscopically when particles make a transition from one state to another.

We shall not carry out the calculation of these corrections here, as the differences between our model and real atomic electrons are too great to make it worthwhile. But nonrelativistic calculations of this sort by Bethe gave the first reasonably accurate theoretical calculation of the Lamb shift, i.e., the difference in the energy levels between the $2 S_{1 / 2}$ and the $2 P_{1 / 2}$ states of the hydrogen atom. The interested reader can consult Bethe's original account [9] or the retelling of it in Sakurai [103], §2.8. Of course, to obtain really accurate and theoretically sound results along these lines, one must proceed from a relativistically correct theory and take account of the contributions of high-energy virtual quanta, so that the results do not depend on the choice of cutoff. We shall say a little more about this in §7.11.

One tends to think of renormalization as a way of "subtracting off infinities," and indeed it plays that role here if we define $H_{I}$ in terms of the original quantum field without cutting off the high frequencies. But the fact that it still has a substantial role to play when the cutoff is employed highlights the fact that it has a deeper significance. The presence of interactions really does change the effective parameters of a system, and those changes must be taken into account whether they be finite or not.

\subsection*{The scattering matrix}
After this warmup, we now turn to the real subject at hand: quantum fields with interactions. The path we follow here is the one we followed in discussing free fields, "canonical quantization." That is, we start out with classical field equations and a relativistically invariant Lagrangian from which they are derived, then replace the classical field variables by quantum fields, which are Fourier integrals of creation and annihilation operators that satisfy suitable commutation or anticommutation relations. The Lagrangian will be a sum of free-field terms and interaction terms, and the same will be true of the corresponding Hamiltonian. We then make the assumption that perturbation theory will yield some meaningful results, so that we can calculate the time evolution by using the Dyson series.

What are we trying to calculate, anyhow? In the typical particle-physics experiment some incoming particles with certain momenta, initially far apart, come together and interact with each other, producing some outgoing particles with certain momenta that again become far apart and eventually arrive at detectors. The phrase "far apart" is meant to indicate that the particles no longer interact and can\\
be treated as free particles, and we are emphasizing that the momenta of the particles are usually a lot more important than their positions. Speaking loosely, there is an "incoming" state |in> consisting of free particles and an "outgoing" state |out consisting of free particles, and we wish to find the transformation $\mid$ in $\rangle \mapsto \mid$ out $\rangle$.

Setting up a precise description of this situation is a rather subtle business. We shall do so quite informally and refer the reader to Weinberg [131], §3.1, for a more careful analysis. Let $U(t)=e^{-i t H}$ be the time evolution operator; we assume that $H$ is the sum of a free-field Hamiltonian $H_{0}$ and an interaction Hamiltonian $H_{I}$. If $|v\rangle$ is the state vector of the system at the present time, one might at first think that $\mid$ in $\rangle$ and $\mid$ out $\rangle$ would be $\lim _{t \rightarrow-\infty} U(t)|v\rangle$ and $\lim _{t \rightarrow+\infty} U(t)|v\rangle$, but these limits are unlikely to exist, as the particles move off to infinity as $t \rightarrow \pm \infty$. But if $U(t)|v\rangle$ is essentially a collection of free particles for $\pm t$ large, we can keep that collection in our field of view by moving it back to the present using the freefield evolution $U_{0}(-t)=e^{i H_{0} t}$. Thus we are led to the interaction-picture time dependence $|v(t)\rangle=U_{0}(-t) U(t)|v\rangle$ as discussed in §6.1, and |in $\rangle$ and |out $\rangle$ should be the limits of $|v(t)\rangle$ as $t \rightarrow-\infty$ or $t \rightarrow+\infty$. In short, the transformation $\mid$ in $\rangle \rightarrow \mid$ out $\rangle$ that we are interested in is the scattering operator

$$
S=\lim _{t_{0} \rightarrow-\infty, t_{1} \rightarrow+\infty} U_{0}\left(-t_{1}\right) U\left(t_{1}-t_{0}\right) U_{0}\left(t_{0}\right)
$$

More practically, we are interested in the matrix elements of $S,\langle$ out $| S \mid$ in $\rangle$, which are collectively known as the scattering matrix or $S$-matrix.

From the point of view of rigorous mathematics, the existence and unitarity of the scattering operator is an interesting and highly nontrivial question. (There is a rigorous scattering theory in the context of the Wightman axioms due to Haag and Ruelle; see Jost $[\mathbf{7 0}]$, Bogolubov et al. $[\mathbf{1 1}],[\mathbf{1 2}]$, and Araki $[\mathbf{3}]$. The scattering theory of nonrelativistic quantum mechanics is developed in Reed and Simon [96].) But we have already made our bargain with the devil, so we shall not stop to worry about this, but rather proceed directly to try to calculate the matrix elements by perturbation theory. By (6.3), the operator $V(t)=U_{0}(-t) U(t)$ is the solution of the initial value problem $i V^{\prime}(t)=H_{I}(t) V(t), V(0)=I$, where $H_{I}(t)=U_{0}(-t) H_{I} U_{0}(t)$ is the interaction-picture form of $H_{I}$. Therefore, the operator

$$
V\left(t, t_{0}\right)=U_{0}(-t) U\left(t-t_{0}\right) U_{0}\left(t_{0}\right)=V(t) V\left(t_{0}\right)^{-1}
$$

is the solution of the initial value problem

$$
i \frac{d}{d t} V\left(t, t_{0}\right)=H_{I}(t) V(t) V\left(t_{0}\right)^{-1}=H_{I}(t) V\left(t, t_{0}\right), \quad V\left(t_{0}, t_{0}\right)=I
$$

so by the same calculation that leads to (6.4), it is formally given by the Dyson series

$$
\begin{aligned}
V\left(t, t_{0}\right) & \sim \mathcal{T} \exp \frac{1}{i} \int_{t_{0}}^{t} H_{I}(t) d t \\
& =I+\sum_{1}^{\infty} \frac{1}{i^{n} n!} \int_{t_{0}}^{t} \cdots \int_{t_{0}}^{t} \mathcal{T}\left[H_{I}\left(\tau_{1}\right) \cdots H_{I}\left(\tau_{n}\right)\right] d \tau_{1} \cdots d \tau_{n}
\end{aligned}
$$

The corresponding formal expansion for the scattering operator (6.23) is then (6.24)

$$
S=V(+\infty,-\infty) \sim I+\sum_{1}^{\infty} \frac{1}{i^{n} n!} \int_{-\infty}^{\infty} \cdots \int_{-\infty}^{\infty} \mathcal{T}\left[H_{I}\left(\tau_{1}\right) \cdots H_{I}\left(t_{n}\right)\right] d \tau_{1} \cdots d \tau_{n}
$$

Proceeding on the faith that formula (6.24) has some useful content, we proceed to calculate its matrix elements $\langle$ out $| S \mid$ in $\rangle$ by calculating the corresponding elements for the operators on the right. In line with the ideas sketched above, we take the in and out states to be idealized states describing particles with definite momenta, namely, states obtained by applying a finite sequence of creation operators $a_{j}\left(\mathbf{p}_{j}\right)$ to the vacuum state. More precisely, the "vacuum state" here is the no-particle state in the tensor product of the Fock spaces for the free fields in question, which we denote by $|0\rangle,{ }^{3}$ and the in and out states are taken to be of the form

$$
a_{1}^{\dagger}\left(\mathbf{p}_{1}\right) \cdots a_{k}^{\dagger}\left(\mathbf{p}_{k}\right)|0\rangle
$$

where the subscripts $1, \ldots, k$ on the $a^{\dagger}$ 's specify the particle species, spin state, and any other relevant parameters.

Now, there are both mathematical and physical drawbacks to this. Mathematically, one must recall that $a_{j}^{\dagger}(\cdot)$ is a distribution rather than a function; applying it to $|0\rangle$ yields what physicists call a "non-normalizable state," that is, a generalized state whose Fourier-transformed wave function is a delta-function in momentum space. Physicists have no a priori objection to such states, but they bear only a vague resemblance to a collection of distinct particles that are supposed to be "far apart" from each other, as they are not localized in position at all. ${ }^{4}$ Nonetheless, we brush all such objections aside for the time being and proceed; once we get some results it will be time to ask what significance they have.

Before proceeding with the quantum field theory, let us briefly examine how these ideas work in a simpler situation: the scattering of a particle by a fixed potential. Let $V(\mathbf{x})$ be a potential function on $\mathbb{R}^{3}$, which we assume to be negligibly small outside some bounded set $\Sigma$; the minimal condition we need is that $V \in L^{1}\left(\mathbb{R}^{3}\right)$. We send in a particle with momentum $\mathbf{p}$ from far outside $\Sigma$; the potential deflects it, and it emerges with some momentum $\mathbf{q}$. What is the amplitude for this process?

Here the Hilbert space is $L^{2}\left(\mathbb{R}^{3}\right)$, the (nonrelativistic) state space for a single particle; the free Hamiltonian $H_{0}$ is $-\nabla^{2} / 2 m$, and the "interaction" Hamiltonian $H_{I}$ is multiplication by the function $V$. The state $|\mathbf{p}\rangle$ of a free particle with momentum $\mathbf{p}$ is given by the wave function $e^{i \mathbf{p} \cdot \mathbf{x}}$, and we wish to calculate $\langle\mathbf{q}| S|\mathbf{p}\rangle$. Let us consider just the first-order approximation to $S$ arising from the series (6.24):

$$
S \approx I+\frac{1}{i} \int_{-\infty}^{\infty} H_{I}(t) d t=I+\frac{1}{i} \int_{-\infty}^{\infty} e^{i H_{0} t} V e^{-i H_{0} t} d t
$$

Since $H_{0}|\mathbf{p}\rangle=\left(|\mathbf{p}|^{2} \mid / 2 m\right)|\mathbf{p}\rangle$, we have

$$
\begin{aligned}
\langle\mathbf{q}| S|\mathbf{p}\rangle & \approx\langle\mathbf{q} \mid \mathbf{p}\rangle+\frac{1}{i} \int_{\infty}^{\infty}\langle\mathbf{q}| e^{i|\mathbf{q}|^{2} t / 2 m} V e^{-i|\mathbf{p}|^{2} t / 2 m}|\mathbf{p}\rangle d t \\
& =\int e^{i(\mathbf{p}-\mathbf{q}) \cdot \mathbf{x}} d^{3} \mathbf{x}+\frac{1}{i} \int_{-\infty}^{\infty} e^{i\left(|\mathbf{q}|^{2}-|\mathbf{p}|^{2}\right) t / 2 m} d t \int_{\mathbb{R}^{3}} e^{-i(\mathbf{q}-\mathbf{p}) \cdot \mathbf{x}} V(\mathbf{x}) d^{3} \mathbf{x}
\end{aligned}
$$

\footnotetext{${ }^{3}$ The other common symbol for the vacuum, $\Omega$, will be used for a different vacuum state in §6.11.\\
${ }^{4}$ Still, given that space is infinite in extent, particles whose position is uniformly distributed over space should be far apart on the average.
}
so by the Fourier inversion formula (1.5),
$$
\langle\mathbf{q}| S|\mathbf{p}\rangle \approx(2 \pi)^{3} \delta(\mathbf{q}-\mathbf{p})-2 \pi i \delta\left(\frac{|\mathbf{q}|^{2}-|\mathbf{p}|^{2}}{2 m}\right) \widehat{V}(\mathbf{q}-\mathbf{p})
$$

This is known as the Born approximation to the scattering amplitude. We shall find it useful later on in comparing the results of quantum field theory with more classical descriptions of various processes.

The formula (6.25) requires some commentary. First, the delta-functions are the price we pay for using the idealized states $\langle\mathbf{q}|$ and $|\mathbf{p}\rangle$. In this situation, of course, it is easy to restate the result in terms of honest $L^{2}$ states with wave functions $g$ and $f$ : just multiply both sides by $\widehat{g}(\mathbf{q}) \widehat{f}(\mathbf{p})$ and integrate over $\mathbf{p}$ and q to obtain

$$
\langle g| S|f\rangle \approx\langle g \mid f\rangle-2 \pi i \int_{\mathbb{R}^{3}} \delta\left(\frac{|\mathbf{q}|^{2}-|\mathbf{p}|^{2}}{2 m}\right) \widehat{g}(\mathbf{q}) \widehat{V}(\mathbf{q}-\mathbf{p}) \widehat{f}(\mathbf{p}) \frac{d^{3} \mathbf{p} d^{3} \mathbf{q}}{(2 \pi)^{6}}
$$

(The integral on the right makes sense without further interpretation if $f$ and $g$ have a little regularity.) The delta-function $\delta(\mathbf{q}-\mathbf{p})$ is always there simply because the perturbation series starts with the identity, but it disappears as soon as $\mathbf{q}$ differs from $\mathbf{p}$ even slightly, and it is the second term in (6.25), the matrix element for $S-I$, that contains the interesting information. The delta-function in it expresses conservation of energy; we shall refer to the rest - that is, $-2 \pi i \widehat{V}(\mathbf{q}-\mathbf{p})$ - as the Born amplitude for scattering by the potential $V$. We refer to Reed and Simon [96], §XI.6, for a detailed mathematical treatment of the Born approximation and the full perturbative series of which it is the first-order part.

We shall find that the S-matrix in quantum field processes has a form similar to (6.25). More specifically, for incoming particles with 4 -momenta $p_{1}, \ldots, p_{m}$ and outgoing particles with 4 -momenta $q_{1}, \ldots, q_{n}$, we will have

$$
\left\langle q_{1}, \ldots, q_{n}\right| S-I\left|p_{1}, \ldots, p_{m}\right\rangle=i(2 \pi)^{4} \delta\left(\sum q_{k}-\sum p_{j}\right) M\left(q_{1}, \ldots, q_{n} ; p_{1}, \ldots, p_{m}\right)
$$

where the quantity of real interest is $M$. Note that here the delta-function expresses conservation of total energy and momentum. In (6.25) there was conservation of energy only, because the fixed potential does not respect conservation of momentum.

Let us now return to quantum fields. We have in mind a field theory with $k$ interacting fields, corresponding to $k$ species of particles and antiparticles. Each field will contribute a free-field term to the classical Lagrangian density. For a real scalar field $\phi$ of mass $m$, this term is $\frac{1}{2}\left[(\partial \phi)^{2}-m^{2} \phi^{2}\right]$; for a Dirac field $\psi$ of mass $m$ it is $\bar{\psi}\left(i \gamma^{\mu} \partial_{\mu}-m\right) \psi$; for the electromagnetic field it is $-\frac{1}{4} F_{\mu \nu} F^{\mu \nu}$, and so forth. In addition, there will be interaction terms, which we assume to be polynomials in the fields and their conjugates. (Here the "conjugate" of a field is the field that becomes the operator adjoint on the quantum level.) Since the Lagrangian must be real, nonreal fields (i.e., non-Hermitian ones - ones whose particles have distinct antiparticles) and their conjugates must appear symmetrically in the Lagrangian. This is true also for Fermion fields. (Fermions that are their own antiparticles are theoretically possible; none are definitely known in nature so far, although they may have something to do with neutrinos.) These interaction terms will contain numerical coefficients, the "coupling constants."

The two examples that we shall keep returning to for most of the rest of this book are as follows:

\begin{itemize}
  \item Quantum electrodynamics (QED). Here we have one Dirac (spin- $\frac{1}{2}$ ) field $\psi$ of mass $m>0$ (representing charged particles of some species and their antiparticles) and one neutral vector field $A_{\mu}$ of mass 0 (representing photons), and the Lagrangian density is
\end{itemize}

$$
\mathcal{L}=\bar{\psi}\left(i \gamma^{\mu} \partial_{\mu}-m\right) \psi-\frac{1}{4} F_{\mu \nu} F^{\mu \nu}-e A_{\mu} \bar{\psi} \gamma^{\mu} \psi
$$

where $F_{\mu \nu}=\partial_{\mu} A_{\nu}-\partial_{\nu} A_{\mu}$ and $e$ (the coupling constant) is the charge of the particles in question. More generally, we can allow several Dirac fields $\psi_{j}$ here, each one representing a different type of charged particle; we simply include one term $\bar{\psi}_{j}\left(i \gamma^{\mu} \partial_{\mu}-m_{j}\right) \psi_{j}$ and one term $-e_{j} A_{\mu} \bar{\psi}_{j} \gamma^{\mu} \psi_{j}$ in the Lagrangian for each such field. The resulting theory describes only the electromagnetic interactions of these various charged particles. (If the particles interact with each other in other ways, one must include additional interaction terms in the Lagrangian, resulting in a more complicated theory.) Most of the time it is enough to consider just one species of charged particle, and the most important case is that of electrons, for which $e$ is conventionally taken to be negative (about $-\sqrt{4 \pi / 137}$ ). When we speak of QED in the sequel, we shall always mean the one-Fermion form arising from (6.27) unless we say otherwise, and we shall take the Fermion in question to be the electron.

\begin{itemize}
  \item The $\phi^{4}$ scalar field theory. Here we have just one self-interacting neutral scalar field $\phi$ of mass $m>0$. The Lagrangian density is
\end{itemize}

$$
\mathcal{L}=\frac{1}{2}\left(\partial_{\mu} \phi\right)\left(\partial^{\mu} \phi\right)-\frac{1}{2} m^{2} \phi^{2}-\frac{\lambda}{4!} \phi^{4}
$$

where the coupling constant $\lambda$ is assumed positive (and small). (The factor of $4!$ is just for convenience.) This corresponds to the nonlinear Klein-Gordon equation $\partial^{2} \phi=m^{2} \phi+\lambda \phi^{3} / 3$ !. Fields of this sort do occur in real-world physical theories; in particular, the Higgs Boson that plays a crucial role in gauge field theories is a neutral scalar particle whose selfinteraction is of a similar type. But the main reason for the ubiquitous appearance of the $\phi^{4}$ scalar field theory in quantum field theory texts is that it is the simplest example of an interacting field theory, which offers a setting in which to learn how quantum fields work without the various complicating factors of QED. In space-time dimensions 2 and 3 it is also one of the few quantum field theories that admits a mathematically rigorous model.\\
We shall meet other physically important fields in Chapter 9. In addition, we shall occasionally make a few remarks about two other field theories:

\begin{itemize}
  \item Yukawa field theory. In the simplest version of this, we have one Dirac field $\psi$ of mass $m_{\psi}>0$ and one neutral scalar field $\phi$ of mass $m_{\phi}>0$, and the Lagrangian density is
\end{itemize}

$$
\mathcal{L}=\bar{\psi}\left(i \gamma^{\mu} \partial_{\mu}-m_{\psi}\right) \psi+\frac{1}{2}(\partial \phi)^{2}-\frac{1}{2} m_{\phi}^{2} \phi^{2}-g \phi \bar{\psi} \psi
$$

The coupling constant is $g$. With a little elaboration, this is the model needed for the first reasonably successful theory of the strong interaction, which accounts for the attraction between nucleons in an atomic nucleus in terms of exchange of (virtual) pions. In the simple Lagrangian (6.29), $\psi$ can represent either protons or neutrons and $\phi$ represents neutral pions.

For a more realistic theory, one needs two Dirac fields $\psi_{p}$ and $\psi_{n}$ for protons and neutrons, as well as two scalar fields $\phi_{0}$ (real) and $\phi_{1}$ (complex) for neutral and charged pions. Moreover, pions have negative parity (they are "pseudoscalars" rather than "scalars"), which entails an extra factor of $\gamma^{5}$ in the interactions. Besides the free-field terms for each particle type, the Lagrangian includes the interactions

$$
-g\left[\phi_{0}\left(\bar{\psi}_{p} \gamma^{5} \psi_{p}+\bar{\psi}_{n} \gamma^{5} \psi_{n}\right)+\phi_{1} \bar{\psi}_{p} \gamma^{5} \psi_{n}+\phi_{1}^{*} \bar{\psi}_{n} \gamma^{5} \psi_{p}\right] .
$$

The first two terms represent a proton or neutron emitting or absorbing a neutral pion; the next one represents a neutron absorbing a positive pion or emitting a negative pion and turning into a proton, and the last one a proton absorbing a negative pion or emitting a positive pion and turning into a neutron. ${ }^{5}$

\begin{itemize}
  \item The Fermi model for beta decay. Here we have four Dirac fields $\psi_{e}, \psi_{p}$, $\psi_{n}$, and $\psi_{\nu}$ representing electrons, protons, neutrons, and neutrinos. The free-field Lagrangian is the sum of the four $\bar{\psi}\left(i \gamma^{\mu} \partial_{\mu}-m_{\psi}\right) \psi$ 's, and the interaction terms are
\end{itemize}

$$
-G\left[\bar{\psi}_{p} \Gamma \psi_{n} \bar{\psi}_{e} \Gamma \psi_{\nu}+\left(\bar{\psi}_{p} \Gamma \psi_{n} \bar{\psi}_{e} \Gamma \psi_{\nu}\right)^{*}\right],
$$

where $G$ is the coupling constant and $\Gamma$ is a suitable combination of Dirac matrices. (There are several possibilities here; we shall be more specific when we discuss this model in § 9.4.) This interaction is meant to model the decay of the neutron, $n \rightarrow e+p+\bar{\nu}$, and related processes.\\
In all of these cases (as well as others of importance), the passage from the Lagrangian to the Hamiltonian on the classical level is quite simple, following the paradigm

$$
L=T-V \quad \longrightarrow \quad H=T+V
$$

with $T$ and $V$ the kinetic and potential energies. The kinetic terms are those involving the time derivatives of the fields in question, and the potential terms are everything else. In particular, the kinetic terms are always part of the free-field terms, so the interaction Hamiltonian - the crucial ingredient for the perturbationtheoretic treatment - is the spatial integral of a Hamiltonian density,

$$
H_{I}(t)=\int_{\mathbb{R}^{3}} \mathcal{H}_{I}(t, \mathbf{x}) d^{3} \mathbf{x}
$$

and $\mathcal{H}_{I}(t, \mathbf{x})$ in turn is simply the negative of the sum of the interaction terms in the Lagrangian density. For the three examples above the interaction Hamiltonian densities are, respectively,

$$
\mathcal{H}_{I}=e \bar{\psi} \gamma^{\mu} \psi A_{\mu}, \quad \mathcal{H}_{I}=\frac{\lambda}{4!} \phi^{4}, \quad \mathcal{H}_{I}=g \bar{\psi} \psi \phi .
$$

Thus in all these cases, as well as more complicated ones that arise in the standard model, the classical interaction Hamiltonian density is a sum of products of field variables.

When we pass to the quantum picture, the field variables become operatorvalued functions (or rather distributions). Suppose the theory involves $k$ fields $\phi_{1}, \ldots, \phi_{k}$, each of which acts on its own Fock space $\mathcal{F}_{j}$. The Hilbert space for the

\footnotetext{${ }^{5}$ The fact that pions are unstable particles that readily decay into leptons is largely irrelevant here; the lifetime of the pion is large in comparison with the characteristic time scale of the strong interaction.
}
interacting theory is then the tensor product $\mathcal{F}=\mathcal{F}_{1} \otimes \cdots \otimes \mathcal{F}_{k}$, and the fields act on this space in the obvious way, with the $j$ th field acting "in the $j$ th variable." Each Fock space $\mathcal{F}_{j}$ comes with its own creation and annihilation operators $a_{j}^{\dagger}$ and $a_{j}$ for particles and $b_{j}^{\dagger}$ and $b_{j}$ for antiparticles (if particles and antiparticles are distinct), and the field $\phi_{j}$ is a Fourier integral of such operators. (As noted in §4.5, this picture needs to be modified slightly so that the creation/annihilation operators for different Fermion fields anticommute rather than commute.)

Our notation for creation and annihilation operators will be as follows. For each particle species, its creation and annihilation operators are functions of momentum $\mathbf{p} \in \mathbb{R}^{3}$ and a discrete variable $\sigma$ that specifies the spin state of the particle. (Actually they are distributions rather than functions in $\mathbf{p}$, but we ignore this distinction here. And of course $\sigma$ is absent for spin-zero fields.) It will be convenient to introduce a label for the particle species as a third variable $\pi$. (If a particle has a distinct antiparticle, these two are considered different species for this purpose. Any other discrete parameters relevant to the problem should also be incorporated into $\pi$.$) Thus, we write a(\mathbf{p}, \sigma, \pi)$ or $a^{\dagger}(\mathbf{p}, \sigma, \pi)$ for the operator that annihilates or creates a particle of species $\pi$ with momentum $\mathbf{p}$ and spin state $\sigma$. The canonical (anti)commutation relations are then

$$
\begin{aligned}
a(\mathbf{p}, \sigma, \pi) a^{\dagger}\left(\mathbf{p}^{\prime}, \sigma^{\prime}, \pi^{\prime}\right) & = \pm a^{\dagger}\left(\mathbf{p}^{\prime} \sigma^{\prime}, \pi^{\prime}\right) a(\mathbf{p}, \sigma, \pi)+(2 \pi)^{3} \delta\left(\mathbf{p}-\mathbf{p}^{\prime}\right) \delta_{\sigma \sigma^{\prime}} \delta_{\pi \pi^{\prime}} \\
a(\mathbf{p}, \sigma, \pi) a\left(\mathbf{p}^{\prime}, \sigma^{\prime}, \pi^{\prime}\right) & = \pm a\left(\mathbf{p}^{\prime}, \sigma^{\prime}, \pi^{\prime}\right) a(\mathbf{p}, \sigma, \pi) \\
a^{\dagger}(\mathbf{p}, \sigma, \pi) a^{\dagger}\left(\mathbf{p}^{\prime}, \sigma^{\prime}, \pi^{\prime}\right) & = \pm a^{\dagger}\left(\mathbf{p}^{\prime}, \sigma^{\prime}, \pi^{\prime}\right) a^{\dagger}(\mathbf{p}, \sigma, \pi)
\end{aligned}
$$

Here the $\pm \operatorname{sign}$ is + if at least one of the species $\pi, \pi^{\prime}$ is Bosonic and - if both species are Fermionic.

What is the quantum interaction Hamiltonian density? The classical $\mathcal{H}_{I}(t, \mathbf{x})$ is a product of field variables, but the corresponding quantum objects are operatorvalued distributions, so it is not clear how to multiply them, or how to integrate the resulting product over $\mathbb{R}^{3}$ to obtain the actual interaction Hamiltonian $H_{I}(t)$ as a well-defined operator on the state space. In fact, a rigorous mathematical construction of $H_{I}(t)$ is generally not available, but fortunately it is not really needed to extract meaningful results from the theory. Rather, one has to regard the formal expression for $\mathcal{H}_{I}(t, \mathbf{x})$ simply as a way of encoding certain information to be used as input for calculations that do have a well-defined meaning. One might draw a parallel with a formal power series, which encodes certain useful algebraic information about its sequence of coefficients that does not depend on the convergence of the series.

Of course, if $H_{I}$ cannot be taken seriously as an operator, then neither can the full Hamiltonian $H$, the unitary operators $e^{-i t H}$ it generates, the scattering operator $S$, or the individual terms in its Dyson series $\sum_{0}^{\infty} S_{n}$. What can be taken seriously, as it turns out, are the S -matrix elements in perturbation theory, that is, the quantities $\sum_{0}^{N}\langle$ out $| S_{n} \mid$ in $\rangle$ for $N$ finite. They are given by sums of integrals of well-defined functions over $\mathbb{R}^{4 k}$ for suitable $k$. Many of these integrals diverge (a reflection of the ill-definedness of the operators just mentioned), but in favorable cases the divergences can be removed by renormalization, and the resulting quantities can then be taken to the lab and compared with experiment. Our goal, therefore, is to compute these matrix elements.

\subsection*{Evaluation of the S-matrix: first steps}
By (6.24) and (6.32), the scattering operator $S$ is formally given by

$$
S \sim I+\sum_{1}^{\infty} \frac{1}{i^{n} n!} \int_{\mathbb{R}^{4}} \cdots \int_{\mathbb{R}^{4}} \mathcal{T}\left[H_{I}\left(x_{1}\right) \cdots H_{I}\left(x_{n}\right)\right] d^{4} x_{1} \cdots d^{4} x_{n}
$$

It is not clear whether this makes any sense at all, but at least on a formal level it is rather nice. In particular, the fact that it is given by integrals over spacetime suggests at least the possibility of relativistic invariance, except for the timeordering inside the integrals. We shall have to examine this point in more detail later; for now, we just forge resolutely ahead. Recall that $\mathcal{H}_{I}$ is a sum of products of field operators. If

$$
|\mathrm{in}\rangle=\prod_{1}^{k_{\mathrm{in}}} a^{\dagger}\left(\mathbf{p}_{j}^{\mathrm{in}}, \sigma_{j}^{\mathrm{in}}, \pi_{j}^{\mathrm{in}}\right)|0\rangle, \quad|\mathrm{out}\rangle=\prod_{1}^{k_{\mathrm{out}}} a^{\dagger}\left(\mathbf{p}_{j}^{\mathrm{out}}, \sigma_{j}^{\mathrm{out}}, \pi_{j}^{\mathrm{out}}\right)|0\rangle,
$$

the S -matrix element $\langle$ out $| S \mid$ in $\rangle$ is given by a sum of integrals of terms of the form

$$
\langle 0| \prod_{1}^{k_{\mathrm{out}}} a\left(\mathbf{p}_{j}^{\mathrm{out}}, \sigma_{j}^{\mathrm{out}}, \pi_{j}^{\mathrm{out}}\right) \mathcal{T}\left[\prod_{1}^{k_{\mathrm{field}}} \phi_{j}\left(x_{j}\right)\right] \prod_{1}^{k_{\mathrm{in}}} a^{\dagger}\left(\mathbf{p}_{j}^{\mathrm{in}}, \sigma_{j}^{\mathrm{in}}, \pi_{j}^{\mathrm{in}}\right)|0\rangle .
$$

One small technical point arises here. In passing from the classical to the quantum interaction Hamiltonian, one is replacing classical functions, which commute, with operators, which may not, so how is one to order the factors? Actually, this is rarely an issue. Since operators pertaining to particles of different species always commute or anticommute, the only possible problem arises when there is a product of a non-Hermitian field with its own adjoint. In the specific theories we are concerned with, the only such terms are of the form $\bar{\psi} \gamma^{\mu} \psi$ where $\psi$ is a Dirac field, and here we take the order dictated by the matrix algebra, with the $\bar{\psi}$ on the left. Another prescription that has some theoretical justification is to declare that all products of fields with themselves or their adjoints should be Wick ordered; see Weinberg [131], p. 200. We shall say more about this later, but it will not be a serious concern for us.

Our first job is to evaluate expressions of the form (6.35). The field operators $\phi_{j}\left(x_{j}\right)$ are themselves Fourier integrals of the operators $a(\mathbf{p}, \sigma, \pi)$ and $a^{\dagger}(\mathbf{p}, \sigma, \pi)$, so we have to evaluate the vacuum expectation values of products of creation and annihilation operators. The method for doing this is simple: use the commutation relations (6.33) to move each creation operator to the left and each annihilation operator to the right. When any annihilation operator reaches the right end and acts on $|0\rangle$ it yields 0 , and when any creation operator reaches the left end and acts on $\langle 0|$ (via its adjoint) it yields 0 . The only nonzero terms will arise when each creation operator (in the initial ket $|\mathrm{in}\rangle$ or in one of the fields) is paired with an annihilation operator (in the final bra <out| or in one of the fields); their commutator will yield a numerical factor (containing delta-functions), and the resulting product of these factors is what must be integrated to get the final result.

Let us work out a very simple example (too simple to be useful all by itself, although it occurs as a part of larger calculations). Consider a single neutral scalar field $\phi$ and its associated creation operator $a^{\dagger}(\mathbf{p})$, and let $\mid$ in $\rangle=a^{\dagger}(\mathbf{p})|0\rangle$ and\\
$\mid$ out $\rangle=|0\rangle$. We have

$$
\langle 0| \phi(t, \mathbf{x}) a^{\dagger}(\mathbf{p})|0\rangle=\int \frac{1}{\sqrt{2 \omega_{\mathbf{q}}}}\langle 0|\left[e^{i q_{\mu} x^{\mu}} a(\mathbf{q})+e^{-i q_{\mu} x^{\mu}} a^{\dagger}(\mathbf{q})\right] a^{\dagger}(\mathbf{p})|0\rangle \frac{d^{3} \mathbf{q}}{(2 \pi)^{3}}
$$

Now $\langle 0| a^{\dagger}(\mathbf{q})=\langle 0| a^{\dagger}(\mathbf{p})=0$, and $a(\mathbf{q}) a^{\dagger}(\mathbf{p})=a^{\dagger}(\mathbf{p}) a(\mathbf{q})+(2 \pi)^{3} \delta(\mathbf{q}-\mathbf{p})$, so

$$
\langle 0| \phi(t, \mathbf{x}) a^{\dagger}(\mathbf{p})|0\rangle=\int \frac{1}{\sqrt{2 \omega_{\mathbf{q}}}} e^{i q_{\mu} x^{\mu}} \delta(\mathbf{q}-\mathbf{p}) d^{3} \mathbf{q}=\frac{e^{i p_{\mu} x^{\mu}}}{\sqrt{2 \omega_{\mathbf{p}}}}
$$

The reader may be reassured to see that the final answer is a perfectly well-defined function!

Returning to the general case, we need a generalization of the notion of Wick ordering that was introduced in $\oint$ 5.3. First, suppose $A_{1}, \ldots, A_{n}$ are operators that are sums or integrals of Bosonic creation and annihilation operators; that is, $A_{j}= A_{j}^{a}+A_{j}^{c}$ where $A_{j}^{a}$ (resp. $A_{j}^{c}$ ) is a sum or integral of annihilation (resp. creation) operators satisfying the commutation relations (6.33) with the ± signs taken to be + . For example, $A_{j}$ could be a field operator $\phi(t, \mathbf{x})$ or a single creation operator $a^{\dagger}(\mathbf{p})$. The Wick ordered product or normally ordered product : $A_{1} \cdots A_{n}$ : is the product $A_{1} \cdots A_{n}$ with the terms rearranged so that all creation operators are to the left of all annihilation operators. For example, when $n=2$,

$$
\begin{aligned}
: A_{1} A_{2}: & =:\left(A_{1}^{a}+A_{1}^{c}\right)\left(A_{2}^{a}+A_{2}^{c}\right):=A_{1}^{a} A_{2}^{a}+A_{1}^{c} A_{2}^{a}+A_{2}^{c} A_{1}^{a}+A_{1}^{c} A_{2}^{c} \\
& =A_{1} A_{2}-\left[A_{1}^{a}, A_{2}^{c}\right]
\end{aligned}
$$

and for $n=3$,

$$
\begin{aligned}
: A_{1} A_{2} A_{3}:= & A_{1}^{a} A_{2}^{a} A_{3}^{a}+A_{3}^{c} A_{1}^{a} A_{2}^{a}+A_{2}^{c} A_{1}^{a} A_{3}^{a}+A_{1}^{c} A_{2}^{a} A_{3}^{a} \\
& +A_{1}^{c} A_{2}^{c} A_{3}^{a}+A_{1}^{c} A_{3}^{c} A_{2}^{a}+A_{2}^{c} A_{3}^{c} A_{1}^{a}+A_{1}^{c} A_{2}^{c} A_{3}^{c} .
\end{aligned}
$$

(The order of the $A^{a}$ 's and $A^{c}$ 's is immaterial, as creation operators commute with each other.)

If some or all of the $A_{n}$ involve Fermionic creation and annihilation operators, the operation of Wick ordering is defined as above except that a factor of -1 is attached to any term when two Fermionic operators are interchanged. Thus, the analogues of (6.37) and (6.38) when all the $A_{j}$ are Fermionic are

$$
\begin{aligned}
: A_{1} A_{2}: & =:\left(A_{1}^{a}+A_{1}^{c}\right)\left(A_{2}^{a}+A_{2}^{c}\right):=A_{1}^{a} A_{2}^{a}+A_{1}^{c} A_{2}^{a}-A_{2}^{c} A_{1}^{a}+A_{1}^{c} A_{2}^{c} \\
& =A_{1} A_{2}-\left\{A_{1}^{a}, A_{2}^{c}\right\}
\end{aligned}
$$

and

$$
\begin{aligned}
: A_{1} A_{2} A_{3}:= & A_{1}^{a} A_{2}^{a} A_{3}^{a}+A_{3}^{c} A_{1}^{a} A_{2}^{a}-A_{2}^{c} A_{1}^{a} A_{3}^{a}+A_{1}^{c} A_{2}^{a} A_{3}^{a} \\
& +A_{1}^{c} A_{2}^{c} A_{3}^{a}-A_{1}^{c} A_{3}^{c} A_{2}^{a}+A_{2}^{c} A_{3}^{c} A_{1}^{a}+A_{1}^{c} A_{2}^{c} A_{3}^{c} .
\end{aligned}
$$

In all cases, the essential feature of Wick ordered products is that their vacuum expectation values vanish:

$$
\langle 0|: A_{1} \cdots A_{n}:|0\rangle=0 .
$$

We shall also modify our definition of time-ordered products where Fermion fields are concerned, again by introducing a minus sign when two Fermionic operators are interchanged. Namely, if $A_{1}(t), \ldots, A_{n}(t)$ are Fermionic operators depending on a time parameter $t$, we set

$$
\mathcal{T}\left[A_{1}\left(t_{1}\right) \cdots A_{n}\left(t_{n}\right)\right]=(\operatorname{sgn} \sigma) A_{\sigma(1)}\left(t_{\sigma(1)}\right) \cdots A_{\sigma(n)}\left(t_{\sigma(n)}\right),
$$

where $\sigma$ is the permutation of $1, \ldots, n$ such that $t_{\sigma(1)}>\cdots>t_{\sigma(n)}$. This convention has no effect on the time-ordered products of interaction Hamiltonians that occur in the Dyson series, because Fermion fields always occur in pairs in such Hamiltonians, so interchanging $\mathcal{H}_{I}\left(t_{1}\right)$ with $\mathcal{H}_{I}\left(t_{2}\right)$ always involves an even number of Fermion interchanges.

Next, suppose $A_{1}$ and $A_{2}$ are sums or integrals of Bosonic or Fermionic creation and annihilation operators as above, and suppose that one or both of them depend on time parameters. If both do, and their time parameters are different, $\mathcal{T}\left(A_{1} A_{2}\right)$ denotes their time-ordered product; otherwise, $\mathcal{T}\left(A_{1} A_{2}\right)=A_{1} A_{2}$. The contraction of $A_{1}$ with $A_{2}$ is the scalar $\overbrace{A_{1} A_{2}}$ defined by

$$
\overbrace{A_{1} A_{2}} I=\mathcal{T}\left(A_{1} A_{2}\right)-: A_{1} A_{2}:
$$

To see that the difference on the right is indeed a scalar multiple of $I$, suppose that at most one of the $A_{j}$ 's is Fermionic. By (6.37), if at most one of the $A$ 's is time-dependent, then

$$
\mathcal{T}\left(A_{1} A_{2}\right)-: A_{1} A_{2}:=\left[A_{1}^{a}, A_{2}^{c}\right]
$$

whereas if $A_{1}=A_{1}\left(t_{1}\right)$ and $A_{2}=A_{2}\left(t_{2}\right)$, then

$$
\mathcal{T}\left(A_{1} A_{2}\right)-: A_{1} A_{2}:= \begin{cases}{\left[A_{1}^{a}, A_{2}^{c}\right]} & \text { if } t_{1}>t_{2} \\ {\left[A_{2}^{a}, A_{1}^{c}\right]} & \text { if } t_{2}>t_{1}\end{cases}
$$

If both $A_{j}$ 's are Fermionic, the result is the same except that commutators are replaced by anticommutators, and there is an extra minus sign from interchanging $A_{1}$ and $A_{2}$ in the case $t_{2}>t_{1}$. In all cases, these commutators or anticommutators are scalars by (6.33). (As usual, we shall be cavalier about identifying the scalar $\overbrace{A_{1} A_{2}}$ with the operator $\overbrace{A_{1} A_{2}} I$.)

Recall that we are trying to evaluate vacuum expectation values of integrals of products of operators $A_{j}$ of the types we have been discussing. More precisely, the products involve annihilation operators on the left, time-ordered field operators in the middle, and creation operators on the right: symbolically, $A^{a} \mathcal{T}\left[A^{f}\right] A^{c}$. As a matter of notational convenience, we agree that the time-ordering on $A^{f}$ may encompass the other operators without changing anything:

$$
\mathcal{T}\left[A^{a} A^{f} A^{c}\right]=A^{a} \mathcal{T}\left[A^{f}\right] A^{c} .
$$

Moreover, the individual fields that make up the Hamiltonian $H_{I}(t)$, which all involve the same time $t$, are taken in the order in which they appear in $H_{I}(t)$. With this understanding, the result of moving all the annihilation operators to the right as described above is known as Wick's theorem:\\
$\mathcal{T}\left(A_{1} \cdots A_{n}\right)$ is equal to the sum of all operators obtained by contracting $k$ pairs of $A_{j}$ 's, $0 \leq k \leq[n / 2]$, and Wick-ordering the product of the remaining $A_{j}$ 's.

The proof is a straightforward but tedious induction on $n$, and we omit it. But to make the meaning clear, we shall write out the result for small $n$. For $n=2$, the theorem states that $\mathcal{T}\left(A_{1} A_{2}\right)=: A_{1} A_{2}:+\overbrace{A_{1} A_{2}}$, which is just the definition of $\overbrace{A_{1} A_{2}}$. For $n=3$, we have

$$
\mathcal{T}\left(A_{1} A_{2} A_{3}\right)=: A_{1} A_{2} A_{3}:+\overbrace{A_{1} A_{2}} A_{3}+\overbrace{A_{1} A_{3}} A_{2}+\overbrace{A_{2} A_{3}} A_{1}
$$

and for $n=4$ we have

$$
\begin{aligned}
\mathcal{T}\left(A_{1} A_{2} A_{3} A_{4}\right)= & : A_{1} A_{2} A_{3} A_{4}: \\
& +\overbrace{A_{1} A_{2}}: A_{3} A_{4}:+\overbrace{A_{1} A_{3}}: A_{2} A_{4}:+\overbrace{A_{1} A_{4}}: A_{2} A_{3}: \\
& +\overbrace{A_{2} A_{3}}: A_{1} A_{4}:+\overbrace{A_{2} A_{4}}: A_{1} A_{3}:+\overbrace{A_{3} A_{4}}: A_{1} A_{2}: \\
& +\overbrace{A_{1} A_{2}} \overbrace{A_{3} A_{4}}+\overbrace{A_{1} A_{3}} \overbrace{A_{2} A_{4}}+\overbrace{A_{1} A_{4}} \overbrace{A_{2} A_{3}} .
\end{aligned}
$$

As an immediate corollary of Wick's theorem and (6.39), we achieve our goal of evaluating the vacuum expectation value of a time-ordered product $\mathcal{T}\left(A_{1}, \ldots, A_{n}\right)$ :

The vacuum expectation value of $\mathcal{T}\left(A_{1} \cdots A_{n}\right)$ is

$$
\langle 0| \mathcal{T}\left(A_{1} \cdots A_{n}\right)|0\rangle= \begin{cases}0 & i_{n} n \text { is odd } \\ \sum \overbrace{A_{j_{1}} A_{j_{2}}} \overbrace{A_{j_{3}} A_{j_{4}}} \cdots \overbrace{A_{j_{n-1}} A_{j_{n}}} & \text { if } n \text { is even },\end{cases}
$$

where the sum in the even case is over all $1 \cdot 3 \cdots(n-1)$ ways of grouping the $A_{j}$ 's into $n / 2$ ordered pairs with the ordering in each pair being the same as in the original product ( $j_{1}<j_{2}, j_{3}<j_{4}$, etc. ).

Our interest with these vacuum expectation values is as ingredients in the Dyson series for the scattering matrix. In that situation the operators $A_{j}$ are either field operators (from the interaction Hamiltonians) or individual creation or annihilation operators (from the in and out states). The former are time-dependent; the latter are not. So we need to compute $\overbrace{A_{1} A_{2}}$ where each $A_{j}$ is either a creation or annihilation operator or a field operator. If both $A_{j}$ 's are individual creation or annihilation operators, the result is simply given by the canonical commutation relations (6.33). If one $A_{j}$ is a field operator, the calculation is still simple; we essentially did it in (6.36). To wit, if the field operator is

$$
\phi_{\pi}(x)=\sum_{\tau} \int f(\mathbf{q})\left[u(\mathbf{q}, \tau, \pi) a(\mathbf{q}, \tau, \pi) e^{-i q_{\mu} x^{\mu}}+v(\mathbf{q}, \tau, \pi) a^{\dagger}(\mathbf{q}, \tau, \bar{\pi}) e^{i q_{\mu} x^{\mu}}\right] \frac{d^{3} \mathbf{q}}{(2 \pi)^{3}}
$$

( $\bar{\pi}$ being the antiparticle of $\pi$, and $f(\mathbf{q})$ incorporating the remaining scalar factors such as $\sqrt{2 \omega_{\mathbf{q}}}$ ), then

$$
\begin{aligned}
& \overbrace{a\left(\mathbf{p}, \sigma, \pi^{\prime}\right) \phi_{\pi}(x)}=f(\mathbf{p}) v(\mathbf{p}, \sigma, \pi) e^{i p_{\mu} x^{\mu}} \delta_{\bar{\pi} \pi^{\prime}} \\
& \overbrace{\phi_{\pi}(x) a^{\dagger}\left(\mathbf{p}, \sigma, \pi^{\prime}\right)}=f(\mathbf{p}) u(\mathbf{p}, \sigma, \pi) e^{-i p_{\mu} x^{\mu}} \delta_{\pi \pi^{\prime}}
\end{aligned}
$$

(The contractions $\overbrace{a^{\dagger}(\mathbf{p}, \sigma, \pi) \phi\left(x_{\mu}\right)}$ and $\overbrace{\phi\left(x_{\mu}\right) a(\mathbf{p}, \sigma, \pi)}$ don't arise in practice, and in any case they both vanish.)

The interesting case is where both $A_{j}$ 's are field operators. Here each $A_{j}$ involves annihilation operators for some species of particle and creation operators for its antiparticle. Unless the particles of one $A_{j}$ are the antiparticles for the other one, these operators all commute or anticommute, with the result that the contraction vanishes. Thus the only nontrivial contractions are those involving a field and its adjoint. These contractions (up to a conventional factor of $i$ ) are known as propagators for the field in question.

The reason for the name is that by (6.39) and (6.41),

$$
\overbrace{\phi(t, \mathbf{x}) \phi^{\dagger}(s, \mathbf{y})}=\langle 0| \mathcal{T}\left[\phi(t, \mathbf{x}) \phi^{\dagger}(s, \mathbf{y})\right]|0\rangle .
$$

The expression on the right is the amplitude for the field to create a particle out of the vacuum at $\mathbf{y}$ at time $s$ and annihilate it at $\mathbf{x}$ at time $t$ if $t>s$, and the amplitude for the field to create an antiparticle out of the vacuum at $\mathbf{x}$ at time $t$ and annihilate it at $\mathbf{y}$ at time $s$ if $s>t$; in either case the particle or antiparticle "propagates" from $\mathbf{y}$ to $\mathbf{x}$ or from $\mathbf{x}$ to $\mathbf{y}$. We shall see that this interpretation remains valid when propagators are incorporated into more realistic processes.

We must now pause to study propagators in some detail.

\subsection*{Propagators}
We begin by considering a neutral scalar field $\phi$ of mass $m$ :

$$
\phi(t, \mathbf{x})=\int \frac{1}{\sqrt{2 \omega_{\mathbf{p}}}}\left[e^{-i \omega_{\mathbf{p}} t+i \mathbf{p} \cdot \mathbf{x}} a(\mathbf{p})+e^{i \omega_{\mathbf{p}} t-i \mathbf{p} \cdot \mathbf{x}} a^{\dagger}(\mathbf{p})\right] \frac{d^{3} \mathbf{p}}{(2 \pi)^{3}}
$$

If $t_{1}>t_{2}, \overbrace{\phi\left(t_{1}, \mathbf{x}_{1}\right) \phi\left(t_{2}, \mathbf{x}_{2}\right)}$ is equal to

$$
\begin{aligned}
& \iint \frac{1}{2 \sqrt{\omega_{\mathbf{p}_{1}} \omega_{\mathbf{p}_{2}}}} e^{-i \omega_{\mathbf{p}_{1}} t_{1}+i \mathbf{p}_{1} \cdot \mathbf{x}_{1}+i \omega_{\mathbf{p}_{2}} t_{2}-i \mathbf{p}_{2} \cdot \mathbf{x}_{2}}\left[a\left(\mathbf{p}_{1}\right), a^{\dagger}\left(\mathbf{p}_{2}\right)\right] \frac{d^{3} \mathbf{p}_{1} d^{3} \mathbf{p}_{2}}{(2 \pi)^{6}} \\
= & \iint \frac{1}{2 \sqrt{\omega_{\mathbf{p}_{1}} \omega_{\mathbf{p}_{2}}}} e^{-i \omega_{\mathbf{p}_{1}} t_{1}+i \mathbf{p}_{1} \cdot \mathbf{x}_{1}+i \omega_{\mathbf{p}_{2}} t_{2}-i \mathbf{p}_{2} \cdot \mathbf{x}_{2}} \delta\left(\mathbf{p}_{1}-\mathbf{p}_{2}\right) \frac{d^{3} \mathbf{p}_{1} d^{3} \mathbf{p}_{2}}{(2 \pi)^{3}} \\
= & \int \frac{1}{2 \omega_{\mathbf{p}}} e^{-i \omega_{\mathbf{p}}\left(t_{1}-t_{2}\right)+i \mathbf{p} \cdot\left(\mathbf{x}_{1}-\mathbf{x}_{2}\right)} \frac{d^{3} \mathbf{p}}{(2 \pi)^{3}} .
\end{aligned}
$$

If $t_{1}<t_{2}$ we must replace $t_{1}-t_{2}$ in this last integral by $t_{2}-t_{1}$, but the $\mathbf{x}_{1}-\mathbf{x}_{2}$ can remain as it is because the substitution $\mathbf{p} \rightarrow-\mathbf{p}$ shows that the integral is even in $\mathbf{x}_{1}-\mathbf{x}_{2}$. In short,

$$
\overbrace{\phi\left(t_{1}, \mathbf{x}_{1}\right) \phi\left(t_{2}, \mathbf{x}_{2}\right)}=-i \Delta_{F}\left(t_{1}-t_{2}, \mathbf{x}_{1}-\mathbf{x}_{2}\right),
$$

where

$$
\Delta_{F}(t, \mathbf{x})=i \int \frac{e^{-i \omega_{\mathbf{p}}|t|+i \mathbf{p} \cdot \mathbf{x}}}{2 \omega_{\mathbf{p}}} \frac{d^{3} \mathbf{p}}{(2 \pi)^{3}}
$$

This is to be interpreted as the tempered distribution on $\mathbb{R}^{4}$ whose Fourier transform in the space variables $\mathbf{x}$ is

$$
\left(\mathcal{F}_{\mathbf{x}} \Delta_{F}\right)(t, \mathbf{p})=\frac{i e^{-i \omega_{\mathbf{p}}|t|}}{2 \omega_{\mathbf{p}}} .
$$

(The foregoing derivation of (6.44) is easily made rigorous by integrating both sides against a test function $f\left(\mathbf{x}_{1}\right) g\left(\mathbf{x}_{2}\right)$.)\\
$\Delta_{F}$ is called the Feynman propagator. It is both physically and mathematically interesting. Here is a list of its most important properties:\\
i. $\Delta_{F}$ is a fundamental solution of the Klein-Gordon operator:

$$
\left(\partial_{t}^{2}-\nabla_{\mathbf{x}}^{2}+m^{2}\right) \Delta_{F}=\delta
$$

ii. The full space-time Fourier transform of $\Delta_{F}$ is

$$
\widehat{\Delta}_{F}(\omega, \mathbf{p})=\lim _{\epsilon \rightarrow 0+} \frac{1}{-\omega^{2}+|\mathbf{p}|^{2}+m^{2}-i \epsilon},
$$

that is,

$$
\widehat{\Delta}_{F}(p)=\lim _{\epsilon \rightarrow 0+} \frac{1}{-p^{2}+m^{2}-i \epsilon}
$$

where the limit is taken in the weak topology of tempered distributions.\\
iii. $\Delta_{F}$ is invariant under the full Lorentz group.\\
iv. $\Delta_{F}$ is a $C^{\infty}$ function away from the light cone, and its support is $\mathbb{R}^{4}$.

Let us prove these assertions. (Mathematical readers should find this a refreshing change of pace!) For (i), by taking the Fourier transform in $\mathbf{x}$ it is equivalent to show that $g(t)=i e^{-i \omega_{\mathbf{p}}|t|}$ satisfies $\left(\partial_{t}^{2}+|\mathbf{p}|^{2}+m^{2}\right) g(t)=2 \omega_{\mathbf{p}} \delta(t)$. This is elementary: $\partial_{t} g=\omega_{\mathbf{p}}(\operatorname{sgn} t) e^{-i \omega_{\mathbf{p}}|t|}$, so

$$
\partial_{t}^{2} g(t)=2 \omega_{\mathbf{p}} \delta(t) e^{-i \omega_{\mathbf{p}}|t|}-i \omega_{\mathbf{p}}^{2} e^{-i \omega_{\mathbf{p}}|t|}=2 \omega_{\mathbf{p}} \delta(t)-\left(|\mathbf{p}|^{2}+m^{2}\right) g(t)
$$

Next, given $\epsilon>0$ and $\mathbf{p} \in \mathbb{R}^{3}$, let $a=a(\epsilon, \mathbf{p})$ be the square root of $i \epsilon-|\mathbf{p}|^{2}-m^{2}$ with positive real part. Then, for $\omega \in \mathbb{R}$,

$$
\begin{aligned}
\frac{1}{-\omega^{2}+|\mathbf{p}|^{2}+m^{2}-i \epsilon} & =-\frac{1}{2 a}\left[\frac{1}{a+i \omega}+\frac{1}{a-i \omega}\right] \\
& =-\frac{1}{2 a}\left[\int_{0}^{\infty} e^{-(a+i \omega) t} d t+\int_{-\infty}^{0} e^{(a-i \omega) t} d t\right]
\end{aligned}
$$

That is, as a function of $\omega,\left(-\omega^{2}+|\mathbf{p}|^{2}+m^{2}-i \epsilon\right)^{-1}$ is the Fourier transform in $t$ of $-(1 / 2 a) e^{-a|t|}$. Now let $\epsilon \rightarrow 0+$ : we have $a \rightarrow i \sqrt{|\mathbf{p}|^{2}+m^{2}}=i \omega_{\mathbf{p}}$, so $-(1 / 2 a) e^{-a|t|} \rightarrow\left(i / 2 \omega_{\mathbf{p}}\right) e^{-i \omega_{\mathbf{p}}|t|}$. This convergence takes place in the topology of tempered distributions on $\mathbb{R}^{4}$, by the dominated convergence theorem, so (ii) follows from (6.46). Also, the function $h(\omega, \mathbf{p})=\left(-\omega^{2}+|\mathbf{p}|^{2}+m^{2}-i \epsilon\right)^{-1}$ is manifestly Lorentz-invariant on $\mathbb{R}^{4}$; hence so is its limit as $\epsilon \rightarrow 0+$, and so is the latter's Fourier transform $\Delta_{F}$, which proves (iii).

One can compute $\Delta_{F}$ explicitly by replacing $|t|$ by $|t|-i \epsilon$ in (6.45), evaluating the resulting absolutely convergent integral, and letting $\epsilon \rightarrow 0+$ at the end of the calculation. Since $e^{-\omega_{\mathbf{p}}(\epsilon+i|t|)} / \omega_{\mathbf{p}}$ is a radial function of $\mathbf{p}$, the integration reduces to a one-dimensional one that is part of the classic lore of special functions, and the result is that

$$
\Delta_{F}(t, \mathbf{x})=\frac{i m}{4 \pi^{2} \sqrt{|\mathbf{x}|^{2}-t^{2}}} K_{1}\left(m \sqrt{|\mathbf{x}|^{2}-t^{2}}\right) \text { on } \mathbb{R}^{4} \backslash\{(t, \mathbf{x}):|t|=|\mathbf{x}|\}
$$

where $K_{1}$ is the modified Bessel function of order 1 . We omit the details since we shall have no use for this explicit formula, but it implies (iv).

The fact that $\Delta_{F}$ is $C^{\infty}$ on the complement of the light cone can also be deduced from general considerations about differential equations (see Folland [46], $\oint 8 \mathrm{G})$. Since $\Delta_{F}$ is a fundamental solution for the Klein-Gordon operator, its wave front set is contained in the characteristic variety of that operator (i.e., the union of the light cones in the cotangent spaces at all points of $\mathbb{R}^{4}$ ) together with the wave front set of $\delta$ (i.e., the cotangent space at the origin). On the other hand, since $\Delta_{F}$ is Lorentz-invariant, its wave front set is contained in the union of the conormal bundles of the orbits of the Lorentz group. Since the normal to the orbit through\\
( $t, \mathbf{x}$ ) does not lie on the light cone unless ( $t, \mathbf{x}$ ) itself does, $\Delta_{F}$ has no wave front set away from the light cone and is therefore $C^{\infty}$ there.

Important remark: It is common in the physics literature to suppress the " $\lim _{\epsilon \rightarrow 0+}$ " in (6.47) and (6.48). In general, in formulas of this sort that contain an unspecified $\epsilon$, a " $\lim _{\epsilon \rightarrow 0+}$ " should be understood, with the limit taken in a suitable sense. We shall often avail ourselves of this bit of shorthand.

Property (i) is the reason for factoring out the $-i$ in the definition of $\Delta_{F}$.\\
By the way, the preceding results apply equally to the case $m=0$, where $\Delta_{F}$ is a fundamental solution of the wave operator $\partial_{t}^{2}-\nabla^{2}$. It is quite different from the fundamental solution that usually arises in classical wave mechanics, namely, the one supported on the forward light cone. Let us denote that solution by $G_{r}$ (for "retarded Green's function"), and let $G_{a}$ ("advanced Green's function") be the corresponding solution supported on the backward light cone, $G_{a}(t, \mathbf{x})=G_{r}(-t, \mathbf{x})$. Then $\Delta_{F}-\frac{1}{2}\left(G_{r}+G_{a}\right)$ is a solution of the homogeneous wave equation; in fact,

$$
\left.\left.\Delta_{F}(t, \mathbf{x})\right|_{m=0}=\frac{1}{2}\left(G_{r}+G_{a}\right)(t, \mathbf{x})+\left.\frac{i}{4 \pi^{2}}\left(\partial_{t}^{2}-\nabla_{\mathbf{x}}^{2}\right) \log | | \mathbf{x}\right|^{2}-t^{2} \right\rvert\,
$$

which agrees off the light cone with the function $i / 4 \pi^{2}\left(|\mathbf{x}|^{2}-t^{2}\right)$. It may be instructive to compare the Fourier transforms of $G_{r}$ and $G_{a}$ with that of $\Delta_{F}$ :

$$
\begin{gathered}
\widehat{G}_{r}(\omega, \mathbf{p})=\lim _{\epsilon \rightarrow 0+} \frac{1}{-(\omega-i \epsilon)^{2}+|\mathbf{p}|^{2}}, \quad \widehat{G}_{a}(\omega, \mathbf{p})=\lim _{\epsilon \rightarrow 0+} \frac{1}{-(\omega+i \epsilon)^{2}+|\mathbf{p}|^{2}} \\
\widehat{\Delta}_{F}(\omega, \mathbf{p})=\lim _{\epsilon \rightarrow 0+} \frac{1}{-\omega^{2}+|\mathbf{p}|^{2}-i \epsilon}
\end{gathered}
$$

See Folland [47] for a complete derivation of these facts.\\
We return to the calculation of contractions of field operators. For a charged scalar field

$$
\phi(x)=\int \frac{1}{\sqrt{2 \omega_{\mathbf{p}}}}\left[e^{-i p_{\mu} x^{\mu}} a(\mathbf{p})+e^{i p_{\mu} x^{\mu}} b^{\dagger}(\mathbf{p})\right] \frac{d^{3} \mathbf{p}}{(2 \pi)^{3}},
$$

where $a$ and $b$ are the annihilation operators for particles and antiparticles, the calculation is similar to the neutral case. Since $\left[a(\mathbf{p}), b^{\dagger}\left(\mathbf{p}^{\prime}\right)\right]=0$ for any $\mathbf{p}, \mathbf{p}^{\prime}$, whereas $\left[a(\mathbf{p}), a^{\dagger}\left(\mathbf{p}^{\prime}\right)\right]=\left[b(\mathbf{p}), b^{\dagger}\left(\mathbf{p}^{\prime}\right)\right]=(2 \pi)^{3} \delta\left(\mathbf{p}-\mathbf{p}^{\prime}\right)$, it follows easily that

$$
\begin{gathered}
\overbrace{\phi\left(t_{1}, \mathbf{x}_{1}\right) \phi^{\dagger}\left(t_{2}, \mathbf{x}_{2}\right)}^{\phi\left(t_{1}, \mathbf{x}_{1}\right)}=\overbrace{\phi\left(t_{2}, \mathbf{x}_{2}\right)}^{\phi^{\dagger}\left(t_{1}, \mathbf{x}_{1}\right)}=\overbrace{\phi\left(t_{2}, \mathbf{x}_{2}\right)}^{\phi^{\dagger}\left(t_{1}, \mathbf{x}_{1}\right) \phi^{\dagger}\left(t_{2}, \mathbf{x}_{2}\right)}=-i \Delta_{F}\left(t_{1}-t_{2}, \mathbf{x}_{1}-\mathbf{x}_{2}\right) \\
.
\end{gathered}
$$

Thus the propagator $-i \Delta_{F}$ is the same as for the neutral field.\\
For Dirac fields, the important quantity is the contraction of a component $\psi_{j}$ of a Dirac field $\psi$ with a component $\bar{\psi}_{k}$ of its Dirac adjoint $\bar{\psi}=\psi^{\dagger} \gamma_{0}(j, k=1, \ldots, 4)$ (in either order). We have

$$
\begin{aligned}
\overbrace{\psi_{j}\left(t_{1}, \mathbf{x}_{1}\right) \bar{\psi}_{k}\left(t_{2}, \mathbf{x}_{2}\right)} & =\mathcal{T}\left[\psi_{j}\left(t_{1}, \mathbf{x}_{1}\right) \bar{\psi}_{k}\left(t_{2}, \mathbf{x}_{2}\right)\right]-: \psi_{j}\left(t_{1}, \mathbf{x}_{1}\right) \bar{\psi}_{k}\left(t_{2}, \mathbf{x}_{2}\right): \\
& = \begin{cases}\left\{\psi_{j}^{a}\left(t_{1}, \mathbf{x}_{1}\right), \bar{\psi}_{k}^{c}\left(t_{2}, \mathbf{x}_{2}\right)\right\} & \text { if } t_{1}>t_{2} \\
-\left\{\psi_{k}^{a}\left(t_{2}, \mathbf{x}_{2}\right), \bar{\psi}_{j}^{c}\left(t_{1}, \mathbf{x}_{1}\right)\right\} & \text { if } t_{1}<t_{2}\end{cases}
\end{aligned}
$$

where the superscripts $a$ and $c$ indicate the parts of the fields involving annihilation and creation operators, respectively. These anticommutators are evaluated by using the Fourier expansions (5.32)-(5.33) of the fields and the spin sums (5.34) in much\\
the same way as in the calculation of (5.35). In fact, setting $t=t_{1}-t_{2}$ and $\mathbf{x}=\mathbf{x}_{1}-\mathbf{x}_{2}$, for $t_{1}>t_{2}$ we have

$$
\begin{aligned}
\overbrace{\psi_{j}\left(t_{1}, \mathbf{x}_{1}\right) \bar{\psi}_{k}\left(t_{2}, \mathbf{x}_{2}\right)} & =\int \sum_{s= \pm} \frac{m}{2 \omega_{\mathbf{p}}} e^{-i \omega_{\mathbf{p}} t+i \mathbf{p} \cdot \mathbf{x}} u(\mathbf{p}, s)_{j} \bar{u}(\mathbf{p}, s)_{k} \frac{d^{3} \mathbf{p}}{(2 \pi)^{3}} \\
& =\int \frac{m}{2 \omega_{\mathbf{p}}} e^{-i \omega_{\mathbf{p}} t+i \mathbf{p} \cdot \mathbf{x}}\left(m^{-1} p_{\mu} \gamma^{\mu}+I\right)_{j k} \frac{d^{3} \mathbf{p}}{(2 \pi)^{3}} \\
& =\left(i \gamma^{\mu} \partial_{\mu}+m I\right)_{j k} \int \frac{e^{-i \omega_{\mathbf{p}} t+i \mathbf{p} \cdot \mathbf{x}}}{2 \omega_{\mathbf{p}}} \frac{d^{3} \mathbf{p}}{(2 \pi)^{3}}
\end{aligned}
$$

and by (6.45), the last integral is equal to $-i \Delta_{F}(t, \mathbf{x})$. On the other hand, for $t_{1}<t_{2}$ we have

$$
\begin{aligned}
\overbrace{\psi_{j}\left(t_{1}, \mathbf{x}_{1}\right) \bar{\psi}_{k}\left(t_{2}, \mathbf{x}_{2}\right)} & =-\int \sum_{s= \pm} \frac{m}{2 \omega_{\mathbf{p}}} e^{i \omega_{\mathbf{p}} t-i \mathbf{p} \cdot \mathbf{x}} \bar{v}(\mathbf{p}, s)_{k} v(\mathbf{p}, s)_{j} \frac{d^{3} \mathbf{p}}{(2 \pi)^{3}} \\
& =-\int \frac{m}{2 \omega_{\mathbf{p}}} e^{i \omega_{\mathbf{p}} t-i \mathbf{p} \cdot \mathbf{x}}\left(m^{-1} p_{\mu} \gamma^{\mu}-m I\right)_{j k} \frac{d^{3} \mathbf{p}}{(2 \pi)^{3}} \\
& =\left(i \gamma^{\mu} \partial_{\mu}+m I\right)_{j k} \int \frac{e^{i \omega_{\mathbf{p}} t-i \mathbf{p} \cdot \mathbf{x}}}{2 \omega_{\mathbf{p}}} \frac{d^{3} \mathbf{p}}{(2 \pi)^{3}}
\end{aligned}
$$

The substitution $\mathbf{p} \rightarrow-\mathbf{p}$ shows that the $e^{-i \mathbf{p} \cdot \mathbf{x}}$ can be replaced by $e^{i \mathbf{p} \cdot \mathbf{x}}$ in this last integral, so it is again equal to $-i \Delta_{F}(t, \mathbf{x})$. The calculation with $\psi$ and $\bar{\psi}$ switched is similar, and the result is the same except for a minus sign.

We have shown that

$$
\overbrace{\psi_{j}\left(t_{1}, \mathbf{x}_{1}\right) \bar{\psi}_{k}\left(t_{2}, \mathbf{x}_{2}\right)}=-\overbrace{\bar{\psi}_{j}\left(t_{1}, \mathbf{x}_{1}\right) \psi_{k}\left(t_{2}, \mathbf{x}_{2}\right)}=-i\left[\Delta_{\mathrm{Dirac}}\left(t_{1}-t_{2}, \mathbf{x}_{1}-\mathbf{x}_{2}\right)\right]_{j k},
$$

where the Dirac propagator $\Delta_{\text {Dirac }}$, a $4 \times 4$ matrix-valued distribution, is

$$
\Delta_{\text {Dirac }}=\left(i \gamma^{\mu} \partial_{\mu}+m\right) \Delta_{F},
$$

that is,

$$
\widehat{\Delta}_{\mathrm{Dirac}}(p)=\frac{\left(p_{\mu} \gamma^{\mu}+m\right)}{-p^{2}+m^{2}-i \epsilon}
$$

In this last formula there is an implicit " $\lim _{\epsilon \rightarrow 0+}$," as we described in connection with $\Delta_{F}$.

The Fourier-transformed propagator (6.51) is a matrix-valued distribution that agrees with the function $\left(p_{\mu} \gamma^{\mu}+m\right) /\left(-p^{2}+m^{2}\right)$ off the mass shell $p^{2}=m^{2}$. Since

$$
\left(p_{\mu} \gamma^{\mu}+m\right)\left(p_{\mu} \gamma^{\mu}-m\right)=p^{2}-m^{2}
$$

one often finds (6.51) rewritten as

$$
-i \widehat{\Delta}_{\text {Dirac }}(p)=\frac{i}{p_{\mu} \gamma^{\mu}-m} \quad\left(p^{2} \neq m^{2}\right)
$$

(Of course, the expression on the right denotes the matrix $i\left(p_{\mu} \gamma^{\mu}-m I\right)^{-1}$.) For the same reason, $\Delta_{\text {Dirac }}$ is a fundamental solution for the Dirac operator $-i \gamma^{\mu} \partial_{\mu}+m$.

Similar calculations yield the propagators for field operators for fields of arbitrary spin; see Weinberg [131], §5.7. Their Fourier transforms all turn out to be of the form $F(p) /\left(-p^{2}+m^{2}-i \epsilon\right)$ where $F$ is a polynomial with values in the appropriate spinor space. The one we will need for quantum electrodynamics is the\\
photon field, which is a spin-1 field of mass 0 . Unfortunately, the absence of mass complicates the picture, and a detailed analysis of the situation is rather difficult. Readers may consult physics texts for more information, but they will not find any unanimity about the best way to explain this matter. For the present we content ourselves with a brief heuristic discussion leading to the final result. We shall give another derivation of it - still nonrigorous but perhaps more convincing - in §8.3.

Let us begin with the easier case of a massive vector field instead. Just as the propagators for spin- 0 and spin- $\frac{1}{2}$ fields are given by fundamental solutions for the corresponding field operators - the Klein-Gordon and Dirac operators - so we expect the propagator for the massive vector field to be given by a fundamental solution of the Proca operator. The Proca equation (5.39) equation can be written

$$
\left(-g^{\mu \nu}\left(\partial^{2}+m^{2}\right)+\partial^{\mu} \partial^{\nu}\right) A_{\mu}=0
$$

so we are looking for a fundamental solution of $-g^{\mu \nu}\left(\partial^{2}+m^{2}\right)+\partial^{\mu} \partial^{\nu}$, that is, a matrix-valued distribution whose Fourier transform evaluated at $p$ is the inverse of the matrix

$$
D^{\mu \nu}=g^{\mu \nu}\left(p^{2}-m^{2}\right)-p^{\mu} p^{\nu}
$$

This inverse is readily computed:

$$
\left(D^{-1}\right)_{\mu \nu}=\frac{-g_{\mu \nu}+p_{\mu} p_{\nu} / m^{2}}{-p^{2}+m^{2}}
$$

and indeed, the propagator for the massive vector field is given by

$$
\Delta_{\mu \nu}^{\text {Proca }}=-\left(g_{\mu \nu}+\partial_{\mu} \partial_{\nu} / m^{2}\right) \Delta_{F}, \text { i.e., } \widehat{\Delta}_{\mu \nu}^{\text {Proca }}(p)=\frac{-g_{\mu \nu}+p_{\nu} p_{\mu} / m^{2}}{-p^{2}+m^{2}-i \epsilon} .
$$

Of course, to verify this, and to see that the overall sign and normalization are correct, one must actually calculate the relevant contractions; the reader may try this as an exercise or consult Weinberg [131], §5.3.

Now, (6.53) blows up as $m \rightarrow 0$; the trouble is that when $m=0$ the matrix $D_{\mu \nu}$ is not invertible, for $D_{\mu \nu} p^{\nu}=0$. This has to do, once again, with the fact that the elecromagnetic field cannot propagate longitudinal waves, and with the gaugeinvariance of the theory. If one imposes the Landau gauge condition $\partial_{\mu} A^{\mu}=0$, however, the term $\partial_{\mu} \partial_{\nu} A^{\mu}$ in Maxwell's equation disappears. In fact, under this condition one can replace the Maxwell equation $\left(-g^{\mu \nu} \partial^{2}+\partial^{\mu} \partial^{\nu}\right) A_{\mu}=0$ by

$$
\left(-g^{\mu \nu} \partial^{2}+(1-b) \partial^{\mu} \partial^{\nu}\right) A_{\mu}=0
$$

for any $b \in \mathbb{R}$, and as long as $b \neq 0$ the corresponding matrix

$$
D_{(b)}^{\mu \nu}=g^{\mu \nu} p^{2}+(b-1) p^{\mu} p^{\nu}
$$

is invertible: its inverse is

$$
\frac{g_{\mu \nu}+(a-1) p_{\mu} p_{\nu} / p^{2}}{p^{2}}, \quad a=b^{-1}
$$

And, in fact, this device gives the right answer: in computing S-matrix elements one can take the contractions of photon fields to be given by

$$
\overbrace{A_{\mu}\left(t_{1}, \mathbf{x}_{1}\right) A_{\nu}\left(t_{2}, \mathbf{x}_{2}\right)} \longrightarrow-i \Delta_{\mu \nu}^{(a)}\left(t_{1}-t_{2}, \mathbf{x}_{1}-\mathbf{x}_{2}\right),
$$

where

$$
\widehat{\Delta}_{\mu \nu}^{(a)}(p)=\frac{g_{\mu \nu}+(a-1) p_{\mu} p_{\nu} / p^{2}}{p^{2}+i \epsilon}
$$

We write an arrow rather than an equal sign in (6.54) because $a$ can be taken to be any real number. This indeterminacy may seem unsettling, but as it turns out, the gauge-invariance of the theory guarantees that the contributions of the $p_{\mu} p_{\nu}$ term in the photon propagators ultimately cancel out when one computes a physically meaningful quantity. (We shall explain this more fully in §7.7.) As a result, one can use the freedom in choosing $a$ to simplify calculations as needed. In particular, the default choice is $a=1$ (sometimes called Feynman gauge), which gives the photon propagator

$$
-i \Delta_{\mu \nu}=\left.i g_{\mu \nu} \Delta_{F}\right|_{m=0}, \quad \text { i.e., } \quad \widehat{\Delta}_{\mu \nu}(p)=\frac{g_{\mu \nu}}{p^{2}+i \epsilon}
$$

Propagators and the spin-statistics theorem. The terms of the Dyson series (6.24) are integrals over space-time of time-ordered products of interaction Hamiltonian densities $\mathcal{H}_{I}\left(x_{j}\right)\left(x_{j} \in \mathbb{R}^{4}\right)$. If these integrals are to have a relativistically invariant meaning, it is necessary to examine the time-ordering in more detail. For $x, y \in \mathbb{R}^{4}$ it makes Lorentz-invariant sense to say that $x^{0}>y^{0}$ if $x-y$ is timelike $\left((x-y)^{2}>0\right)$, but if $x-y$ is spacelike there are Lorentz transformations that reverse the sign of $x^{0}-y^{0}$ or make it vanish. Since the time-ordered products are to be integrated over all space-time, the result will have a relativistically invariant meaning only if the time ordering is irrelevant when the space-time points in question are spacelike separated. In other words, we must require of the Hamiltonian density $\mathcal{H}_{I}$ that

$$
\left[\mathcal{H}_{I}(x), \mathcal{H}_{I}(y)\right]=0 \text { when }(x-y)^{2} \leq 0 .
$$

Because of the way $\mathcal{H}_{I}$ is formed from products of field operators, the way to guarantee this is to require that the field operators $\phi(x)$ and $\phi(y)$ commute or anticommute when $x-y$ is spacelike. More precisely, if $\phi$ is a field that occurs an odd number of times in some term in $\mathcal{H}_{I}$ (which can only happen if $\phi=\phi^{\dagger}$ ), we must have $[\phi(x), \phi(y)]=0$ when $x-y$ is spacelike. On the other hand, if $\phi$ and $\phi^{\dagger}$ together occur an even number of times in each term, it is enough to have either $[\phi(x), \phi(y)]=0$ or $\{\phi(x), \phi(y)\}=0$ for $x-y$ spacelike and likewise with $\phi^{\dagger}(y)$ instead of $\phi(y)$, as two minus signs will cancel in the latter case. (For example, in the interaction for QED, the electromagnetic field occurs once and the Dirac field occurs twice; the field $\phi$ occurs four times in the $\phi^{4}$ interaction.)

Calculating these commutators and anticommutators is easy enough; it is essentially the same as the calculation of contractions that we did earlier, except that the time-ordering is absent. For a neutral scalar field $\phi$ of mass $m$ given by (5.11), with $p=\left(\omega_{\mathbf{p}}, \mathbf{p}\right)$ we have

$$
\begin{aligned}
{[\phi(x), \phi(y)]=} & \int \frac{e^{-i p_{\mu} x^{\mu}+i q_{\mu} y^{\mu}}}{2 \sqrt{\omega_{\mathbf{p}} \omega_{\mathbf{q}}}}\left[a(\mathbf{p}), a^{\dagger}(\mathbf{q})\right] \frac{d^{3} \mathbf{p} d^{3} \mathbf{q}}{(2 \pi)^{6}} \\
& +\int \frac{e^{i p_{\mu} x^{\mu}-i q_{\mu} y^{\mu}}}{2 \sqrt{\omega_{\mathbf{p}} \omega_{\mathbf{q}}}}\left[a^{\dagger}(\mathbf{p}), a(\mathbf{q})\right] \frac{d^{3} \mathbf{p} d^{3} \mathbf{q}}{(2 \pi)^{6}}
\end{aligned}
$$

and since $\left[a(\mathbf{p}), a^{\dagger}(\mathbf{q})\right]=-\left[a^{\dagger}(\mathbf{p}), a(\mathbf{q})\right]=(2 \pi)^{3} \delta(\mathbf{p}-\mathbf{q})$, this boils down to

$$
[\phi(x), \phi(y)]=-i \Delta_{+}(x-y)+i \Delta_{+}(y-x),
$$

where

$$
\Delta_{+}(x)=i \int \frac{e^{-i p_{\mu} x^{\mu}}}{2 \omega_{\mathbf{p}}} \frac{d^{3} \mathbf{p}}{(2 \pi)^{3}}
$$

$\Delta_{+}$is just like the Feynman propagator $\Delta_{F}$ except that one has $x^{0}=t$ in the exponent instead of $|t|$; in particular, the Fourier integral is to be interpreted in the sense of distributions as before. (Rather than being a fundamental solution for the Klein-Gordon operator like $\Delta_{F}$, though, $\Delta_{+}$is a solution of the homogeneous Klein-Gordon equation.)\\
$\Delta_{+}$is just $i$ times the Fourier transform of the invariant measure $d^{3} \mathbf{p} / 2 \omega_{\mathbf{p}}$ on the positive mass shell $p_{0}=\omega_{\mathbf{p}}$, and hence it is invariant under the orthochronous Lorentz group. (It is not invariant under time inversion, which interchanges the mass shell with $p_{0}=-\omega_{\mathbf{p}}$.) Moreover, it agrees with $\Delta_{F}$ in the region $t>0$ (where $|t|=t$ ), and hence - since $\Delta_{+}$and $\Delta_{F}$ are both invariant under the orthochronous Lorentz group - it agrees with $\Delta_{F}$ in the whole region outside the backward light cone, and in particular in the spacelike region. Therefore, if $x-y$ is spacelike,

$$
[\phi(x), \phi(y)]=-i \Delta_{+}(x-y)+i \Delta_{+}(y-x)=-i \Delta_{F}(x-y)+i \Delta_{F}(y-x)=0
$$

because $\Delta_{F}$ is even. This is exactly what we were looking for.\\
By the same calculation, for a charged scalar field we have $\left[\phi(x), \phi^{\dagger}(y)\right]=0$ when $x-y$ is spacelike. (That $[\phi(x), \phi(y)]=0$ is automatic since $\phi$ contains annihilation operators for a particle and creation operators for its antiparticle; these operators commute.)

However, suppose we tried to construct a quantum theory involving Fermions with spin zero. The field describing these Fermions would be constructed out of creation and annihilation operators that satisfy anticommutation relations, so we would have to use anticommutators instead of commutators in this calculation. With this change, the calculation proceeds much as before, except that one crucial minus sign changes to a plus sign, giving

$$
\{\phi(x), \phi(y)\}=-i \Delta_{+}(x-y)-i \Delta_{+}(y-x)=-2 i \Delta_{F}(x-y) \text { for } x-y \text { spacelike. }
$$

As we have observed earlier, $\Delta_{F}$ is not zero in the spacelike region, so relativistic invariance is destroyed. Particles of spin zero must be Bosons.

For a Dirac field $\psi$ of mass $m$, calculations similar to the ones above show that

$$
\left\{\psi_{a}(x), \bar{\psi}_{b}(y)\right\}=-i\left[\left(i \gamma^{\mu} \partial_{\mu}+m\right) \Delta_{+}(x-y)+\left(i \gamma^{\mu} \partial_{\mu}-m\right) \Delta_{+}(y-x)\right]_{a b} .
$$

Since $\Delta_{+}$is even in the spacelike region, the terms involving $m$ cancel out there, and since $\partial_{\mu} \Delta_{+}$is odd, so do the terms involving $\gamma^{\mu} \partial_{\mu}$. However, if we tried to construct a Dirac field out of creation and annihilation operators that satisfy commutation relations, we would have to consider $\left[\psi_{a}(x), \bar{\psi}_{b}(y)\right]$ instead, and it would be given by the right side of (6.57) with the plus sign between the two distributions replaced by a minus sign. Those two terms would add up instead of canceling, and relativistic invariance would be violated again. Particles of spin $\frac{1}{2}$ must be Fermions.

Similar considerations apply to particles of higher spin: for fields of spin $s$, the commutators or anticommutators will involve derivatives of order $2 s$ of $\Delta_{+}$, which will be even or odd according as $2 s$ is even or odd. Hence, in order to obtain relativistic invariance, one must use commutators when $2 s$ is even and anticommutators when $s$ is odd. This gives another way of arriving at the conclusion that particles of integer spin must be Bosons; particles of half-integer spin must be Fermions. See Weinberg [131], §5.7, for a more detailed treatment.

\subsection*{Feynman diagrams}
We now return to the Dyson series for a scattering matrix element between initial and final states consisting of free particles with definite momenta. The Nth term in the series is

$$
\begin{aligned}
\frac{(-i)^{N}}{N!} \int \cdots \int\langle 0| a_{1}\left(\mathbf{p}_{1}^{\text {out }}\right) \cdots a_{J}\left(\mathbf{p}_{J}^{\text {out }}\right) \mathcal{T} & {\left[\mathcal{H}_{I}\left(x_{1}\right) \cdots \mathcal{H}_{I}\left(x_{N}\right)\right] } \\
& \cdot a_{1}^{\dagger}\left(\mathbf{p}_{1}^{\text {in }}\right) \cdots a_{K}^{\dagger}\left(\mathbf{p}_{K}^{\text {in }}\right)|0\rangle d^{4} x_{1} \cdots d^{4} x_{N}
\end{aligned}
$$

Each $\mathcal{H}_{I}\left(x_{n}\right)$ is a sum of products of field operators, and (6.58) is the sum of the integrals obtained from it by replacing each $\mathcal{H}_{I}\left(x_{n}\right)$ by a single one of these products (including the relevant coupling constant). There can be any number of different fields here, and the initial and final creation and annihilation operators can be for any of the corresponding particle or antiparticle types. By the corollary (6.42) of Wick's theorem, when each $\mathcal{H}_{I}\left(x_{n}\right)$ is taken to consist of a single product, (6.58) is equal to $(-i)^{N} / N!$ times the sum of all the terms obtained by pairing up all the creation operators $a_{k}^{\dagger}\left(\mathbf{p}_{k}\right)$, the annihilation operators $a_{j}\left(\mathbf{p}_{j}^{\prime}\right)$, and the field operators in the $\mathcal{H}_{I}\left(x_{m}\right)$ in all possible ways and integrating the resulting product of contractions. Each such term is the integral, with respect to $x_{1}, \ldots, x_{N}$, of a product of contractions of the form

$$
\overbrace{a_{j}\left(\mathbf{p}_{j}^{\text {out }}\right) a_{k}^{\dagger}\left(\mathbf{p}_{k}^{\text {in }}\right)}, \quad \overbrace{a_{j}\left(\mathbf{p}_{j}^{\text {out }}\right) \phi\left(x_{n}\right)}, \quad \overbrace{\phi\left(x_{n}\right) a_{k}^{\dagger}\left(\mathbf{p}_{k}^{\text {in }}\right)}, \quad \overbrace{\phi\left(x_{n}\right) \phi^{\dagger}\left(x_{n^{\prime}}\right)},
$$

where $\phi$ is one of the fields contributing to $\mathcal{H}_{I}$. (The notation here is very strippeddown. If $\phi$ is a field involved in $\mathcal{H}_{I}$ then so is $\phi^{\dagger}$, so contractions of the form $\overbrace{a_{j}\left(\mathbf{p}_{j}^{\text {out }}\right) \phi^{\dagger}\left(x_{n}\right)}$ and $\overbrace{\phi^{\dagger}\left(x_{n}\right) a_{k}^{\dagger}\left(\mathbf{p}_{k}^{\text {in }}\right)}$ are also included in this scheme. Moreover, there may be other parameters indicating components of vectors, spin states, particle species, and other entities such as Dirac matrices, that are being suppressed. Finally, the $a$ 's and $a^{\dagger}$ 's are completely unrelated; for instance, $a_{1}\left(\mathbf{p}_{1}^{\text {out }}\right)$ and $a_{1}^{\dagger}\left(\mathbf{p}_{1}^{\text {in }}\right)$ may pertain to different spin states or different particles.)

We now know how to evaluate all of these:

\begin{itemize}
  \item $\overbrace{a_{j}\left(\mathbf{p}_{j}^{\text {out }}\right) a_{k}^{\dagger}\left(\mathbf{p}_{k}^{\text {in }}\right)}$ is zero if $a_{j}$ and $a_{k}^{\dagger}$ pertain to different spin states or different particles and is $(2 \pi)^{3} \delta\left(\mathbf{p}_{j}^{\text {out }}-\mathbf{p}_{k}^{\text {in }}\right)$ otherwise.
  \item $\overbrace{a_{j}\left(\mathbf{p}_{j}^{\text {out }}\right) \phi\left(x_{n}\right)}$ is zero unless $\phi$ creates particles of the type annihilated by $a_{j}$. In that case\\
$\phi(x)=\int e^{-i p_{\mu} x^{\mu}} u\left(\mathbf{p}_{j}\right) a_{j}^{\dagger}(\mathbf{p}) d^{3} \mathbf{p} /(2 \pi)^{3}+\cdots \quad\left(p=\left(\omega_{\mathbf{p}}, \mathbf{p}\right)\right)$,\\
where any further ingredients such as spin components and $\sqrt{2 \omega_{\mathbf{p}}}$ are incorporated in $u(\mathbf{p})$, and we have
\end{itemize}

$$
\overbrace{a_{j}\left(\mathbf{p}_{j}^{\text {out }}\right) \phi\left(x_{n}\right)}=e^{-i p_{\mu}^{\text {out }} x^{\mu}} u\left(\mathbf{p}_{j}^{\text {out }}\right) .
$$

\begin{itemize}
  \item Similarly, $\overbrace{\phi\left(x_{n}\right) a_{k}^{\dagger}\left(\mathbf{p}_{k}^{\mathrm{in}}\right)}$ is zero unless $\phi$ annihilates particles of the type created by $a_{k}^{\dagger}$. In that case
\end{itemize}

$$
\phi(x)=\int e^{-i p_{\mu} x^{\mu}} u\left(\mathbf{p}_{j}\right) a_{j}(\mathbf{p}) d^{3} \mathbf{p} /(2 \pi)^{3}+\cdots \quad\left(p=\left(\omega_{\mathbf{p}}, \mathbf{p}\right)\right)
$$

and we have

$$
\overbrace{\phi\left(x_{n}\right) a_{k}^{\dagger}\left(\mathbf{p}_{k}^{\mathrm{in}}\right)}=e^{i p_{\mu}^{\mathrm{in}} x^{\mu}} u\left(\mathbf{p}_{k}^{\mathrm{in}}\right) .
$$

\begin{itemize}
  \item Finally, $\overbrace{\phi\left(x_{n}\right) \phi^{\dagger}\left(x_{n^{\prime}}\right)}=-i \Delta\left(x_{n}-x_{n^{\prime}}\right)$ where $\Delta$ is the propagator for the field $\phi$.\\
So now it remains to multiply these together and integrate, but we postpone this task for a bit. First, we present the marvelous graphical interpretation of these integrals due to Richard Feynman, which makes it much easier to see what is going on and keep track of all the different terms.
\end{itemize}

Each product of nonzero contractions of the above types corresponds uniquely to a graph with some additional structure called a Feynman diagram, according to the following rules.

\begin{itemize}
  \item For each initial creation operator $a_{j}^{\dagger}\left(\mathbf{p}_{k}^{\mathrm{in}}\right)$, each final annihilation operator $a^{j}\left(\mathbf{p}_{j}^{\text {out }}\right)$, and each point $x_{n}$ where field operators act there is a vertex. The vertices corresponding to the initial (resp. final) operators are called initial (resp. final); initial and final vertices are both called external; and the vertices corresponding to the field operators are called internal. The external vertices are labeled with the on-mass-shell 4-momenta $p=\left(\omega_{\mathbf{p}}, \mathbf{p}\right)$ corresponding to their 3 -momenta $\mathbf{p}_{j}^{\text {out }}$ and $\mathbf{p}_{k}^{\text {in }}$, and the internal vertices are labeled with their points $x_{n}$; other parameters such as spin states and Lorentz indices may be included too. When drawing a Feynman diagram, it is customary to put the vertices for creation operators at the left or bottom of the picture, the vertices for annihilation operators at the right or top of the picture, and the internal vertices in between. In this book we almost always use the left-right convention.
  \item For each contraction there is an edge of the graph between the appropriate vertices: between the external vertices with labels $p_{j}^{\text {out }}$ and $p_{k}^{\text {in }}$ for $\overbrace{a_{j}\left(\mathbf{p}_{j}^{\text {out }}\right) a_{k}^{\dagger}\left(\mathbf{p}_{k}^{\text {in }}\right)}$, between the internal vertices with labels $x_{n}$ and $x_{n^{\prime}}$ for $\overbrace{\phi\left(x_{n}\right) \phi^{\dagger}\left(x_{n^{\prime}}\right)}$, etc. Edges with one end at an external vertex are called external lines; edges with both ends at internal vertices are called internal lines.
  \item Each line is associated to a particular particle species, namely, the one to which the operators in its contraction pertain. It should be labeled in some way to indicate this species. In practice this is done by using solid lines, dotted lines, wavy lines, etc., for the various species. (In particular, in QED electrons are denoted by solid lines and photons by wavy lines.) However, for this purpose, particles and antiparticles are considered as belonging to the same species; the distinction between then comes in the next item.
  \item At this point, for each particle-antiparticle pair that occurs in the calculation, one must decide once and for all which one is the "particle" and which one is the "antiparticle." (In practice this is a matter of iron-clad convention. In particular, electrons and protons are "particles," and the determination for other leptons and baryons is determined by laws to the effect that in any interaction the number of leptons [or baryons] minus the number of antileptons [or antibaryons] is conserved.) The lines of the\\
graph associated to particles with distinct antipaticles are equipped with arrows as follows. (No arrows are drawn on lines corresponding to a particle without a distinct antiparticle, such as a photon; however, arrows may be drawn adjacent to them to indicate the direction of momentum flow, as we shall explain later.)\\
i. The arrow of a line with one end at an initial vertex $p^{\text {in }}$ points away from that vertex if the entity created there is a particle and toward that vertex if it is an antiparticle.\\
ii. The arrow of a line with one end at a final vertex $p^{\text {out }}$ points toward from that vertex if the entity annihilated there is a particle and away from that vertex if it is an antiparticle. (Note that this is consistent with (i) for lines that join initial and final vertices.)\\
iii. The arrow of an internal line joining the vertices $x$ and $x^{\prime}$ points from $x^{\prime}$ to $x$ if the field $\phi$ in the associated contraction $\overbrace{\phi(x) \phi^{\dagger}\left(x^{\prime}\right)}$ annihilates particles and creates antiparticles, and from $x$ to $x^{\prime}$ if it creates particles and annihilates antiparticles.\\
We observe that the following conditions always hold:
  \item Every external vertex is contained in exactly one line, and the other end of that line is either an internal vertex or an external vertex of the opposite sort (initial or final).
  \item Recall that the integral giving rise to a Feynman diagram is obtained from (6.58) by replacing each $\mathcal{H}_{I}\left(x_{n}\right)$ by one of its constituent products of field operators. The number of lines that meet the vertex $x_{n}$ is the number of terms in that product. More precisely, if the contribution of a particular Hermitian field $\phi$ to the product is $\phi^{m}$, there will be $m$ lines of type $\phi$ at $x_{n}$; if the contribution of non-Hermitian field $\phi$ is $\left(\phi \phi^{\dagger}\right)^{m}$, there will be $m$ lines of type $\phi$ with arrows pointing toward the vertex $x_{n}$ and $m$ with arrows pointing away.\\
Part of what makes this graphical representation useful is that one can go the other way, from Feynman diagrams to terms in the Dyson series. To wit, suppose one is given an interaction Hamiltonian $\mathcal{H}_{I}$ of the sort described above and some initial and final states of free particles with definite momenta. One builds a Feynman diagram by starting with $K$ initial vertices labeled with momenta $p_{k}^{\text {in }}, J$ final vertices labeled with momenta $p_{j}^{\text {out }}$, and $N$ internal vertices labeled with spacetime points $x_{n}$, and choosing one product in $\mathcal{H}_{I}\left(x_{n}\right)$ for each $n$ to be the interaction at the vertex $x_{n}$. One then connects the vertices subject to the conditions detailed above. Each way of doing this yields a product of contractions whose values can be read off from the list compiled above; the integral of this product, multiplied by $(-i)^{N} \prod_{1}^{N} \lambda_{n}$, where $\lambda_{n}$ is the coupling constant for the vertex $x_{n}$, yields one of the terms whose sum is (6.58).
\end{itemize}

This is a very neat way to catalogue all the integrals that need to be calculated! However, there are some combinatorial issues that must be addressed. First, we do not distinguish between two Feynman diagrams that differ only in the relabeling of their internal vertices. With that understanding, the factor of $1 / N!$ in (6.58) disappears. (If one compares the present discussion with the derivation of the Dyson series in §6.1, one sees that this $1 / N$ ! is there precisely to cancel the $N$ ! equal contributions of these permuted diagrams.) Second, if the interaction involves\\
powers higher than one of a single field, there are additional counting problems. For example, in the $\phi^{4}$ scalar field theory, each of the four factors of $\phi(x)$ in an interaction term $\mathcal{H}_{I}(x)$ must be contracted with some other operator, and there are generally 4 ! ways to arrange these contractions; it is precisely to compensate for this fact that the factor of $1 / 4!$ is inserted into (6.28). However, if two or three of these factors are contracted with $\phi(y)$ 's that are all in the same $\mathcal{H}_{I}(y)$, this compensation is off by a factor of 2 ! or 3 !, which must therefore be adjusted accordingly. (The corresponding diagrams have two or three lines connecting the vertices $x$ and $y$, and the point is that the diagram is unchanged by the permutation of these lines.) These symmetry factors, which can also arise in other ways, can be annoying to deal with, but they will cause no difficulty in the calculations we shall perform here. (See Weinberg [131], §6.1, for a more extended discussion of these issues.)

One also has to be careful about keeping track of minus signs when the interactions involve Fermions. For one thing, changing the order of the creation operators that define the initial and final states can introduce factors of -1 , so one has to order the operators in the initial and final states consistently. Moreover, when one expresses an S-matrix element as a product of contractions $\overbrace{A B}$, one must permute the operators involved so that each $A$ is adjacent to its partner $B$, and permuting Fermionic operators introduces a factor of the sign of the permutation. In the next section we shall see an example of the importance of these minus signs.

At this point the reader is undoubtedly hungry for some specific examples, so let us look at the lowest-order terms for an S -matrix element in $\phi^{4}$ scalar field theory in which the initial and final states have two particles each: $\mid$ in $\rangle=a^{\dagger}\left(\mathbf{p}_{1}\right) a^{\dagger}\left(\mathbf{p}_{2}\right)|0\rangle$ and $\mid$ out $\rangle=a^{\dagger}\left(\mathbf{p}_{3}\right) a^{\dagger}\left(\mathbf{p}_{4}\right)|0\rangle$. The zeroth term in the perturbation series is

$$
\begin{aligned}
\langle 0| a\left(\mathbf{p}_{3}\right) a\left(\mathbf{p}_{4}\right) & a^{\dagger}\left(\mathbf{p}_{1}\right) a^{\dagger}\left(\mathbf{p}_{2}\right)|0\rangle \\
& =\overbrace{a\left(\mathbf{p}_{3}\right) a^{\dagger}\left(\mathbf{p}_{1}\right)} \overbrace{a\left(\mathbf{p}_{4}\right) a^{\dagger}\left(\mathbf{p}_{2}\right)}+\overbrace{a\left(\mathbf{p}_{3}\right) a^{\dagger}\left(\mathbf{p}_{2}\right)} \overbrace{a\left(\mathbf{p}_{4}\right) a^{\dagger}\left(\mathbf{p}_{1}\right)} \\
& =(2 \pi)^{6}\left[\delta\left(\mathbf{p}_{3}-\mathbf{p}_{1}\right) \delta\left(\mathbf{p}_{4}-\mathbf{p}_{2}\right)+\delta\left(\mathbf{p}_{3}-\mathbf{p}_{2}\right) \delta\left(\mathbf{p}_{4}-\mathbf{p}_{1}\right)\right]
\end{aligned}
$$

which corresponds to the two Feynman diagrams in Figure 6.2. No surprises here: nothing happens, and the two outgoing particles are the two incoming particles.

\begin{figure}[h]
\begin{center}
  \includegraphics[width=\textwidth]{2025_11_18_7960b2c161913da1e816g-168}

\caption{Figure 6.2. The trivial interaction in \$\textbackslash phi\^{}\{4}\$ theory.\}\end{center}
\end{figure}

Note: We shall always indicate internal vertices of Feynman diagrams by black dots. If two lines cross but there is no black dot at the intersection, there is no vertex there.

The first-order term is

$$
\frac{\lambda}{4!} \int\langle 0| a\left(\mathbf{p}_{3}\right) a\left(\mathbf{p}_{4}\right) \phi(x) \phi(x) \phi(x) \phi(x) a^{\dagger}\left(\mathbf{p}_{1}\right) a^{\dagger}\left(\mathbf{p}_{2}\right)|0\rangle d^{4} x .
$$

The various ways of contracting the eight terms in the integrand yield diagrams of the types shown in Figure 6.3. The first one, representing the most basic scattering process of the $\phi^{4}$ interaction, comes from contracting each of the $\phi(x)$ 's with one of the external creation or annihilation operators. There are $4!$ ways to do this, which all give the same result; this cancels the $4!$ in the coupling constant. By (5.12) we have

$$
\overbrace{a\left(\mathbf{p}_{j}\right) \phi(x)}=\int \frac{e^{i p_{\mu} x^{\mu}}}{\sqrt{2 \omega_{\mathbf{p}}}} \delta\left(\mathbf{p}-\mathbf{p}_{j}\right) d^{3} \mathbf{p}=\frac{e^{i p_{j \mu} x^{\mu}}}{\sqrt{2 \omega_{\mathbf{p}_{j}}}}
$$

and likewise $\overbrace{\phi(x) a^{\dagger}\left(\mathbf{p}_{j}\right)}=e^{-i p_{j \mu} x^{\mu}} / \sqrt{2 \omega_{\mathbf{p}_{j}}}$, so the value of Figure 6.3a is

$$
\lambda \int \frac{e^{i\left(p_{3 \mu}+p_{4 \mu}-p_{1 \mu}-p_{2 \mu}\right) x^{\mu}}}{4 \sqrt{\omega_{\mathbf{p}_{1}} \omega_{\mathbf{p}_{2}} \omega_{\mathbf{p}_{3}} \omega_{\mathbf{p}_{4}}}} d^{4} x=\lambda \frac{(2 \pi)^{4}}{4 \sqrt{\omega_{\mathbf{p}_{1}} \omega_{\mathbf{p}_{2}} \omega_{\mathbf{p}_{3}} \omega_{\mathbf{p}_{4}}}} \delta\left(p_{3}+p_{4}-p_{1}-p_{2}\right) .
$$

The delta-function embodies the overall conservation of energy and momentum in the process.

\begin{figure}[h]
\begin{center}
  \includegraphics[width=\textwidth]{2025_11_18_7960b2c161913da1e816g-169}

\caption{Figure 6.3. First-order interactions in \$\textbackslash phi\^{}\{4}\$ theory.\}\end{center}
\end{figure}

Figure 6.3b corresponds to contracting two of the field operators with external creation or annihilation operators and the other two with each other. There are actually four diagrams of this type, depending on which initial vertices are connected to which final ones and which line gets the extra loop. Each of them has a symmetry factor of 2 , because there is only one way to pair two factors of $\phi(x)$ rather than two. Figure 6.3c corresponds to contracting all of the field operators with each other; there are two such diagrams, depending on the pairing of initial and final vertices, and they have a symmetry factor of 8 since there are only 3 ways of grouping the four $\phi$ 's into two pairs. But there is something highly peculiar about these diagrams, because the contraction that corresponds to the loop is $-i \Delta_{F}(x-x)=-i \Delta_{F}(0)$, and $\Delta_{F}$ has a singularity at the origin.

There are two possible attitudes to take to this problem. One is to declare that the interaction Hamiltonian must be Wick ordered at the outset - that is, to take $\mathcal{H}_{I}(x)$ to be $(\lambda / 4!): \phi(x)^{4}$ : rather than $(\lambda / 4!) \phi(x)^{4}$. This has the effect of removing all contractions of operators within the same $\mathcal{H}_{I}(x)$ from consideration; diagrammatically, it amounts to adding a rule that no line can begin and end at the same vertex. This prescription is always followed in the study of rigorous models in space-time dimensions 2 and 3 , and there is a theoretical justification for it noted by Weinberg [131], p. 200. On the other hand, the Fourier representation (6.47) of $\Delta_{F}$ gives a formula for $\Delta_{F}(0)$ as a divergent integral,

$$
\Delta_{F}(0)=\int \frac{1}{-p^{2}+m^{2}} \frac{d^{4} p}{(2 \pi)^{4}} .
$$

As such it is only one of many divergent integrals that arise from Feynman diagrams containing loops, and one can include it in the general renormalization scheme to be discussed in Chapter 7; it then affects the counterterms that must be added to cancel the divergences. In the end, these procedures lead to physically equivalent results as long as the field $\phi$ is considered in isolation, although one-line loops can have a nontrivial effect when there are external fields present. (See, e.g., Weinberg [131], p. 576.) We shall take the easy way out by henceforth forbidding lines to start and end at the same vertex.

\begin{figure}[h]
\begin{center}
  \includegraphics[width=\textwidth]{2025_11_18_7960b2c161913da1e816g-170}

\caption{Figure 6.4. Second-order interactions in \$\textbackslash phi\^{}\{4}\$ theory\}\end{center}
\end{figure}

With this simplification, the reader can easily verify that the second-order term

$$
\iint\langle 0| a\left(\mathbf{p}_{3}\right) a\left(\mathbf{p}_{4}\right) \mathcal{T}\left[: \phi(x)^{4}:: \phi(y)^{4}:\right] a^{\dagger}\left(\mathbf{p}_{1}\right) a^{\dagger}\left(\mathbf{p}_{2}\right)|0\rangle d^{4} x d^{4} y
$$

gives rise to the three kinds of diagrams in Figure 6.4. They have symmetry factors $2!, 3!$, and $4!$, respectively, from the permutation of the internal lines. We shall analyze the first of these in detail in §7.4 and the second in a more qualitative way in §7.6. The last one contains a "vacuum bubble," one of the quantum fluctuations of the vacuum; we shall say more about such diagrams in §6.11.

For QED the basic ideas are the same, but the algebra is more complicated. There are two fields here, the photon field $A$ and the electron field $\psi$ (although one can substitute another charged spin- $\frac{1}{2}$ particle for the electron); the interaction Hamiltonian is $\mathscr{H}_{I}=e A_{\mu} \bar{\psi} \gamma^{\mu} \psi$ where $e$ is the charge of the electron. These fields are both vector-valued, and the vector spaces for both of them are 4 -dimensional, but they are not the same: for the photon field the space is physical space-time or its dual momentum space, whereas for the electron field it is the space of Dirac spinors. Likewise, the photon propagator and the electron propagator are both $4 \times 4$ matrices, but acting on these two different spaces. The standard way to avoid massive confusion is to display all Lorentz indices but no spinor indices explicitly. Here is a list of the main features to keep in mind:\\
i. Each internal vertex, labeled by a space-time point $x$, has its own Lorentz index $\mu$, which is the Lorentz index of the Dirac matrix $\gamma^{\mu}$ in $\mathcal{H}_{I}(x)$; it forms part of the label of the vertex. If two internal vertices lableled by $(x, \mu)$ and $(y, \nu)$ are joined by a photon line, the corresponding photon propagator is $-i \Delta_{\mu \nu}(x-y)$ (which is symmetric in $\mu$ and $\nu$ and even in $x-y$, so the order doesn't matter).\\
ii. The contraction $\overbrace{A_{\mu}(x) a^{\dagger}(\mathbf{p})}$ or $\overbrace{a(\mathbf{p}) A_{\mu}(x)}$ corresponding to an external photon line has the form $\boldsymbol{\epsilon}_{\mu}(\mathbf{p}) / \sqrt{2|\mathbf{p}|}$ or $\boldsymbol{\epsilon}_{\mu}^{*}(\mathbf{p}) / \sqrt{2|\mathbf{p}|}$ where $\boldsymbol{\epsilon}_{\mu}$ is the polarization\\
vector for the incoming or outgoing photon. ${ }^{6}$ This will be contracted into the Dirac matrix $\gamma^{\mu}$ for the internal vertex $x$.\\
iii. If an electron line joins the vertices labeled by $x$ and $y$ and its arrow points from $x$ to $y$, the corresponding propagator is $-i \Delta_{\text {Dirac }}(y-x)$. (Here the order does matter!)\\
iv. The contraction $\overbrace{\psi(x) a^{\dagger}(\mathbf{p})}$ or $\overbrace{a(\mathbf{p}) \bar{\psi}(x)}$ has the form $\sqrt{m / 2 \omega_{\mathbf{p}}} u(\mathbf{p})$ (a column vector) or $\sqrt{m / 2 \omega_{\mathbf{p}}} \bar{u}(\mathbf{p})$ (a row vector), where $u(\mathbf{p})$ is the appropriate Dirac spinor defined by (5.29) and (5.30).\\
v. The electron lines in any diagram fall into two kinds of groups: chains of lines connecting an incoming electron vertex to an outgoing one, and internal loops. The total contribution of a chain connecting external vertices through $n$ internal vertices with labels $\left(x_{j}, \mu_{j}\right)$ has the form\\
$(-i e)^{n} \bar{u}\left(\mathbf{p}^{\text {out }}\right) \gamma^{\mu_{n}} \Delta_{\text {Dirac }}\left(x_{n}-x_{n-1}\right) \gamma^{\mu_{n-1}} \cdots \gamma^{\mu_{2}} \Delta_{\text {Dirac }}\left(x_{2}-x_{1}\right) \gamma^{\mu_{1}} u\left(\mathbf{p}^{\text {in }}\right)$.\\
Note that this is a product of a row vector, some matrices, and a column vector, and hence is a scalar as far as spinor space is concerned. The total contribution of a loop with vertices labeled $\left(x_{1}, \mu_{1}\right), \ldots,\left(x_{n}, \mu_{n}\right)\left(\right.$ and $\left(x_{n+1}, \mu_{n+1}\right)= \left.\left(x_{1}, \mu_{1}\right)\right)$ has the form\\
$-(-i e)^{n} \operatorname{tr}\left[\gamma^{\mu_{n}} \Delta_{\text {Dirac }}\left(x_{n}-x_{n-1}\right) \gamma^{\mu_{n-1}} \cdots \gamma^{\mu_{2}} \Delta_{\text {Dirac }}\left(x_{2}-x_{1}\right) \gamma^{\mu_{1}} \Delta_{\text {Dirac }}\left(x_{1}-x_{n}\right)\right]$, which is also a scalar vis-a-vis spinor space. The extra minus sign comes from the fact that all the propagators from contractions $\overbrace{\psi\left(x_{j+1}\right) \bar{\psi}\left(x_{j}\right)}$ except the last one, which comes from $\overbrace{\bar{\psi}\left(x_{n}\right) \psi\left(x_{1}\right)}=-\overbrace{\psi\left(x_{1}\right) \bar{\psi}\left(x_{n}\right)}$. The fact that the formula involves a trace comes from the workings of the linear algebra; the point is that the product of two matrices $A$ and $B$ is given by $(A B)_{i k}=\sum_{j} A_{i j} B_{j k}$, but when one completes the loop at the end, one sets $k=i$ and sums to obtain $\sum_{i, j} A_{i j} B_{j i}=\operatorname{tr}(A B)$.\\
vi. The Lorentz indices in the expressions in (v) are all summed out against photon propagators or external polarization vectors. Hence the final result is a scalar.

\begin{figure}[h]
\begin{center}
  \includegraphics[width=\textwidth]{2025_11_18_7960b2c161913da1e816g-171}

\caption{Figure 6.5. Some Feynman diagrams for QED.}
\end{center}
\end{figure}

\footnotetext{${ }^{6}$ This is exactly what one would expect, and exactly what one gets from the quantization of the photon field in Coulomb gauge (5.42). Even if one quantizes in another way, the external photons must be in physically correct states with polarization vectors that have no time component and are orthogonal to the momentum. See Weinberg [131], §5.9, for further discussion.
}The algebra is sufficiently tedious that we shall not write out a list of examples here, but rather consider individual examples as the occasion arises. However, a few simple Feynman diagrams for QED are shown in Figure 6.5. The first one is the main contribution to electron-electron scattering; the second is the main contribution to electron-positron annihilation; the third represents a photon turning into an electron-positron pair and back into a photon. The last one represents the fundamental interaction of QED, but unlike the analogous diagram (Figure 6.3a) for the $\phi^{4}$ interaction, it cannot stand on its own. One doesn't need quantum field theory, merely conservation of energy, to see that the amplitude for the depicted process, a (real) electron emitting or absorbing a (real) photon with nothing else happening, is zero. (Hint: Consider the reference frame in which the electron is initially at rest.)

By the way, if one rotates the diagram in Figure 6.5b through a right angle clockwise, it depicts the scattering of a photon off an electron, a process that we shall study in detail in §6.10. It follows that the corresponding amplitudes for electron-positron annihilation and electron-photon scattering are closely related. As the reader may verify, the precise relation is as follows. The values of these diagrams are functions of the external particle momenta; the value of the original diagram with the incoming positron assigned (4-)momentum $p$ and the bottom photon assigned momentum $k$ is equal to the value of the rotated diagram with the outgoing electron assigned momentum $-p$ and the incoming photon assigned momentum $-k$. Note that $p$ and $-p$, or $k$ and $-k$, cannot both be physically allowed, as $p_{0}>0$ for physical particles, so the relation is a mathematical one rather than a directly physical one, but it is useful nonetheless. An analogous result holds for any pair of diagrams that are obtained from one another by replacing some incoming particles by outgoing antiparticles or vice versa, a phenomenon known as crossing symmetry.

The correspondence between the integrals that make up the Dyson series and Feynman diagrams is perfectly precise and well-defined. However, it is customary to go further and think of the Feynman diagrams as schematic pictures of physical processes, and here the interpretation acquires a more imaginative character. The lines starting at the initial or final vertices can clearly be thought of as the trajectories of the initial and final particles entering or leaving the interaction. Once one has made this connection, the temptation is irresistible to interpret the lines joining internal vertices also as trajectories of particles involved in intermediate steps of the interaction. Some of them form chains pertaining to the same particle species that connect an incoming line to an outgoing one by following the arrows; they can be interpreted as the trajectories of the real particles that enter, interact, and leave. The particles that are created and annihilated within the diagram, however, have a much more tenuous claim to reality. They are not observed directly; they sometimes travel faster than the speed of light; they may violate conservation of energy quite flagrantly. They are, in short, the infamous virtual particles that are so ubiquitous in physicists' discourse. In the final analysis, the only existence they possess for certain is as picturesque ways of thinking about the ingredients of the integrals in the Dyson series.

However, it must be said that the line between virtual particles and real ones is not completely sharp. A real particle may be observed to be a little off mass shell because of the uncertainty principle, and a photon emitted by an electron in the\\
sun that manages to reach the earth before being absorbed by another electron has every right to claim reality even though the photon in Figure 6.5a is presumably virtual.

The arrows in a Feynman diagram indicate the direction of travel of the (virtual) particles in question: they point in the direction of travel for particles and in the opposite direction for antiparticles. The convention for the direction of the arrow on an edge joining two internal vertices is explained as follows. If $\phi$ is a field that destroys particles and creates antiparticles, the contraction $\overbrace{\phi(x) \phi^{\dagger}(y)}$ involves either the commutator of the annihilation operators in $\phi(x)$ with the creation operators in $\phi^{\dagger}(y)$ or the commutator of the creation operators in $\phi(x)$ with the annihilation operators in $\phi^{\dagger}(y)$, depending on the time-ordering of $x$ and $y$. That is, if $x_{0}>y_{0}, \overbrace{\phi(x) \phi^{\dagger}(y)}$ represents the creation of a (virtual) particle at $y$ followed by its destruction at $x$, whereas if $y_{0}>x_{0}$, it represents the creation of a (virtual) antiparticle at $x$ followed by its destruction at $y$. Both processes are represented by the same directed line in the Feynman diagram. In fact, Feynman argued for their physical indistinguishability; he was fond of asserting that antiparticles are just particles traveling backward in time.

\subsection*{Feynman diagrams in momentum space}
What we have described so far is the "position space" correspondence between Feynman diagrams and multiple integrals that contribute to an S -matrix element; that is, the vertices in the diagram are labeled by position-space variables, and these are the variables of integration. However, when one actually performs the calculation, it is easier to use a "momentum space" representation instead, because the most convenient way of expressing the propagators is by their Fourier expansions. The way this works is as follows.

Each integral represented by a Feynman diagram is of the form

$$
\int \cdots \int F_{1}\left(x, \mathbf{p}^{\text {out }}\right) F_{2}(x) F_{3}\left(x, \mathbf{p}^{\text {in }}\right) d^{4} x_{1} \cdots d^{4} x_{N}
$$

where $x=\left(x_{1}, \ldots, x_{n}\right) \in \mathbb{R}^{4 n}$ and $\mathbf{p}^{\text {out }}$ and $\mathbf{p}^{\text {in }}$ are likewise concatenations of the outgoing and incoming momentum vectors. More precisely,\\
i. $F_{1}\left(x, \mathbf{p}^{\text {out }}\right)$ is a product of factors $\exp \left(i p_{j \mu}^{\text {out }} x_{n}^{\mu}\right)$, multiplied by coefficients that depend on $p_{j}^{\text {out }}$ and the field types but not on $x_{n}$, coming from the final lines (i.e., the contractions $\overbrace{a\left(\mathbf{p}_{j}^{\text {out }}\right) \phi\left(x_{n}\right)})$;\\
ii. $F_{2}(x)$ is a product of propagators coming from the internal lines (i.e., the contractions $\overbrace{\phi\left(x_{n}\right) \phi^{\dagger}\left(x_{n^{\prime}}\right)})$;\\
iii. $F_{3}\left(x, \mathbf{p}^{\text {in }}\right)$ is a product of factors $\exp \left(-i p_{k \mu}^{\text {in }} x_{n}^{\mu}\right)$, again multiplied by certain coefficients, coming from the initial lines (i.e., the contractions $\overbrace{\phi\left(x_{n}\right) a^{\dagger}\left(\mathbf{p}_{k}^{\mathrm{in}}\right)})$. (There may also be some combinatorial factors.) We express each of the propagators in $F_{2}(x)$ as a Fourier integral:

$$
-i \Delta\left(x_{n}-x_{n^{\prime}}\right)=-i \int \exp \left[-i q_{\mu}\left(x_{n}-x_{n^{\prime}}\right)^{\mu}\right] \widehat{\Delta}(q) \frac{d^{4} q}{(2 \pi)^{4}}
$$

in which $\widehat{\Delta}(q)$ is a rational function of $q$ with denominator $-q^{2}+m^{2}-i \epsilon$ (and there is an implicit $\lim _{\epsilon \rightarrow 0}$ ). Substituting these integrals for the propagators turns (6.60) into

$$
\int \cdots \iint \cdots \int F_{1}\left(x, \mathbf{p}^{\text {out }}\right) F_{2}^{\prime}(x, q) F_{2}^{\prime \prime}(q) F_{3}\left(x, \mathbf{p}^{\text {in }}\right) d^{4} x_{1} \cdots d^{4} x_{N} \frac{d^{4} q_{1} \cdots d^{4} q_{M}}{(2 \pi)^{4 M}}
$$

where there is one 4 -momentum variable $q_{m}$ for each internal line, and\\
i. $F_{2}^{\prime}(x, q)$ is the product of factors $\exp \left[-i q_{m \mu}\left(x_{n}-x_{n^{\prime}}\right)^{\mu}\right]$;\\
ii. $F_{2}^{\prime \prime}(q)$ is the product of the Fourier transforms of the propagators.

At this point the $x$-dependence in the integrand is simply a product of imaginary exponentials, so the $x$-integrations can be performed by the Fourier inversion formula

$$
\int e^{-i x_{\mu} a^{\mu}} d^{4} x=(2 \pi)^{4} \delta(a)
$$

The result is that the integral (6.62) becomes

$$
G\left(\mathbf{p}^{\text {out }}, \mathbf{p}^{\text {in }}\right) \int \cdots \int(2 \pi)^{4 N} \delta\left(\Sigma_{1}\right) \cdots \delta\left(\Sigma_{N}\right) R_{1}\left(q_{1}\right) \cdots R_{m}\left(q_{M}\right) \frac{d^{4} q_{1} \cdots d^{4} q_{M}}{(2 \pi)^{4 M}}
$$

where\\
i. $G\left(\mathbf{p}^{\text {out }}, \mathbf{p}^{\text {in }}\right)$ is the product of field coefficients coming from the external lines (and any combinatorial factors),\\
ii. the $R_{m}$ are the Fourier-transformed propagators (" $R$ " stands for "rational function"), and\\
iii. $\Sigma_{n}$ is the algebraic sum (see below) of the 4 -momenta ( $p_{j}^{\text {out }}, p_{k}^{\text {in }}$, and/or $q_{m}$ ) associated to the lines that meet at the vertex $x_{n}$.\\
In addition, when any initial vertex $\mathbf{p}_{k}^{\text {in }}$ is connected directly to a final vertex $\mathbf{p}_{j}^{\text {out }}$, there will be an extra factor of $\delta\left(\mathbf{p}_{k}^{\mathrm{in}}-\mathbf{p}_{j}^{\mathrm{out}}\right)$.

The term "algebraic sum" requires some explanation. Some lines that meet at the vertex $x_{n}$ have arrows attached to them; a line is considered "incoming" if its arrow points toward $x_{n}$ and "outgoing" if it points away from $x_{n} \cdot{ }^{7}$ Lines without arrows (corresponding to neutral particles) are arbitrarily assigned arrows for this purpose; one merely has to be consistent in using the same arrow for the calculations at both ends of the line. (These arrows are usually drawn adjacent and parallel to the lines rather than on them.) The algebraic sum $\Sigma_{n}$ of the momenta at $x_{n}$ is then the sum of the outgoing momenta minus the sum of the incoming momenta. Thus, the delta-function $\delta\left(\Sigma_{n}\right)$ expresses "conservation of energy-momentum at the vertex $x_{n}$."

When one thinks of a Feynman diagram as representing the momentum-space integral (6.63), instead of labeling the internal vertices with the position variables, one labels the lines with 4 -momentum variables. The conversion from position to momentum labels is as follows. The line from the initial vertex with momentum $\mathbf{p}_{k}^{\mathrm{in}}$ is labeled by the corresponding 4 -momentum $p_{k}^{\mathrm{in}}=\left(\omega_{\mathbf{p}_{k}^{\mathrm{in}}}, \mathbf{p}_{k}^{\mathrm{in}}\right)$, and likewise for lines to final vertices. The internal line from the vertex $x_{n^{\prime}}$ to the vertex $x_{n}$ (representing the propagator $-i \Delta\left(x_{n}-x_{n^{\prime}}\right)$ ) is labeled by the momentum variable $q$ in the Fourier expansion (6.61). We think of $q$ as flowing in the direction of the arrow associated to

\footnotetext{${ }^{7}$ One can reverse this convention for lines that are considered to represent antiparticles; it doesn't matter as long as the convention is consistent from vertex to vertex.
}\begin{table}[h]
\begin{center}
\begin{tabular}{|l|l|l|}
\hline
External line & $\{\stackrel{x}{\longrightarrow}\}$ & $\frac{1}{\sqrt{2 \omega_{\mathbf{p}}}}$ \\
\hline
Internal line & $p$ & $\frac{-i}{-p^{2}+m^{2}-i \epsilon}$ \\
\hline
Internal vertex & \includegraphics[max width=\textwidth]{2025_11_18_7960b2c161913da1e816g-175}
 & $-i \lambda \delta\left(p_{1}+p_{2}-p_{3}-p_{4}\right)$ \\
\hline
\end{tabular}

\caption{Table 6.1. Momentum-space Feynman rules for \$\textbackslash phi\^{}\{4}\$ scalar field theory. (External vertices are denoted by crosses.)\}
\end{center}
\end{table}

the line (that is, "from $x_{n^{\prime}}$ to $x_{n}$ ") if the line is considered as representing a particle (including neutral particles) and in the opposite direction if the line is considered as representing an antiparticle, as in the preceding paragraph. This distinction is unimportant for internal Boson lines, as their propagators are even functions and $-q_{m}$ is just as good a variable as $q_{m}$ in (6.63); but it is significant for internal Fermion lines (whose propagators are not even) and external lines (whose momenta $p$ must be on mass shell so that $p^{0}>0$ ).

The integral (6.63) corresponding to a diagram labeled in this way is built up as follows:\\
i. For each external line labeled by the momentum $p$ there is a factor $u(p)$ consisting of the appropriate field coefficient arising from the contraction associated to the line.\\
ii. For the internal line labeled with $q_{m}(m=1, \ldots, M)$ there is a factor $-i \widehat{\Delta}\left(q_{m}\right)$, the Fourier-transformed propagator for the field associated to the line.\\
iii. For each vertex there is a factor $-i(2 \pi)^{4} \lambda \delta(\Sigma)$, where $\Sigma$ is the algebraic sum of the 4 -momenta at that vertex and $\lambda$ incorporates the coupling constant and any other algebraic factors such as Dirac matrices for the interaction at that vertex.\\
One now multiplies all these factors together, integrates over all the internal momenta, and also sums over whatever discrete indices (spin, etc.) are present. One must also include combinatorial factors and Fermionic minus signs as needed, just as for the position-space integral.

In this setting the "virtualness" of the particles associated to the internal lines manifests itself in the fact that their 4 -momenta $q_{m}$ are completely arbitrary, not on mass shell.

For future reference, we give a summary of the momentum-space Feynman rules for the $\phi^{4}$ scalar field theory and for QED in Tables 6.1 and 6.2.

We return to the general discussion of the evaluation of integrals associated to Feynman diagrams. Because of the delta-functions in the integrand, some, and

\begin{table}[h]
\begin{center}
\begin{tabular}{|l|l|l|}
\hline
Incoming electron line & \( \times \quad p \rightarrow \) & $\sqrt{m / 2 \omega_{\mathbf{p}}} u(\mathbf{p}, s)$ \\
\hline
Incoming positron line & \includegraphics[max width=\textwidth]{2025_11_18_7960b2c161913da1e816g-176}
 & $\sqrt{m / 2 \omega_{\mathbf{p}}} \bar{v}(\mathbf{p}, s)$ \\
\hline
Incoming photon line & \includegraphics[max width=\textwidth]{2025_11_18_7960b2c161913da1e816g-176(3)}
 & $\epsilon_{\mu}(\mathrm{p})$ \\
\hline
Outgoing electron line & \includegraphics[max width=\textwidth]{2025_11_18_7960b2c161913da1e816g-176(1)}
 & $\sqrt{m / 2 \omega_{\mathbf{p}}} \bar{u}(\mathbf{p}, s)$ \\
\hline
Outgoing positron line & \includegraphics[max width=\textwidth]{2025_11_18_7960b2c161913da1e816g-176(6)}
 & $\sqrt{m / 2 \omega_{\mathbf{p}}} v(\mathbf{p}, s)$ \\
\hline
Outgoing photon line & \includegraphics[max width=\textwidth]{2025_11_18_7960b2c161913da1e816g-176(4)}
 & $\epsilon_{\mu}^{*}(\mathbf{p})$ \\
\hline
Internal electron or positron line & \( \{\xrightarrow[p \rightarrow]{\vec{p}}\} \) &  \\
\hline
Internal photon line & \includegraphics[max width=\textwidth]{2025_11_18_7960b2c161913da1e816g-176(2)}
 & $\frac{-i g_{\mu \nu}}{p^{2}+i \epsilon}$ \\
\hline
Internal vertex & \includegraphics[max width=\textwidth]{2025_11_18_7960b2c161913da1e816g-176(5)}
 & $-i e \gamma^{\mu} \delta\left(p_{1}+q-p_{2}\right)$ \\
\hline
\end{tabular}

\caption{Table 6.2. Momentum-space Feynman rules for QED. (External vertices are denoted by crosses. The Lorentz indices at internal ends of photon lines contract with the Lorentz indices at internal vertices.)}
\end{center}
\end{table}

perhaps all, of the integrations in (6.63) collapse immediately. Each internal momentum variable $q_{m}$ occurs in precisely two of the sums $\Sigma_{n}$ (the ones coming from the vertices at the ends of the line associated to $q$ ), so repeated application of the formula

$$
\int(2 \pi)^{4} \delta\left(q-q^{\prime}\right)(2 \pi)^{4} \delta\left(q^{\prime}-q^{\prime \prime}\right) \frac{d^{4} q^{\prime}}{(2 \pi)^{4}}=(2 \pi)^{4} \delta\left(q-q^{\prime \prime}\right)
$$

(this is honest mathematics: the convolution of two measures) serves to eliminate some or all of the internal momenta, which are then expressed in terms of the external momenta and remaining internal momenta (if any) in the remaining part of the integrand.

The form of the result of performing these trivial integrations depends on the topology of the Feynman diagram. First, if the diagram is disconnected, the integral it represents is just the product of the integrals represented by its connected\\
subdiagrams, so it suffices to consider connected diagrams. For these we quote a basic topological result for one-dimensional cell complexes, the analogue of the Euler formula $V-E+F=2-2 g$ for polyhedra. (A proof can be found in any introductory book on algebraic topology; see, e.g., Croom [19], Theorem 2.5.)

Let $V$ and $E$ be the number of vertices and edges, respectively, in a connected graph $G$, and let $L$ be the number of independent loops in $G$ (i.e., the rank of the homology group $H_{1}(G)$ ). ${ }^{8}$ Then $V-E=1-L$.

In a Feynman diagram, each external vertex lies on only one line, so deletion of external vertices and lines does not affect the connectedness or the numbers $V-E$ and $L$. By applying the preceding theorem to the subgraph consisting of the internal vertices and lines, we obtain:

Let $V$ and $I$ be the number of internal vertices and internal lines, respectively, in a connected Feynman diagram, and let $L$ be the number of independent loops. Then

$$
V-I=1-L
$$

In particular, we always have $I \geq V-1$, so there are always enough $q$-integrals to eat up all but one of the delta-functions $\delta\left(\Sigma_{n}\right)$. As the reader may verify by keeping more careful track of the different internal and external momenta in the sums $\Sigma_{n}$, once these $V-1$ integrations have been performed, the argument of the remaining delta-function is precisely $\sum_{j} p_{j}^{\text {out }}-\sum_{k} p_{k}^{\text {in }}$, the sum of the final 4momenta minus the sum of the initial 4-momenta. Hence this delta-function simply expresses the fact that overall momentum and energy, for the real (nonvirtual) initial and final particles, must be conserved by the interaction. It contains no internal momenta and so plays no role in any further integrations.

If $L=0$, that is, if the diagram is a tree, then $I=V-1$, there are no more integrations after the delta-functions have been reduced, and the calculation is complete. The result is a product of (1) field coefficients evaluated at the external momenta, (2) propagators whose internal momentum variable has been expressed in terms of the external momenta via the conditions $\Sigma_{n}=0$; (3) the overall momentum conservation delta-function $(2 \pi)^{4} \delta\left(\sum p_{j}^{\text {out }}-\sum p_{k}^{\text {in }}\right)$.

If $L>0$, there are $L$ additional integrations to be performed, and the integrand is a rational function of the remaining internal momenta (a product of propagators). These integrals have an unfortunate propensity for diverging, so their analysis must be postponed until the next chapter.

There are a couple of modifications to the rules for evaluating external lines in a Feynman diagram that are frequently useful. First, one may wish to consider a part of a Feynman diagram in isolation, so that its external lines are actually internal lines of a larger diagram. In that case the external particles are not on mass shell and the external lines correspond to propagators rather than field coefficients. (The same idea is relevant to evaluation of vacuum expectation values of products of fields, as we shall explain in §6.11.) On the other hand, sometimes one wants to concentrate on the integrations associated to the internal lines without worrying about the quantities associated to one or more external lines at all, so one omits the

\footnotetext{${ }^{8}$ For example, the diagram in Figure 6.4b has three loops, one formed by the two curved arcs and the other two formed by one of the arcs and the line segment in between, of which any two are independent; the third is the sum or the difference of the other two in $H_{1}(G)$.
}
corresponding field coefficients or propagators entirely. This is known as amputating the external line.

The integral associated to any Feynman diagram will be called the value of the Feynman diagram. However, it is often convenient to take a short-cut and identify a diagram with its value. Thus, we may speak of a "sum of Feynman diagrams," a "divergent Feynman diagram," etc.

\subsection*{Cross sections and decay rates}
At this point, let us see what we have accomplished. We have shown how to compute the scattering matrix elements between initial and final states consisting of collections of particles with definite momenta in terms of quantities represented by Feynman diagrams; more precisely, our expectation is that the sum of finitely many of these quantities will yield a good approximation to the matrix element. Some of the quantities are divergent, but there are methods to be discussed in the next chapter for removing the divergence. Let us assume either that this has been done or that we are content to use only diagrams without loops where the divergence problem does not occur (the "tree level approximation"), so that we have obtained a formula for the scattering matrix element. Now what? That is, what does this matrix element have to do with quantities that the experimentalists can measure?

The question of interpretation is clearly nontrivial. A scattering matrix element is a transition amplitude between the initial and final states, and the actual transition probability is obtained by taking the absolute square of this amplitude. But the matrix element always contains a momentum-conservation delta-function, so how is one supposed to interpret its square? Of course, some such problem is not unexpected, since the initial and final states with definite momenta are not actual elements of the appropriate Hilbert space. To get more rigorous results, one should replace them with honest wave packets. However, if the final answer is to have any experimental relevance, it had better not depend in any serious way on the precise structure of the wave packets. In the typical situation where an interaction results from the collision of two initial particles prepared in a particle accelerator, one usually has very little precise information about the wave functions of these particles except that their momenta are well localized about certain values determined by the experiment - so that, on some level, the assumption of definite momenta has to be a good approximation after all.

In short, the bad news is that mathematical precision is difficult to achieve, but the good news is that it is also somewhat irrelevant: the end result needs to be insensitive to a certain amount of sloppiness. As Weinberg [131], p. 134, says, "(as far as I know) no interesting open problems in physics hinge on getting the fine points right regarding these matters." We therefore follow him in taking a quick and dirty route to the main results, using the same device we employed in §5.1 and §6.2: putting the particles in a box. This also has the advantage of clarifying the normalization constants. (See Peskin and Schroeder [89], §4.5, or Itzykson and Zuber [66], §5-1-1, for alternate derivations using wave packets. However, one should be aware that different authors use different normalization conventions; in particular, factors of $2 \pi$ are likely not to be invariant under change of reference text.)

Before getting started on this, we need to make a general observation about scattering amplitudes. If one of the initial particles is in exactly the same state as\\
one of the final particles (same species, same momentum, same spin, etc.), there is a possibility - in fact, an overwhelming likelihood - that one incoming particle simply goes in one side and out the other without interacting with anything else. (This corresponds to Feynman diagrams in which one initial vertex is connected to one final vertex by a single line that is disconnected from the rest of the diagram.) More generally, if the sets of initial and final particles can each be partitioned into $k$ subsets such that the total momentum in each incoming subset is equal to the total momentum in the corresponding outgoing subset, the main contribution to the total scattering process will be $k$ separate interactions involving the various subsets. (This corresponds to Feynman diagrams with at least $k$ connected components, each component carrying its own momentum-conservation delta-function.)

We assume at the outset that we are only interested in "irreducible" scattering processes in which all the incoming particles participate, that is, for which the Feynman diagrams are connected. We therefore disallow outgoing states in which some subset of the outgoing particles has exactly the same total momentum as some subset of the incoming particles. This is hardly a serious restriction, as the forbidden subset of the joint momentum space of the outgoing particles is a subvariety of positive codimension (and hence measure zero). Under this condition, if $\mid$ in $\rangle$ and $\mid$ out $\rangle$ are the initial and final states, and $p^{\text {in }}$ and $p^{\text {out }}$ are their total 4 -momenta, the S-matrix element $\langle$ out $| S \mid$ in $\rangle$ has the form

$$
S(\text { in } \rightarrow \text { out })=\langle\text { out }| S \mid \text { in }\rangle=i(2 \pi)^{4} \delta\left(p_{\beta}-p_{\alpha}\right) M(\text { in } \rightarrow \text { out })
$$

where $M$ (in → out) is free of delta-functions. (The $i$ is conventional.) Here $\mid$ in $\rangle$ and $\mid$ out $\rangle$ are of the form $a_{1}^{\dagger}\left(\mathbf{p}_{1}\right) \cdots a_{N}^{\dagger}\left(\mathbf{p}_{N}\right)|0\rangle$ where the subscripts on the $a$ 's encode spin states, particle species, and any other relevant parameters. $N$ will always denote the number of particles in such a state, $\mathbf{p}=\sum_{1}^{N} \mathbf{p}_{j}$ will denote its total 3-momentum, $E_{\mathbf{p}}=\sum \sqrt{\left|\mathbf{p}_{j}\right|^{2}+m_{j}^{2}}$ its total energy, and $p=\left(E_{\mathbf{p}}, \mathbf{p}\right)$ its total 4-momentum; all of these will carry superscripts "in" or "out."

We now replace $\mathbb{R}^{3}$ by a box $\mathcal{B}=\left[-\frac{1}{2} L, \frac{1}{2} L\right]^{3}$ with volume $V=L^{3}$. (One may envision $L$ as a length scale that is extremely large in comparison to subatomic particles but perhaps not in comparison to the experimental apparatus: just a region large enough so that all events of interest happen well inside it.) For simplicity (and here again, specifics are unimportant), we impose periodic boundary conditions on the sides of the box, so that the 3 -momenta of the particles are restricted to the lattice $[(2 \pi / L) \mathbb{Z}]^{3}$, and the Dirac delta $\delta\left(\mathbf{p}_{1}-\mathbf{p}_{2}\right)$ is replaced by the rescaled Kronecker delta

$$
\delta_{\mathcal{B}}\left(\mathbf{p}_{1}-\mathbf{p}_{2}\right)=\int_{\mathcal{B}} e^{i\left(\mathbf{p}_{1}-\mathbf{p}_{2}\right) \cdot \mathbf{x}} \frac{d^{3} \mathbf{x}}{(2 \pi)^{3}}=\frac{V}{(2 \pi)^{3}} \delta_{\mathbf{p}_{1} \mathbf{p}_{2}}
$$

Having cut down space to a finite size, we need also to restrict the time during which the interaction can take place to a large but finite interval $\left[-\frac{1}{2} T, \frac{1}{2} T\right]$, since the particles can no longer escape to infinity at large positive and negative times. Here again, in practice $T$ should be only large enough to encompass the interaction in question comfortably; it need not be large on a human or cosmological time scale. The Dirac delta $\delta\left(E_{1}-E_{2}\right)$ is then replaced by a smooth bump function,

$$
\delta_{T}\left(E_{1}-E_{2}\right)=\int_{-T / 2}^{T / 2} e^{i\left(E_{1}-E_{2}\right) t} \frac{d t}{2 \pi}=\frac{1}{\pi} \frac{\sin \frac{1}{2}\left(E_{1}-E_{2}\right) T}{E_{1}-E_{2}}
$$

We must now address the question of the normalization of the initial and final states. In order to obtain a transition probability, we must use state vectors with norm one. In $\mathbb{R}^{3}$ this is problematic because the states $\mid$ in $\rangle$ and $\mid$ out $\rangle$ have definite momenta and are therefore not normalizable. But in the box, we use the creation operators $a_{\mathcal{B}}^{\dagger}(\mathbf{p})\left(\mathbf{p} \in[(2 \pi / L) \mathbb{Z}]^{3}\right)$ appropriate to the box instead, and then it is easy: the vectors $a_{\mathcal{B}}^{\dagger}(\mathbf{p})|0\rangle$ always have norm one, and hence so do all states of the form $a_{\mathcal{B} 1}^{\dagger}\left(\mathbf{p}_{1}\right) \cdots a_{\mathcal{B} N}^{\dagger}\left(\mathbf{p}_{N}\right)|0\rangle$ provided only that no two of the operators $a_{\mathcal{B} j}^{\dagger}\left(\mathbf{p}_{j}\right)$ are identical (creating particles of the same species in the same spin state with exactly the same momentum). ${ }^{9}$ This is an entirely harmless restriction, and we adopt it henceforth. However, there is a rescaling implicit in the replacement of $a^{\dagger}$ by $a_{\mathcal{B}}^{\dagger}$, because as we saw in $\oint 5.1$ (in the calculations leading to (5.11), where we denoted $a_{\mathcal{B}}^{\dagger}(\mathbf{p})$ by $A_{\mathbf{p}}^{\dagger}$ ), it is not $a_{\mathcal{B}}^{\dagger}$ but $L^{3 / 2} a_{\mathcal{B}}^{\dagger}=V^{1 / 2} a_{\mathcal{B}}^{\dagger}$ that turns into $a^{\dagger}$ as the box is removed. We see this again in the rescaling of delta-functions in (6.66): the inner product of $V^{1 / 2} a_{\mathcal{B}}^{\dagger}\left(\mathbf{p}_{1}\right)|0\rangle$ and $V^{1 / 2} a_{\mathcal{B}}^{\dagger}\left(\mathbf{p}_{2}\right)|0\rangle$ is $V \delta_{\mathbf{p}_{1} \mathbf{p}_{2}}=(2 \pi)^{3} \delta_{\mathcal{B}}\left(\mathbf{p}_{1}-\mathbf{p}_{2}\right)$, which turns into $(2 \pi)^{3} \delta\left(\mathbf{p}_{1}-\mathbf{p}_{2}\right)=\langle 0| a\left(\mathbf{p}_{1}\right) a^{\dagger}\left(\mathbf{p}_{2}\right)|0\rangle$ when the box is removed.

We can therefore adapt the S-matrix element (6.65) to the box by replacing the initial and final states $\mid$ in $\rangle$ and $\mid$ out $\rangle$ with

$$
\left.V^{N^{\text {in }} / 2}|\operatorname{in}\rangle_{\mathcal{B}}=V^{N^{\text {in }} / 2} \prod_{1}^{N^{\text {in }}} a^{\dagger}\left(\mathbf{p}_{j}^{\text {in }}\right)|0\rangle, \quad V^{N^{\text {out }} / 2} \mid \text { out }\right\rangle_{\mathcal{B}}=V^{N^{\text {out }} / 2} \prod_{1}^{N^{\text {out }}} a^{\dagger}\left(\mathbf{p}_{j}^{\text {out }}\right)|0\rangle
$$

(we are suppressing indices that indicate spin states, particle species, etc.) and the delta-function $\delta\left(p^{\text {out }}-p^{\text {in }}\right)$ with box delta-functions:

$$
\begin{aligned}
\left.\left.S\left(V^{N^{\text {in }} / 2} \mid \text { in }\right\rangle_{\mathcal{B}} \rightarrow V^{N^{\text {out }} / 2} \mid \text { out }\right\rangle_{\mathcal{B}}\right) & \\
& \quad=i(2 \pi)^{4} \delta_{\mathcal{B}}\left(\mathbf{p}^{\text {out }}-\mathbf{p}^{\text {in }}\right) \delta_{T}\left(E^{\text {out }}-E^{\text {in }}\right) M(\text { in } \rightarrow \text { out })
\end{aligned}
$$

(The factor $M($ in $\rightarrow$ out), which contains the hard-won information from the quantum field theory, remains the same!) The transition amplitude $S_{\mathcal{B}}$ (in → out) for the normalized states $u_{\mathcal{B}}^{\text {in }}$ and $u_{\mathcal{B}}^{\text {out }}$ is then

$$
S_{\mathcal{B}}(\text { in } \rightarrow \text { out })=i V^{\left(-N^{\text {in }}-N^{\text {out }}\right) / 2}(2 \pi)^{4} \delta_{\mathcal{B}}\left(\mathbf{p}^{\text {out }}-\mathbf{p}^{\text {in }}\right) \delta_{T}\left(E^{\text {out }}-E^{\text {in }}\right) M(\text { in } \rightarrow \text { out })
$$

and the transition probability is the absolute square of this quantity:

$$
P_{\mathcal{B}}(\text { in } \rightarrow \text { out })=V^{-N^{\text {in }}-N^{\text {out }}}(2 \pi)^{8} \delta_{\mathcal{B}}\left(\mathbf{p}^{\text {out }}-\mathbf{p}^{\text {in }}\right)^{2} \delta_{T}\left(E^{\text {out }}-E^{\text {in }}\right)^{2} \mid\left. M(\text { in } \rightarrow \text { out })\right|^{2}
$$

We can now deal with the squared delta-functions, as they no longer have infinite singularities. Indeed, by (6.66) we have

$$
\delta_{\mathcal{B}}\left(\mathbf{p}_{\beta}-\mathbf{p}_{\alpha}\right)^{2}=\delta_{\mathcal{B}}(0) \delta_{\mathcal{B}}\left(\mathbf{p}_{\beta}-\mathbf{p}_{\alpha}\right)=\frac{V}{(2 \pi)^{3}} \delta_{\mathcal{B}}\left(p_{\beta}-p_{\alpha}\right)
$$

There is no corresponding exact formula for $\delta_{T}^{2}$, but we can make an analogous approximation from (6.67):

$$
\delta_{T}\left(E_{\beta}-E_{\alpha}\right)^{2} \approx \delta_{T}(0) \delta_{T}\left(E_{\beta}-E_{\alpha}\right)=\frac{T}{2 \pi} \delta_{T}\left(E_{\beta}-E_{\alpha}\right)
$$

Using this approximation, we obtain\\
$P_{\mathcal{B}}($ in $\rightarrow$ out $)=V^{-N^{\text {in }}-N^{\text {out }}+1} T(2 \pi)^{4} \delta_{\mathcal{B}}\left(\mathbf{p}^{\text {out }}-\mathbf{p}^{\text {in }}\right) \delta_{T}\left(E^{\text {out }}-E^{\text {in }}\right) \mid . M($ in $\rightarrow$ out $)|^{2}$.

\footnotetext{${ }^{9}$ If $m$ of the operators are identical, the norm of $u_{\mathcal{B}}$ has a factor of $\sqrt{m!}$.
}At this point we can drop the pretense of the box as far as the discreteness of possible momenta is concerned. In practice the probability of finding each outgoing particle with any particular momentum is infinitesimal, so we are interested instead in particles whose momentum lies in a small bit $d^{3} \mathbf{p}$ of momentum space. The number of points of the lattice of box momenta in this small region is $\left[V /(2 \pi)^{3}\right] d^{3} \mathbf{p}$. Hence, if $\mathbf{p}_{1}, \ldots, \mathbf{p}_{N^{\text {out }}}$ are the momenta of the outgoing particles (so that $\mathbf{p}^{\text {out }}= \mathbf{p}_{1}+\cdots+\mathbf{p}_{N^{\text {out }}}$ ) and we write

$$
d^{3 N^{\text {out }}} \mathbf{p}^{\text {out }}=d^{3} \mathbf{p}_{1} \cdots d^{3} \mathbf{p}_{N^{\text {out }}}
$$

for short, the differential transition probability for outgoing particles having momenta in the range $d^{3 N^{\text {out }}} \mathbf{p}^{\text {out }}$ is

$$
\begin{aligned}
& d P(\text { in } \rightarrow \text { out }) \\
& =V^{1-N^{\text {in }}} T(2 \pi)^{4} \delta\left(\mathbf{p}^{\text {out }}-\mathbf{p}^{\text {in }}\right) \delta_{T}\left(E^{\text {out }}-E^{\text {in }}\right) \left\lvert\,\left. M(\text { in } \rightarrow \text { out })\right|^{2} \frac{d^{3 N^{\text {out }}} \mathbf{p}^{\text {out }}}{(2 \pi)^{3 N^{\text {out }}}}\right.
\end{aligned}
$$

Here, having passed to the continuum limit in describing the momentum states, we have done the same for the momentum-conservation factor and replaced $\delta_{\mathcal{B}}$ by $\delta$.

This is the general formula for interpreting the S -matrix elements. What remains is to explain the role of $V$ and $T$, and for this it is necessary to consider the specific cases of various numbers of initial particles.

For experiments in particle accelerators, by far the most common situation is that there are two initial particles that are made to collide. More precisely, one manufactures a beam of particles of one species that is directed at a target consisting of a bunch of particles of another (or the same) species. (In some experiments, two particle beams are directed at each other; the only difference is in whether one group of particles is at rest in the laboratory frame or not. For the moment we assume that the target is at rest.) In experiments of this sort, the quantity that is usually measured is the scattering cross section $\sigma$, which is the number of scattering events per unit time, unit volume, unit density of target particles, and unit flux of beam particles. (Here density means number of particles per unit volume, and the flux of the beam is its density times the speed of its particles.) Since time, volume, density, and flux have dimensions $[t],\left[l^{3}\right],\left[1 / l^{3}\right]$, and $\left[l / t l^{3}\right]=\left[1 / t l^{2}\right], \sigma$ has dimensions $\left[l^{2}\right]$, i.e., it is an area. One can think of it as the effective cross-sectional area of the target that is scattered by each particle in the beam.

Now, returning to our model, we take the box $\mathcal{B}$ to be the (macroscopic) region where the interaction of the beam and target takes place, and $T$ to be the (correspondingly long) time during which the interaction takes place. Since the wave functions $\psi$ of particles with definite momenta are evenly spread out in space, i.e., $|\psi|^{2}=$ constant, $1 / V$ can be interpreted as the density of each particle, so normalizing a quantity to be "per unit density" means multiplying by $V$. Thus, from the definition of $\sigma$ in the preceding paragraph, we see that the differential cross-section for an initial state $|\mathrm{in}\rangle$ to be scattered into a range of states with momenta in the infinitesimal region $d^{3 N^{\text {out }}} \mathbf{p}^{\text {out }}$ is

$$
d \sigma(\text { in } \rightarrow \text { out })=d P(\text { in } \rightarrow \text { out }) \cdot \frac{1}{T} \cdot \frac{1}{V} \cdot V \cdot \frac{V}{|\mathbf{v}|}=\frac{V d P(\text { in } \rightarrow \text { out })}{T|\mathbf{v}|}
$$

where $\mathbf{v}$ is the velocity of the beam. Plugging (6.68) into this, noting that $N^{\mathrm{in}}=2$, and replacing the smeared-out $\delta_{T}$ by the plain $\delta$, we obtain

$$
d \sigma(\text { in } \rightarrow \text { out })=(2 \pi)^{4} \delta\left(p^{\text {out }}-p^{\text {in }}\right)|\mathbf{v}|^{-1} \left\lvert\,\left. M(\text { in } \rightarrow \text { out })\right|^{2} \frac{d^{3 N^{\text {out }}} \mathbf{p}^{\text {out }}}{(2 \pi)^{3 N_{\text {out }}}}\right.
$$

Finally the $V$ and $T$ have disappeared, and we have a useful result! Despite the considerable amount of handwaving in its derivation, this is actually the correct formula for computing cross-sections.

Let us say a few words about the Lorentz transformation properties of $d \sigma$, with an eye to removing the hypothesis that the target is stationary. The states $\mid$ in $\rangle$ and |out describe particles in definite spin states, so the transformation law for the matrix elements $S(\mathrm{in} \rightarrow$ out $)$ or $M(\mathrm{in} \rightarrow$ out $)$ involves the transformation of these spin states - in general, a complicated mess. However, the particles are not normally prepared or detected in any particular spin state, so the experimentally measured quantity is obtained by averaging $d \sigma($ in $\rightarrow$ out) over an orthonormal basis of spin states for the incoming particles and summing it over an orthonormal basis of spin states for the outgoing particles (with the given momenta); we denote this process simply by $\sum_{\text {spins }}$. A close examination of the definitions of the variables that enter into $S($ in $\rightarrow$ out $)$ and $M($ in $\rightarrow$ out $)$, which we omit (see Weinberg [131], §3.4), shows that the quantity

$$
R(\text { in } \rightarrow \text { out })=\sum_{\text {spins }} \mid\left. M(\text { in } \rightarrow \text { out })\right|^{2} \prod E_{j}^{\text {in }} \prod E_{k}^{\text {out }}
$$

is Lorentz-invariant, where the $E_{j}^{\text {in }}$ and $E_{k}^{\text {out }}$ are the energies of the incoming and outgoing particles, respectively. On the other hand,

$$
\frac{d^{3 N^{\text {out }}} p^{\text {out }}}{(2 \pi)^{3 N^{\text {out }}} \prod E_{k}^{\text {out }}}=\prod \frac{d \mathbf{p}_{k}^{\text {out }}}{(2 \pi)^{3} E_{k}^{\text {out }}}
$$

is a Lorentz-invariant measure on the momentum space of the outgoing particles. Therefore, if $E_{1}$ and $E_{2}$ are the energies of the incoming particles, the quantity

$$
\sum_{\text {spins }} d \sigma(\text { in } \rightarrow \text { out })=\frac{(2 \pi)^{4} \delta\left(p^{\text {out }}-p^{\text {in }}\right)}{E_{1}^{\text {in }} E_{2}^{\text {in }}|\mathbf{v}|} R(\text { in } \rightarrow \text { out }) \prod \frac{d \mathbf{p}_{k}^{\text {out }}}{(2 \pi)^{3} E_{k}^{\text {out }}}
$$

will be Lorentz-invariant provided we replace the term $E_{1} E_{2}|\mathbf{v}|$ by the Lorentzinvariant quantity that agrees with it when the second particle is at rest, namely,

$$
U=\sqrt{\left(p_{1 \mu} p_{2}^{\mu}\right)^{2}-m_{1}^{2} m_{2}^{2}}
$$

$p_{1}, p_{2}$ and $m_{1}, m_{2}$ being the 4 -momenta and masses of the two initial particles. (Indeed, $U$ is clearly Lorentz-invariant; and if $\mathbf{p}_{2}=\mathbf{0}$, then $m_{2}=E_{2}$ and $p_{1 \mu} p_{2}^{\mu}= E_{1} E_{2}$, so $U=E_{2} \sqrt{E_{1}^{2}-m_{1}^{2}}=E_{2}\left|\mathbf{p}_{1}\right|=E_{1} E_{2}\left|\mathbf{v}_{1}\right|$.)

The quantity $U / E_{1} E_{2}$ can still be interpreted as a "relative velocity" when both incoming particles are in motion provided that they are aimed directly at each other so that $\mathbf{p}_{1} \cdot \mathbf{p}_{2}=-\left|\mathbf{p}_{1}\right|\left|\mathbf{p}_{2}\right|$. Indeed, in this case a short calculation shows that $U=E_{1}\left|\mathbf{p}_{2}\right|+E_{2}\left|\mathbf{p}_{1}\right|$, so

$$
\frac{U}{E_{1} E_{2}}=\frac{\left|\mathbf{p}_{1}\right|}{E_{1}}+\frac{\left|\mathbf{p}_{2}\right|}{E_{2}}=\left|\mathbf{v}_{1}-\mathbf{v}_{2}\right|,
$$

the absolute difference in velocities of the two particles as viewed from the laboratory frame. (It is not, however, the speed of one particle as viewed from the rest frame of the other.)

After the case of two initial particles, the most important situation is that of one initial particle. The "scattering" process here is the decay of an unstable particle. When $N^{\mathrm{in}}=1$, the $V$ factor in (6.68) disappears, so dividing by $T$ (again, the time during which the decay process occurs) yields the differential transition rate

$$
d \Gamma(\text { in } \rightarrow \text { out })=\frac{d P(\text { in } \rightarrow \text { out })}{T}
$$

If one replaces $\delta_{T}$ by $\delta$ (a step that requires some comment - see below), integrates over momenta, and sums over spins and possible decay modes, one obtains a formula for the total decay rate of the particle:

$$
\Gamma=\sum_{\text {spins, modes }} \frac{1}{E^{\text {in }}} \int R(\text { in } \rightarrow \text { out })(2 \pi)^{4} \delta\left(p^{\text {out }}-p^{\text {in }}\right) \prod \frac{d^{3} \mathbf{p}_{k}^{\text {out }}}{(2 \pi)^{3} E_{k}^{\text {out }}},
$$

where $R$ (in → out) is as in (6.70). This expression is Lorentz-invariant except for the energy $E^{\text {in }}$ of the initial particle; that factor incorporates the relativistic time dilation, according to which more rapidly moving particles (with respect to the laboratory frame) decay more slowly (with respect to the laboratory clock). The "half-life" of the particle in the usual sense is $(\log 2) / \Gamma_{0}$, where $\Gamma_{0}$ is the decay rate in the rest frame of the particle.

However, one cannot let $T \rightarrow \infty$ with impunity to replace $\delta_{T}$ by $\delta$ in (6.71): unstable particles are unstable, after all, and one cannot assume that they come from the infinitely distant past. This is the same problem that we encountered in studying the decay of an excited state in §6.2: for the approximation $\delta_{T} \approx \delta$ to be valid, $T$ must be large enough so that $1 / T$ is much less than the energies $E^{\text {in }}, E_{k}^{\text {out }}$ involved in the process, but $d P($ in $\rightarrow$ out $) / T$ can be interpreted as a transition rate only if $T$ is much less than the mean lifetime of the particle, so there needs to be a gap between these two quantities. The energies involved in typical particle decays correspond to times of $10^{-20}$ second or less, so this analysis is valid for all but the extremely short-lived particles. (See Peskin and Schroeder [89], §7.3, for an analysis of the latter situation.) As with the decay of excited states, the fact that $\delta_{T}$ is not exactly $\delta$ manifests itself in the fact that there is an unremovable uncertainty in the experimentally observed energy difference $E^{\text {out }}-E^{\text {in }}$.

Processes with three or more incoming particles are occasionally encountered, and the general formula (6.68) is again the bridge from field theory to experimental or observational results concerning them. However, we shall say no more about them here.

\subsection*{QED, the Coulomb potential, and the Yukawa potential}
It is time to see how to connect our machinery to the real world. To begin with, the description of electromagnetism provided by QED, in terms of emission and absorption of photons by electrons or other Fermions, looks very different from the description given by classical field theory, so we had better be able to relate the latter to the former. For this purpose, let us examine the basic Feynman diagram for the interaction of two electrons (Figure 6.6). We shall calculate the value of this\\
diagram in excruciating detail; later calculations of a similar sort can then be done somewhat more tersely. The momentum-space integral for this diagram is

$$
\begin{array}{r}
\int \frac{m^{2}}{4 \sqrt{\omega_{\mathbf{p}_{1}} \omega_{\mathbf{p}_{2}} \omega_{\mathbf{p}_{1}^{\prime}} \omega_{\mathbf{p}_{2}^{\prime}}}} \bar{u}\left(\mathbf{p}_{1}^{\prime}, s_{1}^{\prime}\right)\left(-i e \gamma^{\mu}\right) u\left(\mathbf{p}_{1}, s_{1}\right) \frac{-i g_{\mu \nu}}{q^{2}} \bar{u}\left(\mathbf{p}_{2}^{\prime}, s_{2}^{\prime}\right)\left(-i e \gamma^{\nu}\right) u\left(\mathbf{p}_{2}, s_{2}\right) \\
\times(2 \pi)^{4} \delta\left(p_{1}^{\prime}+q-p_{1}\right)(2 \pi)^{4} \delta\left(p_{2}^{\prime}-q-p_{2}\right) \frac{d^{4} q}{(2 \pi)^{4}}
\end{array}
$$

\begin{figure}[h]
\begin{center}
  \includegraphics[width=\textwidth]{2025_11_18_7960b2c161913da1e816g-184}

\caption{Figure 6.6. The basic Feynman diagram for electron-electron scattering.}
\end{center}
\end{figure}

Some explanations:\\
i. For the external (on-mass-shell) momenta we have $p=\left(\omega_{\mathbf{p}}, \mathbf{p}\right)$ as usual.\\
ii. The internal momentum $q$ is taken to flow from the first particle to the second one. It could equally well be taken to flow the other way; the only change would be in the signs of $q$ in the delta-functions.\\
iii. The Fourier-transformed photon propagator is the limit of $-i g_{\mu \nu} /\left(q^{2}+i \epsilon\right)$ as $\epsilon \rightarrow 0+$, and we have passed to the limit immediately since the $i \epsilon$ serves no purpose here.\\
iv. The factors of $m$ and $\omega_{\mathbf{p}}$ come from the $\sqrt{m / 2 \omega_{\mathbf{p}}}$ in the formula (5.32) for the Dirac field.

The $q$-integration in (6.72) is easy to perform:

$$
\begin{aligned}
\int \frac{-i g_{\mu \nu}}{q^{2}}(2 \pi)^{4} \delta\left(p_{1}^{\prime}+q-p_{1}\right)(2 \pi)^{4} \delta\left(p_{2}^{\prime}\right. & \left.-q-p_{2}\right) \frac{d^{4} q}{(2 \pi)^{4}} \\
& =\frac{-i g_{\mu \nu}}{\left(p_{1}^{\prime}-p_{1}\right)^{2}}(2 \pi)^{4} \delta\left(p_{1}^{\prime}+p_{2}^{\prime}-p_{1}-p_{2}\right)
\end{aligned}
$$

(We could equally well write $\left(p_{2}^{\prime}-p_{2}\right)^{2}$ instead of $\left(p_{1}^{\prime}-p_{1}\right)^{2}$.) The overall momentumconservation delta-function (with its faithful $(2 \pi)^{4}$ ) can safely remain in the background, so we factor it out and consider only the remaining part, which we denote by $i M$ as in (6.65):

$$
i M=\frac{i m^{2} e^{2}}{4 \sqrt{\omega_{\mathbf{p}_{1}} \omega_{\mathbf{p}_{2}} \omega_{\mathbf{p}_{1}^{\prime}} \omega_{\mathbf{p}_{2}^{\prime}}}} \bar{u}\left(\mathbf{p}_{1}^{\prime}, s_{1}^{\prime}\right) \gamma^{\mu} u\left(\mathbf{p}_{1}, s_{1}\right) \frac{g_{\mu \nu}}{\left(p_{1}^{\prime}-p_{1}\right)^{2}} \bar{u}\left(\mathbf{p}_{2}^{\prime}, s_{2}^{\prime}\right) \gamma^{\nu} u\left(\mathbf{p}_{2}, s_{2}\right) .
$$

Since Fermions are involved, we had better check that we have not misplaced a minus sign. The position-space integral corresponding to (6.72) comes from

$$
\begin{aligned}
\frac{1}{i^{2}} \iint\langle 0| a\left(\mathbf{p}_{1}^{\prime}, s_{1}^{\prime}\right) a\left(\mathbf{p}_{2}^{\prime}, s_{2}^{\prime}\right) A_{\mu}(x) \bar{\psi}(x) \gamma^{\mu} \psi(x) A_{\nu}(y) \bar{\psi}(y) \gamma^{\nu} \psi(y) & \\
& \times a^{\dagger}\left(\mathbf{p}_{2}, s_{2}\right) a^{\dagger}\left(\mathbf{p}_{1}, s_{1}\right)|0\rangle d^{4} x d^{4} y
\end{aligned}
$$

by pairing up the operators into the contractions

$$
\frac{1}{i^{2}} \iint \overbrace{a\left(\mathbf{p}_{1}^{\prime}, s_{1}^{\prime}\right) \bar{\psi}(x)} \gamma^{\mu} \overbrace{\psi(x) a^{\dagger}\left(\mathbf{p}_{1}, s_{1}\right)} \overbrace{A_{\mu}(x) A_{\nu}(y)} \overbrace{\times a\left(\mathbf{p}_{2}^{\prime}, s_{2}^{\prime}\right) \bar{\psi}(y)} \gamma^{\nu} \overbrace{\psi(y) a^{\dagger}\left(\mathbf{p}_{2}, s_{2}\right)} d^{4} x d^{4} y
$$

(Note that the orders of the creation operators for the initial and final states $a^{\dagger}\left(\mathbf{p}_{2}, s_{2}\right) a^{\dagger}\left(\mathbf{p}_{1}, s_{1}\right)|0\rangle$ and $a^{\dagger}\left(\mathbf{p}_{2}^{\prime}, s_{2}^{\prime}\right) a^{\dagger}\left(\mathbf{p}_{1}^{\prime}, s_{1}^{\prime}\right)|0\rangle$ are consistent!) One can easily check that the permutation of Fermion operators needed to change the order from that in (6.74) to that in (6.75) is even, so the sign in (6.73) is indeed correct. However, Figure 6.6 is only one of two diagrams that arise from the integral (6.74) and contribute to the basic photon exchange process; the other one is Figure 6.7, with the final particles reversed so that $a\left(\mathbf{p}_{1}^{\prime}, s_{1}^{\prime}\right)$ is contracted with $\bar{\psi}(y)$ and $a\left(\mathbf{p}_{2}^{\prime}, s_{2}^{\prime}\right)$ with $\bar{\psi}(x)$. The permutation needed to effect this arrangement is odd, so the value of Figure 6.7 is obtained from (6.73) by switching the $\bar{u}\left(\mathbf{p}_{j}^{\prime}, s_{j}^{\prime}\right)$ 's and inserting a factor of -1 . This -1 is important because it is the sum of the two diagrams that gives the S -matrix element. (Permuting the operators to go from (6.74) to (6.75) may seem to do violence to the matrix algebra, but it does not, as one sees by writing out all the components of the spinor fields explicitly; it is these individual components that are contracted with one another.)

\begin{figure}[h]
\begin{center}
  \includegraphics[width=\textwidth]{2025_11_18_7960b2c161913da1e816g-185}

\caption{Figure 6.7. Electron-electron scattering with outgoing particles interchanged.}
\end{center}
\end{figure}

By the way, we could equally well have used one of the alternative photon propagators (6.55). The result is exactly the same, because

$$
\bar{u}\left(\mathbf{p}_{1}^{\prime}, s_{1}^{\prime}\right) \gamma^{\mu} p_{1 \mu}^{\prime}=m \bar{u}\left(\mathbf{p}_{1}^{\prime}, s_{1}^{\prime}\right), \quad p_{1 \mu} \gamma^{\mu} u\left(\mathbf{p}_{1}, s_{1}\right)=m u\left(\mathbf{p}_{1}, s_{1}\right),
$$

so addition of any term proportional to $\left(p_{1}^{\prime}-p_{1}\right)_{\mu}$ to the $g_{\mu \nu}$ in (6.73) yields two terms that cancel out. Thus our claim that the $q_{\mu} q_{\nu}$ term in (6.55) does not contribute to physically meaningful quantities is verified in this simple case.

The fact that the diagrams of both Figures 6.6 and 6.7 contribute to the S matrix element complicates things a bit, so to make the connection with classical physics as simply as possible we shall consider the scattering of two different Fermions with possibly different masses $m_{1}$ and $m_{2}$ and possibly different charges\\
$e_{1}$ and $e_{2}$. Then Figure 6.6, where the subscripts 1 and 2 remain unmixed, is the only relevant one. Its value is still given by (6.73) except that the factor $m^{2} e^{2}$ must be replaced by $m_{1} m_{2} e_{1} e_{2}$.

Let us examine what happens to (6.73) in the nonrelativistic limit where the external momenta (and hence also the internal momentum $\mathbf{q}=\mathbf{p}_{1}-\mathbf{p}_{1}^{\prime}=\mathbf{p}_{2}^{\prime}-\mathbf{p}_{2}$ ) are all very small by replacing all the quantities in (6.73) by their lowest-order terms as functions of the momenta:

$$
\begin{aligned}
\omega_{\mathbf{p}_{j}} \rightarrow m_{j}, \quad \omega_{\mathbf{p}_{j}^{\prime}} \rightarrow m_{j}, \quad\left(p_{1}^{\prime}-p_{1}\right)^{2} \rightarrow-\left|\mathbf{p}_{1}^{\prime}-\mathbf{p}_{1}\right|^{2} \\
u\left(\mathbf{p}_{j}, s\right) \rightarrow u(\mathbf{0}, s), \quad u\left(\mathbf{p}_{j}^{\prime}, s_{j}^{\prime}\right) \rightarrow u\left(\mathbf{0}, s^{\prime}\right)
\end{aligned}
$$

Here the $u(\mathbf{0}, s)$, with $s= \pm$, are given by (5.29). We have

$$
\bar{u}(\mathbf{0}, s) \gamma^{0} u\left(\mathbf{0}, s^{\prime}\right)=u(\mathbf{0}, s)^{\dagger} u\left(\mathbf{0}, s^{\prime}\right)=2 \delta_{s s^{\prime}}
$$

by (5.29), and for $j=1,2,3$,

$$
\bar{u}(\mathbf{0}, s) \gamma^{j} u\left(\mathbf{0}, s^{\prime}\right)=u_{R}(\mathbf{0}, s)^{\dagger} \sigma_{j} u_{R}\left(\mathbf{0}, s^{\prime}\right)-u_{L}(\mathbf{0}, s)^{\dagger} \sigma_{j} u_{L}\left(\mathbf{0}, s^{\prime}\right)=0
$$

where $u_{R}$ and $u_{L}$ are as in (4.27) and the $\sigma_{j}$ are the Pauli matrices, simply because $u_{L}(\mathbf{0}, s)=u_{R}(\mathbf{0}, s)$ for all $s$. Therefore, since $g_{00}=1$, in this approximation (6.73) boils down to

$$
i M=\frac{-i e_{1} e_{2}}{\left|\mathbf{p}_{1}^{\prime}-\mathbf{p}_{1}\right|^{2}} \delta_{s_{1} s_{1}^{\prime}} \delta_{s_{2} s_{2}^{\prime}}
$$

The first conclusion is that the spins of the two particles are separately conserved by the interaction. To get to the more interesting electrostatic aspect, concentrate on the first particle and compare (6.78) with the Born approximation (6.25) to the S -matrix element for the scattering of the particle by a potential $V(\mathbf{x})$. Putting aside the energy-conservation delta-function and its $2 \pi$ as we did for the 4 -momentum-conservation delta-function here, ${ }^{10}$ we see that (6.78) gives (in the first-order approximation) the amplitude for scattering by the potential $V$ such that $\widehat{V}(\mathbf{q})=e_{1} e_{2} /|\mathbf{q}|^{2}$. The condition $|\mathbf{q}|^{2} \widehat{V}(\mathbf{q})=e_{1} e_{2}$ is equivalent to $-\nabla^{2} V(\mathbf{x})=e_{1} e_{2} \delta(\mathbf{x})$, so $V$ is just what it should be: the Coulomb potential $V(\mathbf{x})=e_{1} e_{2} / 4 \pi|\mathbf{x}|$. (If this result isn't familiar, see the calculation at the end of this section and take the limit as $m_{\phi} \rightarrow 0$.) Thus we recover the classical physics in the low-energy limit.

What happens if we replace one of the particles - say, the second one - by its antiparticle? The $a^{\dagger}\left(\mathbf{p}_{2}, s_{2}\right)$ and $a\left(\mathbf{p}_{2}^{\prime}, s_{2}^{\prime}\right)$ in (6.74) must be replaced by $b^{\dagger}\left(\mathbf{p}_{2}, s_{2}\right)$ and $b\left(\mathbf{p}_{2}^{\prime}, s_{2}^{\prime}\right)$, the creation and annihilation operators for the antiparticle. This has the effect that the $u\left(\mathbf{p}_{2}, s_{2}\right)$ and $\bar{u}\left(\mathbf{p}_{2}^{\prime}, s_{2}^{\prime}\right)$ in (6.73) are replaced by $v\left(\mathbf{p}_{2}, s_{2}\right)$ and $\bar{v}\left(\mathbf{p}_{2}^{\prime}, s_{2}^{\prime}\right)$. In the nonrelativistic approximation, these $v$ 's still satisfy (6.76) and (6.77) (this time because $v_{L}=-v_{R}$ ), so it would seem that (6.78) is unchanged. But no: there is an extra minus sign in going from (6.74) to (6.75) because $b$ and $b^{\dagger}$ contract with $\psi$ and $\bar{\psi}$ rather than the other way around, and this minus sign survives in (6.78). Thus this calculation is consistent with the fact that particles and antiparticles have opposite charges.

Let us carry this line of thought further by considering the corresponding calculations for the Yukawa interaction (6.29): the exchange of a quantum of the

\footnotetext{${ }^{10}$ In the center-of-mass frame in which $\mathbf{p}_{1}+\mathbf{p}_{2}=\mathbf{0}$, the spatial momentum-conservation delta-function disappears on integration over $\mathbf{p}_{2}$; all that is left is the energy-conservation deltafunction.
}
scalar field $\phi$ between two distinguishable Fermions (say, a proton and a neutron). The Feynman diagram in Figure 6.6 still serves the purpose, with the wavy line now representing a quantum of the field $\phi$. The corresponding quantities (6.73), (6.74), and (6.75) are the same as before except that (i) the field $A_{\mu}$ is replaced by $\phi$ and the photon propagator $-i g_{\mu \nu} / q^{2}$ is replaced by the scalar field propagator $-i /\left(-q^{2}+m_{\phi}^{2}\right)$, (ii) the $e_{1}$ and $e_{2}$ are replaced by the coupling constant $g$ for the Yukawa interaction (which we now assume to be the same for both particles, as it is for the proton and neutron), and (iii) the Dirac matrices $\gamma^{\mu}$ and $\gamma^{\nu}$ are omitted. ${ }^{11}$

We pass to the nonrelativistic approximation. The analogue of (6.76) is

$$
\bar{u}(\mathbf{0}, s) u\left(\mathbf{0}, s^{\prime}\right)=u(\mathbf{0}, s)^{\dagger} u\left(\mathbf{0}, s^{\prime}\right)=2 \delta_{s s^{\prime}}
$$

as before, because $\gamma_{0} u=u$. (There are no terms corresponding to (6.77).) The analogue of (6.78) is therefore

$$
i M=\frac{i g^{2}}{\left|\mathbf{p}_{1}^{\prime}-\mathbf{p}_{1}\right|^{2}+m_{\phi}^{2}} \delta_{s_{1} s_{1}^{\prime}} \delta_{s_{2} s_{2}^{\prime}}
$$

which corresponds to the potential $V$ such that $\widehat{V}(\mathbf{q})=-g^{2} /\left(|\mathbf{q}|^{2}+m_{\phi}^{2}\right)$. By a calculation that we shall sketch at the end of this section, we find that

$$
V(\mathbf{x})=-\frac{g^{2} e^{-m_{\phi}|\mathbf{x}|}}{4 \pi|\mathbf{x}|}
$$

This is called the Yukawa potential. It resembles the Coulomb potential for $\mathbf{x}$ small but decays rapidly for $\mathbf{x}$ large, and it is negative. Hence it describes an attractive short-range force whose range is essentially $1 / m_{\phi}$. The fact that the Yukawa interaction produces an attractive force whereas the electromagnetic interaction between particles of the same charge produces a repulsive force comes out of the opposite signs of $q^{2}$ in the photon and scalar field propagators.

What happens if we replace the second particle by its antiparticle in this situation? There is an extra minus sign from the permutation of Fermion operators as with the electromagnetic interaction, but here there is yet another minus sign from the replacement of $u$ 's by $v$ 's for the second particle, because $\gamma^{0} v=-v$ and hence

$$
\bar{v}(\mathbf{0}, s) v\left(\mathbf{0}, s^{\prime}\right)=-v(\mathbf{0}, s)^{\dagger} v\left(\mathbf{0}, s^{\prime}\right)=-2 \delta_{s s^{\prime}}
$$

Thus the force remains attractive! If we also replace the first particle by its antiparticle, we get another pair of minus signs, so (6.79) again remains unchanged. The Yukawa force is universally attractive as long as all Fermions in question interact with the scalar field with the same coupling constant $g$.

Yukawa proposed his interaction in 1935 as a model for the strong force that binds atomic nuclei. The existence of a particle representing the quantum of the $\phi$ field was one of the main predictions of the theory, and Yukawa estimated its mass at around 150 MeV from experimental data about the range of the strong force. The first particle to be discovered that seemed like a reasonable candidate to fit Yukawa's theory was the muon, but it didn't fill the bill for various reasons; Yukawa was finally vindicated with the discovery of pions (with masses around 140 MeV ) in 1947. His theory doesn't have many more notable successes, however, because the coupling constant $g$ is too large to allow the effective use of perturbation theory.

\footnotetext{${ }^{11}$ If one uses the pseudoscalar interaction (6.30) with $\gamma^{5}$ in place of $\gamma^{\mu}$ instead, the effect is to replace $u$ spinors with $v$ spinors and vice versa, but the final result concerning the Yukawa potential is unchanged.
}We conclude this section by sketching the derivation of (6.80), which is a bit delicate because $\widehat{V} \notin L^{1}$. One can begin by restricting the integration to the ball $|\mathbf{q}| \leq R$. Since the result is a radial function of $\mathbf{x}$, it suffices to take $\mathbf{x}=(0,0, a)$ with $a>0$; then integration in spherical coordinates gives

$$
\begin{aligned}
\int_{|\mathbf{q}| \leq R} \frac{e^{i \mathbf{q} \cdot \mathbf{x}}}{|\mathbf{q}|^{2}+m_{\phi}^{2}} \frac{d^{3} \mathbf{q}}{(2 \pi)^{3}} & =\int_{0}^{R} \int_{0}^{\pi} \frac{e^{i a r \cos \theta}}{r^{2}+m_{\phi}^{2}} r^{2} \sin \theta \frac{d \theta d r}{(2 \pi)^{2}} \\
& =\int_{0}^{R} \frac{r\left(e^{i a r}-e^{-i a r}\right)}{4 \pi^{2} i a\left(r^{2}+m_{\phi}^{2}\right)} d r=\int_{-R}^{R} \frac{r e^{i a r}}{4 \pi^{2} i a\left(r^{2}+m_{\phi}^{2}\right)} d r
\end{aligned}
$$

The limit of this as $R \rightarrow \infty$ can be evaluated via the residue theorem to give $e^{-m_{\phi} a} / 4 \pi a$. To certify the validity of (6.80) from this, one needs to verify that the convergence take place not just pointwise but in a sense at least as strong as weak convergence of distributions. Alternatively, once one has used this calculation to guess (6.80), it is a completely elementary exercise (using spherical coordinates as above) to check that the Fourier transform of the $L^{1}$ function $e^{-m_{\phi}|\mathbf{x}|} / 4 \pi|\mathbf{x}|$ is $1 /\left(|\mathbf{q}|^{2}+m_{\phi}^{2}\right)$.

\subsection*{Compton scattering}
As another illustration of the basic calculations of QED, we sketch the calculation of the cross-section for the scattering of a photon by an electron, known as Compton scattering, in the first-order approximation. The initial state $|\mathrm{in}\rangle$ consists of an electron with momentum $\mathbf{p}$ and spin state $s$ and a photon with momentum $\mathbf{k}$ and polarization 4-vector $\epsilon=(0, \boldsymbol{\epsilon})$; the final state $\mid$ out $\rangle$ consists of an electron with momentum $\mathbf{p}^{\prime}$ and spin state $s^{\prime}$ and a photon with momentum $\mathbf{k}^{\prime}$ and polarization 4-vector $\epsilon^{\prime}=\left(0, \boldsymbol{\epsilon}^{\prime}\right)$. (Here $\boldsymbol{\epsilon}$ and $\boldsymbol{\epsilon}^{\prime}$ must be orthogonal to $\mathbf{k}$ and $\mathbf{k}^{\prime}$, respectively.) We shall assume the electron is initially at rest, so that $\mathbf{p}=\mathbf{0}$ and $p_{0}=m$, where $m$ is the mass of the electron. We denote by $\omega=|\mathbf{k}|$ and $\omega^{\prime}=\left|\mathbf{k}^{\prime}\right|$ the energies of the incoming and outgoing photons, and by $p, k, p^{\prime}, k^{\prime}$ the 4 -momenta of the electrons and photons (i.e., $p=(m, \mathbf{0}), k=(\omega, \mathbf{k})$, etc.).

Conservation of energy and momentum means that

$$
m+\omega=\sqrt{m^{2}+\left|\mathbf{p}^{\prime}\right|^{2}}+\omega^{\prime}, \quad \mathbf{k}=\mathbf{p}^{\prime}+\mathbf{k}^{\prime},
$$

so that\\
(6.81) $\left(m+\omega-\omega^{\prime}\right)^{2}=m^{2}+\left|\mathbf{p}^{\prime}\right|^{2}=m^{2}+\left|\mathbf{k}-\mathbf{k}^{\prime}\right|^{2}=m^{2}+\omega^{2}+\omega^{\prime 2}-2 \omega \omega^{\prime} \cos \theta$,\\
where $\theta$ is the angle between the incoming and outgoing photon momenta. A little algebraic manipulation with this yields the relation between the incoming and outgoing photon energies (or wavelengths):

$$
\frac{1}{\omega^{\prime}}=\frac{1}{\omega}+\frac{1-\cos \theta}{m}
$$

which was confirmed in Compton's experiments on the scattering of x-rays by electrons.\\
(At the end of the nineteenth century the wave theory of light appeared to have won a decisive victory over Newton's corpuscular theory. In 1900 and 1905, Planck and Einstein revived the hypothesis that electromagnetic radiation is quantized to explain the spectrum of black-body radiation and the photoelectric effect, but many physicists remained dubious. Compton's experiments in 1923 clinched the matter: the relation (6.82) is almost obvious if x-rays consist of photons, as we have just\\[0pt]
seen, but it is hard to account for if one uses a wave model. See Pais [88] for a fuller account of this history.)

\begin{figure}[h]
\begin{center}
  \includegraphics[width=\textwidth]{2025_11_18_7960b2c161913da1e816g-189}

\caption{Figure 6.8. Feynman diagrams for Compton scattering.}
\end{center}
\end{figure}

In the first-order approximation, there are two Feynman diagrams that contribute to this process, shown in Figure 6.8. The sum of these diagrams is

$$
\begin{gathered}
S=\frac{(-i e)^{2}(2 \pi)^{4} \sqrt{m}}{4 \sqrt{\omega_{\mathbf{p}^{\prime}} \omega \omega^{\prime}}} \int \bar{u}\left(\mathbf{p}^{\prime}, s^{\prime}\right)\left[\epsilon_{\mu}^{\prime *} \gamma^{\mu} \frac{-i\left(q_{\rho} \gamma^{\rho}+m\right)}{-q^{2}+m^{2}} \epsilon_{\nu} \gamma^{\nu} \delta\left(q-p^{\prime}-k^{\prime}\right) \delta(q-p-k)\right. \\
\left.+\epsilon_{\mu} \gamma^{\mu} \frac{-i\left(q_{\rho} \gamma^{\rho}+m\right)}{-q^{2}+m^{2}} \epsilon_{\nu}^{\prime *} \gamma^{\nu} \delta\left(q+k-p^{\prime}\right) \delta\left(q+k^{\prime}-p\right)\right] u(\mathbf{0}, s) d^{4} q
\end{gathered}
$$

(The $(2 \pi)^{4}$ comes from the $(2 \pi)^{4}$ 's attached to the delta-functions and the $1 /(2 \pi)^{4}$ attached to the $d^{4} q$, and we have used the fact that $\omega_{\mathbf{p}}=\omega_{\mathbf{0}}=m$.)

After performing the integration (merely a convolution of delta-functions) and using the facts that $p^{2}=m^{2}$ and $k^{2}=k^{\prime 2}=0$ to simplify the resulting denominators $-(p+k)^{2}+m^{2}$ and $-\left(p-k^{\prime}\right)^{2}+m^{2}$, we find that

$$
S=(2 \pi)^{4} \delta\left(p^{\prime}+k^{\prime}-p-k\right) i M,
$$

where

$$
\begin{aligned}
& M=\frac{-e^{2} \sqrt{m}}{8 \sqrt{\omega_{\mathbf{p}^{\prime}} \omega \omega^{\prime}}} \\
& \times \bar{u}\left(\mathbf{p}^{\prime}, s^{\prime}\right)\left[\frac{\epsilon_{\mu}^{\prime *} \gamma^{\mu}\left[(p+k)_{\rho} \gamma^{\rho}+m\right] \epsilon_{\nu} \gamma^{\nu}}{p_{\rho} k^{\rho}}-\frac{\epsilon_{\mu} \gamma^{\mu}\left[\left(p-k^{\prime}\right)_{\rho} \gamma^{\rho}+m\right] \epsilon_{\nu}^{\prime *} \gamma^{\nu}}{p_{\rho} k^{\prime \rho}}\right] u(\mathbf{0}, s)
\end{aligned}
$$

This $M$ is the delta-function-free matrix element that is used to calculate crosssections. Namely, according to (6.68), the differential cross-section for scattering into the region $d^{3} \mathbf{p}^{\prime} d^{3} \mathbf{k}^{\prime}$ of momentum space is

$$
d \sigma=(2 \pi)^{4} \delta\left(p^{\prime}+k^{\prime}-p-k\right)|M|^{2} \frac{d^{3} \mathbf{p}^{\prime} d^{3} \mathbf{k}^{\prime}}{(2 \pi)^{6}} .
$$

The factor of $|\mathbf{v}|$ has disappeared because we have assumed that the electron is at rest, so $|\mathbf{v}|$ is just the speed of the incoming photon, namely, 1. Now,

$$
(2 \pi)^{4} \delta\left(p^{\prime}+k^{\prime}-p-k\right) \frac{d^{3} \mathbf{p}^{\prime} d^{3} \mathbf{k}^{\prime}}{(2 \pi)^{6}}=\delta\left(\mathbf{p}^{\prime}+\mathbf{k}^{\prime}-\mathbf{k}\right) d^{3} \mathbf{p}^{\prime} \cdot \delta\left(p_{0}^{\prime}+\omega^{\prime}-m-\omega\right) \frac{d^{3} \mathbf{k}^{\prime}}{(2 \pi)^{2}} .
$$

The first delta-function on the right just means that in calculating the cross-section we must set $\mathbf{p}^{\prime}=\mathbf{k}-\mathbf{k}^{\prime}$ and drop the $d^{3} \mathbf{p}^{\prime}$. Likewise, if we write the element of volume in $\mathbf{k}^{\prime}$-space in spherical coordinates,

$$
d^{3} \mathbf{k}^{\prime}=\omega^{\prime 2} d \omega^{\prime} d \Omega,
$$

where $d \Omega$ is surface measure on the unit sphere, the second delta-function fixes $\omega^{\prime}$ at the value given by (6.82) and eliminates the $d \omega^{\prime}$, but an extra factor arises from the general fact that $\delta(f(t))=\delta\left(t-t_{0}\right) / f^{\prime}\left(t_{0}\right)$ when $f$ is a smooth function that vanishes at $t_{0}$. In our case, by (6.81),

$$
f\left(\omega^{\prime}\right)=p_{0}^{\prime}+\omega^{\prime}-m-\omega=\sqrt{\omega^{2}-2 \omega \omega^{\prime} \cos \theta+\omega^{\prime 2}+m^{2}}+\omega^{\prime}-m-\omega,
$$

so with $p_{0}^{\prime}=m+\omega-\omega^{\prime}$ and $\omega^{\prime}$ given by (6.82),

$$
f^{\prime}\left(\omega^{\prime}\right)=\frac{\omega^{\prime}-\omega \cos \theta}{p_{0}^{\prime}}+1=\frac{-\omega \cos \theta+m+\omega}{m+\omega-\omega^{\prime}}=\frac{m \omega}{\omega^{\prime}\left(m+\omega-\omega^{\prime}\right)} .
$$

The upshot is that the differential cross-section for the direction of the outgoing photon to be in the solid angle $d \Omega$ located at an angle $\theta$ from the direction of the incoming photon is

$$
d \sigma=|M|^{2} \frac{\left(m+\omega-\omega^{\prime}\right) \omega^{\prime 3}}{(2 \pi)^{2} m \omega} d \Omega
$$

where $\omega^{\prime}$ is determined by $\omega$ according to (6.82).\\
This is all very well, but the simple notation $|M|^{2}$ conceals a rather formidable expression (see (6.83)) to calculate. Moreover, $|M|^{2}$ depends on the spin and helicity states of the incoming and outgoing electrons and photons. Normally one does not wish to specify or measure the electron spin state, so one must average over the two orthonormal states $s= \pm$ for the incoming electron and sum over the two states for the outgoing electron. Plugging in the definition of $M$ and reducing the resulting expression to something manageable takes several pages of tedious algebra involving products of Dirac matrices, which the mathematical tourist would probably prefer to skip. We refer to Weinberg [131], §8.7, or Peskin and Schroeder [89], §5.5, for the details, with the warning that their normalizations differ from ours. However, the final result is quite simple:

$$
\frac{1}{2} \sum_{s, s^{\prime}} d \sigma=\frac{e^{4} \omega^{\prime 2}}{64 \pi^{2} m^{2} \omega^{2}}\left[\frac{\omega}{\omega^{\prime}}+\frac{\omega^{\prime}}{\omega}-2+4\left(\boldsymbol{\epsilon} \cdot \boldsymbol{\epsilon}^{\prime}\right)^{2}\right] d \Omega
$$

This formula was originally derived in 1929 by Klein and Nishina by more oldfashioned methods.

Furthermore, the polarization $\boldsymbol{\epsilon}$ of the incoming photon is also usually unspecified, so one should average over the two orthonormal possibilities for it. The result is

$$
\frac{1}{4} \sum_{s, s^{\prime}, \boldsymbol{\epsilon}} d \sigma=\frac{e^{4} \omega^{\prime 2}}{32 \pi^{2} m^{2} \omega^{2}}\left[\frac{\omega}{\omega^{\prime}}+\frac{\omega^{\prime}}{\omega}-2\left(\boldsymbol{\epsilon}^{\prime} \cdot \widehat{\mathbf{k}}\right)^{2}\right] d \Omega
$$

where $\widehat{\mathbf{k}}=\mathbf{k} /|\mathbf{k}|$. It follows that $d \sigma$ is maximized when $\boldsymbol{\epsilon}^{\prime}$ is orthogonal to $\mathbf{k}$, and $\boldsymbol{\epsilon}^{\prime}$ is always orthogonal to $\mathbf{k}^{\prime}$, so the outgoing photon is preferentially polarized in the direction orthogonal to the plane of the scattering.

The cross-section for outgoing photons of arbitrary polarization is obtained by summing (6.84) over two orthonormal values of $\boldsymbol{\epsilon}^{\prime}$, and the result is

$$
\frac{1}{4} \sum_{s, s^{\prime}, \boldsymbol{\epsilon}, \boldsymbol{\epsilon}^{\prime}} d \sigma=\frac{e^{4} \omega^{\prime 2}}{32 \pi^{2} m^{2} \omega^{2}}\left[\frac{\omega}{\omega^{\prime}}+\frac{\omega^{\prime}}{\omega}-1+\cos ^{2} \theta\right] d \Omega
$$

(Recall that $\theta$ is the angle between the incoming and outgoing photon momenta.) If we assume that the incoming photon energy $\omega$ is much less than the mass of the\\
electron (as it is for any photon less energetic than a gamma ray), (6.82) shows that $\omega^{\prime} / \omega \approx 1$. With this approximation, (6.85) simplifies to

$$
\frac{1}{4} \sum_{s, s^{\prime}, \boldsymbol{\epsilon}, \boldsymbol{\epsilon}^{\prime}} d \sigma=\frac{e^{4}}{32 \pi^{2} m^{2}}\left(1+\cos ^{2} \theta\right) d \Omega
$$

and since

$$
\int\left(1+\cos ^{2} \theta\right) d \Omega=\int_{0}^{2 \pi} \int_{0}^{\pi}\left(1+\cos ^{2} \theta\right) \sin \theta d \theta d \phi=\frac{16 \pi}{3}
$$

the total cross-section is

$$
\sigma_{T}=\frac{e^{4}}{6 \pi m^{2}}=\frac{8 \pi r_{0}^{2}}{3}
$$

where $r_{0}$ is the Compton radius (1.7). ${ }^{12}$ This result was originally derived in the context of classical electrodynamics as a formula relating to the scattering of lowenergy radiation by a static charge.

\subsection*{The Gell-Mann-Low and LSZ formulas}
In the preceding sections we have taken the most direct path to the integrals corresponding to Feynman diagrams that are used to compute quantities that can be tested against experiment. All of our calculations have been based on free field operators; we have not so much as mentioned operators representing interacting fields. ${ }^{13}$ Moreover, the full interacting Hamiltonian was quickly swept under the rug after making a brief appearance at the beginning of the chapter.

From a mathematical point of view, this is all to the good. A precise mathematical construction of the interacting fields that describe actual fundamental physical processes in 4-dimensional space-time is still lacking and may not be feasible without serious modifications to the theory. Similarly, we have no way to define the Hamiltonian $H$ in a mathematically rigorous way as a self-adjoint operator. It was presented as the sum of the free Hamiltonian $H_{0}$, which is well-defined, and the interaction Hamiltonian $H_{I}$; but the latter was presented as the integral of a density consisting of products of fields, which are operator-valued distributions rather than functions. Fortunately, this formal characterization of $H_{I}$ is sufficient to lead to well-defined calculations in perturbation theory, but we have not really defined $H_{I}$ as an operator; and if we had, we would still not have shown how to define the appropriate domain to make the sum $H_{0}+H_{I}$ self-adjoint.

The devil now offers us another bargain: If we are willing to accept the existence of the interacting fields and the full Hamiltonian associated to them (but without any expectation of finding a mathematically rigorous way to construct them) and to make a few arguments on the level of pure hand-waving about their structure, we can use them to obtain a more complete picture of the theory.

From a practical point of view, it is possible to do a lot of quantum field theory without accepting this bargain. The perturbation theory of this chapter together with the renormalization techniques of the next one suffice to calculate cross sections, anomalous magnetic moments, energy level shifts, and other such items of experimental interest. In this section, however, we shall accept the bargain to the extent of making the connection between scattering processes and Feynman

\footnotetext{${ }^{12}$ The subscript $T$ can stand for either "total" or "Thomson."\\
${ }^{13}$ Hence the carefully chosen title of the chapter.
}
diagrams, on the one hand, and vacuum expectation values of time-ordered products of interacting fields, on the other. The latter are widely used in the physics literature and are of fundamental importance in the study of rigorous models that satisfy the Wightman axioms. But from a strict mathematical point of view, everything in this section should be taken on the level of heuristics and plausibility arguments. We will add some credibility to main results by rederiving them in an entirely different, still nonrigorous, but intuitively very appealing way in Chapter 8.

To keep the notation manageable, we consider the simple case of the $\phi^{4}$ scalar field theory. The generalization to several fields with arbitrary spin and interactions given by suitable products of field operators is straightforward and will be sketched in due course.

The setup is as follows: We begin with a Hilbert space of states on which we are given a self-adjoint operator $H$, the full Hamiltonian, whose spectrum is bounded below by some number $E_{0}$ that is a simple eigenvalue with eigenvector $|\Omega\rangle$. Furthermore, we are given an operator-valued distribution $\phi$ on $\mathbb{R}^{4}$, and as usual, we shall pretend that it is a function for notational purposes. The time dependence of $\phi$ is "Heisenberg-picture," that is,

$$
\phi(t, \mathbf{x})=e^{i H t} \phi(0, \mathbf{x}) e^{-i H t}
$$

and $H$ is formally defined in terms of $\phi$ by

$$
H=\int\left[\frac{1}{2}\left[\left(\partial_{t} \phi(t, \mathbf{x})\right)^{2}+\left|\nabla_{\mathbf{x}} \phi(t, \mathbf{x})\right|^{2}+m^{2} \phi(t, \mathbf{x})^{2}\right]+\frac{\lambda}{4!} \phi(t, \mathbf{x})^{4}\right] d^{3} \mathbf{x}
$$

(The expression on the right apparently depends on $t$, but in fact it does not, in view of (6.86) and the fact that $H$ commutes with $e^{i H t}$.) Associated to $\phi$ and $H$ is the classical Lagrangian density

$$
\mathcal{L}=\frac{1}{2}\left[\left(\partial_{\mu} \phi\right)\left(\partial^{\mu} \phi\right)-m^{2} \phi^{2}\right]-\frac{\lambda}{4!} \phi^{4}
$$

and with respect to this Lagrangian there is a canonically conjugate field $\pi$, namely, $\pi=\partial \mathcal{L} / \partial\left(\partial_{t} \phi\right)=\partial_{t} \phi$ just as for the free field. We assume that $\phi$ and $\pi$ satisfy the canonical equal-time commutation relations

$$
[\phi(t, \mathbf{x}), \pi(t, \mathbf{y})]=i \delta(\mathbf{x}-\mathbf{y}), \quad[\phi(t, \mathbf{x}), \phi(t, \mathbf{y})]=[\pi(t, \mathbf{x}), \pi(t, \mathbf{y})]=0
$$

We fix a reference time, which may as well be $t_{0}=0$, and expand $\phi(0, \mathbf{x})$ and $\pi(0, \mathbf{x})$ as Fourier integrals,

$$
\phi(0, \mathbf{x})=\int e^{i \mathbf{p} \cdot \mathbf{x}} \widehat{\phi}(\mathbf{p}) \frac{d^{3} \mathbf{p}}{(2 \pi)^{3}}, \quad \pi(0, \mathbf{x})=\int e^{i \mathbf{p} \cdot \mathbf{x}} \widehat{\pi}(\mathbf{p}) \frac{d^{3} \mathbf{p}}{(2 \pi)^{3}}
$$

where $\widehat{\phi}$ and $\widehat{\pi}$ are operator-valued distributions on $\mathbb{R}^{3}$. Since $\phi(0, \mathbf{x})$ and $\pi(0, \mathbf{x})$ are Hermitian, we have $\widehat{\phi}(\mathbf{p})^{\dagger}=\widehat{\phi}(-\mathbf{p})$ and $\widehat{\pi}(\mathbf{p})^{\dagger}=\widehat{\pi}(-\mathbf{p})$. We next set

$$
a(\mathbf{p})=\frac{\omega_{\mathbf{p}} \widehat{\phi}(\mathbf{p})+i \widehat{\pi}(\mathbf{p})}{\sqrt{2 \omega_{\mathbf{p}}}} \quad\left(\omega_{\mathbf{p}}=\sqrt{m^{2}+|\mathbf{p}|^{2}}\right)
$$

Then $a^{\dagger}(-\mathbf{p})=\left[\omega_{\mathbf{p}} \widehat{\phi}(\mathbf{p})-i \widehat{\pi}(\mathbf{p})\right] / \sqrt{2 \omega_{\mathbf{p}}}$, so

$$
\widehat{\phi}(\mathbf{p})=\frac{a(\mathbf{p})+a^{\dagger}(-\mathbf{p})}{\sqrt{2 \omega_{\mathbf{p}}}}, \quad \widehat{\pi}(\mathbf{p})=-i \sqrt{\frac{\omega_{\mathbf{p}}}{2}}\left[a(\mathbf{p})-a^{\dagger}(-\mathbf{p})\right]
$$

and hence $\phi(0, \mathbf{x})$ and $\pi(0, \mathbf{x})$ are given by the same integrals (5.11) and (5.19) as for free fields. A short calculation then shows that the commutation relations (6.88)\\
for the fields imply the canonical commutation relations (5.10) for the operators $a(\mathbf{p})$ and $a^{\dagger}(\mathbf{p})$. In short, the field $\phi(0, \mathbf{x})$ has the same formal structure as the free field studied in Chapter 5; what is different is the time evolution.

We wish to describe things in the "interaction picture," obtained by decomposing the Hamiltonian (6.87), with $t=0$, into the free part and the interaction part:

$$
\begin{aligned}
H & =H_{0}+H_{I} \\
H_{0} & =\int \frac{1}{2}\left[\left(\partial_{t} \phi(0, \mathbf{x})\right)^{2}+\left|\nabla_{\mathbf{x}} \phi(0, \mathbf{x})\right|^{2}+m^{2} \phi(0, \mathbf{x})^{2}\right] d^{3} \mathbf{x} \\
H_{I} & =\frac{\lambda}{4!} \int \phi(0, \mathbf{x})^{4} d^{3} \mathbf{x}
\end{aligned}
$$

As in §5.3, $H_{0}$ needs to be renormalized by subtracting off an infinite constant, or by Wick ordering, so that the lowest eigenvalue of $H_{0}$ is 0 . When this is done, the calculations in §5.3 following (5.23) show that $H_{0}$ is given in terms of $a(\mathbf{p})$ and $a^{\dagger}(\mathbf{p})$ by (5.22). We define the interaction picture field $\phi_{0}$ by

$$
\phi_{0}(t, \mathbf{x})=e^{i H_{0} t} \phi(0, \mathbf{x}) e^{-i H_{0} t}
$$

The quantities on the right are defined in terms of $a(\mathbf{p})$ and $a^{\dagger}(\mathbf{p})$ just as in the free field case, so $\phi_{0}$ is indeed structurally identical to a free field and is given by (5.12):

$$
\phi_{0}(x)=\int \frac{1}{\sqrt{2 \omega_{\mathbf{p}}}}\left(e^{-i p_{\mu} x^{\mu}} a(\mathbf{p})+e^{i p_{\mu} x^{\mu}} a^{\dagger}(\mathbf{p})\right) \frac{d^{3} \mathbf{p}}{(2 \pi)^{3}}
$$

In the interaction picture, the interaction Hamiltonian is the time-dependent operator $H_{I}(t)$ defined by

$$
H_{I}(t)=e^{i H_{0} t} H_{I} e^{-i H_{0} t}=\frac{\lambda}{4!} \int \phi_{0}(t, \mathbf{x})^{4} d^{3} \mathbf{x}
$$

As in §6.1 and §6.3, we define

$$
V\left(t, t^{\prime}\right)=e^{i H_{0} t} e^{-i H\left(t-t^{\prime}\right)} e^{-i H_{0} t^{\prime}}
$$

so that

$$
V\left(t, t^{\prime}\right) V\left(t^{\prime}, t^{\prime \prime}\right)=V\left(t, t^{\prime \prime}\right), \quad \phi(t, \mathbf{x})=V(0, t) \phi_{0}(t, \mathbf{x}) V(t, 0)
$$

and, for $t>t^{\prime}, V\left(t, t^{\prime}\right)$ is given perturbation-theoretically by the Dyson series:

$$
V\left(t, t^{\prime}\right)=\mathcal{T} \exp \left[\frac{1}{i} \int_{t^{\prime}}^{t} H_{I}(\tau)\right] d \tau
$$

This is the point where the mathematical credibility of the argument falls to the infinitesimal level. One has to face the fact that the Stone-von Neumann theorem is false in infinite dimensions, so there are many inequivalent representations of the canonical commutation relations (6.33) (see Segal [111]). We therefore have no right to expect that the field $\phi_{0}$ is just the free field studied in Chapter 5. In fact, a theorem of Haag (see Streater and Wightman [116] or Bogolubov et al. [11], [12]) says in effect that it cannot be unless the interaction is trivial. Nonetheless, we shall proceed: after a bit we shall arrive at some results that beg to be interpreted in perturbation theory, in terms of Feynman diagrams, and on that level they seem to be perfectly correct. There are some unresolved mysteries here about the radical restructuring of a quantum system that is produced by any nontrivial interaction of fields, but this is not the place to explore them.

The next step is to examine the vacuum states. As mentioned earlier, we denote the vacuum for $H$ (i.e., the unique state whose eigenvalue $E_{0}$ is the infimum of the spectrum of $H)$ by $|\Omega\rangle$. The free Hamiltonian $H_{0}$ has its own vacuum state, i.e., its eigenstate with eigenvalue 0 , which we denote as before by $|0\rangle$. There is no reason for $|\Omega\rangle$ and $|0\rangle$ to coincide, but to extricate ourselves from the wilderness of the preceding paragraph we must take on faith that they have some overlap, i.e., that $\langle\Omega \mid 0\rangle \neq 0$. Granted this, there is a neat device for expressing $|\Omega\rangle$ in terms of $|0\rangle$. Let $P$ be the projection-valued measure associated to $H$, so that $H$ has the spectral resolution

$$
H=\int_{\left[E_{0}, \infty\right)} E d P(E)=E_{0} P\left(\left\{E_{0}\right\}\right)+\int_{\left(E_{0}, \infty\right)} E d P(E)
$$

Then for any $T \in \mathbb{R}$,

$$
e^{-i T\left(H-E_{0}\right)}|0\rangle=|\Omega\rangle\langle\Omega \mid 0\rangle+\int_{\left(E_{0}, \infty\right)} e^{-i T\left(E-E_{0}\right)} d P(E)|0\rangle
$$

If we now give $T$ a negative imaginary part, $T \rightarrow T(1-i \epsilon)$, and send it to infinity, the second term vanishes and we obtain

$$
|\Omega\rangle=\lim _{T \rightarrow \infty(1-i \epsilon)} \frac{e^{-i T\left(H-E_{0}\right)}|0\rangle}{\langle\Omega \mid 0\rangle}=\lim _{T \rightarrow \infty(1-i \epsilon)} \frac{e^{i E_{0} T} V(0,-T)|0\rangle}{\langle\Omega \mid 0\rangle},
$$

since $e^{i H_{0} T}|0\rangle=|0\rangle$. At this point $\epsilon$ is unrestricted, but in what follows we will need to let it tend to zero, and one should think of it intuitively as being infinitesimal.

We are now ready to express vacuum expectation values of products of interacting fields in terms of quantities defined in terms of free fields. For any $x_{1}, \ldots, x_{n} \in \mathbb{R}^{4}$, let $t_{j}=\left(x_{j}\right)^{0}$; then by (6.91) we have

$$
\phi\left(x_{1}\right) \cdots \phi\left(x_{n}\right)=V\left(0, t_{1}\right) \phi_{0}\left(x_{1}\right) V\left(t_{1}, t_{2}\right) \phi_{0}\left(x_{2}\right) \cdots V\left(t_{n-1}, t_{n}\right) \phi_{0}\left(x_{n}\right) V\left(t_{n}, 0\right)
$$

and hence by (6.93),

$$
\begin{array}{r}
\langle\Omega| \phi\left(x_{1}\right) \cdots \phi\left(x_{n}\right)|\Omega\rangle=\lim _{T \rightarrow \infty(1-i \epsilon)} \frac{e^{2 i E_{0} T}}{|\langle\Omega \mid 0\rangle|^{2}}\langle 0| V\left(T, t_{1}\right) \phi_{0}\left(x_{1}\right) V\left(t_{1}, t_{2}\right) \phi_{0}\left(x_{2}\right) \cdots \\
\cdots V\left(t_{n-1}, t_{n}\right) \phi_{0}\left(x_{n}\right) V\left(t_{n},-T\right)|0\rangle
\end{array}
$$

We can get rid of the unpleasant numerical factor in front by observing that the same equation holds with $n=0$ :

$$
1=\langle\Omega \mid \Omega\rangle=\lim _{T \rightarrow \infty(1-i \epsilon)} \frac{e^{2 i E_{0} T}}{|\langle\Omega \mid 0\rangle|^{2}}\langle 0| V(T,-T)|0\rangle
$$

Taking the quotient of these two equalities yields

$$
\begin{aligned}
& \langle\Omega| \phi\left(x_{1}\right) \cdots \phi\left(x_{n}\right)|\Omega\rangle \\
= & \lim _{T \rightarrow \infty(1-i \epsilon)} \frac{\langle 0| V\left(T, t_{1}\right) \phi_{0}\left(x_{1}\right) V\left(t_{1}, t_{2}\right) \phi_{0}\left(x_{2}\right) \cdots V\left(t_{n-1}, t_{n}\right) \phi_{0}\left(x_{n}\right) V\left(t_{n},-T\right)|0\rangle}{\langle 0| V(T,-T)|0\rangle}
\end{aligned}
$$

Finally, suppose that $x_{1}, \ldots, x_{n}$ are time-ordered: $t_{1} \geq \cdots \geq t_{n}$. Then, for $T$ sufficiently large, the time parameters in this expression all decrease from left to right (where the ordering refers to the real parts of $T$ and $-T$ ), so we can expand\\
the $V$ 's in their Dyson series (6.92). Since $\left(\exp \int_{a}^{b}\right)\left(\exp \int_{b}^{c}\right)=\exp \int_{a}^{c}$, the integrals in these series combine to yield an elegant final result:

$$
\begin{aligned}
& \langle\Omega| \mathcal{T}\left[\phi\left(x_{1}\right) \cdots \phi\left(x_{n}\right)\right]|\Omega\rangle \\
& \quad=\lim _{T \rightarrow \infty(1-i \epsilon)} \frac{\langle 0| \mathcal{T}\left[\phi_{0}\left(x_{1}\right) \cdots \phi_{0}\left(x_{n}\right) \exp \left(-i \int_{-T}^{T} H_{I}(\tau) d \tau\right)\right]|0\rangle}{\langle 0| \mathcal{T}\left[\exp \left(-i \int_{-T}^{T} H_{I}(\tau) d \tau\right)\right]|0\rangle}
\end{aligned}
$$

To be perfectly clear about the meaning of this: the numerator on the right side of (6.94) is the series whose $k$ th term is

$$
\langle 0| \frac{(-i)^{k}}{k!} \int_{-T}^{T} \cdots \int_{-T}^{T} \mathcal{T}\left[\phi_{0}\left(x_{1}\right) \cdots \phi_{0}\left(x_{n}\right) H_{I}\left(\tau_{1}\right) \cdots H_{I}\left(\tau_{k}\right)\right] d \tau_{1} \cdots d \tau_{k}|0\rangle,
$$

and similarly for the denominator. These series are to be interpreted perturbationtheoretically: that is, one uses the finite partial sums to obtain approximate results, in a way that we shall explain shortly, without worrying about the convergence of the entire series. The only slightly mysterious thing is the $i \epsilon$, and we shall shed some more light on it in due course.

We considered the simplest self-interacting scalar field theory in deriving (6.94) only in order not to clutter the picture up with complicated notation, but the result is valid in great generality for quantum field theories involving particles of arbitrary spin whose interaction is given by a sum of products of field operators, as our calculations did not depend on the specific properties of the fields or the interaction in any essential way. The $\phi\left(x_{j}\right)$ need not all come from the same field but can be various components of the various different fields in the theory, and $H_{I}(\tau)$ is the interaction-picture Hamiltonian given by the first equation in (6.90), the expression on the right of the second equation being replaced by the appropriate interaction for the theory in question. One merely has to exercise care, as always, to introduce appropriate minus signs when time-ordering the Fermion fields. (It is always legitimate to assemble all the $H_{I}$ factors on the right, however, as they each contain even numbers of Fermionic operators.) Let us restate (6.94) with the minor notational change necessary to indicate this generality:

$$
\begin{aligned}
& \langle\Omega| \mathcal{T}\left[\phi_{1}\left(x_{1}\right) \cdots \phi_{n}\left(x_{n}\right)\right]|\Omega\rangle \\
& \quad=\lim _{T \rightarrow \infty(1-i \epsilon)} \frac{\langle 0| \mathcal{T}\left[\phi_{01}\left(x_{1}\right) \cdots \phi_{0 n}\left(x_{n}\right) \exp \left(-i \int_{-T}^{T} H_{I}(\tau) d \tau\right)\right]|0\rangle}{\langle 0| \mathcal{T}\left[\exp \left(-i \int_{-T}^{T} H_{I}(\tau) d \tau\right)\right]|0\rangle}
\end{aligned}
$$

This result is known as the Gell-Mann-Low formula; it was first derived in [53].\\
There are several things that need to be said about the Gell-Mann-Low formula. First, if one writes the interaction Hamiltonian as the integral of a Hamiltonian density, $H_{I}(t)=\int \mathcal{H}_{I}(t, \mathbf{x}) d^{3} \mathbf{x}$, and then lets $T \rightarrow \infty$ (ignoring the $i \epsilon$ for the time being), the typical term (6.95) in the numerator becomes

$$
\begin{aligned}
& (6.97) \quad I_{k}\left(x_{1}, \ldots, x_{n}\right) \\
& =\langle 0| \frac{(-i)^{k}}{k!} \int_{\mathbb{R}^{4}} \cdots \int_{\mathbb{R}^{4}} \mathcal{T}\left[\phi_{01}\left(x_{1}\right) \cdots \phi_{0 n}\left(x_{n}\right) \mathcal{H}_{I}\left(y_{1}\right) \cdots \mathcal{H}_{I}\left(y_{k}\right)\right] d^{4} y_{1} \cdots d^{4} y_{k}|0\rangle
\end{aligned}
$$

This is exactly the sort of thing we had to evaluate in computing S-matrix elements in §6.4, except that the creation and annihilation operators for initial and final\\
particles have been replaced by the free-field operators $\phi_{0 j}\left(x_{j}\right)$. Hence the value of (6.97) is computed just as before, as an integral of a sum of products of contractions of free-field operators, and it can be represented by a position-space Feynman diagram just as before. The only difference is that the external vertices are labeled by the points $x_{i}$ and the external lines represent propagators. Conversely, given such a diagram with internal vertices labeled by $y_{j}$ and external vertices labeled by $x_{i}$, one obtains a contribution to (6.97) that is an integral over the internal variables $y_{j}$ of a product of propagators $-i \Delta_{j_{1} j_{2}}\left(y_{j_{1}}-y_{j_{2}}\right)$ and $-i \Delta_{k}\left(x_{k}-y_{j}\right)$. (Note for the external lines one indeed has $\Delta_{k}\left(x_{k}-y_{j}\right)$ and not $\Delta_{k}\left(y_{j}-x_{k}\right)$, because the $\phi\left(x_{k}\right)$ in (6.96) are all on the left.)

It is more convenient to pass to the momentum-space representation, obtained by writing each propagator as a Fourier integral. However, since the external vertices here also involve propagators, one must pass to their momentum-space representation too, that is, to consider the Fourier transform

$$
\widehat{I}_{k}\left(p_{1}, \ldots, p_{n}\right)=\int \cdots \int I_{k}\left(x_{1}, \ldots, x_{n}\right) e^{i \sum p_{j \mu} x_{j}^{\mu}} d^{4} x_{1} \cdots d^{4} x_{n}
$$

Writing each propagator in the form

$$
-i \Delta(z)=-i \int \widehat{\Delta}(q) e^{-i q_{\mu} z^{\mu}} \frac{d^{4} q}{(2 \pi)^{4}} \quad\left(z=x_{k}-y_{j} \text { or } y_{j_{1}}-y_{j_{2}}\right)
$$

we obtain an integral over all position variables $x_{i}$ and $y_{j}$ and all internal momenta $q$ of exponentials times Fourier-transformed propagators. The position variables appear only in the exponentials, so the integration over them reduces to integrals of the form $\int e^{-i q_{\mu} z^{\mu}} d^{4} q /(2 \pi)^{4}$, where $z$ is one of the vertices and $q$ is the algebraic sum of the momenta at that vertex; these give the momentum-conservation deltafunctions as before. The $p_{j}$ are the momenta associated to the external lines, which here are arbitrary (not on mass shell), and the Fourier variables $q$ are the momenta associated to the internal lines. Note that according to the remark at the end of the preceding paragraph and the conventions we established in §6.7, the external momenta always flow toward the external vertices.

This is the right place to account for the mysterious $i \epsilon$ in (6.96). It will be convenient to write $e^{-i \epsilon}$ instead of $1-i \epsilon$; then the position-space integrals of the preceding paragraph are not really $\int e^{-i q_{\mu} z^{\mu}} d^{4} q /(2 \pi)^{4}$ but

$$
\lim _{T \rightarrow \infty e^{-i \epsilon}} \int_{\mathbb{R}^{3}} \int_{-T}^{T} e^{-i p_{0} z^{0}+i \mathbf{p} \cdot \mathbf{z}} \frac{d z^{0} d \mathbf{z}}{(2 \pi)^{4}}
$$

If $p_{0}$ is real, this is not good: the integrand blows up exponentially at one end or the other, and even when one is doing informal distribution theory this is a disaster. But if we replace $p_{0}$ by $e^{i \epsilon} p_{0}$, the $e^{i \epsilon}$ 's cancel and we obtain the purely oscillatory integral whose value is the delta-function we want. This modification of $p_{0}$ is one way to produce the correct $i \epsilon$ in the denominator of the Fourier-transformed propagator - we shall explain this in detail in $\oint 7.1$ - so in fact the $i \epsilon$ fits into the picture perfectly.

We shall say more shortly about the relation between the Gell-Mann-Low formula and the calculation of S-matrix elements, but first we pause to explain the role of the denominator. The denominator in (6.96) is of the same form as the numerator, but without the factors $\phi_{0 j}\left(x_{j}\right)$; thus, it is the sum of the values of all the Feynman diagrams with no external vertices. Such diagrams represent fluctuations\\
of the vacuum, and we shall refer to them as vacuum bubbles. We encountered one such bubble for the $\phi^{4}$ scalar field theory in Figure 6.4c; a few of them for QED are shown in Figure 6.9.

\begin{figure}[h]
\begin{center}
  \includegraphics[width=\textwidth]{2025_11_18_7960b2c161913da1e816g-197}

\caption{Figure 6.9. Some vacuum bubbles in QED.}
\end{center}
\end{figure}

Let $\left\{B_{i}\right\}_{1}^{\infty}$ be a list of all the connected vacuum bubbles for the theory in question. Then every vacuum bubble has the form $B=\sum_{i=1}^{\infty} n_{i} B_{i}$ for some nonnegative integers $n_{i}$ of which all but finitely many are zero; that is, $B$ is the union of $n_{1}$ copies of $B_{1}, n_{2}$ copies of $B_{2}$, etc. If $V$ and $V_{i}$ denote the values of $B$ and $B_{i}$, respectively, we have

$$
V=\prod_{i=1}^{\infty} \frac{1}{n_{i}!} V_{i}^{n_{i}} .
$$

The $\prod n_{i}!$ is the symmetry factor of the bubble; it is there because permutation of the $n_{i}$ copies of $B_{i}$ does not change $B$. The denominator of (6.96) is the sum of all these quantities, namely,

$$
\langle 0| \mathcal{T}\left[\exp \left(-i \int \mathcal{H}_{I}(y) d^{4} y\right)\right]|0\rangle=\sum_{n_{1}, n_{2}, \ldots \geq 0, n_{1}+n_{2}+\cdots<\infty} \prod_{i=1}^{\infty} \frac{1}{n_{i}!} V_{i}^{n_{i}} .
$$

If we provisionally assume that $\sum V_{i}<\infty$, so that $V_{i} \rightarrow 0$ in particular, we can drop the restriction that $n_{1}+n_{2}+\cdots<\infty$, since $\prod_{i=1}^{\infty} V_{i}^{n_{i}} / n_{i}!=0$ whenever infinitely many $n_{i}$ are nonzero. In this case we have

$$
\langle 0| \mathcal{T}\left[\exp \left(-i \int \mathcal{H}_{I}(y) d^{4} y\right)\right]|0\rangle=\prod_{i=1}^{\infty} \sum_{n_{i}=0}^{\infty} \frac{1}{n_{i}!} V_{i}^{n_{i}}=\exp \left(\sum_{i=1}^{\infty} V_{i}\right) .
$$

Now, in fact, the values $V_{i}$ are given by integrals that are generally divergent. However, there are regularization procedures, which will be discussed in the next chapter, for expressing the $V_{i}$ in a systematic way as limits of finite integrals, and one should replace the $V_{i}$ by these regularized finite values. Having done this, however, there is still no reason for the series $\sum V_{i}$ to converge. Rather, both sides of (6.98) should be interpreted in perturbation theory. One thinks of $\mathcal{H}_{I}$ as containing a coefficient $\lambda$; then the $V_{i}$ also contain various powers of $\lambda$, and one drops all terms containing powers of $\lambda$ higher than $N$ for a fixed but arbitrary $N$. Then (6.98) is valid with this interpretation, and the extra terms we inserted (with infinitely many $n_{i} \neq 0$ ) have no effect because they are of infinite order in $\lambda$. (This may seem very murky, but don't give up just yet!)

Now consider the Feynman diagrams contributing to the numerator of (6.96). In general they are disconnected, and many of them contain vacuum bubbles. Indeed, each diagram can be written in the form $D+\sum n_{i} B_{i}$ where $D$ is a diagram\\
with no vacuum bubbles (i.e., such that each connected component contains some external vertices). The value of this diagram is the value of $D$ times the value of the bubble $\sum n_{i} B_{i}$, so by the preceding calculation, the sum of these values, over all possible bubbles, is the value of $D$ times $\exp \left(\sum V_{i}\right)$. Hence, when one divides by the denominator (6.98), the values of all the vacuum bubbles (whatever devices one uses to interpret them) simply cancel out. In short:

The right side of the Gell-Mann-Low formula (6.96) is the sum of all the position-space integrals represented by Feynman diagrams that contain no vacuum bubbles and that possess $n$ external lines corresponding to the fields $\phi_{1}, \ldots, \phi_{n}$, in which the $j$ th external line $L_{j}$ represents the propagator $-i \Delta_{j}\left(x_{j}-y_{j}\right)$. (Here $\Delta_{j}$ is the propagator for $\phi_{j}$ and $y_{j}$ is the variable of integration that labels the vertex at the internal end of $L_{j}$ ).

It is frequently more convenient to pass to the momentum space representation by taking the Fourier transform. Thus, we set

$$
W\left(p_{1}, \ldots, p_{n}\right)=\int \cdots \int\langle\Omega| \mathcal{T}\left[\phi_{1}\left(x_{1}\right) \cdots \phi_{n}\left(x_{n}\right)\right]|\Omega\rangle e^{i p_{1 \mu} x_{1}^{\mu}+\cdots+i p_{n \mu} x_{n}^{\mu}} d^{4} x_{1} \cdots d^{4} x_{n}
$$

In view of the remarks in the second paragraph following (6.96), the preceding result can then be restated as follows:\\
$W\left(p_{1}, \ldots, p_{n}\right)$ is the sum of all the momentum-space integrals represented by Feynman diagrams that contain no vacuum bubbles and that possess $n$ external lines corresponding to the fields $\phi_{1}, \ldots, \phi_{n}$, in which the $j$ th external line represents the Fourier-transformed propagator $-i \widehat{\Delta}_{j}\left(p_{j}\right)$ and $p_{j}$ flows toward the external vertex.

The overall momentum-conservation delta-function that appears in all Feynman diagrams comes out of the translation invariance of the vacuum expectation values. That is, the translation-invariance of the vacuum $|\Omega\rangle$ implies that

$$
\langle\Omega| \mathcal{T}\left[\phi_{1}\left(x_{1}\right) \cdots \phi_{n}\left(x_{n}\right)\right]|\Omega\rangle=\langle\Omega| \mathcal{T}\left[\phi_{1}(0) \phi_{2}\left(x_{2}-x_{1}\right) \cdots \phi_{n}\left(x_{n}-x_{1}\right)\right]|\Omega\rangle .
$$

Substituting this into (6.99) and then making the change of variable $x_{j}=x_{j}^{\prime}+x_{1}$ turns the integrand of (6.99) into a function whose only dependence on $x_{1}$ is a factor $e^{i\left(\sum p_{j}\right)_{\mu} x_{1}^{\mu}}$; integration over $x_{1}$ then gives $(2 \pi)^{4} \delta\left(\sum p_{j}\right)$.

We now return to the relation between vacuum expectation values of products of fields and S -matrix elements, that is, the problem of calculating

$$
\langle 0| a_{1}\left(\mathbf{p}_{1}^{\text {out }}\right) \cdots a_{J}\left(\mathbf{p}_{J}^{\text {out }}\right) S a_{1}^{\dagger}\left(\mathbf{p}_{1}^{\text {in }}\right) \cdots a_{K}^{\dagger}\left(\mathbf{p}_{K}^{\text {in }}\right)|0\rangle
$$

from

$$
\langle\Omega| \mathcal{T}\left[\phi_{1}\left(x_{1}\right) \cdots \phi_{n}\left(x_{n}\right)\right]|\Omega\rangle .
$$

We shall assume, as we did in §6.8, that no proper subset of $\left\{\mathbf{p}_{1}^{\text {in }}, \ldots, \mathbf{p}_{K}^{\text {in }}\right\}$ has the same sum as any proper subset of $\left\{\mathbf{p}_{1}^{\text {out }}, \ldots, \mathbf{p}_{J}^{\text {out }}\right\}$, so that the only Feynman diagrams that contribute to (6.101) are connected ones. (The analysis in the general case is similar but a little more complicated.) The first step is to convert the position dependence of (6.102) into momentum dependence by taking the Fourier transform (6.99). We are going to take $n=J+K$ and the $p_{j}$ to be, essentially, the on-massshell 4 -momenta corresponding to the 3 -momenta in (6.101), but there has to be a twist because the momenta in (6.99) all flow toward the external vertices, whereas the momentum $\mathbf{p}_{k}^{\text {in }}$ in (6.101) should flow away from its initial vertex. (Recall (6.43),\\
which shows that an incoming momentum $\mathbf{p}$ contributes a factor $e^{-i p_{\mu} x^{\mu}}$ to the position-space integral for the S -matrix element, whereas an ougtoing momentum contributes a factor of $e^{i p_{\mu} x^{\mu}}$; the sign is important because $p_{0}=\omega_{\mathbf{p}}>0$ in either case.) Thus, to get the right momentum dependence we must consider

$$
W\left(p_{1}, \ldots, p_{J},-p_{J+1}, \ldots,-p_{J+K}\right) \quad\left(p_{l}=\left(\omega_{\mathbf{p}_{l}}, \mathbf{p}_{l}\right)\right),
$$

where

$$
\mathbf{p}_{j}=\mathbf{p}_{j}^{\text {out }} \text { for } 1 \leq j \leq J, \quad \mathbf{p}_{J+k}=\mathbf{p}_{k}^{\text {in }} \text { for } 1 \leq k \leq K .
$$

(Note that the momentum-conservation delta-function herein is $\delta\left(\sum_{j} p_{j}-\sum_{k} p_{J+k}\right)$, as it should be.)

The preceding discussion shows that this quantity is the sum of the values of all connected Feynman diagrams with external lines labeled by the momenta $\mathbf{p}_{k}^{\text {out }}, \mathbf{p}_{j}^{\text {in }}$ and representing the Fourier-transformed propagators $-i \widehat{\Delta}_{j}\left(p_{j}^{\text {out }}\right),-i \widehat{\Delta}_{k}\left(-p_{k}^{\text {in }}\right)$. On the other hand, (6.101) is the sum of the values of all Feynman diagrams of the same sort in which external lines represent the appropriate coefficient functions $u_{k}\left(\mathbf{p}_{k}^{\text {in }}\right)$ and $u_{j}^{*}\left(\mathbf{p}_{j}^{\text {out }}\right)$. Thus, to get from (6.99) to (6.101) one has merely to replace the external propagators by coefficient functions. The result is

$$
\begin{aligned}
& \langle 0| a_{1}\left(\mathbf{p}_{1}^{\text {out }}\right) \cdots a_{J}\left(\mathbf{p}_{J}^{\text {out }}\right) S a_{1}^{\dagger}\left(\mathbf{p}_{1}^{\text {in }}\right) \cdots a_{K}^{\dagger}\left(\mathbf{p}_{K}^{\text {in }}\right)|0\rangle \\
& =i^{J+K} \prod_{j=1}^{J} \frac{u_{j}^{*}\left(\mathbf{p}_{j}^{\text {out }}\right)}{\widehat{\Delta}_{j}\left(p_{j}^{\text {out }}\right)} \prod_{k=1}^{K} \frac{u_{k}\left(\mathbf{p}_{k}^{\text {in }}\right)}{\widehat{\Delta}_{k}\left(-p_{k}^{\text {in }}\right)} W\left(p_{1}^{\text {out }}, \ldots, p_{J}^{\text {out }},-p_{1}^{\text {in }}, \ldots,-p_{K}^{\text {in }}\right),
\end{aligned}
$$

where $W$ is as in (6.99). This is known as the Lehmann-Symanzik-Zimmermann reduction formula, or LSZ formula for short.

For a scalar field, the function $\Delta$ is even, and it is a fundamental solution for the Klein-Gordon operator $\partial^{2}+m^{2}$, so dividing by $\widehat{\Delta}$ in momentum space is the same as applying the operator $\partial^{2}+m^{2}$ in position space. Thus, for scalar fields the LSZ formula can be restated as

$$
\begin{aligned}
& \langle 0| a_{1}\left(\mathbf{p}_{1}^{\text {out }}\right) \cdots a_{J}\left(\mathbf{p}_{J}^{\text {out }}\right) S a_{1}^{\dagger}\left(\mathbf{p}_{1}^{\text {in }}\right) \cdots a_{K}^{\dagger}\left(\mathbf{p}_{K}^{\text {in }}\right)|0\rangle \\
= & i^{J+K} \prod_{j, k} u_{j}^{*}\left(p_{j}^{\text {out }}\right) u_{k}\left(p_{k}^{\text {in }}\right) \int \cdots \int \exp \left[i \sum_{j} p_{j \mu}^{\text {out }} x^{j \mu}-i \sum_{k} p_{k \mu}^{\text {in }} x_{J+k}^{\mu}\right] \times \\
& \left(\partial_{x_{1}}^{2}+m^{2}\right) \cdots\left(\partial_{x_{J+K}}^{2}+m^{2}\right)\langle\Omega| \mathcal{T}\left[\phi_{1}\left(x_{1}\right) \ldots \phi_{J+K}\left(x_{J+K}\right)\right]|\Omega\rangle d^{4} x_{1} \cdots d^{4} x_{J+K} .
\end{aligned}
$$

Similarly, the Dirac propagator is the fundamental solution for the Dirac operator $-i \gamma^{\mu} \partial_{\mu}+m$, so when spin- $\frac{1}{2}$ fields are involved one obtains a similar formula with $\partial_{x_{j}}^{2}+m^{2}$ replaced by $-i \gamma^{\mu}\left(\partial_{x_{j}}\right)_{\mu}+m$ and $\partial_{x_{J+k}}^{2}+m^{2}$ replaced by $i \gamma^{\mu}\left(\partial_{x_{J+k}}\right)_{\mu}+m$.

There is one technicality that must be mentioned here. As we shall discuss in the next chapter, to carry out the evaluation of Feynman diagrams involving loops, it is necessary to renormalize the field operators, that is, to multiply the field $\phi$ by a constant that is conventionally called $Z^{-1 / 2}$. (If the theory involves several different fields, each one has its own renormalization factor.) The fields $\phi_{j}$ in (6.99) and hence in the LSZ formula (6.103) must be taken to be renormalized fields. Otherwise, the appropriate $Z$ factors must be included in the formula.

We shall content ourselves with an informal qualitative explanation of this fact. The Feynman-diagram picture suggests that particles propagate from one vertex to the next without anything happening in between, but in fact they have encounters with virtual particles along the way. As we shall see in § 7.6, this effectively turns\\
the $1 /\left(-p^{2}+m^{2}\right)$ in the propagator into $Z /\left(-p^{2}+(m+\delta m)^{2}\right)$ for some constants $Z$ and $\delta m$. The $\delta m$ is a mass shift of the sort we encountered in §6.2, but the $Z$ must be compensated for by renormalizing the field. The point is that the fields in the contractions that give rise to the external propagators in (6.102) must also be adjusted in this way, and this gives rise to the extra factors of $Z$.

We derived the LSZ formula on a diagram-by-diagram basis, that is, in perturbation theory. The physics literature also contains a number of (nonrigorous) nonperturbative derivations of it; for example, Itzykson and Zuber [66], §5-1-3, Peskin and Schroeder [89], §7.2, and Weinberg [131], §10.3. In addition, it is possible to prove the LSZ formula rigorously within the setting of the Wightman axioms and the Haag-Ruelle scattering theory; see Bogolubov et al. [11], [12], and Araki [3].



\pagebreak
\printbibliography
\end{document}